\documentclass{report}

\usepackage[utf8]{inputenc}
\usepackage{eumat}
\usepackage[Conny]{fncychap}
\usepackage[bahasa]{babel}

% Rename Contents
\addto\captionsenglish{\renewcommand{\contentsname}{\vspace{-0.5cm} \textbf{Daftar Isi} \vspace{-2cm}}}

\begin{document}

% Cover Page
\begin{titlepage}
    \begin{center}
        \vspace*{0,2cm}

        \Huge
        \textbf{TUGAS PEKAN 16 APLIKASI KOMPUTER}
        
        \vspace{1cm}
        
        \LARGE
        BAB LaTeX dan Markdown  
        
        \vspace{1cm}
        
        \includegraphics[width=0.5\textwidth]{Logo UNY.png}

        \vspace{1cm}
        
        \textbf{Shalih Abdillah}\\
        22305144009\\
        Matematika B 2022
        
        \vspace{2cm}
        
        \Large
        \textbf{PRODI MATEMATIKA}\\
        \textbf{DEPARTEMEN PENDIDIKAN MATEMATIKA}\\
        \textbf{FAKULTAS MATEMATIKA DAN ILMU PENGETAHUAN ALAM}
        \textbf{UNIVERSITAS NEGERI YOGYAKARTA}\\
        \textbf{2023}
        
    \end{center}
\end{titlepage}

\newpage
\tableofcontents

\chapter{KB Pekan 2 : Belajar Menggunakan Software EMT}
\input{Pertemuan 2/Shalih Abdillah_22305144009_ApKom1}

\newpage
\chapter{KB Pekan 3 : Menggunakan EMT untuk menyelesaikan masalah-masalah Aljabar}
\input{Pertemuan 4/Shalih Abdillah_22305144009_ApKom2}

\newpage
\chapter{KB Pekan 4: Menggunakan EMT untuk mengambar grafik 2 dimensi (2D)}
\input{Pertemuan 6/Shalih Abdillah_22305144009_ApKom3}

\newpage
\chapter{KB Pekan 5: Menggunakan EMT untuk mengambar grafik 3 dimensi (3D)}
\documentclass{article}

\usepackage{eumat}

\begin{eulernotebook}
\eulersubheading{Shalih Abdillah(22305144009) }
\eulerheading{Menggambar Plot 3D dengan EMT}
\begin{eulercomment}
Bagian ini merupakan pengenalan plot 3d pada Euler. Kita memerlukan
plot 3D untuk memvisualisasikan dari suatu fungsi dua variabel.

Euler menggambar fungsi tersebut menggunakan algoritma pengurutan
untuk menyembunyikan bagian di latar belakang. Secara umum, Euler
menggunakan proyeksi pusat. Standarnya adalah dari kuadran x-y positif
menuju titik asal x=y=z=0, tetapi sudut=0° terlihat dari arah sumbu y.
Sudut pandang dan tinggi dapat diubah.

Euler dapat merencanakan

- permukaan dengan bayangan dan garis level atau rentang level,\\
- awan poin,\\
- kurva parametrik,\\
- permukaan implisit.

Plot 3D dari suatu fungsi menggunakan plot3d. Cara termudah adalah
dengan memplot ekspresi dalam x dan y. Parameter r mengatur kisaran
plot di sekitar (0,0).
\end{eulercomment}
\begin{eulerprompt}
>aspect(1); plot3d("x^2+sin(y)",r=pi):
\end{eulerprompt}
\eulerimg{27}{images/Shalih Abdillah_22305144009_ApKom4-001.png}
\begin{eulerprompt}
>aspect(3,2); plot3d("2*x^3+sin(y)",r=pi):
\end{eulerprompt}
\eulerimg{17}{images/Shalih Abdillah_22305144009_ApKom4-002.png}
\begin{eulercomment}
Dibawah ini percobaan dari beberapa perintah dengan mengganti salah
satu dari variabelnya.
\end{eulercomment}
\begin{eulerprompt}
>aspect(1.5); plot3d("x^2+sin(y)",r=pi):
\end{eulerprompt}
\eulerimg{17}{images/Shalih Abdillah_22305144009_ApKom4-003.png}
\begin{eulerprompt}
>aspect(2.5); plot3d("x^2+sin(y)",r=pi):
\end{eulerprompt}
\eulerimg{10}{images/Shalih Abdillah_22305144009_ApKom4-004.png}
\begin{eulercomment}
Dari berbagai percobaan diatas, dapat dilihat bahwa semakin besar
nilai aspect-nya, hasil plot 3D semakin kecil.\\
Aspect ialah proporsi antara sumbu x dan sumbu y. Artinya jika aspect
1 maka proporsi x dan y sama.

Selanjutnya akan dilihat perbedaan hasil grafik dengan parameter yang
berubah adalah r.
\end{eulercomment}
\begin{eulerprompt}
>aspect(1); plot3d("x^2+sin(y)",r=100):
\end{eulerprompt}
\eulerimg{27}{images/Shalih Abdillah_22305144009_ApKom4-005.png}
\begin{eulerprompt}
>aspect(1); plot3d("x^2+sin(y)",r=pi):
\end{eulerprompt}
\eulerimg{27}{images/Shalih Abdillah_22305144009_ApKom4-006.png}
\begin{eulerprompt}
>aspect(1); plot3d("x^2+sin(y)",r=2):
\end{eulerprompt}
\eulerimg{27}{images/Shalih Abdillah_22305144009_ApKom4-007.png}
\begin{eulerprompt}
>aspect(1); plot3d("x^2+sin(y)",r=1):
\end{eulerprompt}
\eulerimg{27}{images/Shalih Abdillah_22305144009_ApKom4-008.png}
\begin{eulercomment}
Berbagai perbedaan hasil yang diperoleh ketika nilai r nya dirubah.
Semakin besar nilai r nya plot yag dihasilkan semakin kecil sehingga
ketelitiannya semakin kecil. Di sisi lain jika nilai r yang ditetapkan
terlalu kecil, plot yang dihasilkan juga kurang baik karena belum
dapat mewakili plot secara garis beras. Darisini diperoleh informasi
bahwasanya pemilihan nilai r haruslah tepat agar plot yag dihasilkan
lebih akurat dan memahamkan, sesuai yag kita butuhkan.


\begin{eulercomment}
\eulerheading{Fungsi Dua Variabel}
\begin{eulercomment}
Untuk grafik fungsi, gunakan

- ekspresi sederhana dalam x dan y,\\
- nama fungsi dari dua variabell\\
- atau matriks data.

Standarnya adalah kotak kawat yang diisi dengan warna berbeda di kedua
sisi. Perhatikan bahwa jumlah default interval grid adalah 10, tetapi
plot menggunakan jumlah default 40x40 persegi panjang untuk membangun
permukaan. Ini bisa diubah.

- n=40, n=[40,40]: jumlah garis grid di setiap arah\\
- grid=10, grid=[10,10]: jumlah garis grid di setiap arah.

Kami menggunakan default n=40 dan grid=10.
\end{eulercomment}
\begin{eulerprompt}
>plot3d("x^2+y^2"):
\end{eulerprompt}
\eulerimg{27}{images/Shalih Abdillah_22305144009_ApKom4-009.png}
\begin{eulerprompt}
>plot3d("x^2+cos(x)^5"):
\end{eulerprompt}
\eulerimg{17}{images/Shalih Abdillah_22305144009_ApKom4-010.png}
\begin{eulercomment}
Interaksi pengguna dimungkinkan dengan \textgreater{}parameter pengguna. Pengguna
dapat menekan tombol berikut.

- kiri, kanan, atas, bawah: putar sudut pandang\\
- +,-: memperbesar atau memperkecil\\
- a: menghasilkan anaglyph (lihat di bawah)\\
- l: beralih memutar sumber cahaya (lihat di bawah)\\
- spasi: reset ke default\\
- kembali: akhiri interaksi
\end{eulercomment}
\begin{eulerprompt}
>plot3d("exp(-x^2+y^2)",>user, ...
>  title="Turn with the vector keys (press return to finish)"):
\end{eulerprompt}
\eulerimg{27}{images/Shalih Abdillah_22305144009_ApKom4-011.png}
\begin{eulercomment}
Rentang plot untuk fungsi dapat ditentukan dengan

- a,b: rentang-x\\
- c,d: rentang-y\\
- r: persegi simetris di sekitar (0,0).\\
- n: jumlah subinterval untuk plot.

Ada beberapa parameter untuk menskalakan fungsi atau mengubah tampilan
grafik.

fscale: skala ke nilai fungsi (defaultnya adalah \textless{}fscale).\\
skala: angka atau vektor 1x2 untuk skala ke arah x dan y.\\
bingkai: jenis bingkai (default 1).

Catatan: Jika sudah menggunakan atau mendifinisikan nilai r maka tidak
perlyu mendifinisikan lagi nilai a,b,c,dan d nya.
\end{eulercomment}
\begin{eulerprompt}
>plot3d("exp(-(x^2+y^2)/5)",r=10,n=90,fscale=4,scale=1.5,frame=1):
\end{eulerprompt}
\eulerimg{27}{images/Shalih Abdillah_22305144009_ApKom4-012.png}
\begin{eulercomment}
Tampilan dapat diubah dengan berbagai cara.

- jarak: jarak pandang ke plot.\\
- zoom: nilai zoom.\\
- sudut: sudut terhadap sumbu y negatif dalam radian.\\
- tinggi: ketinggian tampilan dalam radian.

Nilai default dapat diperiksa atau diubah dengan fungsi view(). Ini
mengembalikan parameter dalam urutan di atas.
\end{eulercomment}
\begin{eulerprompt}
>view
\end{eulerprompt}
\begin{euleroutput}
  [5,  2.6,  2,  0.4]
\end{euleroutput}
\begin{eulercomment}
5 menunjukkan jarak pandang ke plot,\\
2.6 menunjukkan nilai zoom,\\
2 menunjukkan nilai sudut terhadap y negatif dalam radian, dan\\
0.4 menunjukkan ketinggian tampilan dalam radian.

Jarak yang lebih dekat membutuhkan lebih sedikit zoom. Efeknya lebih
seperti lensa sudut lebar.

Dalam contoh berikut, sudut=0 dan tinggi=0 terlihat dari sumbu y
negatif. Label sumbu untuk y disembunyikan dalam kasus ini.
\end{eulercomment}
\begin{eulerprompt}
>plot3d("x^2+y",distance=3,zoom=2,angle=0,height=0):
\end{eulerprompt}
\eulerimg{17}{images/Shalih Abdillah_22305144009_ApKom4-013.png}
\begin{eulerprompt}
>plot3d("x^5+y",distance=3,zoom=1.4,angle=1,height=10):
\end{eulerprompt}
\eulerimg{17}{images/Shalih Abdillah_22305144009_ApKom4-014.png}
\begin{eulercomment}
Plot terlihat selalu ke pusat kubus plot. Anda dapat memindahkan pusat
dengan parameter tengah.
\end{eulercomment}
\begin{eulerprompt}
>plot3d("x^4+y^2",a=0,b=1,c=-1,d=1,angle=-20°,height=20°, ...
>  center=[0.4,0,0],zoom=5):
\end{eulerprompt}
\eulerimg{27}{images/Shalih Abdillah_22305144009_ApKom4-015.png}
\begin{eulercomment}
Plot diskalakan agar sesuai dengan kubus satuan untuk dilihat. Jadi
tidak perlu mengubah jarak atau zoom tergantung pada ukuran plot.
Namun, label mengacu pada ukuran sebenarnya.

Jika Anda mematikannya dengan scale=false, Anda perlu berhati-hati,
bahwa plot masih cocok dengan jendela plot, dengan mengubah jarak
pandang atau zoom, dan memindahkan pusat.
\end{eulercomment}
\begin{eulerprompt}
>plot3d("5*exp(-x^2-y^2)",r=2,<fscale,<scale,distance=13,height=50°, ...
>  center=[0,0,-2],frame=3):
\end{eulerprompt}
\eulerimg{27}{images/Shalih Abdillah_22305144009_ApKom4-016.png}
\begin{eulerprompt}
>plot3d("5*exp(cos(x)^2-y^2)",r=2,<fscale,<scale,distance=13,height=50°, ...
>  center=[0,0,-2],frame=3):
\end{eulerprompt}
\eulerimg{27}{images/Shalih Abdillah_22305144009_ApKom4-017.png}
\begin{eulercomment}
Sebuah plot kutub juga tersedia. Parameter polar=true menggambar plot
polar. Fungsi tersebut harus tetap merupakan fungsi dari x dan y.
Parameter "fscale" menskalakan fungsi dengan skala sendiri. Jika
tidak, fungsi diskalakan agar sesuai dengan kubus.
\end{eulercomment}
\begin{eulerprompt}
>plot3d("1/(x^2+y^2+1)",r=5,>polar, ...
>fscale=2,>hue,n=100,zoom=4,>contour,color=blue):
\end{eulerprompt}
\eulerimg{27}{images/Shalih Abdillah_22305144009_ApKom4-018.png}
\begin{eulerprompt}
>plot3d("1/(x^2+y^2-30)",r=5,>polar, ...
>fscale=2,>hue,n=100,zoom=4,>contour,color=white):
\end{eulerprompt}
\eulerimg{27}{images/Shalih Abdillah_22305144009_ApKom4-019.png}
\begin{eulerprompt}
>function f(r) := exp(-r/2)*cos(r); ...
>plot3d("f(x^2+y^2)",>polar,scale=[1,1,0.4],r=2pi,frame=3,zoom=4):
\end{eulerprompt}
\eulerimg{27}{images/Shalih Abdillah_22305144009_ApKom4-020.png}
\begin{eulerprompt}
>function f(r) := exp(-r/2)*cos(r); ...
>plot3d("f(x^2+y^2)",>polar,scale=[1,1,0.4],r=2pi,frame=3,zoom=5,>contour,color=white):
\end{eulerprompt}
\eulerimg{27}{images/Shalih Abdillah_22305144009_ApKom4-021.png}
\begin{eulercomment}
Rotasi parameter memutar fungsi dalam x di sekitar sumbu x.

- rotate=1: Menggunakan sumbu x\\
- rotate=2: Menggunakan sumbu z
\end{eulercomment}
\begin{eulerprompt}
>plot3d("x^2+1",a=-1,b=1,rotate=true,grid=5):
\end{eulerprompt}
\eulerimg{27}{images/Shalih Abdillah_22305144009_ApKom4-022.png}
\begin{eulerprompt}
>plot3d("x^2+1",a=-1,b=1,rotate=false,grid=5):
\end{eulerprompt}
\eulerimg{27}{images/Shalih Abdillah_22305144009_ApKom4-023.png}
\begin{eulercomment}
Perintah rotate = false menghasilkan plot yang sama dengan perintah
tanpa rotate.
\end{eulercomment}
\begin{eulerprompt}
>plot3d("x^2+1",a=-1,b=1,grid=5):
\end{eulerprompt}
\eulerimg{27}{images/Shalih Abdillah_22305144009_ApKom4-024.png}
\begin{eulercomment}
DIbawah ini merupakan hasil plot dari dngsi yang di rotasi dengan
perintah rotate = 1 artinya dirotasi terhadap sumbu y, yang
menghasilkan dengan nilai yang sama dengan perintah rotate = true.
\end{eulercomment}
\begin{eulerprompt}
>plot3d("x^2+1",a=-1,b=1,rotate=1,grid=5):
\end{eulerprompt}
\eulerimg{27}{images/Shalih Abdillah_22305144009_ApKom4-025.png}
\begin{eulercomment}
DIbawah ini merupakan hasil plot dari dngsi yang di rotasi dengan
perintah rotate = 2 artinya dirotasi terhadap sumbu z.
\end{eulercomment}
\begin{eulerprompt}
>plot3d("x^2+1",a=-1,b=1,rotate=2,grid=5):
\end{eulerprompt}
\eulerimg{27}{images/Shalih Abdillah_22305144009_ApKom4-026.png}
\begin{eulercomment}
Dibawah ini merupakan plot dengan tiga fungsi.
\end{eulercomment}
\begin{eulerprompt}
>plot3d("x","x^2+y^2","y",r=2,zoom=3.5,frame=3):
\end{eulerprompt}
\eulerimg{27}{images/Shalih Abdillah_22305144009_ApKom4-027.png}
\begin{eulercomment}
Dibawah ini dilakukan percobaan untuk plot tiga fungsi, ternyata tidak
bisa dikenai perintah rotate (rotasi).
\end{eulercomment}
\begin{eulerprompt}
>plot3d("x","x^2+y^2","y",r=2,zoom=3.5,frame=3):
\end{eulerprompt}
\eulerimg{27}{images/Shalih Abdillah_22305144009_ApKom4-028.png}
\eulerheading{Plot Kontur}
\begin{eulercomment}
Untuk plot, Euler menambahkan garis grid. Sebagai gantinya
dimungkinkan untuk menggunakan garis level dan rona satu warna atau
rona berwarna spektral. Euler dapat menggambar tinggi fungsi pada plot
dengan bayangan. Di semua plot 3D, Euler dapat menghasilkan anaglyph
merah/sian.

-\textgreater{} hue: Menyalakan bayangan cahaya alih-alih kabel.\\
-\textgreater{} kontur: Memplot garis kontur otomatis pada plot.\\
- level=... (atau level): Sebuah vektor nilai untuk garis kontur.

Standarnya adalah level="auto", yang menghitung beberapa garis level
secara otomatis. Seperti yang Anda lihat di plot, level sebenarnya
adalah rentang level.

Gaya default dapat diubah. Untuk plot kontur berikut, kami menggunakan
grid yang lebih halus untuk 100x100 poin, skala fungsi dan plot, dan
menggunakan sudut pandang yang berbeda.
\end{eulercomment}
\begin{eulerprompt}
>plot3d("exp(-x^2-y^2)",r=2,n=100,level="thin", ...
> >contour,>spectral,fscale=1,scale=1.1,angle=45°,height=20°):
\end{eulerprompt}
\eulerimg{27}{images/Shalih Abdillah_22305144009_ApKom4-029.png}
\begin{eulerprompt}
>plot3d("exp(x*y)",angle=100°,>contour,color=green):
\end{eulerprompt}
\eulerimg{27}{images/Shalih Abdillah_22305144009_ApKom4-030.png}
\begin{eulercomment}
Bayangan default menggunakan warna abu-abu. Tetapi rentang warna
spektral juga tersedia.

-\textgreater{} spektral: Menggunakan skema spektral default\\
- color=...: Menggunakan warna khusus atau skema spektral

Untuk plot berikut, kami menggunakan skema spektral default dan
menambah jumlah titik untuk mendapatkan tampilan yang sangat halus.
\end{eulercomment}
\begin{eulerprompt}
>plot3d("x^2+y^2",>spectral,>contour,n=100):
\end{eulerprompt}
\eulerimg{27}{images/Shalih Abdillah_22305144009_ApKom4-031.png}
\begin{eulercomment}
Alih-alih garis level otomatis, kita juga dapat mengatur nilai garis
level. Ini akan menghasilkan garis level tipis alih-alih rentang
level.
\end{eulercomment}
\begin{eulerprompt}
>plot3d("x^2-y^2",0,1,0,1,angle=220°,level=-1:0.2:1,color=redgreen):
\end{eulerprompt}
\eulerimg{27}{images/Shalih Abdillah_22305144009_ApKom4-032.png}
\begin{eulercomment}
Dalam plot berikut, kami menggunakan dua pita level yang sangat luas
dari -0,1 hingga 1, dan dari 0,9 hingga 1. Ini dimasukkan sebagai
matriks dengan batas level sebagai kolom.

Selain itu, kami melapisi kisi dengan 10 interval di setiap arah.
\end{eulercomment}
\begin{eulerprompt}
>plot3d("x^2+y^3",level=[-0.1,0.9;0,1], ...
>  >spectral,angle=30°,grid=10,contourcolor=gray):
\end{eulerprompt}
\eulerimg{27}{images/Shalih Abdillah_22305144009_ApKom4-033.png}
\begin{eulercomment}
Dalam contoh berikut, kami memplot himpunan, di mana

\end{eulercomment}
\begin{eulerformula}
\[
f(x,y) = x^y-y^x = 0
\]
\end{eulerformula}
\begin{eulercomment}
Kami menggunakan satu garis tipis untuk garis level.
\end{eulercomment}
\begin{eulerprompt}
>plot3d("x^y-y^x",level=0,a=0,b=6,c=0,d=6,contourcolor=red,n=100):
\end{eulerprompt}
\eulerimg{27}{images/Shalih Abdillah_22305144009_ApKom4-034.png}
\begin{eulercomment}
Dimungkinkan untuk menunjukkan bidang kontur di bawah plot. Warna dan
jarak ke plot dapat ditentukan.
\end{eulercomment}
\begin{eulerprompt}
>plot3d("x^2+y^4",>cp,cpcolor=green,cpdelta=0.2):
\end{eulerprompt}
\eulerimg{27}{images/Shalih Abdillah_22305144009_ApKom4-035.png}
\begin{eulercomment}
Berikut adalah beberapa gaya lagi. Kami selalu mematikan frame, dan
menggunakan berbagai skema warna untuk plot dan grid.
\end{eulercomment}
\begin{eulerprompt}
>figure(2,2); ...
>expr="y^3-x^2"; ...
>figure(1);  ...
>  plot3d(expr,<frame,>cp,cpcolor=spectral); ...
>figure(2);  ...
>  plot3d(expr,<frame,>spectral,grid=10,cp=2); ...
>figure(3);  ...
>  plot3d(expr,<frame,>contour,color=gray,nc=5,cp=3,cpcolor=greenred); ...
>figure(4);  ...
>  plot3d(expr,<frame,>hue,grid=10,>transparent,>cp,cpcolor=gray); ...
>figure(0):
\end{eulerprompt}
\eulerimg{27}{images/Shalih Abdillah_22305144009_ApKom4-036.png}
\begin{eulercomment}
Ada beberapa skema spektral lainnya, bernomor dari 1 hingga 9. Tetapi
Anda juga dapat menggunakan warna=nilai, di mana nilai

- spektral: untuk rentang dari biru ke merah\\
- putih: untuk rentang yang lebih redup\\
-kuningbiru,ungu hijau,birukuning,hijaumerah\\
- biru kuning, hijau ungu, kuning biru, merah hijau
\end{eulercomment}
\begin{eulerprompt}
>figure(3,3); ...
>for i=1:9;  ...
>  figure(i); plot3d("x^2+y^2",spectral=i,>contour,>cp,<frame,zoom=4);  ...
>end; ...
>figure(0):
\end{eulerprompt}
\eulerimg{27}{images/Shalih Abdillah_22305144009_ApKom4-037.png}
\begin{eulercomment}
Sumber cahaya dapat diubah dengan l dan tombol kursor selama interaksi
pengguna. Itu juga dapat diatur dengan parameter.

- cahaya: arah untuk cahaya\\
- amb: cahaya sekitar antara 0 dan 1

Perhatikan bahwa program tidak membuat perbedaan antara sisi plot.
Tidak ada bayangan. Untuk ini, Anda perlu Povray.
\end{eulercomment}
\begin{eulerprompt}
>plot3d("-x^2-y^2", ...
>  hue=true,light=[0,1,1],amb=0,user=true, ...
>  title="Press l and cursor keys (return to exit)"):
\end{eulerprompt}
\eulerimg{27}{images/Shalih Abdillah_22305144009_ApKom4-038.png}
\begin{eulercomment}
Parameter warna mengubah warna permukaan. Warna garis level juga dapat
diubah.
\end{eulercomment}
\begin{eulerprompt}
>plot3d("-x^2-y^2",color=rgb(0.2,0.2,0),hue=true,frame=false, ...
>  zoom=3,contourcolor=red,level=-2:0.1:1,dl=0.01):
\end{eulerprompt}
\eulerimg{27}{images/Shalih Abdillah_22305144009_ApKom4-039.png}
\begin{eulercomment}
Penggantian perintah hue=true menjadi hue=false memberikan efek cahaya
pada bagian bawah plot. Plot dengan perintah hue=false memiliki efek
gradasi warna yang cenderung semakin gelap pada bagian bawah plot.
\end{eulercomment}
\begin{eulerprompt}
>plot3d("-x^2-y^2",color=rgb(0.2,0.2,0),hue=false,frame=false, ...
>  zoom=3,contourcolor=red,level=-2:0.1:1,dl=0.01):
\end{eulerprompt}
\eulerimg{27}{images/Shalih Abdillah_22305144009_ApKom4-040.png}
\begin{eulercomment}
Dibawah ini hasil dari penambahan perintah frame=true, sehingga plot
memiliki bingkai ketika diprint.
\end{eulercomment}
\begin{eulerprompt}
>plot3d("-x^2-y^2",color=rgb(0.2,0.2,0),hue=false,frame=true, ...
>  zoom=3,contourcolor=red,level=-2:0.1:1,dl=0.01):
\end{eulerprompt}
\eulerimg{27}{images/Shalih Abdillah_22305144009_ApKom4-041.png}
\begin{eulerprompt}
>plot3d("-x^2-y^2",color=rgb(0.2,0.2,0),hue=false,frame=false, ...
>  zoom=3,contourcolor=red,level=-2:0.1:1,dl=0.01):
\end{eulerprompt}
\eulerimg{27}{images/Shalih Abdillah_22305144009_ApKom4-042.png}
\begin{eulercomment}
Warna 0 memberikan efek pelangi khusus.
\end{eulercomment}
\begin{eulerprompt}
>plot3d("x^2/(x^2+y^2+1)",color=0,hue=true,grid=10):
\end{eulerprompt}
\eulerimg{27}{images/Shalih Abdillah_22305144009_ApKom4-043.png}
\begin{eulercomment}
Permukaannya juga bisa transparan.
\end{eulercomment}
\begin{eulerprompt}
>plot3d("x^2+y^2",>transparent,grid=10,wirecolor=red):
\end{eulerprompt}
\eulerimg{27}{images/Shalih Abdillah_22305144009_ApKom4-044.png}
\eulerheading{Plot Implisit}
\begin{eulercomment}
Ada juga plot implisit dalam tiga dimensi. Euler menghasilkan
pemotongan melalui objek. Fitur plot3d termasuk plot implisit.
Plot-plot ini menunjukkan himpunan nol dari suatu fungsi dalam tiga
variabel.\\
Solusi dari

\end{eulercomment}
\begin{eulerformula}
\[
f(x,y,z) = 0
\]
\end{eulerformula}
\begin{eulercomment}
dapat divisualisasikan dalam potongan sejajar dengan bidang x-y-, x-z-
dan y-z.

- implisit=1: potong sejajar dengan bidang y-z\\
- implisit=2: potong sejajar dengan bidang x-z\\
- implisit=4: potong sejajar dengan bidang x-y

Tambahkan nilai-nilai ini, jika Anda suka. Dalam contoh kita plot

\end{eulercomment}
\begin{eulerformula}
\[
M = \{ (x,y,z) : x^2+y^3+zy=1 \}
\]
\end{eulerformula}
\begin{eulerprompt}
>plot3d("x^2+y^3+z*y-1",r=5,implicit=3):
\end{eulerprompt}
\eulerimg{27}{images/Shalih Abdillah_22305144009_ApKom4-045.png}
\begin{eulerprompt}
>plot3d("x^2+y^2+4*x*z+z^3",>implicit,r=2,zoom=2.5):
\end{eulerprompt}
\eulerimg{27}{images/Shalih Abdillah_22305144009_ApKom4-046.png}
\eulerheading{Merencanakan Data 3D}
\begin{eulercomment}
Sama seperti plot2d, plot3d menerima data. Untuk objek 3D, Anda perlu
menyediakan matriks nilai x-, y- dan z, atau tiga fungsi atau ekspresi
fx(x,y), fy(x,y), fz(x,y).

\end{eulercomment}
\begin{eulerformula}
\[
\gamma(t,s) = (x(t,s),y(t,s),z(t,s))
\]
\end{eulerformula}
\begin{eulercomment}
Karena x,y,z adalah matriks, kita asumsikan bahwa (t,s) melalui sebuah
kotak persegi. Hasilnya, Anda dapat memplot gambar persegi panjang di
ruang angkasa.

Anda dapat menggunakan bahasa matriks Euler untuk menghasilkan
koordinat secara efektif.

Dalam contoh berikut, kami menggunakan vektor nilai t dan vektor kolom
nilai s untuk membuat parameter permukaan bola. Dalam gambar kita
dapat menandai daerah, dalam kasus kita daerah kutub.
\end{eulercomment}
\begin{eulerprompt}
>t=linspace(0,2pi,180); s=linspace(-pi/2,pi/2,90)'; ...
>x=cos(s)*cos(t); y=cos(s)*sin(t); z=sin(s); ...
>plot3d(x,y,z,>hue, ...
>color=blue,<frame,grid=[10,20], ...
>values=s,contourcolor=red,level=[90°-24°;90°-22°], ...
>scale=1.4,height=50°):
\end{eulerprompt}
\eulerimg{27}{images/Shalih Abdillah_22305144009_ApKom4-047.png}
\begin{eulerprompt}
>t=linspace(0,2pi,180); s=linspace(-pi/2,pi/2,90)'; ...
>x=cos(s)*cos(t); y=cos(s)*sin(t); z=sin(s); ...
>plot3d(x,y,z,>hue, ...
>color=white,<frame,grid=[100,100], ...
>values=s,contourcolor=red,level=[90°-24°;90°-22°], ...
>scale=1.4,height=50°):
\end{eulerprompt}
\eulerimg{27}{images/Shalih Abdillah_22305144009_ApKom4-048.png}
\begin{eulercomment}
Berikut adalah contoh, yang merupakan grafik fungsi.
\end{eulercomment}
\begin{eulerprompt}
>t=-1:0.1:1; s=(-1:0.1:1)'; plot3d(t,s,t*s,grid=110):
\end{eulerprompt}
\eulerimg{27}{images/Shalih Abdillah_22305144009_ApKom4-049.png}
\begin{eulercomment}
Namun, kita bisa membuat segala macam permukaan. Berikut adalah
permukaan yang sama dengan fungsi.

\end{eulercomment}
\begin{eulerformula}
\[
x = y \, z
\]
\end{eulerformula}
\begin{eulerprompt}
>plot3d(t*s,t,s,angle=180°,grid=10):
\end{eulerprompt}
\eulerimg{27}{images/Shalih Abdillah_22305144009_ApKom4-050.png}
\begin{eulercomment}
Dengan lebih banyak usaha, kami dapat menghasilkan banyak permukaan.

Dalam contoh berikut, kita membuat tampilan bayangan dari bola yang
terdistorsi. Koordinat biasa untuk bola adalah

\end{eulercomment}
\begin{eulerformula}
\[
\gamma(t,s) = (\cos(t)\cos(s),\sin(t)\sin(s),\cos(s))
\]
\end{eulerformula}
\begin{eulercomment}
dengan,

\end{eulercomment}
\begin{eulerformula}
\[
0 \le t \le 2\pi, \quad \frac{-\pi}{2} \le s \le \frac{\pi}{2}.
\]
\end{eulerformula}
\begin{eulercomment}
Kami mendistorsi ini dengan sebuah faktor

\end{eulercomment}
\begin{eulerformula}
\[
d(t,s) = \frac{\cos(4t)+\cos(8s)}{4}.
\]
\end{eulerformula}
\begin{eulerprompt}
>t=linspace(0,2pi,320); s=linspace(-pi/2,pi/2,160)'; ...
>d=1+0.2*(cos(4*t)+cos(8*s)); ...
>plot3d(cos(t)*cos(s)*d,sin(t)*cos(s)*d,sin(s)*d,hue=1, ...
>  light=[1,0,1],frame=0,zoom=5):
\end{eulerprompt}
\eulerimg{27}{images/Shalih Abdillah_22305144009_ApKom4-051.png}
\begin{eulercomment}
Tentu saja, titik cloud juga dimungkinkan. Untuk memplot data titik
dalam ruang, kita membutuhkan tiga vektor untuk koordinat titik-titik
tersebut.

Gayanya sama seperti di plot2d dengan points=true;
\end{eulercomment}
\begin{eulerprompt}
>n=500;  ...
>  plot3d(normal(1,n),normal(1,n),normal(1,n),points=true,style="."):
\end{eulerprompt}
\eulerimg{27}{images/Shalih Abdillah_22305144009_ApKom4-052.png}
\begin{eulercomment}
Dimungkinkan juga untuk memplot kurva dalam 3D. Dalam hal ini, lebih
mudah untuk menghitung titik-titik kurva. Untuk kurva di pesawat kami
menggunakan urutan koordinat dan parameter wire=true.
\end{eulercomment}
\begin{eulerprompt}
>t=linspace(0,8pi,500); ...
>plot3d(sin(t),cos(t),t/10,>wire,zoom=3):
\end{eulerprompt}
\eulerimg{27}{images/Shalih Abdillah_22305144009_ApKom4-053.png}
\begin{eulerprompt}
>t=linspace(0,4pi,1000); plot3d(cos(t),sin(t),t/2pi,>wire, ...
>linewidth=3,wirecolor=blue):
\end{eulerprompt}
\eulerimg{27}{images/Shalih Abdillah_22305144009_ApKom4-054.png}
\begin{eulerprompt}
>t=linspace(0,10pi,1000); plot3d(cos(t),sin(t),t/2pi,>wire, ...
>linewidth=3,wirecolor=green,disconnect=true):
\end{eulerprompt}
\eulerimg{27}{images/Shalih Abdillah_22305144009_ApKom4-055.png}
\begin{eulerprompt}
>X=cumsum(normal(3,100)); ...
> plot3d(X[1],X[2],X[3],>anaglyph,>wire):
\end{eulerprompt}
\eulerimg{27}{images/Shalih Abdillah_22305144009_ApKom4-056.png}
\begin{eulercomment}
EMT juga dapat memplot dalam mode anaglyph. Untuk melihat plot seperti
itu, Anda memerlukan kacamata merah/sian.
\end{eulercomment}
\begin{eulerprompt}
> plot3d("x^2+y^3",>anaglyph,>contour,angle=30°):
\end{eulerprompt}
\eulerimg{27}{images/Shalih Abdillah_22305144009_ApKom4-057.png}
\begin{eulercomment}
Seringkali, skema warna spektral digunakan untuk plot. Ini menekankan
ketinggian fungsi.
\end{eulercomment}
\begin{eulerprompt}
>plot3d("x^2*y^3-y",>spectral,>contour,zoom=3.2):
\end{eulerprompt}
\eulerimg{27}{images/Shalih Abdillah_22305144009_ApKom4-058.png}
\begin{eulercomment}
Hasil tampilan plot dengan tambahn perintah frame=false dan \textgreater{}anaglyph
\end{eulercomment}
\begin{eulerprompt}
>plot3d("x^2*y^3-y",>spectral,>anaglyph,>contour,zoom=3.2,frame=false):
\end{eulerprompt}
\eulerimg{27}{images/Shalih Abdillah_22305144009_ApKom4-059.png}
\begin{eulercomment}
Euler juga dapat memplot permukaan berparameter, ketika parameternya
adalah nilai x-, y-, dan z dari gambar kotak persegi panjang dalam
ruang.

Untuk demo berikut, kami mengatur parameter u- dan v-, dan
menghasilkan koordinat ruang dari ini.
\end{eulercomment}
\begin{eulerprompt}
>u=linspace(-1,1,10); v=linspace(0,2*pi,50)'; ...
>X=(3+u*cos(v/2))*cos(v); Y=(3+u*cos(v/2))*sin(v); Z=u*sin(v/2); ...
>plot3d(X,Y,Z,>anaglyph,<frame,>wire,scale=2.3):
\end{eulerprompt}
\eulerimg{27}{images/Shalih Abdillah_22305144009_ApKom4-060.png}
\begin{eulercomment}
Berikut adalah contoh yang lebih rumit, yang megah dengan kacamata
merah/sian.
\end{eulercomment}
\begin{eulerprompt}
>u:=linspace(-pi,pi,160); v:=linspace(-pi,pi,400)';  ...
>x:=(4*(1+.25*sin(3*v))+cos(u))*cos(2*v); ...
>y:=(4*(1+.25*sin(3*v))+cos(u))*sin(2*v); ...
> z=sin(u)+2*cos(3*v); ...
>plot3d(x,y,z,frame=0,scale=1.5,hue=1,light=[1,0,-1],zoom=2.8,>anaglyph):
\end{eulerprompt}
\eulerimg{27}{images/Shalih Abdillah_22305144009_ApKom4-061.png}
\eulerheading{Plot Statistik}
\begin{eulercomment}
Plot bar juga dimungkinkan. Untuk ini, kita harus menyediakan

- x: vektor baris dengan n+1 elemen\\
- y: vektor kolom dengan n+1 elemen\\
- z: matriks nilai nxn.

z bisa lebih besar, tetapi hanya nilai nxn yang akan digunakan.

Dalam contoh, pertama-tama kita menghitung nilainya. Kemudian kita
sesuaikan x dan y, sehingga vektor berpusat pada nilai yang digunakan.
\end{eulercomment}
\begin{eulerprompt}
>x=-1:0.1:1; y=x'; z=x^2+y^2; ...
>xa=(x|1.1)-0.05; ya=(y_1.1)-0.05; ...
>plot3d(xa,ya,z,bar=true):
\end{eulerprompt}
\eulerimg{27}{images/Shalih Abdillah_22305144009_ApKom4-062.png}
\begin{eulerprompt}
>x=-1:0.1:1; y=x'; z=x^2+y^2; ...
>xa=(x|1.1)-0.05; ya=(y_1.1)-0.05; ...
>plot3d(xa,ya,z,bar=true,>anaglyph):
\end{eulerprompt}
\eulerimg{27}{images/Shalih Abdillah_22305144009_ApKom4-063.png}
\begin{eulercomment}
Dimungkinkan untuk membagi plot permukaan menjadi dua atau lebih
bagian.
\end{eulercomment}
\begin{eulerprompt}
>x=-1:0.1:1; y=x'; z=x+y; d=zeros(size(x)); ...
>plot3d(x,y,z,disconnect=2:2:20):
\end{eulerprompt}
\eulerimg{27}{images/Shalih Abdillah_22305144009_ApKom4-064.png}
\begin{eulercomment}
Jika memuat atau menghasilkan matriks data M dari file dan perlu
memplotnya dalam 3D, Anda dapat menskalakan matriks ke [-1,1] dengan
scale(M), atau menskalakan matriks dengan \textgreater{}zscale. Ini dapat
dikombinasikan dengan faktor penskalaan individu yang diterapkan
sebagai tambahan.
\end{eulercomment}
\begin{eulerprompt}
>i=1:20; j=i'; ...
>plot3d(i*j^2+100*normal(20,20),>zscale,scale=[1,1,1.5],angle=-40°,zoom=1.8):
\end{eulerprompt}
\eulerimg{27}{images/Shalih Abdillah_22305144009_ApKom4-065.png}
\begin{eulerprompt}
>Z=intrandom(5,100,6); v=zeros(5,6); ...
>loop 1 to 5; v[#]=getmultiplicities(1:6,Z[#]); end; ...
>columnsplot3d(v',scols=1:5,ccols=[1:5]):
\end{eulerprompt}
\eulerimg{27}{images/Shalih Abdillah_22305144009_ApKom4-066.png}
\eulerheading{Permukaan Benda Putar}
\begin{eulerprompt}
>plot2d("(x^2+y^2-1)^3-x^2*y^3",r=1.3, ...
>style="#",color=red,<outline, ...
>level=[-2;0],n=100):
\end{eulerprompt}
\eulerimg{27}{images/Shalih Abdillah_22305144009_ApKom4-067.png}
\begin{eulerprompt}
>ekspresi &= (x^2+y^2-1)^3-x^2*y^3; $ekspresi
\end{eulerprompt}
\begin{eulerformula}
\[
\left(y^2+x^2-1\right)^3-x^2\,y^3
\]
\end{eulerformula}
\begin{eulercomment}
Kami ingin memutar kurva jantung di sekitar sumbu y. Berikut adalah
ungkapan, yang mendefinisikan hati:

\end{eulercomment}
\begin{eulerformula}
\[
f(x,y)=(x^2+y^2-1)^3-x^2.y^3.
\]
\end{eulerformula}
\begin{eulercomment}
Selanjutnya kita atur

\end{eulercomment}
\begin{eulerformula}
\[
x=r.cos(a),\quad y=r.sin(a).
\]
\end{eulerformula}
\begin{eulerprompt}
>function fr(r,a) &= ekspresi with [x=r*cos(a),y=r*sin(a)] | trigreduce; $fr(r,a)
\end{eulerprompt}
\begin{eulerformula}
\[
\left(r^2-1\right)^3+\frac{\left(\sin \left(5\,a\right)-\sin \left(  3\,a\right)-2\,\sin a\right)\,r^5}{16}
\]
\end{eulerformula}
\begin{eulercomment}
Hal ini memungkinkan untuk mendefinisikan fungsi numerik, yang
memecahkan r, jika a diberikan. Dengan fungsi itu kita dapat memplot
jantung yang diputar sebagai permukaan parametrik.
\end{eulercomment}
\begin{eulerprompt}
>function map f(a) := bisect("fr",0,2;a); ...
>t=linspace(-pi/2,pi/2,100); r=f(t);  ...
>s=linspace(pi,2pi,100)'; ...
>plot3d(r*cos(t)*sin(s),r*cos(t)*cos(s),r*sin(t), ...
>>hue,<frame,color=red,zoom=4,amb=0,max=0.7,grid=12,height=50°):
\end{eulerprompt}
\eulerimg{27}{images/Shalih Abdillah_22305144009_ApKom4-070.png}
\begin{eulercomment}
Berikut ini adalah plot 3D dari gambar di atas yang diputar di sekitar
sumbu z. Kami mendefinisikan fungsi, yang menggambarkan objek.
\end{eulercomment}
\begin{eulerprompt}
>function f(x,y,z) ...
\end{eulerprompt}
\begin{eulerudf}
  r=x^2+y^2;
  return (r+z^2-1)^3-r*z^3;
   endfunction
\end{eulerudf}
\begin{eulerprompt}
>plot3d("f(x,y,z)", ...
>xmin=0,xmax=1.2,ymin=-1.2,ymax=1.2,zmin=-1.2,zmax=1.4, ...
>implicit=1,angle=-30°,zoom=2.5,n=[10,60,60],>anaglyph):
\end{eulerprompt}
\eulerimg{27}{images/Shalih Abdillah_22305144009_ApKom4-071.png}
\eulerheading{Plot 3D Khusus}
\begin{eulercomment}
Fungsi plot3d bagus untuk dimiliki, tetapi tidak memenuhi semua
kebutuhan. Selain rutinitas yang lebih mendasar, dimungkinkan untuk
mendapatkan plot berbingkai dari objek apa pun yang Anda suka.

Meskipun Euler bukan program 3D, ia dapat menggabungkan beberapa objek
dasar. Kami mencoba memvisualisasikan paraboloid dan garis
singgungnya.
\end{eulercomment}
\begin{eulerprompt}
>function myplot ...
\end{eulerprompt}
\begin{eulerudf}
    y=0:0.01:1; x=(0.1:0.01:1)';
    plot3d(x,y,0.2*(x-0.1)/2,<scale,<frame,>hue, ..
      hues=0.5,>contour,color=orange);
    h=holding(1);
    plot3d(x,y,(x^2+y^2)/2,<scale,<frame,>contour,>hue);
    holding(h);
  endfunction
\end{eulerudf}
\begin{eulercomment}
Sekarang framedplot() menyediakan frame, dan mengatur tampilan.
\end{eulercomment}
\begin{eulerprompt}
>framedplot("myplot",[0.1,1,0,1,0,1],angle=-45°, ...
>  center=[0,0,-0.7],zoom=6):
\end{eulerprompt}
\eulerimg{27}{images/Shalih Abdillah_22305144009_ApKom4-072.png}
\begin{eulercomment}
Dengan cara yang sama, Anda dapat memplot bidang kontur secara manual.
Perhatikan bahwa plot3d() menyetel jendela ke fullwindow() secara
default, tetapi plotcontourplane() mengasumsikan itu.
\end{eulercomment}
\begin{eulerprompt}
>x=-1:0.02:1.1; y=x'; z=x^2-y^4;
>function myplot (x,y,z) ...
\end{eulerprompt}
\begin{eulerudf}
    zoom(2);
    wi=fullwindow();
    plotcontourplane(x,y,z,level="auto",<scale);
    plot3d(x,y,z,>hue,<scale,>add,color=white,level="thin");
    window(wi);
    reset();
  endfunction
\end{eulerudf}
\begin{eulerprompt}
>myplot(x,y,z):
\end{eulerprompt}
\eulerimg{27}{images/Shalih Abdillah_22305144009_ApKom4-073.png}
\eulerheading{Animasi}
\begin{eulercomment}
Euler dapat menggunakan frame untuk menghitung animasi terlebih
dahulu.

Salah satu fungsi yang memanfaatkan teknik ini adalah rotate. Itu
dapat mengubah sudut pandang dan menggambar ulang plot 3D. Fungsi
memanggil addpage() untuk setiap plot baru. Akhirnya itu menjiwai
plot.

Silakan pelajari sumber rotasi untuk melihat lebih detail.
\end{eulercomment}
\begin{eulerprompt}
>function testplot () := plot3d("x^2+y^3"); ...
>rotate("testplot"); testplot():
\end{eulerprompt}
\eulerimg{27}{images/Shalih Abdillah_22305144009_ApKom4-074.png}
\eulerheading{Menggambar Povray}
\begin{eulercomment}
Dengan bantuan file Euler povray.e, Euler dapat menghasilkan file
Povray. Hasilnya sangat bagus untuk dilihat.

Anda perlu menginstal Povray (32bit atau 64bit) dari
http://www.povray.org/, dan meletakkan sub-direktori "bin" dari Povray ke jalur lingkungan, atau mengatur variabel "defaultpovray" dengan path lengkap yang menunjuk ke "pvengine.exe".

Antarmuka Povray dari Euler menghasilkan file Povray di direktori home
pengguna, dan memanggil Povray untuk mengurai file-file ini. Nama file
default adalah current.pov, dan direktori default adalah eulerhome(),
biasanya c:\textbackslash{}Users\textbackslash{}Username\textbackslash{}Euler. Povray menghasilkan file PNG, yang
dapat dimuat oleh Euler ke dalam buku catatan. Untuk membersihkan
file-file ini, gunakan povclear().

Fungsi pov3d memiliki semangat yang sama dengan plot3d. Ini dapat
menghasilkan grafik fungsi f(x,y), atau permukaan dengan koordinat
X,Y,Z dalam matriks, termasuk garis level opsional. Fungsi ini memulai
raytracer secara otomatis, dan memuat adegan ke dalam notebook Euler.

Selain pov3d(), ada banyak fungsi yang menghasilkan objek Povray.
Fungsi-fungsi ini mengembalikan string, yang berisi kode Povray untuk
objek. Untuk menggunakan fungsi ini, mulai file Povray dengan
povstart(). Kemudian gunakan writeln(...) untuk menulis objek ke file
adegan. Terakhir, akhiri file dengan povend(). Secara default,
raytracer akan dimulai, dan PNG akan dimasukkan ke dalam notebook
Euler.

Fungsi objek memiliki parameter yang disebut "look", yang membutuhkan
string dengan kode Povray untuk tekstur dan hasil akhir objek. Fungsi
povlook() dapat digunakan untuk menghasilkan string ini. Ini memiliki
parameter untuk warna, transparansi, Phong Shading dll.

Perhatikan bahwa alam semesta Povray memiliki sistem koordinat lain.
Antarmuka ini menerjemahkan semua koordinat ke sistem Povray. Jadi
Anda dapat terus berpikir dalam sistem koordinat Euler dengan z
menunjuk vertikal ke atas, a nd x,y,z sumbu dalam arti tangan kanan.\\
Anda perlu memuat file povray.
\end{eulercomment}
\begin{eulerprompt}
>load povray;
\end{eulerprompt}
\begin{eulercomment}
Pastikan, direktori bin Povray ada di jalur. Jika tidak, edit variabel
berikut sehingga berisi path ke povray yang dapat dieksekusi.
\end{eulercomment}
\begin{eulerprompt}
>defaultpovray="C:\(\backslash\)Program Files\(\backslash\)POV-Ray\(\backslash\)v3.7\(\backslash\)bin\(\backslash\)pvengine.exe"
\end{eulerprompt}
\begin{euleroutput}
  C:\(\backslash\)Program Files\(\backslash\)POV-Ray\(\backslash\)v3.7\(\backslash\)bin\(\backslash\)pvengine.exe
\end{euleroutput}
\begin{eulercomment}
Untuk kesan pertama, kami memplot fungsi sederhana. Perintah berikut
menghasilkan file povray di direktori pengguna Anda, dan menjalankan
Povray untuk ray tracing file ini.

Jika Anda memulai perintah berikut, GUI Povray akan terbuka,
menjalankan file, dan menutup secara otomatis. Karena alasan keamanan,
Anda akan ditanya, apakah Anda ingin mengizinkan file exe untuk
dijalankan. Anda dapat menekan batal untuk menghentikan pertanyaan
lebih lanjut. Anda mungkin harus menekan OK di jendela Povray untuk
mengakui dialog awal Povray.
\end{eulercomment}
\begin{eulerprompt}
>pov3d("x^2+y^2",zoom=3);
\end{eulerprompt}
\eulerimg{27}{images/Shalih Abdillah_22305144009_ApKom4-075.png}
\begin{eulercomment}
Kita dapat membuat fungsi menjadi transparan dan menambahkan hasil
akhir lainnya. Kami juga dapat menambahkan garis level ke plot fungsi.
\end{eulercomment}
\begin{eulerprompt}
>pov3d("x^2+y^3",axiscolor=red,angle=20°, ...
>  look=povlook(blue,0.2),level=-1:0.5:1,zoom=3.8);
\end{eulerprompt}
\eulerimg{27}{images/Shalih Abdillah_22305144009_ApKom4-076.png}
\begin{eulercomment}
Terkadang perlu untuk mencegah penskalaan fungsi, dan menskalakan
fungsi dengan tangan.

Kami memplot himpunan titik di bidang kompleks, di mana produk dari
jarak ke 1 dan -1 sama dengan 1.
\end{eulercomment}
\begin{eulerprompt}
>pov3d("((x-1)^2+y^2)*((x+1)^2+y^2)/40",r=1.5, ...
>  angle=-120°,level=1/40,dlevel=0.005,light=[-1,1,1],height=45°,n=50, ...
>  <fscale,zoom=3.8);
\end{eulerprompt}
\eulerimg{27}{images/Shalih Abdillah_22305144009_ApKom4-077.png}
\eulerheading{Merencanakan dengan Koordinat}
\begin{eulercomment}
Alih-alih fungsi, kita dapat memplot dengan koordinat. Seperti pada
plot3d, kita membutuhkan tiga matriks untuk mendefinisikan objek.

Dalam contoh kita memutar fungsi di sekitar sumbu z.
\end{eulercomment}
\begin{eulerprompt}
>function f(x) := x^3-x+1; ...
>x=-1:0.01:1; t=linspace(0,2pi,8)'; ...
>Z=x; X=cos(t)*f(x); Y=sin(t)*f(x); ...
>pov3d(X,Y,Z,angle=40°,height=20°,axis=0,zoom=4,light=[10,-5,5]);
\end{eulerprompt}
\eulerimg{27}{images/Shalih Abdillah_22305144009_ApKom4-078.png}
\begin{eulercomment}
Dalam contoh berikut, kami memplot gelombang teredam. Kami
menghasilkan gelombang dengan bahasa matriks Euler.

Kami juga menunjukkan, bagaimana objek tambahan dapat ditambahkan ke
adegan pov3d. Untuk pembuatan objek, lihat contoh berikut. Perhatikan
bahwa plot3d menskalakan plot, sehingga cocok dengan kubus satuan.
\end{eulercomment}
\begin{eulerprompt}
>r=linspace(0,1,80); phi=linspace(0,2pi,80)'; ...
>x=r*cos(phi); y=r*sin(phi); z=exp(-5*r)*cos(8*pi*r)/3;  ...
>pov3d(x,y,z,zoom=5,axis=0,add=povsphere([0,0,0.5],0.1,povlook(green)), ...
>  w=500,h=300);
\end{eulerprompt}
\eulerimg{16}{images/Shalih Abdillah_22305144009_ApKom4-079.png}
\begin{eulercomment}
Dengan metode bayangan canggih dari Povray, sangat sedikit titik yang
dapat menghasilkan permukaan yang sangat halus. Hanya di perbatasan
dan dalam bayang-bayang triknya mungkin menjadi jelas.

Untuk ini, kita perlu menambahkan vektor normal di setiap titik
matriks.
\end{eulercomment}
\begin{eulerprompt}
>Z &= x^2*y^3
\end{eulerprompt}
\begin{euleroutput}
  
                                   2  3
                                  x  y
  
\end{euleroutput}
\begin{eulercomment}
Persamaan permukaannya adalah [x,y,Z]. Kami menghitung dua turunan ke
x dan y ini dan mengambil produk silang sebagai normal.
\end{eulercomment}
\begin{eulerprompt}
>dx &= diff([x,y,Z],x); dy &= diff([x,y,Z],y);
\end{eulerprompt}
\begin{eulercomment}
Kami mendefinisikan normal sebagai produk silang dari turunan ini, dan
mendefinisikan fungsi koordinat.
\end{eulercomment}
\begin{eulerprompt}
>N &= crossproduct(dx,dy); NX &= N[1]; NY &= N[2]; NZ &= N[3]; N,
\end{eulerprompt}
\begin{euleroutput}
  
                                 3       2  2
                         [- 2 x y , - 3 x  y , 1]
  
\end{euleroutput}
\begin{eulercomment}
Kita gunakan 25 poin/ titik saja.
\end{eulercomment}
\begin{eulerprompt}
>x=-1:0.5:1; y=x';
>pov3d(x,y,Z(x,y),angle=10°, ...
>  xv=NX(x,y),yv=NY(x,y),zv=NZ(x,y),<shadow);
\end{eulerprompt}
\eulerimg{27}{images/Shalih Abdillah_22305144009_ApKom4-080.png}
\begin{eulercomment}
Berikut ini adalah simpul Trefoil yang dilakukan oleh A. Busser di
Povray. Ada versi yang ditingkatkan dari ini dalam contoh.

Lihat: Contoh\textbackslash{}Trefoil Simpul \textbar{} Simpul trefoil

Untuk tampilan yang bagus dengan tidak terlalu banyak titik, kami
menambahkan vektor normal di sini. Kami menggunakan Maxima untuk
menghitung normal bagi kami. Pertama, ketiga fungsi koordinat sebagai
ekspresi simbolik.
\end{eulercomment}
\begin{eulerprompt}
>X &= ((4+sin(3*y))+cos(x))*cos(2*y); ...
>Y &= ((4+sin(3*y))+cos(x))*sin(2*y); ...
>Z &= sin(x)+2*cos(3*y);
\end{eulerprompt}
\begin{eulercomment}
Kemudian kedua vektor turunan ke x dan y.
\end{eulercomment}
\begin{eulerprompt}
>dx &= diff([X,Y,Z],x); dy &= diff([X,Y,Z],y);
\end{eulerprompt}
\begin{eulercomment}
Sekarang normal, yang merupakan produk silang dari dua turunan.
\end{eulercomment}
\begin{eulerprompt}
>dn &= crossproduct(dx,dy);
\end{eulerprompt}
\begin{eulercomment}
Kami sekarang mengevaluasi semua ini secara numerik.
\end{eulercomment}
\begin{eulerprompt}
>x:=linspace(-%pi,%pi,40); y:=linspace(-%pi,%pi,100)';
\end{eulerprompt}
\begin{eulercomment}
Vektor normal adalah evaluasi dari ekspresi simbolik dn[i] untuk
i=1,2,3. Sintaks untuk ini adalah \&"expression"(parameters). Ini
adalah alternatif dari metode pada contoh sebelumnya, di mana kita
mendefinisikan ekspresi simbolik NX, NY, NZ terlebih dahulu.
\end{eulercomment}
\begin{eulerprompt}
>pov3d(X(x,y),Y(x,y),Z(x,y),axis=0,zoom=5,w=450,h=350, ...
>  <shadow,look=povlook(gray), ...
>  xv=&"dn[1]"(x,y), yv=&"dn[2]"(x,y), zv=&"dn[3]"(x,y));
\end{eulerprompt}
\eulerimg{21}{images/Shalih Abdillah_22305144009_ApKom4-081.png}
\begin{eulercomment}
Kami juga dapat menghasilkan grid dalam 3D.
\end{eulercomment}
\begin{eulerprompt}
>povstart(zoom=4); ...
>x=-1:0.5:1; r=1-(x+1)^2/6; ...
>t=(0°:30°:360°)'; y=r*cos(t); z=r*sin(t); ...
>writeln(povgrid(x,y,z,d=0.02,dballs=0.05)); ...
>povend();
\end{eulerprompt}
\eulerimg{27}{images/Shalih Abdillah_22305144009_ApKom4-082.png}
\begin{eulercomment}
Dengan povgrid(), kurva dimungkinkan.
\end{eulercomment}
\begin{eulerprompt}
>povstart(center=[0,0,1],zoom=3.6); ...
>t=linspace(0,2,1000); r=exp(-t); ...
>x=cos(2*pi*10*t)*r; y=sin(2*pi*10*t)*r; z=t; ...
>writeln(povgrid(x,y,z,povlook(red))); ...
>writeAxis(0,2,axis=3); ...
>povend();
\end{eulerprompt}
\eulerimg{27}{images/Shalih Abdillah_22305144009_ApKom4-083.png}
\eulerheading{Objek Povray}
\begin{eulercomment}
Di atas, kami menggunakan pov3d untuk memplot permukaan. Antarmuka
povray di Euler juga dapat menghasilkan objek Povray. Objek-objek ini
disimpan sebagai string di Euler, dan perlu ditulis ke file Povray.

Kami memulai output dengan povstart().
\end{eulercomment}
\begin{eulerprompt}
>povstart(zoom=4);
\end{eulerprompt}
\begin{eulercomment}
Pertama kita mendefinisikan tiga silinder, dan menyimpannya dalam
string di Euler.

Fungsi povx() dll. hanya mengembalikan vektor [1,0,0], yang dapat
digunakan sebagai gantinya.
\end{eulercomment}
\begin{eulerprompt}
>c1=povcylinder(-povx,povx,1,povlook(red)); ...
>c2=povcylinder(-povy,povy,1,povlook(green)); ...
>c3=povcylinder(-povz,povz,1,povlook(blue)); ...
\end{eulerprompt}
\begin{eulercomment}
String berisi kode Povray, yang tidak perlu kita pahami pada saat itu.
\end{eulercomment}
\begin{eulerprompt}
>c1
\end{eulerprompt}
\begin{euleroutput}
  cylinder \{ <-1,0,0>, <1,0,0>, 1
   texture \{ pigment \{ color rgb <0.564706,0.0627451,0.0627451> \}  \} 
   finish \{ ambient 0.2 \} 
   \}
\end{euleroutput}
\begin{eulercomment}
Seperti yang Anda lihat, kami menambahkan tekstur ke objek dalam tiga
warna berbeda.

Itu dilakukan oleh povlook(), yang mengembalikan string dengan kode
Povray yang relevan. Kita dapat menggunakan warna Euler default, atau
menentukan warna kita sendiri. Kami juga dapat menambahkan
transparansi, atau mengubah cahaya sekitar.
\end{eulercomment}
\begin{eulerprompt}
>povlook(rgb(0.1,0.2,0.3),0.1,0.5);
\end{eulerprompt}
\begin{eulercomment}
Sekarang kita mendefinisikan objek persimpangan, dan menulis hasilnya
ke file.
\end{eulercomment}
\begin{eulerprompt}
>writeln(povintersection([c1,c2,c3]));
\end{eulerprompt}
\begin{eulercomment}
Persimpangan tiga silinder sulit untuk divisualisasikan, jika Anda
belum pernah melihatnya sebelumnya.
\end{eulercomment}
\begin{eulerprompt}
>povend;
\end{eulerprompt}
\eulerimg{27}{images/Shalih Abdillah_22305144009_ApKom4-084.png}
\begin{eulercomment}
Fungsi berikut menghasilkan fraktal secara rekursif.

Fungsi pertama menunjukkan, bagaimana Euler menangani objek Povray
sederhana. Fungsi povbox() mengembalikan string, yang berisi koordinat
kotak, tekstur, dan hasil akhir.
\end{eulercomment}
\begin{eulerprompt}
>function onebox(x,y,z,d) := povbox([x,y,z],[x+d,y+d,z+d],povlook());
>function fractal (x,y,z,h,n) ...
\end{eulerprompt}
\begin{eulerudf}
   if n==1 then writeln(onebox(x,y,z,h));
   else
     h=h/3;
     fractal(x,y,z,h,n-1);
     fractal(x+2*h,y,z,h,n-1);
     fractal(x,y+2*h,z,h,n-1);
     fractal(x,y,z+2*h,h,n-1);
     fractal(x+2*h,y+2*h,z,h,n-1);
     fractal(x+2*h,y,z+2*h,h,n-1);
     fractal(x,y+2*h,z+2*h,h,n-1);
     fractal(x+2*h,y+2*h,z+2*h,h,n-1);
     fractal(x+h,y+h,z+h,h,n-1);
   endif;
  endfunction
\end{eulerudf}
\begin{eulerprompt}
>povstart(fade=10,<shadow):
\end{eulerprompt}
\eulerimg{27}{images/Shalih Abdillah_22305144009_ApKom4-085.png}
\begin{eulerprompt}
>fractal(-1,-1,-1,2,4);
>povend();
\end{eulerprompt}
\eulerimg{27}{images/Shalih Abdillah_22305144009_ApKom4-086.png}
\begin{eulercomment}
Differences allow cutting off one object from another. Like
intersections, there are part of the CSG objects of Povray.
\end{eulercomment}
\begin{eulerprompt}
>povstart(light=[5,-5,5],fade=10);
\end{eulerprompt}
\begin{eulercomment}
Untuk demonstrasi ini, kami mendefinisikan objek di Povray, alih-alih
menggunakan string di Euler. Definisi ditulis ke file segera.

Koordinat kotak -1 berarti [-1,-1,-1].
\end{eulercomment}
\begin{eulerprompt}
>povdefine("mycube",povbox(-1,1));
\end{eulerprompt}
\begin{eulercomment}
Kita dapat menggunakan objek ini di povobject(), yang mengembalikan
string seperti biasa.
\end{eulercomment}
\begin{eulerprompt}
>c1=povobject("mycube",povlook(red));
\end{eulerprompt}
\begin{eulercomment}
Kami menghasilkan kubus kedua, dan memutar dan menskalakannya sedikit.
\end{eulercomment}
\begin{eulerprompt}
>c2=povobject("mycube",povlook(yellow),translate=[1,1,1], ...
>  rotate=xrotate(10°)+yrotate(10°), scale=1.2);
\end{eulerprompt}
\begin{eulercomment}
Kemudian kita ambil selisih kedua benda tersebut.
\end{eulercomment}
\begin{eulerprompt}
>writeln(povdifference(c1,c2));
\end{eulerprompt}
\begin{eulercomment}
Sekarang tambahkan tiga sumbu.
\end{eulercomment}
\begin{eulerprompt}
>writeAxis(-1.2,1.2,axis=1); ...
>writeAxis(-1.2,1.2,axis=2); ...
>writeAxis(-1.2,1.2,axis=4); ...
>povend();
\end{eulerprompt}
\eulerimg{27}{images/Shalih Abdillah_22305144009_ApKom4-087.png}
\eulerheading{Fungsi Implisit}
\begin{eulercomment}
Povray dapat memplot himpunan di mana f(x,y,z)=0, seperti parameter
implisit di plot3d. Namun, hasilnya terlihat jauh lebih baik.

Sintaks untuk fungsinya sedikit berbeda. Anda tidak dapat menggunakan
output dari ekspresi Maxima atau Euler.
\end{eulercomment}
\begin{eulerprompt}
>povstart(angle=70°,height=50°,zoom=4);
\end{eulerprompt}
\begin{eulercomment}
Buat permukaan implisit. Perhatikan sintaks yang berbeda dalam
ekspresi.
\end{eulercomment}
\begin{eulerprompt}
>writeln(povsurface("pow(x,2)*y-pow(y,3)-pow(z,2)",povlook(green))); ...
>writeAxes(); ...
>povend();
\end{eulerprompt}
\eulerimg{27}{images/Shalih Abdillah_22305144009_ApKom4-088.png}
\eulerheading{Objek Jala}
\begin{eulercomment}
Dalam contoh ini, kami menunjukkan cara membuat objek mesh, dan
menggambarnya dengan informasi tambahan.

Kami ingin memaksimalkan xy di bawah kondisi x+y=1 dan menunjukkan
sentuhan tangensial dari garis level.
\end{eulercomment}
\begin{eulerprompt}
>\(\backslash\)povstart(angle=-10°,center=[0.5,0.5,0.5],zoom=7);
\end{eulerprompt}
\begin{euleroutput}
  Syntax error in expression, or unfinished expression!
  Error in:
  \(\backslash\)povstart(angle=-10°,center=[0.5,0.5,0.5],zoom=7); ...
  ^
\end{euleroutput}
\begin{eulercomment}
Kami tidak dapat menyimpan objek dalam string seperti sebelumnya,
karena terlalu besar. Jadi kita mendefinisikan objek dalam file Povray
menggunakan #declare. Fungsi povtriangle() melakukan ini secara
otomatis. Itu dapat menerima vektor normal seperti pov3d().

Berikut ini mendefinisikan objek mesh, dan langsung menulisnya ke
dalam file.
\end{eulercomment}
\begin{eulerprompt}
>x=0:0.02:1; y=x'; z=x*y; vx=-y; vy=-x; vz=1;
>mesh=povtriangles(x,y,z,"",vx,vy,vz);
\end{eulerprompt}
\begin{euleroutput}
  #declare mesh1 =
  mesh2 \{
  vertex_vectors \{ 2601,
  <0,0,0>,
  <0.02,0,0>,
  <0.04,0,0>,
  <0.06,0,0>,
  <0.08,0,0>,
  <0.1,0,0>,
  <0.12,0,0>,
  <0.14,0,0>,
  <0.16,0,0>,
  <0.18,0,0>,
  <0.2,0,0>,
  <0.22,0,0>,
  <0.24,0,0>,
  <0.26,0,0>,
  <0.28,0,0>,
  <0.3,0,0>,
  <0.32,0,0>,
  <0.34,0,0>,
  <0.36,0,0>,
  <0.38,0,0>,
  <0.4,0,0>,
  <0.42,0,0>,
  <0.44,0,0>,
  <0.46,0,0>,
  <0.48,0,0>,
  <0.5,0,0>,
  <0.52,0,0>,
  <0.54,0,0>,
  <0.56,0,0>,
  <0.58,0,0>,
  <0.6,0,0>,
  <0.62,0,0>,
  <0.64,0,0>,
  <0.66,0,0>,
  <0.68,0,0>,
  <0.7,0,0>,
  <0.72,0,0>,
  <0.74,0,0>,
  <0.76,0,0>,
  <0.78,0,0>,
  <0.8,0,0>,
  <0.82,0,0>,
  <0.84,0,0>,
  <0.86,0,0>,
  <0.88,0,0>,
  <0.9,0,0>,
  <0.92,0,0>,
  <0.94,0,0>,
  <0.96,0,0>,
  <0.98,0,0>,
  <1,0,0>,
  <0,0.02,0>,
  <0.02,0.02,0.0004>,
  <0.04,0.02,0.0008>,
  <0.06,0.02,0.0012>,
  <0.08,0.02,0.0016>,
  <0.1,0.02,0.002>,
  <0.12,0.02,0.0024>,
  <0.14,0.02,0.0028>,
  <0.16,0.02,0.0032>,
  <0.18,0.02,0.0036>,
  <0.2,0.02,0.004>,
  <0.22,0.02,0.0044>,
  <0.24,0.02,0.0048>,
  <0.26,0.02,0.0052>,
  <0.28,0.02,0.0056>,
  <0.3,0.02,0.006>,
  <0.32,0.02,0.0064>,
  <0.34,0.02,0.0068>,
  <0.36,0.02,0.0072>,
  <0.38,0.02,0.0076>,
  <0.4,0.02,0.008>,
  <0.42,0.02,0.0084>,
  <0.44,0.02,0.0088>,
  <0.46,0.02,0.0092>,
  <0.48,0.02,0.0096>,
  <0.5,0.02,0.01>,
  <0.52,0.02,0.0104>,
  <0.54,0.02,0.0108>,
  <0.56,0.02,0.0112>,
  <0.58,0.02,0.0116>,
  <0.6,0.02,0.012>,
  <0.62,0.02,0.0124>,
  <0.64,0.02,0.0128>,
  <0.66,0.02,0.0132>,
  <0.68,0.02,0.0136>,
  <0.7,0.02,0.014>,
  <0.72,0.02,0.0144>,
  <0.74,0.02,0.0148>,
  <0.76,0.02,0.0152>,
  <0.78,0.02,0.0156>,
  <0.8,0.02,0.016>,
  <0.82,0.02,0.0164>,
  <0.84,0.02,0.0168>,
  <0.86,0.02,0.0172>,
  <0.88,0.02,0.0176>,
  <0.9,0.02,0.018>,
  <0.92,0.02,0.0184>,
  <0.94,0.02,0.0188>,
  <0.96,0.02,0.0192>,
  <0.98,0.02,0.0196>,
  <1,0.02,0.02>,
  <0,0.04,0>,
  <0.02,0.04,0.0008>,
  <0.04,0.04,0.0016>,
  <0.06,0.04,0.0024>,
  <0.08,0.04,0.0032>,
  <0.1,0.04,0.004>,
  <0.12,0.04,0.0048>,
  <0.14,0.04,0.0056>,
  <0.16,0.04,0.0064>,
  <0.18,0.04,0.0072>,
  <0.2,0.04,0.008>,
  <0.22,0.04,0.0088>,
  <0.24,0.04,0.0096>,
  <0.26,0.04,0.0104>,
  <0.28,0.04,0.0112>,
  <0.3,0.04,0.012>,
  <0.32,0.04,0.0128>,
  <0.34,0.04,0.0136>,
  <0.36,0.04,0.0144>,
  <0.38,0.04,0.0152>,
  <0.4,0.04,0.016>,
  <0.42,0.04,0.0168>,
  <0.44,0.04,0.0176>,
  <0.46,0.04,0.0184>,
  <0.48,0.04,0.0192>,
  <0.5,0.04,0.02>,
  <0.52,0.04,0.0208>,
  <0.54,0.04,0.0216>,
  <0.56,0.04,0.0224>,
  <0.58,0.04,0.0232>,
  <0.6,0.04,0.024>,
  <0.62,0.04,0.0248>,
  <0.64,0.04,0.0256>,
  <0.66,0.04,0.0264>,
  <0.68,0.04,0.0272>,
  <0.7,0.04,0.028>,
  <0.72,0.04,0.0288>,
  <0.74,0.04,0.0296>,
  <0.76,0.04,0.0304>,
  <0.78,0.04,0.0312>,
  <0.8,0.04,0.032>,
  <0.82,0.04,0.0328>,
  <0.84,0.04,0.0336>,
  <0.86,0.04,0.0344>,
  <0.88,0.04,0.0352>,
  <0.9,0.04,0.036>,
  <0.92,0.04,0.0368>,
  <0.94,0.04,0.0376>,
  <0.96,0.04,0.0384>,
  <0.98,0.04,0.0392>,
  <1,0.04,0.04>,
  <0,0.06,0>,
  <0.02,0.06,0.0012>,
  <0.04,0.06,0.0024>,
  <0.06,0.06,0.0036>,
  <0.08,0.06,0.0048>,
  <0.1,0.06,0.006>,
  <0.12,0.06,0.0072>,
  <0.14,0.06,0.0084>,
  <0.16,0.06,0.0096>,
  <0.18,0.06,0.0108>,
  <0.2,0.06,0.012>,
  <0.22,0.06,0.0132>,
  <0.24,0.06,0.0144>,
  <0.26,0.06,0.0156>,
  <0.28,0.06,0.0168>,
  <0.3,0.06,0.018>,
  <0.32,0.06,0.0192>,
  <0.34,0.06,0.0204>,
  <0.36,0.06,0.0216>,
  <0.38,0.06,0.0228>,
  <0.4,0.06,0.024>,
  <0.42,0.06,0.0252>,
  <0.44,0.06,0.0264>,
  <0.46,0.06,0.0276>,
  <0.48,0.06,0.0288>,
  <0.5,0.06,0.03>,
  <0.52,0.06,0.0312>,
  <0.54,0.06,0.0324>,
  <0.56,0.06,0.0336>,
  <0.58,0.06,0.0348>,
  <0.6,0.06,0.036>,
  <0.62,0.06,0.0372>,
  <0.64,0.06,0.0384>,
  <0.66,0.06,0.0396>,
  <0.68,0.06,0.0408>,
  <0.7,0.06,0.042>,
  <0.72,0.06,0.0432>,
  <0.74,0.06,0.0444>,
  <0.76,0.06,0.0456>,
  <0.78,0.06,0.0468>,
  <0.8,0.06,0.048>,
  <0.82,0.06,0.0492>,
  <0.84,0.06,0.0504>,
  <0.86,0.06,0.0516>,
  <0.88,0.06,0.0528>,
  <0.9,0.06,0.054>,
  <0.92,0.06,0.0552>,
  <0.94,0.06,0.0564>,
  <0.96,0.06,0.0576>,
  <0.98,0.06,0.0588>,
  <1,0.06,0.06>,
  <0,0.08,0>,
  <0.02,0.08,0.0016>,
  <0.04,0.08,0.0032>,
  <0.06,0.08,0.0048>,
  <0.08,0.08,0.0064>,
  <0.1,0.08,0.008>,
  <0.12,0.08,0.0096>,
  <0.14,0.08,0.0112>,
  <0.16,0.08,0.0128>,
  <0.18,0.08,0.0144>,
  <0.2,0.08,0.016>,
  <0.22,0.08,0.0176>,
  <0.24,0.08,0.0192>,
  <0.26,0.08,0.0208>,
  <0.28,0.08,0.0224>,
  <0.3,0.08,0.024>,
  <0.32,0.08,0.0256>,
  <0.34,0.08,0.0272>,
  <0.36,0.08,0.0288>,
  <0.38,0.08,0.0304>,
  <0.4,0.08,0.032>,
  <0.42,0.08,0.0336>,
  <0.44,0.08,0.0352>,
  <0.46,0.08,0.0368>,
  <0.48,0.08,0.0384>,
  <0.5,0.08,0.04>,
  <0.52,0.08,0.0416>,
  <0.54,0.08,0.0432>,
  <0.56,0.08,0.0448>,
  <0.58,0.08,0.0464>,
  <0.6,0.08,0.048>,
  <0.62,0.08,0.0496>,
  <0.64,0.08,0.0512>,
  <0.66,0.08,0.0528>,
  <0.68,0.08,0.0544>,
  <0.7,0.08,0.056>,
  <0.72,0.08,0.0576>,
  <0.74,0.08,0.0592>,
  <0.76,0.08,0.0608>,
  <0.78,0.08,0.0624>,
  <0.8,0.08,0.064>,
  <0.82,0.08,0.0656>,
  <0.84,0.08,0.0672>,
  <0.86,0.08,0.0688>,
  <0.88,0.08,0.0704>,
  <0.9,0.08,0.072>,
  <0.92,0.08,0.0736>,
  <0.94,0.08,0.0752>,
  <0.96,0.08,0.0768>,
  <0.98,0.08,0.0784>,
  <1,0.08,0.08>,
  <0,0.1,0>,
  <0.02,0.1,0.002>,
  <0.04,0.1,0.004>,
  <0.06,0.1,0.006>,
  <0.08,0.1,0.008>,
  <0.1,0.1,0.01>,
  <0.12,0.1,0.012>,
  <0.14,0.1,0.014>,
  <0.16,0.1,0.016>,
  <0.18,0.1,0.018>,
  <0.2,0.1,0.02>,
  <0.22,0.1,0.022>,
  <0.24,0.1,0.024>,
  <0.26,0.1,0.026>,
  <0.28,0.1,0.028>,
  <0.3,0.1,0.03>,
  <0.32,0.1,0.032>,
  <0.34,0.1,0.034>,
  <0.36,0.1,0.036>,
  <0.38,0.1,0.038>,
  <0.4,0.1,0.04>,
  <0.42,0.1,0.042>,
  <0.44,0.1,0.044>,
  <0.46,0.1,0.046>,
  <0.48,0.1,0.048>,
  <0.5,0.1,0.05>,
  <0.52,0.1,0.052>,
  <0.54,0.1,0.054>,
  <0.56,0.1,0.056>,
  <0.58,0.1,0.058>,
  <0.6,0.1,0.06>,
  <0.62,0.1,0.062>,
  <0.64,0.1,0.064>,
  <0.66,0.1,0.066>,
  <0.68,0.1,0.068>,
  <0.7,0.1,0.07>,
  <0.72,0.1,0.072>,
  <0.74,0.1,0.074>,
  <0.76,0.1,0.076>,
  <0.78,0.1,0.078>,
  <0.8,0.1,0.08>,
  <0.82,0.1,0.082>,
  <0.84,0.1,0.084>,
  <0.86,0.1,0.086>,
  <0.88,0.1,0.088>,
  <0.9,0.1,0.09>,
  <0.92,0.1,0.092>,
  <0.94,0.1,0.094>,
  <0.96,0.1,0.096>,
  <0.98,0.1,0.098>,
  <1,0.1,0.1>,
  <0,0.12,0>,
  <0.02,0.12,0.0024>,
  <0.04,0.12,0.0048>,
  <0.06,0.12,0.0072>,
  <0.08,0.12,0.0096>,
  <0.1,0.12,0.012>,
  <0.12,0.12,0.0144>,
  <0.14,0.12,0.0168>,
  <0.16,0.12,0.0192>,
  <0.18,0.12,0.0216>,
  <0.2,0.12,0.024>,
  <0.22,0.12,0.0264>,
  <0.24,0.12,0.0288>,
  <0.26,0.12,0.0312>,
  <0.28,0.12,0.0336>,
  <0.3,0.12,0.036>,
  <0.32,0.12,0.0384>,
  <0.34,0.12,0.0408>,
  <0.36,0.12,0.0432>,
  <0.38,0.12,0.0456>,
  <0.4,0.12,0.048>,
  <0.42,0.12,0.0504>,
  <0.44,0.12,0.0528>,
  <0.46,0.12,0.0552>,
  <0.48,0.12,0.0576>,
  <0.5,0.12,0.06>,
  <0.52,0.12,0.0624>,
  <0.54,0.12,0.0648>,
  <0.56,0.12,0.0672>,
  <0.58,0.12,0.0696>,
  <0.6,0.12,0.072>,
  <0.62,0.12,0.0744>,
  <0.64,0.12,0.0768>,
  <0.66,0.12,0.0792>,
  <0.68,0.12,0.0816>,
  <0.7,0.12,0.084>,
  <0.72,0.12,0.0864>,
  <0.74,0.12,0.0888>,
  <0.76,0.12,0.0912>,
  <0.78,0.12,0.0936>,
  <0.8,0.12,0.096>,
  <0.82,0.12,0.0984>,
  <0.84,0.12,0.1008>,
  <0.86,0.12,0.1032>,
  <0.88,0.12,0.1056>,
  <0.9,0.12,0.108>,
  <0.92,0.12,0.1104>,
  <0.94,0.12,0.1128>,
  <0.96,0.12,0.1152>,
  <0.98,0.12,0.1176>,
  <1,0.12,0.12>,
  <0,0.14,0>,
  <0.02,0.14,0.0028>,
  <0.04,0.14,0.0056>,
  <0.06,0.14,0.0084>,
  <0.08,0.14,0.0112>,
  <0.1,0.14,0.014>,
  <0.12,0.14,0.0168>,
  <0.14,0.14,0.0196>,
  <0.16,0.14,0.0224>,
  <0.18,0.14,0.0252>,
  <0.2,0.14,0.028>,
  <0.22,0.14,0.0308>,
  <0.24,0.14,0.0336>,
  <0.26,0.14,0.0364>,
  <0.28,0.14,0.0392>,
  <0.3,0.14,0.042>,
  <0.32,0.14,0.0448>,
  <0.34,0.14,0.0476>,
  <0.36,0.14,0.0504>,
  <0.38,0.14,0.0532>,
  <0.4,0.14,0.056>,
  <0.42,0.14,0.0588>,
  <0.44,0.14,0.0616>,
  <0.46,0.14,0.0644>,
  <0.48,0.14,0.0672>,
  <0.5,0.14,0.07>,
  <0.52,0.14,0.0728>,
  <0.54,0.14,0.0756>,
  <0.56,0.14,0.0784>,
  <0.58,0.14,0.0812>,
  <0.6,0.14,0.084>,
  <0.62,0.14,0.0868>,
  <0.64,0.14,0.0896>,
  <0.66,0.14,0.0924>,
  <0.68,0.14,0.0952>,
  <0.7,0.14,0.098>,
  <0.72,0.14,0.1008>,
  <0.74,0.14,0.1036>,
  <0.76,0.14,0.1064>,
  <0.78,0.14,0.1092>,
  <0.8,0.14,0.112>,
  <0.82,0.14,0.1148>,
  <0.84,0.14,0.1176>,
  <0.86,0.14,0.1204>,
  <0.88,0.14,0.1232>,
  <0.9,0.14,0.126>,
  <0.92,0.14,0.1288>,
  <0.94,0.14,0.1316>,
  <0.96,0.14,0.1344>,
  <0.98,0.14,0.1372>,
  <1,0.14,0.14>,
  <0,0.16,0>,
  <0.02,0.16,0.0032>,
  <0.04,0.16,0.0064>,
  <0.06,0.16,0.0096>,
  <0.08,0.16,0.0128>,
  <0.1,0.16,0.016>,
  <0.12,0.16,0.0192>,
  <0.14,0.16,0.0224>,
  <0.16,0.16,0.0256>,
  <0.18,0.16,0.0288>,
  <0.2,0.16,0.032>,
  <0.22,0.16,0.0352>,
  <0.24,0.16,0.0384>,
  <0.26,0.16,0.0416>,
  <0.28,0.16,0.0448>,
  <0.3,0.16,0.048>,
  <0.32,0.16,0.0512>,
  <0.34,0.16,0.0544>,
  <0.36,0.16,0.0576>,
  <0.38,0.16,0.0608>,
  <0.4,0.16,0.064>,
  <0.42,0.16,0.0672>,
  <0.44,0.16,0.0704>,
  <0.46,0.16,0.0736>,
  <0.48,0.16,0.0768>,
  <0.5,0.16,0.08>,
  <0.52,0.16,0.0832>,
  <0.54,0.16,0.0864>,
  <0.56,0.16,0.0896>,
  <0.58,0.16,0.0928>,
  <0.6,0.16,0.096>,
  <0.62,0.16,0.0992>,
  <0.64,0.16,0.1024>,
  <0.66,0.16,0.1056>,
  <0.68,0.16,0.1088>,
  <0.7,0.16,0.112>,
  <0.72,0.16,0.1152>,
  <0.74,0.16,0.1184>,
  <0.76,0.16,0.1216>,
  <0.78,0.16,0.1248>,
  <0.8,0.16,0.128>,
  <0.82,0.16,0.1312>,
  <0.84,0.16,0.1344>,
  <0.86,0.16,0.1376>,
  <0.88,0.16,0.1408>,
  <0.9,0.16,0.144>,
  <0.92,0.16,0.1472>,
  <0.94,0.16,0.1504>,
  <0.96,0.16,0.1536>,
  <0.98,0.16,0.1568>,
  <1,0.16,0.16>,
  <0,0.18,0>,
  <0.02,0.18,0.0036>,
  <0.04,0.18,0.0072>,
  <0.06,0.18,0.0108>,
  <0.08,0.18,0.0144>,
  <0.1,0.18,0.018>,
  <0.12,0.18,0.0216>,
  <0.14,0.18,0.0252>,
  <0.16,0.18,0.0288>,
  <0.18,0.18,0.0324>,
  <0.2,0.18,0.036>,
  <0.22,0.18,0.0396>,
  <0.24,0.18,0.0432>,
  <0.26,0.18,0.0468>,
  <0.28,0.18,0.0504>,
  <0.3,0.18,0.054>,
  <0.32,0.18,0.0576>,
  <0.34,0.18,0.0612>,
  <0.36,0.18,0.0648>,
  <0.38,0.18,0.0684>,
  <0.4,0.18,0.072>,
  <0.42,0.18,0.0756>,
  <0.44,0.18,0.0792>,
  <0.46,0.18,0.0828>,
  <0.48,0.18,0.0864>,
  <0.5,0.18,0.09>,
  <0.52,0.18,0.0936>,
  <0.54,0.18,0.0972>,
  <0.56,0.18,0.1008>,
  <0.58,0.18,0.1044>,
  <0.6,0.18,0.108>,
  <0.62,0.18,0.1116>,
  <0.64,0.18,0.1152>,
  <0.66,0.18,0.1188>,
  <0.68,0.18,0.1224>,
  <0.7,0.18,0.126>,
  <0.72,0.18,0.1296>,
  <0.74,0.18,0.1332>,
  <0.76,0.18,0.1368>,
  <0.78,0.18,0.1404>,
  <0.8,0.18,0.144>,
  <0.82,0.18,0.1476>,
  <0.84,0.18,0.1512>,
  <0.86,0.18,0.1548>,
  <0.88,0.18,0.1584>,
  <0.9,0.18,0.162>,
  <0.92,0.18,0.1656>,
  <0.94,0.18,0.1692>,
  <0.96,0.18,0.1728>,
  <0.98,0.18,0.1764>,
  <1,0.18,0.18>,
  <0,0.2,0>,
  <0.02,0.2,0.004>,
  <0.04,0.2,0.008>,
  <0.06,0.2,0.012>,
  <0.08,0.2,0.016>,
  <0.1,0.2,0.02>,
  <0.12,0.2,0.024>,
  <0.14,0.2,0.028>,
  <0.16,0.2,0.032>,
  <0.18,0.2,0.036>,
  <0.2,0.2,0.04>,
  <0.22,0.2,0.044>,
  <0.24,0.2,0.048>,
  <0.26,0.2,0.052>,
  <0.28,0.2,0.056>,
  <0.3,0.2,0.06>,
  <0.32,0.2,0.064>,
  <0.34,0.2,0.068>,
  <0.36,0.2,0.072>,
  <0.38,0.2,0.076>,
  <0.4,0.2,0.08>,
  <0.42,0.2,0.084>,
  <0.44,0.2,0.088>,
  <0.46,0.2,0.092>,
  <0.48,0.2,0.096>,
  <0.5,0.2,0.1>,
  <0.52,0.2,0.104>,
  <0.54,0.2,0.108>,
  <0.56,0.2,0.112>,
  <0.58,0.2,0.116>,
  <0.6,0.2,0.12>,
  <0.62,0.2,0.124>,
  <0.64,0.2,0.128>,
  <0.66,0.2,0.132>,
  <0.68,0.2,0.136>,
  <0.7,0.2,0.14>,
  <0.72,0.2,0.144>,
  <0.74,0.2,0.148>,
  <0.76,0.2,0.152>,
  <0.78,0.2,0.156>,
  <0.8,0.2,0.16>,
  <0.82,0.2,0.164>,
  <0.84,0.2,0.168>,
  <0.86,0.2,0.172>,
  <0.88,0.2,0.176>,
  <0.9,0.2,0.18>,
  <0.92,0.2,0.184>,
  <0.94,0.2,0.188>,
  <0.96,0.2,0.192>,
  <0.98,0.2,0.196>,
  <1,0.2,0.2>,
  <0,0.22,0>,
  <0.02,0.22,0.0044>,
  <0.04,0.22,0.0088>,
  <0.06,0.22,0.0132>,
  <0.08,0.22,0.0176>,
  <0.1,0.22,0.022>,
  <0.12,0.22,0.0264>,
  <0.14,0.22,0.0308>,
  <0.16,0.22,0.0352>,
  <0.18,0.22,0.0396>,
  <0.2,0.22,0.044>,
  <0.22,0.22,0.0484>,
  <0.24,0.22,0.0528>,
  <0.26,0.22,0.0572>,
  <0.28,0.22,0.0616>,
  <0.3,0.22,0.066>,
  <0.32,0.22,0.0704>,
  <0.34,0.22,0.0748>,
  <0.36,0.22,0.0792>,
  <0.38,0.22,0.0836>,
  <0.4,0.22,0.088>,
  <0.42,0.22,0.0924>,
  <0.44,0.22,0.0968>,
  <0.46,0.22,0.1012>,
  <0.48,0.22,0.1056>,
  <0.5,0.22,0.11>,
  <0.52,0.22,0.1144>,
  <0.54,0.22,0.1188>,
  <0.56,0.22,0.1232>,
  <0.58,0.22,0.1276>,
  <0.6,0.22,0.132>,
  <0.62,0.22,0.1364>,
  <0.64,0.22,0.1408>,
  <0.66,0.22,0.1452>,
  <0.68,0.22,0.1496>,
  <0.7,0.22,0.154>,
  <0.72,0.22,0.1584>,
  <0.74,0.22,0.1628>,
  <0.76,0.22,0.1672>,
  <0.78,0.22,0.1716>,
  <0.8,0.22,0.176>,
  <0.82,0.22,0.1804>,
  <0.84,0.22,0.1848>,
  <0.86,0.22,0.1892>,
  <0.88,0.22,0.1936>,
  <0.9,0.22,0.198>,
  <0.92,0.22,0.2024>,
  <0.94,0.22,0.2068>,
  <0.96,0.22,0.2112>,
  <0.98,0.22,0.2156>,
  <1,0.22,0.22>,
  <0,0.24,0>,
  <0.02,0.24,0.0048>,
  <0.04,0.24,0.0096>,
  <0.06,0.24,0.0144>,
  <0.08,0.24,0.0192>,
  <0.1,0.24,0.024>,
  <0.12,0.24,0.0288>,
  <0.14,0.24,0.0336>,
  <0.16,0.24,0.0384>,
  <0.18,0.24,0.0432>,
  <0.2,0.24,0.048>,
  <0.22,0.24,0.0528>,
  <0.24,0.24,0.0576>,
  <0.26,0.24,0.0624>,
  <0.28,0.24,0.0672>,
  <0.3,0.24,0.072>,
  <0.32,0.24,0.0768>,
  <0.34,0.24,0.0816>,
  <0.36,0.24,0.0864>,
  <0.38,0.24,0.0912>,
  <0.4,0.24,0.096>,
  <0.42,0.24,0.1008>,
  <0.44,0.24,0.1056>,
  <0.46,0.24,0.1104>,
  <0.48,0.24,0.1152>,
  <0.5,0.24,0.12>,
  <0.52,0.24,0.1248>,
  <0.54,0.24,0.1296>,
  <0.56,0.24,0.1344>,
  <0.58,0.24,0.1392>,
  <0.6,0.24,0.144>,
  <0.62,0.24,0.1488>,
  <0.64,0.24,0.1536>,
  <0.66,0.24,0.1584>,
  <0.68,0.24,0.1632>,
  <0.7,0.24,0.168>,
  <0.72,0.24,0.1728>,
  <0.74,0.24,0.1776>,
  <0.76,0.24,0.1824>,
  <0.78,0.24,0.1872>,
  <0.8,0.24,0.192>,
  <0.82,0.24,0.1968>,
  <0.84,0.24,0.2016>,
  <0.86,0.24,0.2064>,
  <0.88,0.24,0.2112>,
  <0.9,0.24,0.216>,
  <0.92,0.24,0.2208>,
  <0.94,0.24,0.2256>,
  <0.96,0.24,0.2304>,
  <0.98,0.24,0.2352>,
  <1,0.24,0.24>,
  <0,0.26,0>,
  <0.02,0.26,0.0052>,
  <0.04,0.26,0.0104>,
  <0.06,0.26,0.0156>,
  <0.08,0.26,0.0208>,
  <0.1,0.26,0.026>,
  <0.12,0.26,0.0312>,
  <0.14,0.26,0.0364>,
  <0.16,0.26,0.0416>,
  <0.18,0.26,0.0468>,
  <0.2,0.26,0.052>,
  <0.22,0.26,0.0572>,
  <0.24,0.26,0.0624>,
  <0.26,0.26,0.0676>,
  <0.28,0.26,0.0728>,
  <0.3,0.26,0.078>,
  <0.32,0.26,0.0832>,
  <0.34,0.26,0.0884>,
  <0.36,0.26,0.0936>,
  <0.38,0.26,0.0988>,
  <0.4,0.26,0.104>,
  <0.42,0.26,0.1092>,
  <0.44,0.26,0.1144>,
  <0.46,0.26,0.1196>,
  <0.48,0.26,0.1248>,
  <0.5,0.26,0.13>,
  <0.52,0.26,0.1352>,
  <0.54,0.26,0.1404>,
  <0.56,0.26,0.1456>,
  <0.58,0.26,0.1508>,
  <0.6,0.26,0.156>,
  <0.62,0.26,0.1612>,
  <0.64,0.26,0.1664>,
  <0.66,0.26,0.1716>,
  <0.68,0.26,0.1768>,
  <0.7,0.26,0.182>,
  <0.72,0.26,0.1872>,
  <0.74,0.26,0.1924>,
  <0.76,0.26,0.1976>,
  <0.78,0.26,0.2028>,
  <0.8,0.26,0.208>,
  <0.82,0.26,0.2132>,
  <0.84,0.26,0.2184>,
  <0.86,0.26,0.2236>,
  <0.88,0.26,0.2288>,
  <0.9,0.26,0.234>,
  <0.92,0.26,0.2392>,
  <0.94,0.26,0.2444>,
  <0.96,0.26,0.2496>,
  <0.98,0.26,0.2548>,
  <1,0.26,0.26>,
  <0,0.28,0>,
  <0.02,0.28,0.0056>,
  <0.04,0.28,0.0112>,
  <0.06,0.28,0.0168>,
  <0.08,0.28,0.0224>,
  <0.1,0.28,0.028>,
  <0.12,0.28,0.0336>,
  <0.14,0.28,0.0392>,
  <0.16,0.28,0.0448>,
  <0.18,0.28,0.0504>,
  <0.2,0.28,0.056>,
  <0.22,0.28,0.0616>,
  <0.24,0.28,0.0672>,
  <0.26,0.28,0.0728>,
  <0.28,0.28,0.0784>,
  <0.3,0.28,0.084>,
  <0.32,0.28,0.0896>,
  <0.34,0.28,0.0952>,
  <0.36,0.28,0.1008>,
  <0.38,0.28,0.1064>,
  <0.4,0.28,0.112>,
  <0.42,0.28,0.1176>,
  <0.44,0.28,0.1232>,
  <0.46,0.28,0.1288>,
  <0.48,0.28,0.1344>,
  <0.5,0.28,0.14>,
  <0.52,0.28,0.1456>,
  <0.54,0.28,0.1512>,
  <0.56,0.28,0.1568>,
  <0.58,0.28,0.1624>,
  <0.6,0.28,0.168>,
  <0.62,0.28,0.1736>,
  <0.64,0.28,0.1792>,
  <0.66,0.28,0.1848>,
  <0.68,0.28,0.1904>,
  <0.7,0.28,0.196>,
  <0.72,0.28,0.2016>,
  <0.74,0.28,0.2072>,
  <0.76,0.28,0.2128>,
  <0.78,0.28,0.2184>,
  <0.8,0.28,0.224>,
  <0.82,0.28,0.2296>,
  <0.84,0.28,0.2352>,
  <0.86,0.28,0.2408>,
  <0.88,0.28,0.2464>,
  <0.9,0.28,0.252>,
  <0.92,0.28,0.2576>,
  <0.94,0.28,0.2632>,
  <0.96,0.28,0.2688>,
  <0.98,0.28,0.2744>,
  <1,0.28,0.28>,
  <0,0.3,0>,
  <0.02,0.3,0.006>,
  <0.04,0.3,0.012>,
  <0.06,0.3,0.018>,
  <0.08,0.3,0.024>,
  <0.1,0.3,0.03>,
  <0.12,0.3,0.036>,
  <0.14,0.3,0.042>,
  <0.16,0.3,0.048>,
  <0.18,0.3,0.054>,
  <0.2,0.3,0.06>,
  <0.22,0.3,0.066>,
  <0.24,0.3,0.072>,
  <0.26,0.3,0.078>,
  <0.28,0.3,0.084>,
  <0.3,0.3,0.09>,
  <0.32,0.3,0.096>,
  <0.34,0.3,0.102>,
  <0.36,0.3,0.108>,
  <0.38,0.3,0.114>,
  <0.4,0.3,0.12>,
  <0.42,0.3,0.126>,
  <0.44,0.3,0.132>,
  <0.46,0.3,0.138>,
  <0.48,0.3,0.144>,
  <0.5,0.3,0.15>,
  <0.52,0.3,0.156>,
  <0.54,0.3,0.162>,
  <0.56,0.3,0.168>,
  <0.58,0.3,0.174>,
  <0.6,0.3,0.18>,
  <0.62,0.3,0.186>,
  <0.64,0.3,0.192>,
  <0.66,0.3,0.198>,
  <0.68,0.3,0.204>,
  <0.7,0.3,0.21>,
  <0.72,0.3,0.216>,
  <0.74,0.3,0.222>,
  <0.76,0.3,0.228>,
  <0.78,0.3,0.234>,
  <0.8,0.3,0.24>,
  <0.82,0.3,0.246>,
  <0.84,0.3,0.252>,
  <0.86,0.3,0.258>,
  <0.88,0.3,0.264>,
  <0.9,0.3,0.27>,
  <0.92,0.3,0.276>,
  <0.94,0.3,0.282>,
  <0.96,0.3,0.288>,
  <0.98,0.3,0.294>,
  <1,0.3,0.3>,
  <0,0.32,0>,
  <0.02,0.32,0.0064>,
  <0.04,0.32,0.0128>,
  <0.06,0.32,0.0192>,
  <0.08,0.32,0.0256>,
  <0.1,0.32,0.032>,
  <0.12,0.32,0.0384>,
  <0.14,0.32,0.0448>,
  <0.16,0.32,0.0512>,
  <0.18,0.32,0.0576>,
  <0.2,0.32,0.064>,
  <0.22,0.32,0.0704>,
  <0.24,0.32,0.0768>,
  <0.26,0.32,0.0832>,
  <0.28,0.32,0.0896>,
  <0.3,0.32,0.096>,
  <0.32,0.32,0.1024>,
  <0.34,0.32,0.1088>,
  <0.36,0.32,0.1152>,
  <0.38,0.32,0.1216>,
  <0.4,0.32,0.128>,
  <0.42,0.32,0.1344>,
  <0.44,0.32,0.1408>,
  <0.46,0.32,0.1472>,
  <0.48,0.32,0.1536>,
  <0.5,0.32,0.16>,
  <0.52,0.32,0.1664>,
  <0.54,0.32,0.1728>,
  <0.56,0.32,0.1792>,
  <0.58,0.32,0.1856>,
  <0.6,0.32,0.192>,
  <0.62,0.32,0.1984>,
  <0.64,0.32,0.2048>,
  <0.66,0.32,0.2112>,
  <0.68,0.32,0.2176>,
  <0.7,0.32,0.224>,
  <0.72,0.32,0.2304>,
  <0.74,0.32,0.2368>,
  <0.76,0.32,0.2432>,
  <0.78,0.32,0.2496>,
  <0.8,0.32,0.256>,
  <0.82,0.32,0.2624>,
  <0.84,0.32,0.2688>,
  <0.86,0.32,0.2752>,
  <0.88,0.32,0.2816>,
  <0.9,0.32,0.288>,
  <0.92,0.32,0.2944>,
  <0.94,0.32,0.3008>,
  <0.96,0.32,0.3072>,
  <0.98,0.32,0.3136>,
  <1,0.32,0.32>,
  <0,0.34,0>,
  <0.02,0.34,0.0068>,
  <0.04,0.34,0.0136>,
  <0.06,0.34,0.0204>,
  <0.08,0.34,0.0272>,
  <0.1,0.34,0.034>,
  <0.12,0.34,0.0408>,
  <0.14,0.34,0.0476>,
  <0.16,0.34,0.0544>,
  <0.18,0.34,0.0612>,
  <0.2,0.34,0.068>,
  <0.22,0.34,0.0748>,
  <0.24,0.34,0.0816>,
  <0.26,0.34,0.0884>,
  <0.28,0.34,0.0952>,
  <0.3,0.34,0.102>,
  <0.32,0.34,0.1088>,
  <0.34,0.34,0.1156>,
  <0.36,0.34,0.1224>,
  <0.38,0.34,0.1292>,
  <0.4,0.34,0.136>,
  <0.42,0.34,0.1428>,
  <0.44,0.34,0.1496>,
  <0.46,0.34,0.1564>,
  <0.48,0.34,0.1632>,
  <0.5,0.34,0.17>,
  <0.52,0.34,0.1768>,
  <0.54,0.34,0.1836>,
  <0.56,0.34,0.1904>,
  <0.58,0.34,0.1972>,
  <0.6,0.34,0.204>,
  <0.62,0.34,0.2108>,
  <0.64,0.34,0.2176>,
  <0.66,0.34,0.2244>,
  <0.68,0.34,0.2312>,
  <0.7,0.34,0.238>,
  <0.72,0.34,0.2448>,
  <0.74,0.34,0.2516>,
  <0.76,0.34,0.2584>,
  <0.78,0.34,0.2652>,
  <0.8,0.34,0.272>,
  <0.82,0.34,0.2788>,
  <0.84,0.34,0.2856>,
  <0.86,0.34,0.2924>,
  <0.88,0.34,0.2992>,
  <0.9,0.34,0.306>,
  <0.92,0.34,0.3128>,
  <0.94,0.34,0.3196>,
  <0.96,0.34,0.3264>,
  <0.98,0.34,0.3332>,
  <1,0.34,0.34>,
  <0,0.36,0>,
  <0.02,0.36,0.0072>,
  <0.04,0.36,0.0144>,
  <0.06,0.36,0.0216>,
  <0.08,0.36,0.0288>,
  <0.1,0.36,0.036>,
  <0.12,0.36,0.0432>,
  <0.14,0.36,0.0504>,
  <0.16,0.36,0.0576>,
  <0.18,0.36,0.0648>,
  <0.2,0.36,0.072>,
  <0.22,0.36,0.0792>,
  <0.24,0.36,0.0864>,
  <0.26,0.36,0.0936>,
  <0.28,0.36,0.1008>,
  <0.3,0.36,0.108>,
  <0.32,0.36,0.1152>,
  <0.34,0.36,0.1224>,
  <0.36,0.36,0.1296>,
  <0.38,0.36,0.1368>,
  <0.4,0.36,0.144>,
  <0.42,0.36,0.1512>,
  <0.44,0.36,0.1584>,
  <0.46,0.36,0.1656>,
  <0.48,0.36,0.1728>,
  <0.5,0.36,0.18>,
  <0.52,0.36,0.1872>,
  <0.54,0.36,0.1944>,
  <0.56,0.36,0.2016>,
  <0.58,0.36,0.2088>,
  <0.6,0.36,0.216>,
  <0.62,0.36,0.2232>,
  <0.64,0.36,0.2304>,
  <0.66,0.36,0.2376>,
  <0.68,0.36,0.2448>,
  <0.7,0.36,0.252>,
  <0.72,0.36,0.2592>,
  <0.74,0.36,0.2664>,
  <0.76,0.36,0.2736>,
  <0.78,0.36,0.2808>,
  <0.8,0.36,0.288>,
  <0.82,0.36,0.2952>,
  <0.84,0.36,0.3024>,
  <0.86,0.36,0.3096>,
  <0.88,0.36,0.3168>,
  <0.9,0.36,0.324>,
  <0.92,0.36,0.3312>,
  <0.94,0.36,0.3384>,
  <0.96,0.36,0.3456>,
  <0.98,0.36,0.3528>,
  <1,0.36,0.36>,
  <0,0.38,0>,
  <0.02,0.38,0.0076>,
  <0.04,0.38,0.0152>,
  <0.06,0.38,0.0228>,
  <0.08,0.38,0.0304>,
  <0.1,0.38,0.038>,
  <0.12,0.38,0.0456>,
  <0.14,0.38,0.0532>,
  <0.16,0.38,0.0608>,
  <0.18,0.38,0.0684>,
  <0.2,0.38,0.076>,
  <0.22,0.38,0.0836>,
  <0.24,0.38,0.0912>,
  <0.26,0.38,0.0988>,
  <0.28,0.38,0.1064>,
  <0.3,0.38,0.114>,
  <0.32,0.38,0.1216>,
  <0.34,0.38,0.1292>,
  <0.36,0.38,0.1368>,
  <0.38,0.38,0.1444>,
  <0.4,0.38,0.152>,
  <0.42,0.38,0.1596>,
  <0.44,0.38,0.1672>,
  <0.46,0.38,0.1748>,
  <0.48,0.38,0.1824>,
  <0.5,0.38,0.19>,
  <0.52,0.38,0.1976>,
  <0.54,0.38,0.2052>,
  <0.56,0.38,0.2128>,
  <0.58,0.38,0.2204>,
  <0.6,0.38,0.228>,
  <0.62,0.38,0.2356>,
  <0.64,0.38,0.2432>,
  <0.66,0.38,0.2508>,
  <0.68,0.38,0.2584>,
  <0.7,0.38,0.266>,
  <0.72,0.38,0.2736>,
  <0.74,0.38,0.2812>,
  <0.76,0.38,0.2888>,
  <0.78,0.38,0.2964>,
  <0.8,0.38,0.304>,
  <0.82,0.38,0.3116>,
  <0.84,0.38,0.3192>,
  <0.86,0.38,0.3268>,
  <0.88,0.38,0.3344>,
  <0.9,0.38,0.342>,
  <0.92,0.38,0.3496>,
  <0.94,0.38,0.3572>,
  <0.96,0.38,0.3648>,
  <0.98,0.38,0.3724>,
  <1,0.38,0.38>,
  <0,0.4,0>,
  <0.02,0.4,0.008>,
  <0.04,0.4,0.016>,
  <0.06,0.4,0.024>,
  <0.08,0.4,0.032>,
  <0.1,0.4,0.04>,
  <0.12,0.4,0.048>,
  <0.14,0.4,0.056>,
  <0.16,0.4,0.064>,
  <0.18,0.4,0.072>,
  <0.2,0.4,0.08>,
  <0.22,0.4,0.088>,
  <0.24,0.4,0.096>,
  <0.26,0.4,0.104>,
  <0.28,0.4,0.112>,
  <0.3,0.4,0.12>,
  <0.32,0.4,0.128>,
  <0.34,0.4,0.136>,
  <0.36,0.4,0.144>,
  <0.38,0.4,0.152>,
  <0.4,0.4,0.16>,
  <0.42,0.4,0.168>,
  <0.44,0.4,0.176>,
  <0.46,0.4,0.184>,
  <0.48,0.4,0.192>,
  <0.5,0.4,0.2>,
  <0.52,0.4,0.208>,
  <0.54,0.4,0.216>,
  <0.56,0.4,0.224>,
  <0.58,0.4,0.232>,
  <0.6,0.4,0.24>,
  <0.62,0.4,0.248>,
  <0.64,0.4,0.256>,
  <0.66,0.4,0.264>,
  <0.68,0.4,0.272>,
  <0.7,0.4,0.28>,
  <0.72,0.4,0.288>,
  <0.74,0.4,0.296>,
  <0.76,0.4,0.304>,
  <0.78,0.4,0.312>,
  <0.8,0.4,0.32>,
  <0.82,0.4,0.328>,
  <0.84,0.4,0.336>,
  <0.86,0.4,0.344>,
  <0.88,0.4,0.352>,
  <0.9,0.4,0.36>,
  <0.92,0.4,0.368>,
  <0.94,0.4,0.376>,
  <0.96,0.4,0.384>,
  <0.98,0.4,0.392>,
  <1,0.4,0.4>,
  <0,0.42,0>,
  <0.02,0.42,0.0084>,
  <0.04,0.42,0.0168>,
  <0.06,0.42,0.0252>,
  <0.08,0.42,0.0336>,
  <0.1,0.42,0.042>,
  <0.12,0.42,0.0504>,
  <0.14,0.42,0.0588>,
  <0.16,0.42,0.0672>,
  <0.18,0.42,0.0756>,
  <0.2,0.42,0.084>,
  <0.22,0.42,0.0924>,
  <0.24,0.42,0.1008>,
  <0.26,0.42,0.1092>,
  <0.28,0.42,0.1176>,
  <0.3,0.42,0.126>,
  <0.32,0.42,0.1344>,
  <0.34,0.42,0.1428>,
  <0.36,0.42,0.1512>,
  <0.38,0.42,0.1596>,
  <0.4,0.42,0.168>,
  <0.42,0.42,0.1764>,
  <0.44,0.42,0.1848>,
  <0.46,0.42,0.1932>,
  <0.48,0.42,0.2016>,
  <0.5,0.42,0.21>,
  <0.52,0.42,0.2184>,
  <0.54,0.42,0.2268>,
  <0.56,0.42,0.2352>,
  <0.58,0.42,0.2436>,
  <0.6,0.42,0.252>,
  <0.62,0.42,0.2604>,
  <0.64,0.42,0.2688>,
  <0.66,0.42,0.2772>,
  <0.68,0.42,0.2856>,
  <0.7,0.42,0.294>,
  <0.72,0.42,0.3024>,
  <0.74,0.42,0.3108>,
  <0.76,0.42,0.3192>,
  <0.78,0.42,0.3276>,
  <0.8,0.42,0.336>,
  <0.82,0.42,0.3444>,
  <0.84,0.42,0.3528>,
  <0.86,0.42,0.3612>,
  <0.88,0.42,0.3696>,
  <0.9,0.42,0.378>,
  <0.92,0.42,0.3864>,
  <0.94,0.42,0.3948>,
  <0.96,0.42,0.4032>,
  <0.98,0.42,0.4116>,
  <1,0.42,0.42>,
  <0,0.44,0>,
  <0.02,0.44,0.0088>,
  <0.04,0.44,0.0176>,
  <0.06,0.44,0.0264>,
  <0.08,0.44,0.0352>,
  <0.1,0.44,0.044>,
  <0.12,0.44,0.0528>,
  <0.14,0.44,0.0616>,
  <0.16,0.44,0.0704>,
  <0.18,0.44,0.0792>,
  <0.2,0.44,0.088>,
  <0.22,0.44,0.0968>,
  <0.24,0.44,0.1056>,
  <0.26,0.44,0.1144>,
  <0.28,0.44,0.1232>,
  <0.3,0.44,0.132>,
  <0.32,0.44,0.1408>,
  <0.34,0.44,0.1496>,
  <0.36,0.44,0.1584>,
  <0.38,0.44,0.1672>,
  <0.4,0.44,0.176>,
  <0.42,0.44,0.1848>,
  <0.44,0.44,0.1936>,
  <0.46,0.44,0.2024>,
  <0.48,0.44,0.2112>,
  <0.5,0.44,0.22>,
  <0.52,0.44,0.2288>,
  <0.54,0.44,0.2376>,
  <0.56,0.44,0.2464>,
  <0.58,0.44,0.2552>,
  <0.6,0.44,0.264>,
  <0.62,0.44,0.2728>,
  <0.64,0.44,0.2816>,
  <0.66,0.44,0.2904>,
  <0.68,0.44,0.2992>,
  <0.7,0.44,0.308>,
  <0.72,0.44,0.3168>,
  <0.74,0.44,0.3256>,
  <0.76,0.44,0.3344>,
  <0.78,0.44,0.3432>,
  <0.8,0.44,0.352>,
  <0.82,0.44,0.3608>,
  <0.84,0.44,0.3696>,
  <0.86,0.44,0.3784>,
  <0.88,0.44,0.3872>,
  <0.9,0.44,0.396>,
  <0.92,0.44,0.4048>,
  <0.94,0.44,0.4136>,
  <0.96,0.44,0.4224>,
  <0.98,0.44,0.4312>,
  <1,0.44,0.44>,
  <0,0.46,0>,
  <0.02,0.46,0.0092>,
  <0.04,0.46,0.0184>,
  <0.06,0.46,0.0276>,
  <0.08,0.46,0.0368>,
  <0.1,0.46,0.046>,
  <0.12,0.46,0.0552>,
  <0.14,0.46,0.0644>,
  <0.16,0.46,0.0736>,
  <0.18,0.46,0.0828>,
  <0.2,0.46,0.092>,
  <0.22,0.46,0.1012>,
  <0.24,0.46,0.1104>,
  <0.26,0.46,0.1196>,
  <0.28,0.46,0.1288>,
  <0.3,0.46,0.138>,
  <0.32,0.46,0.1472>,
  <0.34,0.46,0.1564>,
  <0.36,0.46,0.1656>,
  <0.38,0.46,0.1748>,
  <0.4,0.46,0.184>,
  <0.42,0.46,0.1932>,
  <0.44,0.46,0.2024>,
  <0.46,0.46,0.2116>,
  <0.48,0.46,0.2208>,
  <0.5,0.46,0.23>,
  <0.52,0.46,0.2392>,
  <0.54,0.46,0.2484>,
  <0.56,0.46,0.2576>,
  <0.58,0.46,0.2668>,
  <0.6,0.46,0.276>,
  <0.62,0.46,0.2852>,
  <0.64,0.46,0.2944>,
  <0.66,0.46,0.3036>,
  <0.68,0.46,0.3128>,
  <0.7,0.46,0.322>,
  <0.72,0.46,0.3312>,
  <0.74,0.46,0.3404>,
  <0.76,0.46,0.3496>,
  <0.78,0.46,0.3588>,
  <0.8,0.46,0.368>,
  <0.82,0.46,0.3772>,
  <0.84,0.46,0.3864>,
  <0.86,0.46,0.3956>,
  <0.88,0.46,0.4048>,
  <0.9,0.46,0.414>,
  <0.92,0.46,0.4232>,
  <0.94,0.46,0.4324>,
  <0.96,0.46,0.4416>,
  <0.98,0.46,0.4508>,
  <1,0.46,0.46>,
  <0,0.48,0>,
  <0.02,0.48,0.0096>,
  <0.04,0.48,0.0192>,
  <0.06,0.48,0.0288>,
  <0.08,0.48,0.0384>,
  <0.1,0.48,0.048>,
  <0.12,0.48,0.0576>,
  <0.14,0.48,0.0672>,
  <0.16,0.48,0.0768>,
  <0.18,0.48,0.0864>,
  <0.2,0.48,0.096>,
  <0.22,0.48,0.1056>,
  <0.24,0.48,0.1152>,
  <0.26,0.48,0.1248>,
  <0.28,0.48,0.1344>,
  <0.3,0.48,0.144>,
  <0.32,0.48,0.1536>,
  <0.34,0.48,0.1632>,
  <0.36,0.48,0.1728>,
  <0.38,0.48,0.1824>,
  <0.4,0.48,0.192>,
  <0.42,0.48,0.2016>,
  <0.44,0.48,0.2112>,
  <0.46,0.48,0.2208>,
  <0.48,0.48,0.2304>,
  <0.5,0.48,0.24>,
  <0.52,0.48,0.2496>,
  <0.54,0.48,0.2592>,
  <0.56,0.48,0.2688>,
  <0.58,0.48,0.2784>,
  <0.6,0.48,0.288>,
  <0.62,0.48,0.2976>,
  <0.64,0.48,0.3072>,
  <0.66,0.48,0.3168>,
  <0.68,0.48,0.3264>,
  <0.7,0.48,0.336>,
  <0.72,0.48,0.3456>,
  <0.74,0.48,0.3552>,
  <0.76,0.48,0.3648>,
  <0.78,0.48,0.3744>,
  <0.8,0.48,0.384>,
  <0.82,0.48,0.3936>,
  <0.84,0.48,0.4032>,
  <0.86,0.48,0.4128>,
  <0.88,0.48,0.4224>,
  <0.9,0.48,0.432>,
  <0.92,0.48,0.4416>,
  <0.94,0.48,0.4512>,
  <0.96,0.48,0.4608>,
  <0.98,0.48,0.4704>,
  <1,0.48,0.48>,
  <0,0.5,0>,
  <0.02,0.5,0.01>,
  <0.04,0.5,0.02>,
  <0.06,0.5,0.03>,
  <0.08,0.5,0.04>,
  <0.1,0.5,0.05>,
  <0.12,0.5,0.06>,
  <0.14,0.5,0.07>,
  <0.16,0.5,0.08>,
  <0.18,0.5,0.09>,
  <0.2,0.5,0.1>,
  <0.22,0.5,0.11>,
  <0.24,0.5,0.12>,
  <0.26,0.5,0.13>,
  <0.28,0.5,0.14>,
  <0.3,0.5,0.15>,
  <0.32,0.5,0.16>,
  <0.34,0.5,0.17>,
  <0.36,0.5,0.18>,
  <0.38,0.5,0.19>,
  <0.4,0.5,0.2>,
  <0.42,0.5,0.21>,
  <0.44,0.5,0.22>,
  <0.46,0.5,0.23>,
  <0.48,0.5,0.24>,
  <0.5,0.5,0.25>,
  <0.52,0.5,0.26>,
  <0.54,0.5,0.27>,
  <0.56,0.5,0.28>,
  <0.58,0.5,0.29>,
  <0.6,0.5,0.3>,
  <0.62,0.5,0.31>,
  <0.64,0.5,0.32>,
  <0.66,0.5,0.33>,
  <0.68,0.5,0.34>,
  <0.7,0.5,0.35>,
  <0.72,0.5,0.36>,
  <0.74,0.5,0.37>,
  <0.76,0.5,0.38>,
  <0.78,0.5,0.39>,
  <0.8,0.5,0.4>,
  <0.82,0.5,0.41>,
  <0.84,0.5,0.42>,
  <0.86,0.5,0.43>,
  <0.88,0.5,0.44>,
  <0.9,0.5,0.45>,
  <0.92,0.5,0.46>,
  <0.94,0.5,0.47>,
  <0.96,0.5,0.48>,
  <0.98,0.5,0.49>,
  <1,0.5,0.5>,
  <0,0.52,0>,
  <0.02,0.52,0.0104>,
  <0.04,0.52,0.0208>,
  <0.06,0.52,0.0312>,
  <0.08,0.52,0.0416>,
  <0.1,0.52,0.052>,
  <0.12,0.52,0.0624>,
  <0.14,0.52,0.0728>,
  <0.16,0.52,0.0832>,
  <0.18,0.52,0.0936>,
  <0.2,0.52,0.104>,
  <0.22,0.52,0.1144>,
  <0.24,0.52,0.1248>,
  <0.26,0.52,0.1352>,
  <0.28,0.52,0.1456>,
  <0.3,0.52,0.156>,
  <0.32,0.52,0.1664>,
  <0.34,0.52,0.1768>,
  <0.36,0.52,0.1872>,
  <0.38,0.52,0.1976>,
  <0.4,0.52,0.208>,
  <0.42,0.52,0.2184>,
  <0.44,0.52,0.2288>,
  <0.46,0.52,0.2392>,
  <0.48,0.52,0.2496>,
  <0.5,0.52,0.26>,
  <0.52,0.52,0.2704>,
  <0.54,0.52,0.2808>,
  <0.56,0.52,0.2912>,
  <0.58,0.52,0.3016>,
  <0.6,0.52,0.312>,
  <0.62,0.52,0.3224>,
  <0.64,0.52,0.3328>,
  <0.66,0.52,0.3432>,
  <0.68,0.52,0.3536>,
  <0.7,0.52,0.364>,
  <0.72,0.52,0.3744>,
  <0.74,0.52,0.3848>,
  <0.76,0.52,0.3952>,
  <0.78,0.52,0.4056>,
  <0.8,0.52,0.416>,
  <0.82,0.52,0.4264>,
  <0.84,0.52,0.4368>,
  <0.86,0.52,0.4472>,
  <0.88,0.52,0.4576>,
  <0.9,0.52,0.468>,
  <0.92,0.52,0.4784>,
  <0.94,0.52,0.4888>,
  <0.96,0.52,0.4992>,
  <0.98,0.52,0.5096>,
  <1,0.52,0.52>,
  <0,0.54,0>,
  <0.02,0.54,0.0108>,
  <0.04,0.54,0.0216>,
  <0.06,0.54,0.0324>,
  <0.08,0.54,0.0432>,
  <0.1,0.54,0.054>,
  <0.12,0.54,0.0648>,
  <0.14,0.54,0.0756>,
  <0.16,0.54,0.0864>,
  <0.18,0.54,0.0972>,
  <0.2,0.54,0.108>,
  <0.22,0.54,0.1188>,
  <0.24,0.54,0.1296>,
  <0.26,0.54,0.1404>,
  <0.28,0.54,0.1512>,
  <0.3,0.54,0.162>,
  <0.32,0.54,0.1728>,
  <0.34,0.54,0.1836>,
  <0.36,0.54,0.1944>,
  <0.38,0.54,0.2052>,
  <0.4,0.54,0.216>,
  <0.42,0.54,0.2268>,
  <0.44,0.54,0.2376>,
  <0.46,0.54,0.2484>,
  <0.48,0.54,0.2592>,
  <0.5,0.54,0.27>,
  <0.52,0.54,0.2808>,
  <0.54,0.54,0.2916>,
  <0.56,0.54,0.3024>,
  <0.58,0.54,0.3132>,
  <0.6,0.54,0.324>,
  <0.62,0.54,0.3348>,
  <0.64,0.54,0.3456>,
  <0.66,0.54,0.3564>,
  <0.68,0.54,0.3672>,
  <0.7,0.54,0.378>,
  <0.72,0.54,0.3888>,
  <0.74,0.54,0.3996>,
  <0.76,0.54,0.4104>,
  <0.78,0.54,0.4212>,
  <0.8,0.54,0.432>,
  <0.82,0.54,0.4428>,
  <0.84,0.54,0.4536>,
  <0.86,0.54,0.4644>,
  <0.88,0.54,0.4752>,
  <0.9,0.54,0.486>,
  <0.92,0.54,0.4968>,
  <0.94,0.54,0.5076>,
  <0.96,0.54,0.5184>,
  <0.98,0.54,0.5292>,
  <1,0.54,0.54>,
  <0,0.56,0>,
  <0.02,0.56,0.0112>,
  <0.04,0.56,0.0224>,
  <0.06,0.56,0.0336>,
  <0.08,0.56,0.0448>,
  <0.1,0.56,0.056>,
  <0.12,0.56,0.0672>,
  <0.14,0.56,0.0784>,
  <0.16,0.56,0.0896>,
  <0.18,0.56,0.1008>,
  <0.2,0.56,0.112>,
  <0.22,0.56,0.1232>,
  <0.24,0.56,0.1344>,
  <0.26,0.56,0.1456>,
  <0.28,0.56,0.1568>,
  <0.3,0.56,0.168>,
  <0.32,0.56,0.1792>,
  <0.34,0.56,0.1904>,
  <0.36,0.56,0.2016>,
  <0.38,0.56,0.2128>,
  <0.4,0.56,0.224>,
  <0.42,0.56,0.2352>,
  <0.44,0.56,0.2464>,
  <0.46,0.56,0.2576>,
  <0.48,0.56,0.2688>,
  <0.5,0.56,0.28>,
  <0.52,0.56,0.2912>,
  <0.54,0.56,0.3024>,
  <0.56,0.56,0.3136>,
  <0.58,0.56,0.3248>,
  <0.6,0.56,0.336>,
  <0.62,0.56,0.3472>,
  <0.64,0.56,0.3584>,
  <0.66,0.56,0.3696>,
  <0.68,0.56,0.3808>,
  <0.7,0.56,0.392>,
  <0.72,0.56,0.4032>,
  <0.74,0.56,0.4144>,
  <0.76,0.56,0.4256>,
  <0.78,0.56,0.4368>,
  <0.8,0.56,0.448>,
  <0.82,0.56,0.4592>,
  <0.84,0.56,0.4704>,
  <0.86,0.56,0.4816>,
  <0.88,0.56,0.4928>,
  <0.9,0.56,0.504>,
  <0.92,0.56,0.5152>,
  <0.94,0.56,0.5264>,
  <0.96,0.56,0.5376>,
  <0.98,0.56,0.5488>,
  <1,0.56,0.56>,
  <0,0.58,0>,
  <0.02,0.58,0.0116>,
  <0.04,0.58,0.0232>,
  <0.06,0.58,0.0348>,
  <0.08,0.58,0.0464>,
  <0.1,0.58,0.058>,
  <0.12,0.58,0.0696>,
  <0.14,0.58,0.0812>,
  <0.16,0.58,0.0928>,
  <0.18,0.58,0.1044>,
  <0.2,0.58,0.116>,
  <0.22,0.58,0.1276>,
  <0.24,0.58,0.1392>,
  <0.26,0.58,0.1508>,
  <0.28,0.58,0.1624>,
  <0.3,0.58,0.174>,
  <0.32,0.58,0.1856>,
  <0.34,0.58,0.1972>,
  <0.36,0.58,0.2088>,
  <0.38,0.58,0.2204>,
  <0.4,0.58,0.232>,
  <0.42,0.58,0.2436>,
  <0.44,0.58,0.2552>,
  <0.46,0.58,0.2668>,
  <0.48,0.58,0.2784>,
  <0.5,0.58,0.29>,
  <0.52,0.58,0.3016>,
  <0.54,0.58,0.3132>,
  <0.56,0.58,0.3248>,
  <0.58,0.58,0.3364>,
  <0.6,0.58,0.348>,
  <0.62,0.58,0.3596>,
  <0.64,0.58,0.3712>,
  <0.66,0.58,0.3828>,
  <0.68,0.58,0.3944>,
  <0.7,0.58,0.406>,
  <0.72,0.58,0.4176>,
  <0.74,0.58,0.4292>,
  <0.76,0.58,0.4408>,
  <0.78,0.58,0.4524>,
  <0.8,0.58,0.464>,
  <0.82,0.58,0.4756>,
  <0.84,0.58,0.4872>,
  <0.86,0.58,0.4988>,
  <0.88,0.58,0.5104>,
  <0.9,0.58,0.522>,
  <0.92,0.58,0.5336>,
  <0.94,0.58,0.5452>,
  <0.96,0.58,0.5568>,
  <0.98,0.58,0.5684>,
  <1,0.58,0.58>,
  <0,0.6,0>,
  <0.02,0.6,0.012>,
  <0.04,0.6,0.024>,
  <0.06,0.6,0.036>,
  <0.08,0.6,0.048>,
  <0.1,0.6,0.06>,
  <0.12,0.6,0.072>,
  <0.14,0.6,0.084>,
  <0.16,0.6,0.096>,
  <0.18,0.6,0.108>,
  <0.2,0.6,0.12>,
  <0.22,0.6,0.132>,
  <0.24,0.6,0.144>,
  <0.26,0.6,0.156>,
  <0.28,0.6,0.168>,
  <0.3,0.6,0.18>,
  <0.32,0.6,0.192>,
  <0.34,0.6,0.204>,
  <0.36,0.6,0.216>,
  <0.38,0.6,0.228>,
  <0.4,0.6,0.24>,
  <0.42,0.6,0.252>,
  <0.44,0.6,0.264>,
  <0.46,0.6,0.276>,
  <0.48,0.6,0.288>,
  <0.5,0.6,0.3>,
  <0.52,0.6,0.312>,
  <0.54,0.6,0.324>,
  <0.56,0.6,0.336>,
  <0.58,0.6,0.348>,
  <0.6,0.6,0.36>,
  <0.62,0.6,0.372>,
  <0.64,0.6,0.384>,
  <0.66,0.6,0.396>,
  <0.68,0.6,0.408>,
  <0.7,0.6,0.42>,
  <0.72,0.6,0.432>,
  <0.74,0.6,0.444>,
  <0.76,0.6,0.456>,
  <0.78,0.6,0.468>,
  <0.8,0.6,0.48>,
  <0.82,0.6,0.492>,
  <0.84,0.6,0.504>,
  <0.86,0.6,0.516>,
  <0.88,0.6,0.528>,
  <0.9,0.6,0.54>,
  <0.92,0.6,0.552>,
  <0.94,0.6,0.564>,
  <0.96,0.6,0.576>,
  <0.98,0.6,0.588>,
  <1,0.6,0.6>,
  <0,0.62,0>,
  <0.02,0.62,0.0124>,
  <0.04,0.62,0.0248>,
  <0.06,0.62,0.0372>,
  <0.08,0.62,0.0496>,
  <0.1,0.62,0.062>,
  <0.12,0.62,0.0744>,
  <0.14,0.62,0.0868>,
  <0.16,0.62,0.0992>,
  <0.18,0.62,0.1116>,
  <0.2,0.62,0.124>,
  <0.22,0.62,0.1364>,
  <0.24,0.62,0.1488>,
  <0.26,0.62,0.1612>,
  <0.28,0.62,0.1736>,
  <0.3,0.62,0.186>,
  <0.32,0.62,0.1984>,
  <0.34,0.62,0.2108>,
  <0.36,0.62,0.2232>,
  <0.38,0.62,0.2356>,
  <0.4,0.62,0.248>,
  <0.42,0.62,0.2604>,
  <0.44,0.62,0.2728>,
  <0.46,0.62,0.2852>,
  <0.48,0.62,0.2976>,
  <0.5,0.62,0.31>,
  <0.52,0.62,0.3224>,
  <0.54,0.62,0.3348>,
  <0.56,0.62,0.3472>,
  <0.58,0.62,0.3596>,
  <0.6,0.62,0.372>,
  <0.62,0.62,0.3844>,
  <0.64,0.62,0.3968>,
  <0.66,0.62,0.4092>,
  <0.68,0.62,0.4216>,
  <0.7,0.62,0.434>,
  <0.72,0.62,0.4464>,
  <0.74,0.62,0.4588>,
  <0.76,0.62,0.4712>,
  <0.78,0.62,0.4836>,
  <0.8,0.62,0.496>,
  <0.82,0.62,0.5084>,
  <0.84,0.62,0.5208>,
  <0.86,0.62,0.5332>,
  <0.88,0.62,0.5456>,
  <0.9,0.62,0.558>,
  <0.92,0.62,0.5704>,
  <0.94,0.62,0.5828>,
  <0.96,0.62,0.5952>,
  <0.98,0.62,0.6076>,
  <1,0.62,0.62>,
  <0,0.64,0>,
  <0.02,0.64,0.0128>,
  <0.04,0.64,0.0256>,
  <0.06,0.64,0.0384>,
  <0.08,0.64,0.0512>,
  <0.1,0.64,0.064>,
  <0.12,0.64,0.0768>,
  <0.14,0.64,0.0896>,
  <0.16,0.64,0.1024>,
  <0.18,0.64,0.1152>,
  <0.2,0.64,0.128>,
  <0.22,0.64,0.1408>,
  <0.24,0.64,0.1536>,
  <0.26,0.64,0.1664>,
  <0.28,0.64,0.1792>,
  <0.3,0.64,0.192>,
  <0.32,0.64,0.2048>,
  <0.34,0.64,0.2176>,
  <0.36,0.64,0.2304>,
  <0.38,0.64,0.2432>,
  <0.4,0.64,0.256>,
  <0.42,0.64,0.2688>,
  <0.44,0.64,0.2816>,
  <0.46,0.64,0.2944>,
  <0.48,0.64,0.3072>,
  <0.5,0.64,0.32>,
  <0.52,0.64,0.3328>,
  <0.54,0.64,0.3456>,
  <0.56,0.64,0.3584>,
  <0.58,0.64,0.3712>,
  <0.6,0.64,0.384>,
  <0.62,0.64,0.3968>,
  <0.64,0.64,0.4096>,
  <0.66,0.64,0.4224>,
  <0.68,0.64,0.4352>,
  <0.7,0.64,0.448>,
  <0.72,0.64,0.4608>,
  <0.74,0.64,0.4736>,
  <0.76,0.64,0.4864>,
  <0.78,0.64,0.4992>,
  <0.8,0.64,0.512>,
  <0.82,0.64,0.5248>,
  <0.84,0.64,0.5376>,
  <0.86,0.64,0.5504>,
  <0.88,0.64,0.5632>,
  <0.9,0.64,0.576>,
  <0.92,0.64,0.5888>,
  <0.94,0.64,0.6016>,
  <0.96,0.64,0.6144>,
  <0.98,0.64,0.6272>,
  <1,0.64,0.64>,
  <0,0.66,0>,
  <0.02,0.66,0.0132>,
  <0.04,0.66,0.0264>,
  <0.06,0.66,0.0396>,
  <0.08,0.66,0.0528>,
  <0.1,0.66,0.066>,
  <0.12,0.66,0.0792>,
  <0.14,0.66,0.0924>,
  <0.16,0.66,0.1056>,
  <0.18,0.66,0.1188>,
  <0.2,0.66,0.132>,
  <0.22,0.66,0.1452>,
  <0.24,0.66,0.1584>,
  <0.26,0.66,0.1716>,
  <0.28,0.66,0.1848>,
  <0.3,0.66,0.198>,
  <0.32,0.66,0.2112>,
  <0.34,0.66,0.2244>,
  <0.36,0.66,0.2376>,
  <0.38,0.66,0.2508>,
  <0.4,0.66,0.264>,
  <0.42,0.66,0.2772>,
  <0.44,0.66,0.2904>,
  <0.46,0.66,0.3036>,
  <0.48,0.66,0.3168>,
  <0.5,0.66,0.33>,
  <0.52,0.66,0.3432>,
  <0.54,0.66,0.3564>,
  <0.56,0.66,0.3696>,
  <0.58,0.66,0.3828>,
  <0.6,0.66,0.396>,
  <0.62,0.66,0.4092>,
  <0.64,0.66,0.4224>,
  <0.66,0.66,0.4356>,
  <0.68,0.66,0.4488>,
  <0.7,0.66,0.462>,
  <0.72,0.66,0.4752>,
  <0.74,0.66,0.4884>,
  <0.76,0.66,0.5016>,
  <0.78,0.66,0.5148>,
  <0.8,0.66,0.528>,
  <0.82,0.66,0.5412>,
  <0.84,0.66,0.5544>,
  <0.86,0.66,0.5676>,
  <0.88,0.66,0.5808>,
  <0.9,0.66,0.594>,
  <0.92,0.66,0.6072>,
  <0.94,0.66,0.6204>,
  <0.96,0.66,0.6336>,
  <0.98,0.66,0.6468>,
  <1,0.66,0.66>,
  <0,0.68,0>,
  <0.02,0.68,0.0136>,
  <0.04,0.68,0.0272>,
  <0.06,0.68,0.0408>,
  <0.08,0.68,0.0544>,
  <0.1,0.68,0.068>,
  <0.12,0.68,0.0816>,
  <0.14,0.68,0.0952>,
  <0.16,0.68,0.1088>,
  <0.18,0.68,0.1224>,
  <0.2,0.68,0.136>,
  <0.22,0.68,0.1496>,
  <0.24,0.68,0.1632>,
  <0.26,0.68,0.1768>,
  <0.28,0.68,0.1904>,
  <0.3,0.68,0.204>,
  <0.32,0.68,0.2176>,
  <0.34,0.68,0.2312>,
  <0.36,0.68,0.2448>,
  <0.38,0.68,0.2584>,
  <0.4,0.68,0.272>,
  <0.42,0.68,0.2856>,
  <0.44,0.68,0.2992>,
  <0.46,0.68,0.3128>,
  <0.48,0.68,0.3264>,
  <0.5,0.68,0.34>,
  <0.52,0.68,0.3536>,
  <0.54,0.68,0.3672>,
  <0.56,0.68,0.3808>,
  <0.58,0.68,0.3944>,
  <0.6,0.68,0.408>,
  <0.62,0.68,0.4216>,
  <0.64,0.68,0.4352>,
  <0.66,0.68,0.4488>,
  <0.68,0.68,0.4624>,
  <0.7,0.68,0.476>,
  <0.72,0.68,0.4896>,
  <0.74,0.68,0.5032>,
  <0.76,0.68,0.5168>,
  <0.78,0.68,0.5304>,
  <0.8,0.68,0.544>,
  <0.82,0.68,0.5576>,
  <0.84,0.68,0.5712>,
  <0.86,0.68,0.5848>,
  <0.88,0.68,0.5984>,
  <0.9,0.68,0.612>,
  <0.92,0.68,0.6256>,
  <0.94,0.68,0.6392>,
  <0.96,0.68,0.6528>,
  <0.98,0.68,0.6664>,
  <1,0.68,0.68>,
  <0,0.7,0>,
  <0.02,0.7,0.014>,
  <0.04,0.7,0.028>,
  <0.06,0.7,0.042>,
  <0.08,0.7,0.056>,
  <0.1,0.7,0.07>,
  <0.12,0.7,0.084>,
  <0.14,0.7,0.098>,
  <0.16,0.7,0.112>,
  <0.18,0.7,0.126>,
  <0.2,0.7,0.14>,
  <0.22,0.7,0.154>,
  <0.24,0.7,0.168>,
  <0.26,0.7,0.182>,
  <0.28,0.7,0.196>,
  <0.3,0.7,0.21>,
  <0.32,0.7,0.224>,
  <0.34,0.7,0.238>,
  <0.36,0.7,0.252>,
  <0.38,0.7,0.266>,
  <0.4,0.7,0.28>,
  <0.42,0.7,0.294>,
  <0.44,0.7,0.308>,
  <0.46,0.7,0.322>,
  <0.48,0.7,0.336>,
  <0.5,0.7,0.35>,
  <0.52,0.7,0.364>,
  <0.54,0.7,0.378>,
  <0.56,0.7,0.392>,
  <0.58,0.7,0.406>,
  <0.6,0.7,0.42>,
  <0.62,0.7,0.434>,
  <0.64,0.7,0.448>,
  <0.66,0.7,0.462>,
  <0.68,0.7,0.476>,
  <0.7,0.7,0.49>,
  <0.72,0.7,0.504>,
  <0.74,0.7,0.518>,
  <0.76,0.7,0.532>,
  <0.78,0.7,0.546>,
  <0.8,0.7,0.56>,
  <0.82,0.7,0.574>,
  <0.84,0.7,0.588>,
  <0.86,0.7,0.602>,
  <0.88,0.7,0.616>,
  <0.9,0.7,0.63>,
  <0.92,0.7,0.644>,
  <0.94,0.7,0.658>,
  <0.96,0.7,0.672>,
  <0.98,0.7,0.686>,
  <1,0.7,0.7>,
  <0,0.72,0>,
  <0.02,0.72,0.0144>,
  <0.04,0.72,0.0288>,
  <0.06,0.72,0.0432>,
  <0.08,0.72,0.0576>,
  <0.1,0.72,0.072>,
  <0.12,0.72,0.0864>,
  <0.14,0.72,0.1008>,
  <0.16,0.72,0.1152>,
  <0.18,0.72,0.1296>,
  <0.2,0.72,0.144>,
  <0.22,0.72,0.1584>,
  <0.24,0.72,0.1728>,
  <0.26,0.72,0.1872>,
  <0.28,0.72,0.2016>,
  <0.3,0.72,0.216>,
  <0.32,0.72,0.2304>,
  <0.34,0.72,0.2448>,
  <0.36,0.72,0.2592>,
  <0.38,0.72,0.2736>,
  <0.4,0.72,0.288>,
  <0.42,0.72,0.3024>,
  <0.44,0.72,0.3168>,
  <0.46,0.72,0.3312>,
  <0.48,0.72,0.3456>,
  <0.5,0.72,0.36>,
  <0.52,0.72,0.3744>,
  <0.54,0.72,0.3888>,
  <0.56,0.72,0.4032>,
  <0.58,0.72,0.4176>,
  <0.6,0.72,0.432>,
  <0.62,0.72,0.4464>,
  <0.64,0.72,0.4608>,
  <0.66,0.72,0.4752>,
  <0.68,0.72,0.4896>,
  <0.7,0.72,0.504>,
  <0.72,0.72,0.5184>,
  <0.74,0.72,0.5328>,
  <0.76,0.72,0.5472>,
  <0.78,0.72,0.5616>,
  <0.8,0.72,0.576>,
  <0.82,0.72,0.5904>,
  <0.84,0.72,0.6048>,
  <0.86,0.72,0.6192>,
  <0.88,0.72,0.6336>,
  <0.9,0.72,0.648>,
  <0.92,0.72,0.6624>,
  <0.94,0.72,0.6768>,
  <0.96,0.72,0.6912>,
  <0.98,0.72,0.7056>,
  <1,0.72,0.72>,
  <0,0.74,0>,
  <0.02,0.74,0.0148>,
  <0.04,0.74,0.0296>,
  <0.06,0.74,0.0444>,
  <0.08,0.74,0.0592>,
  <0.1,0.74,0.074>,
  <0.12,0.74,0.0888>,
  <0.14,0.74,0.1036>,
  <0.16,0.74,0.1184>,
  <0.18,0.74,0.1332>,
  <0.2,0.74,0.148>,
  <0.22,0.74,0.1628>,
  <0.24,0.74,0.1776>,
  <0.26,0.74,0.1924>,
  <0.28,0.74,0.2072>,
  <0.3,0.74,0.222>,
  <0.32,0.74,0.2368>,
  <0.34,0.74,0.2516>,
  <0.36,0.74,0.2664>,
  <0.38,0.74,0.2812>,
  <0.4,0.74,0.296>,
  <0.42,0.74,0.3108>,
  <0.44,0.74,0.3256>,
  <0.46,0.74,0.3404>,
  <0.48,0.74,0.3552>,
  <0.5,0.74,0.37>,
  <0.52,0.74,0.3848>,
  <0.54,0.74,0.3996>,
  <0.56,0.74,0.4144>,
  <0.58,0.74,0.4292>,
  <0.6,0.74,0.444>,
  <0.62,0.74,0.4588>,
  <0.64,0.74,0.4736>,
  <0.66,0.74,0.4884>,
  <0.68,0.74,0.5032>,
  <0.7,0.74,0.518>,
  <0.72,0.74,0.5328>,
  <0.74,0.74,0.5476>,
  <0.76,0.74,0.5624>,
  <0.78,0.74,0.5772>,
  <0.8,0.74,0.592>,
  <0.82,0.74,0.6068>,
  <0.84,0.74,0.6216>,
  <0.86,0.74,0.6364>,
  <0.88,0.74,0.6512>,
  <0.9,0.74,0.666>,
  <0.92,0.74,0.6808>,
  <0.94,0.74,0.6956>,
  <0.96,0.74,0.7104>,
  <0.98,0.74,0.7252>,
  <1,0.74,0.74>,
  <0,0.76,0>,
  <0.02,0.76,0.0152>,
  <0.04,0.76,0.0304>,
  <0.06,0.76,0.0456>,
  <0.08,0.76,0.0608>,
  <0.1,0.76,0.076>,
  <0.12,0.76,0.0912>,
  <0.14,0.76,0.1064>,
  <0.16,0.76,0.1216>,
  <0.18,0.76,0.1368>,
  <0.2,0.76,0.152>,
  <0.22,0.76,0.1672>,
  <0.24,0.76,0.1824>,
  <0.26,0.76,0.1976>,
  <0.28,0.76,0.2128>,
  <0.3,0.76,0.228>,
  <0.32,0.76,0.2432>,
  <0.34,0.76,0.2584>,
  <0.36,0.76,0.2736>,
  <0.38,0.76,0.2888>,
  <0.4,0.76,0.304>,
  <0.42,0.76,0.3192>,
  <0.44,0.76,0.3344>,
  <0.46,0.76,0.3496>,
  <0.48,0.76,0.3648>,
  <0.5,0.76,0.38>,
  <0.52,0.76,0.3952>,
  <0.54,0.76,0.4104>,
  <0.56,0.76,0.4256>,
  <0.58,0.76,0.4408>,
  <0.6,0.76,0.456>,
  <0.62,0.76,0.4712>,
  <0.64,0.76,0.4864>,
  <0.66,0.76,0.5016>,
  <0.68,0.76,0.5168>,
  <0.7,0.76,0.532>,
  <0.72,0.76,0.5472>,
  <0.74,0.76,0.5624>,
  <0.76,0.76,0.5776>,
  <0.78,0.76,0.5928>,
  <0.8,0.76,0.608>,
  <0.82,0.76,0.6232>,
  <0.84,0.76,0.6384>,
  <0.86,0.76,0.6536>,
  <0.88,0.76,0.6688>,
  <0.9,0.76,0.684>,
  <0.92,0.76,0.6992>,
  <0.94,0.76,0.7144>,
  <0.96,0.76,0.7296>,
  <0.98,0.76,0.7448>,
  <1,0.76,0.76>,
  <0,0.78,0>,
  <0.02,0.78,0.0156>,
  <0.04,0.78,0.0312>,
  <0.06,0.78,0.0468>,
  <0.08,0.78,0.0624>,
  <0.1,0.78,0.078>,
  <0.12,0.78,0.0936>,
  <0.14,0.78,0.1092>,
  <0.16,0.78,0.1248>,
  <0.18,0.78,0.1404>,
  <0.2,0.78,0.156>,
  <0.22,0.78,0.1716>,
  <0.24,0.78,0.1872>,
  <0.26,0.78,0.2028>,
  <0.28,0.78,0.2184>,
  <0.3,0.78,0.234>,
  <0.32,0.78,0.2496>,
  <0.34,0.78,0.2652>,
  <0.36,0.78,0.2808>,
  <0.38,0.78,0.2964>,
  <0.4,0.78,0.312>,
  <0.42,0.78,0.3276>,
  <0.44,0.78,0.3432>,
  <0.46,0.78,0.3588>,
  <0.48,0.78,0.3744>,
  <0.5,0.78,0.39>,
  <0.52,0.78,0.4056>,
  <0.54,0.78,0.4212>,
  <0.56,0.78,0.4368>,
  <0.58,0.78,0.4524>,
  <0.6,0.78,0.468>,
  <0.62,0.78,0.4836>,
  <0.64,0.78,0.4992>,
  <0.66,0.78,0.5148>,
  <0.68,0.78,0.5304>,
  <0.7,0.78,0.546>,
  <0.72,0.78,0.5616>,
  <0.74,0.78,0.5772>,
  <0.76,0.78,0.5928>,
  <0.78,0.78,0.6084>,
  <0.8,0.78,0.624>,
  <0.82,0.78,0.6396>,
  <0.84,0.78,0.6552>,
  <0.86,0.78,0.6708>,
  <0.88,0.78,0.6864>,
  <0.9,0.78,0.702>,
  <0.92,0.78,0.7176>,
  <0.94,0.78,0.7332>,
  <0.96,0.78,0.7488>,
  <0.98,0.78,0.7644>,
  <1,0.78,0.78>,
  <0,0.8,0>,
  <0.02,0.8,0.016>,
  <0.04,0.8,0.032>,
  <0.06,0.8,0.048>,
  <0.08,0.8,0.064>,
  <0.1,0.8,0.08>,
  <0.12,0.8,0.096>,
  <0.14,0.8,0.112>,
  <0.16,0.8,0.128>,
  <0.18,0.8,0.144>,
  <0.2,0.8,0.16>,
  <0.22,0.8,0.176>,
  <0.24,0.8,0.192>,
  <0.26,0.8,0.208>,
  <0.28,0.8,0.224>,
  <0.3,0.8,0.24>,
  <0.32,0.8,0.256>,
  <0.34,0.8,0.272>,
  <0.36,0.8,0.288>,
  <0.38,0.8,0.304>,
  <0.4,0.8,0.32>,
  <0.42,0.8,0.336>,
  <0.44,0.8,0.352>,
  <0.46,0.8,0.368>,
  <0.48,0.8,0.384>,
  <0.5,0.8,0.4>,
  <0.52,0.8,0.416>,
  <0.54,0.8,0.432>,
  <0.56,0.8,0.448>,
  <0.58,0.8,0.464>,
  <0.6,0.8,0.48>,
  <0.62,0.8,0.496>,
  <0.64,0.8,0.512>,
  <0.66,0.8,0.528>,
  <0.68,0.8,0.544>,
  <0.7,0.8,0.56>,
  <0.72,0.8,0.576>,
  <0.74,0.8,0.592>,
  <0.76,0.8,0.608>,
  <0.78,0.8,0.624>,
  <0.8,0.8,0.64>,
  <0.82,0.8,0.656>,
  <0.84,0.8,0.672>,
  <0.86,0.8,0.688>,
  <0.88,0.8,0.704>,
  <0.9,0.8,0.72>,
  <0.92,0.8,0.736>,
  <0.94,0.8,0.752>,
  <0.96,0.8,0.768>,
  <0.98,0.8,0.784>,
  <1,0.8,0.8>,
  <0,0.82,0>,
  <0.02,0.82,0.0164>,
  <0.04,0.82,0.0328>,
  <0.06,0.82,0.0492>,
  <0.08,0.82,0.0656>,
  <0.1,0.82,0.082>,
  <0.12,0.82,0.0984>,
  <0.14,0.82,0.1148>,
  <0.16,0.82,0.1312>,
  <0.18,0.82,0.1476>,
  <0.2,0.82,0.164>,
  <0.22,0.82,0.1804>,
  <0.24,0.82,0.1968>,
  <0.26,0.82,0.2132>,
  <0.28,0.82,0.2296>,
  <0.3,0.82,0.246>,
  <0.32,0.82,0.2624>,
  <0.34,0.82,0.2788>,
  <0.36,0.82,0.2952>,
  <0.38,0.82,0.3116>,
  <0.4,0.82,0.328>,
  <0.42,0.82,0.3444>,
  <0.44,0.82,0.3608>,
  <0.46,0.82,0.3772>,
  <0.48,0.82,0.3936>,
  <0.5,0.82,0.41>,
  <0.52,0.82,0.4264>,
  <0.54,0.82,0.4428>,
  <0.56,0.82,0.4592>,
  <0.58,0.82,0.4756>,
  <0.6,0.82,0.492>,
  <0.62,0.82,0.5084>,
  <0.64,0.82,0.5248>,
  <0.66,0.82,0.5412>,
  <0.68,0.82,0.5576>,
  <0.7,0.82,0.574>,
  <0.72,0.82,0.5904>,
  <0.74,0.82,0.6068>,
  <0.76,0.82,0.6232>,
  <0.78,0.82,0.6396>,
  <0.8,0.82,0.656>,
  <0.82,0.82,0.6724>,
  <0.84,0.82,0.6888>,
  <0.86,0.82,0.7052>,
  <0.88,0.82,0.7216>,
  <0.9,0.82,0.738>,
  <0.92,0.82,0.7544>,
  <0.94,0.82,0.7708>,
  <0.96,0.82,0.7872>,
  <0.98,0.82,0.8036>,
  <1,0.82,0.82>,
  <0,0.84,0>,
  <0.02,0.84,0.0168>,
  <0.04,0.84,0.0336>,
  <0.06,0.84,0.0504>,
  <0.08,0.84,0.0672>,
  <0.1,0.84,0.084>,
  <0.12,0.84,0.1008>,
  <0.14,0.84,0.1176>,
  <0.16,0.84,0.1344>,
  <0.18,0.84,0.1512>,
  <0.2,0.84,0.168>,
  <0.22,0.84,0.1848>,
  <0.24,0.84,0.2016>,
  <0.26,0.84,0.2184>,
  <0.28,0.84,0.2352>,
  <0.3,0.84,0.252>,
  <0.32,0.84,0.2688>,
  <0.34,0.84,0.2856>,
  <0.36,0.84,0.3024>,
  <0.38,0.84,0.3192>,
  <0.4,0.84,0.336>,
  <0.42,0.84,0.3528>,
  <0.44,0.84,0.3696>,
  <0.46,0.84,0.3864>,
  <0.48,0.84,0.4032>,
  <0.5,0.84,0.42>,
  <0.52,0.84,0.4368>,
  <0.54,0.84,0.4536>,
  <0.56,0.84,0.4704>,
  <0.58,0.84,0.4872>,
  <0.6,0.84,0.504>,
  <0.62,0.84,0.5208>,
  <0.64,0.84,0.5376>,
  <0.66,0.84,0.5544>,
  <0.68,0.84,0.5712>,
  <0.7,0.84,0.588>,
  <0.72,0.84,0.6048>,
  <0.74,0.84,0.6216>,
  <0.76,0.84,0.6384>,
  <0.78,0.84,0.6552>,
  <0.8,0.84,0.672>,
  <0.82,0.84,0.6888>,
  <0.84,0.84,0.7056>,
  <0.86,0.84,0.7224>,
  <0.88,0.84,0.7392>,
  <0.9,0.84,0.756>,
  <0.92,0.84,0.7728>,
  <0.94,0.84,0.7896>,
  <0.96,0.84,0.8064>,
  <0.98,0.84,0.8232>,
  <1,0.84,0.84>,
  <0,0.86,0>,
  <0.02,0.86,0.0172>,
  <0.04,0.86,0.0344>,
  <0.06,0.86,0.0516>,
  <0.08,0.86,0.0688>,
  <0.1,0.86,0.086>,
  <0.12,0.86,0.1032>,
  <0.14,0.86,0.1204>,
  <0.16,0.86,0.1376>,
  <0.18,0.86,0.1548>,
  <0.2,0.86,0.172>,
  <0.22,0.86,0.1892>,
  <0.24,0.86,0.2064>,
  <0.26,0.86,0.2236>,
  <0.28,0.86,0.2408>,
  <0.3,0.86,0.258>,
  <0.32,0.86,0.2752>,
  <0.34,0.86,0.2924>,
  <0.36,0.86,0.3096>,
  <0.38,0.86,0.3268>,
  <0.4,0.86,0.344>,
  <0.42,0.86,0.3612>,
  <0.44,0.86,0.3784>,
  <0.46,0.86,0.3956>,
  <0.48,0.86,0.4128>,
  <0.5,0.86,0.43>,
  <0.52,0.86,0.4472>,
  <0.54,0.86,0.4644>,
  <0.56,0.86,0.4816>,
  <0.58,0.86,0.4988>,
  <0.6,0.86,0.516>,
  <0.62,0.86,0.5332>,
  <0.64,0.86,0.5504>,
  <0.66,0.86,0.5676>,
  <0.68,0.86,0.5848>,
  <0.7,0.86,0.602>,
  <0.72,0.86,0.6192>,
  <0.74,0.86,0.6364>,
  <0.76,0.86,0.6536>,
  <0.78,0.86,0.6708>,
  <0.8,0.86,0.688>,
  <0.82,0.86,0.7052>,
  <0.84,0.86,0.7224>,
  <0.86,0.86,0.7396>,
  <0.88,0.86,0.7568>,
  <0.9,0.86,0.774>,
  <0.92,0.86,0.7912>,
  <0.94,0.86,0.8084>,
  <0.96,0.86,0.8256>,
  <0.98,0.86,0.8428>,
  <1,0.86,0.86>,
  <0,0.88,0>,
  <0.02,0.88,0.0176>,
  <0.04,0.88,0.0352>,
  <0.06,0.88,0.0528>,
  <0.08,0.88,0.0704>,
  <0.1,0.88,0.088>,
  <0.12,0.88,0.1056>,
  <0.14,0.88,0.1232>,
  <0.16,0.88,0.1408>,
  <0.18,0.88,0.1584>,
  <0.2,0.88,0.176>,
  <0.22,0.88,0.1936>,
  <0.24,0.88,0.2112>,
  <0.26,0.88,0.2288>,
  <0.28,0.88,0.2464>,
  <0.3,0.88,0.264>,
  <0.32,0.88,0.2816>,
  <0.34,0.88,0.2992>,
  <0.36,0.88,0.3168>,
  <0.38,0.88,0.3344>,
  <0.4,0.88,0.352>,
  <0.42,0.88,0.3696>,
  <0.44,0.88,0.3872>,
  <0.46,0.88,0.4048>,
  <0.48,0.88,0.4224>,
  <0.5,0.88,0.44>,
  <0.52,0.88,0.4576>,
  <0.54,0.88,0.4752>,
  <0.56,0.88,0.4928>,
  <0.58,0.88,0.5104>,
  <0.6,0.88,0.528>,
  <0.62,0.88,0.5456>,
  <0.64,0.88,0.5632>,
  <0.66,0.88,0.5808>,
  <0.68,0.88,0.5984>,
  <0.7,0.88,0.616>,
  <0.72,0.88,0.6336>,
  <0.74,0.88,0.6512>,
  <0.76,0.88,0.6688>,
  <0.78,0.88,0.6864>,
  <0.8,0.88,0.704>,
  <0.82,0.88,0.7216>,
  <0.84,0.88,0.7392>,
  <0.86,0.88,0.7568>,
  <0.88,0.88,0.7744>,
  <0.9,0.88,0.792>,
  <0.92,0.88,0.8096>,
  <0.94,0.88,0.8272>,
  <0.96,0.88,0.8448>,
  <0.98,0.88,0.8624>,
  <1,0.88,0.88>,
  <0,0.9,0>,
  <0.02,0.9,0.018>,
  <0.04,0.9,0.036>,
  <0.06,0.9,0.054>,
  <0.08,0.9,0.072>,
  <0.1,0.9,0.09>,
  <0.12,0.9,0.108>,
  <0.14,0.9,0.126>,
  <0.16,0.9,0.144>,
  <0.18,0.9,0.162>,
  <0.2,0.9,0.18>,
  <0.22,0.9,0.198>,
  <0.24,0.9,0.216>,
  <0.26,0.9,0.234>,
  <0.28,0.9,0.252>,
  <0.3,0.9,0.27>,
  <0.32,0.9,0.288>,
  <0.34,0.9,0.306>,
  <0.36,0.9,0.324>,
  <0.38,0.9,0.342>,
  <0.4,0.9,0.36>,
  <0.42,0.9,0.378>,
  <0.44,0.9,0.396>,
  <0.46,0.9,0.414>,
  <0.48,0.9,0.432>,
  <0.5,0.9,0.45>,
  <0.52,0.9,0.468>,
  <0.54,0.9,0.486>,
  <0.56,0.9,0.504>,
  <0.58,0.9,0.522>,
  <0.6,0.9,0.54>,
  <0.62,0.9,0.558>,
  <0.64,0.9,0.576>,
  <0.66,0.9,0.594>,
  <0.68,0.9,0.612>,
  <0.7,0.9,0.63>,
  <0.72,0.9,0.648>,
  <0.74,0.9,0.666>,
  <0.76,0.9,0.684>,
  <0.78,0.9,0.702>,
  <0.8,0.9,0.72>,
  <0.82,0.9,0.738>,
  <0.84,0.9,0.756>,
  <0.86,0.9,0.774>,
  <0.88,0.9,0.792>,
  <0.9,0.9,0.81>,
  <0.92,0.9,0.828>,
  <0.94,0.9,0.846>,
  <0.96,0.9,0.864>,
  <0.98,0.9,0.882>,
  <1,0.9,0.9>,
  <0,0.92,0>,
  <0.02,0.92,0.0184>,
  <0.04,0.92,0.0368>,
  <0.06,0.92,0.0552>,
  <0.08,0.92,0.0736>,
  <0.1,0.92,0.092>,
  <0.12,0.92,0.1104>,
  <0.14,0.92,0.1288>,
  <0.16,0.92,0.1472>,
  <0.18,0.92,0.1656>,
  <0.2,0.92,0.184>,
  <0.22,0.92,0.2024>,
  <0.24,0.92,0.2208>,
  <0.26,0.92,0.2392>,
  <0.28,0.92,0.2576>,
  <0.3,0.92,0.276>,
  <0.32,0.92,0.2944>,
  <0.34,0.92,0.3128>,
  <0.36,0.92,0.3312>,
  <0.38,0.92,0.3496>,
  <0.4,0.92,0.368>,
  <0.42,0.92,0.3864>,
  <0.44,0.92,0.4048>,
  <0.46,0.92,0.4232>,
  <0.48,0.92,0.4416>,
  <0.5,0.92,0.46>,
  <0.52,0.92,0.4784>,
  <0.54,0.92,0.4968>,
  <0.56,0.92,0.5152>,
  <0.58,0.92,0.5336>,
  <0.6,0.92,0.552>,
  <0.62,0.92,0.5704>,
  <0.64,0.92,0.5888>,
  <0.66,0.92,0.6072>,
  <0.68,0.92,0.6256>,
  <0.7,0.92,0.644>,
  <0.72,0.92,0.6624>,
  <0.74,0.92,0.6808>,
  <0.76,0.92,0.6992>,
  <0.78,0.92,0.7176>,
  <0.8,0.92,0.736>,
  <0.82,0.92,0.7544>,
  <0.84,0.92,0.7728>,
  <0.86,0.92,0.7912>,
  <0.88,0.92,0.8096>,
  <0.9,0.92,0.828>,
  <0.92,0.92,0.8464>,
  <0.94,0.92,0.8648>,
  <0.96,0.92,0.8832>,
  <0.98,0.92,0.9016>,
  <1,0.92,0.92>,
  <0,0.94,0>,
  <0.02,0.94,0.0188>,
  <0.04,0.94,0.0376>,
  <0.06,0.94,0.0564>,
  <0.08,0.94,0.0752>,
  <0.1,0.94,0.094>,
  <0.12,0.94,0.1128>,
  <0.14,0.94,0.1316>,
  <0.16,0.94,0.1504>,
  <0.18,0.94,0.1692>,
  <0.2,0.94,0.188>,
  <0.22,0.94,0.2068>,
  <0.24,0.94,0.2256>,
  <0.26,0.94,0.2444>,
  <0.28,0.94,0.2632>,
  <0.3,0.94,0.282>,
  <0.32,0.94,0.3008>,
  <0.34,0.94,0.3196>,
  <0.36,0.94,0.3384>,
  <0.38,0.94,0.3572>,
  <0.4,0.94,0.376>,
  <0.42,0.94,0.3948>,
  <0.44,0.94,0.4136>,
  <0.46,0.94,0.4324>,
  <0.48,0.94,0.4512>,
  <0.5,0.94,0.47>,
  <0.52,0.94,0.4888>,
  <0.54,0.94,0.5076>,
  <0.56,0.94,0.5264>,
  <0.58,0.94,0.5452>,
  <0.6,0.94,0.564>,
  <0.62,0.94,0.5828>,
  <0.64,0.94,0.6016>,
  <0.66,0.94,0.6204>,
  <0.68,0.94,0.6392>,
  <0.7,0.94,0.658>,
  <0.72,0.94,0.6768>,
  <0.74,0.94,0.6956>,
  <0.76,0.94,0.7144>,
  <0.78,0.94,0.7332>,
  <0.8,0.94,0.752>,
  <0.82,0.94,0.7708>,
  <0.84,0.94,0.7896>,
  <0.86,0.94,0.8084>,
  <0.88,0.94,0.8272>,
  <0.9,0.94,0.846>,
  <0.92,0.94,0.8648>,
  <0.94,0.94,0.8836>,
  <0.96,0.94,0.9024>,
  <0.98,0.94,0.9212>,
  <1,0.94,0.94>,
  <0,0.96,0>,
  <0.02,0.96,0.0192>,
  <0.04,0.96,0.0384>,
  <0.06,0.96,0.0576>,
  <0.08,0.96,0.0768>,
  <0.1,0.96,0.096>,
  <0.12,0.96,0.1152>,
  <0.14,0.96,0.1344>,
  <0.16,0.96,0.1536>,
  <0.18,0.96,0.1728>,
  <0.2,0.96,0.192>,
  <0.22,0.96,0.2112>,
  <0.24,0.96,0.2304>,
  <0.26,0.96,0.2496>,
  <0.28,0.96,0.2688>,
  <0.3,0.96,0.288>,
  <0.32,0.96,0.3072>,
  <0.34,0.96,0.3264>,
  <0.36,0.96,0.3456>,
  <0.38,0.96,0.3648>,
  <0.4,0.96,0.384>,
  <0.42,0.96,0.4032>,
  <0.44,0.96,0.4224>,
  <0.46,0.96,0.4416>,
  <0.48,0.96,0.4608>,
  <0.5,0.96,0.48>,
  <0.52,0.96,0.4992>,
  <0.54,0.96,0.5184>,
  <0.56,0.96,0.5376>,
  <0.58,0.96,0.5568>,
  <0.6,0.96,0.576>,
  <0.62,0.96,0.5952>,
  <0.64,0.96,0.6144>,
  <0.66,0.96,0.6336>,
  <0.68,0.96,0.6528>,
  <0.7,0.96,0.672>,
  <0.72,0.96,0.6912>,
  <0.74,0.96,0.7104>,
  <0.76,0.96,0.7296>,
  <0.78,0.96,0.7488>,
  <0.8,0.96,0.768>,
  <0.82,0.96,0.7872>,
  <0.84,0.96,0.8064>,
  <0.86,0.96,0.8256>,
  <0.88,0.96,0.8448>,
  <0.9,0.96,0.864>,
  <0.92,0.96,0.8832>,
  <0.94,0.96,0.9024>,
  <0.96,0.96,0.9216>,
  <0.98,0.96,0.9408>,
  <1,0.96,0.96>,
  <0,0.98,0>,
  <0.02,0.98,0.0196>,
  <0.04,0.98,0.0392>,
  <0.06,0.98,0.0588>,
  <0.08,0.98,0.0784>,
  <0.1,0.98,0.098>,
  <0.12,0.98,0.1176>,
  <0.14,0.98,0.1372>,
  <0.16,0.98,0.1568>,
  <0.18,0.98,0.1764>,
  <0.2,0.98,0.196>,
  <0.22,0.98,0.2156>,
  <0.24,0.98,0.2352>,
  <0.26,0.98,0.2548>,
  <0.28,0.98,0.2744>,
  <0.3,0.98,0.294>,
  <0.32,0.98,0.3136>,
  <0.34,0.98,0.3332>,
  <0.36,0.98,0.3528>,
  <0.38,0.98,0.3724>,
  <0.4,0.98,0.392>,
  <0.42,0.98,0.4116>,
  <0.44,0.98,0.4312>,
  <0.46,0.98,0.4508>,
  <0.48,0.98,0.4704>,
  <0.5,0.98,0.49>,
  <0.52,0.98,0.5096>,
  <0.54,0.98,0.5292>,
  <0.56,0.98,0.5488>,
  <0.58,0.98,0.5684>,
  <0.6,0.98,0.588>,
  <0.62,0.98,0.6076>,
  <0.64,0.98,0.6272>,
  <0.66,0.98,0.6468>,
  <0.68,0.98,0.6664>,
  <0.7,0.98,0.686>,
  <0.72,0.98,0.7056>,
  <0.74,0.98,0.7252>,
  <0.76,0.98,0.7448>,
  <0.78,0.98,0.7644>,
  <0.8,0.98,0.784>,
  <0.82,0.98,0.8036>,
  <0.84,0.98,0.8232>,
  <0.86,0.98,0.8428>,
  <0.88,0.98,0.8624>,
  <0.9,0.98,0.882>,
  <0.92,0.98,0.9016>,
  <0.94,0.98,0.9212>,
  <0.96,0.98,0.9408>,
  <0.98,0.98,0.9604>,
  <1,0.98,0.98>,
  <0,1,0>,
  <0.02,1,0.02>,
  <0.04,1,0.04>,
  <0.06,1,0.06>,
  <0.08,1,0.08>,
  <0.1,1,0.1>,
  <0.12,1,0.12>,
  <0.14,1,0.14>,
  <0.16,1,0.16>,
  <0.18,1,0.18>,
  <0.2,1,0.2>,
  <0.22,1,0.22>,
  <0.24,1,0.24>,
  <0.26,1,0.26>,
  <0.28,1,0.28>,
  <0.3,1,0.3>,
  <0.32,1,0.32>,
  <0.34,1,0.34>,
  <0.36,1,0.36>,
  <0.38,1,0.38>,
  <0.4,1,0.4>,
  <0.42,1,0.42>,
  <0.44,1,0.44>,
  <0.46,1,0.46>,
  <0.48,1,0.48>,
  <0.5,1,0.5>,
  <0.52,1,0.52>,
  <0.54,1,0.54>,
  <0.56,1,0.56>,
  <0.58,1,0.58>,
  <0.6,1,0.6>,
  <0.62,1,0.62>,
  <0.64,1,0.64>,
  <0.66,1,0.66>,
  <0.68,1,0.68>,
  <0.7,1,0.7>,
  <0.72,1,0.72>,
  <0.74,1,0.74>,
  <0.76,1,0.76>,
  <0.78,1,0.78>,
  <0.8,1,0.8>,
  <0.82,1,0.82>,
  <0.84,1,0.84>,
  <0.86,1,0.86>,
  <0.88,1,0.88>,
  <0.9,1,0.9>,
  <0.92,1,0.92>,
  <0.94,1,0.94>,
  <0.96,1,0.96>,
  <0.98,1,0.98>,
  <1,1,1>,
  \}
  normal_vectors \{ 2601,
  <0,0,1>,
  <0,-0.02,1>,
  <0,-0.04,1>,
  <0,-0.06,1>,
  <0,-0.08,1>,
  <0,-0.1,1>,
  <0,-0.12,1>,
  <0,-0.14,1>,
  <0,-0.16,1>,
  <0,-0.18,1>,
  <0,-0.2,1>,
  <0,-0.22,1>,
  <0,-0.24,1>,
  <0,-0.26,1>,
  <0,-0.28,1>,
  <0,-0.3,1>,
  <0,-0.32,1>,
  <0,-0.34,1>,
  <0,-0.36,1>,
  <0,-0.38,1>,
  <0,-0.4,1>,
  <0,-0.42,1>,
  <0,-0.44,1>,
  <0,-0.46,1>,
  <0,-0.48,1>,
  <0,-0.5,1>,
  <0,-0.52,1>,
  <0,-0.54,1>,
  <0,-0.56,1>,
  <0,-0.58,1>,
  <0,-0.6,1>,
  <0,-0.62,1>,
  <0,-0.64,1>,
  <0,-0.66,1>,
  <0,-0.68,1>,
  <0,-0.7,1>,
  <0,-0.72,1>,
  <0,-0.74,1>,
  <0,-0.76,1>,
  <0,-0.78,1>,
  <0,-0.8,1>,
  <0,-0.82,1>,
  <0,-0.84,1>,
  <0,-0.86,1>,
  <0,-0.88,1>,
  <0,-0.9,1>,
  <0,-0.92,1>,
  <0,-0.94,1>,
  <0,-0.96,1>,
  <0,-0.98,1>,
  <0,-1,1>,
  <-0.02,0,1>,
  <-0.02,-0.02,1>,
  <-0.02,-0.04,1>,
  <-0.02,-0.06,1>,
  <-0.02,-0.08,1>,
  <-0.02,-0.1,1>,
  <-0.02,-0.12,1>,
  <-0.02,-0.14,1>,
  <-0.02,-0.16,1>,
  <-0.02,-0.18,1>,
  <-0.02,-0.2,1>,
  <-0.02,-0.22,1>,
  <-0.02,-0.24,1>,
  <-0.02,-0.26,1>,
  <-0.02,-0.28,1>,
  <-0.02,-0.3,1>,
  <-0.02,-0.32,1>,
  <-0.02,-0.34,1>,
  <-0.02,-0.36,1>,
  <-0.02,-0.38,1>,
  <-0.02,-0.4,1>,
  <-0.02,-0.42,1>,
  <-0.02,-0.44,1>,
  <-0.02,-0.46,1>,
  <-0.02,-0.48,1>,
  <-0.02,-0.5,1>,
  <-0.02,-0.52,1>,
  <-0.02,-0.54,1>,
  <-0.02,-0.56,1>,
  <-0.02,-0.58,1>,
  <-0.02,-0.6,1>,
  <-0.02,-0.62,1>,
  <-0.02,-0.64,1>,
  <-0.02,-0.66,1>,
  <-0.02,-0.68,1>,
  <-0.02,-0.7,1>,
  <-0.02,-0.72,1>,
  <-0.02,-0.74,1>,
  <-0.02,-0.76,1>,
  <-0.02,-0.78,1>,
  <-0.02,-0.8,1>,
  <-0.02,-0.82,1>,
  <-0.02,-0.84,1>,
  <-0.02,-0.86,1>,
  <-0.02,-0.88,1>,
  <-0.02,-0.9,1>,
  <-0.02,-0.92,1>,
  <-0.02,-0.94,1>,
  <-0.02,-0.96,1>,
  <-0.02,-0.98,1>,
  <-0.02,-1,1>,
  <-0.04,0,1>,
  <-0.04,-0.02,1>,
  <-0.04,-0.04,1>,
  <-0.04,-0.06,1>,
  <-0.04,-0.08,1>,
  <-0.04,-0.1,1>,
  <-0.04,-0.12,1>,
  <-0.04,-0.14,1>,
  <-0.04,-0.16,1>,
  <-0.04,-0.18,1>,
  <-0.04,-0.2,1>,
  <-0.04,-0.22,1>,
  <-0.04,-0.24,1>,
  <-0.04,-0.26,1>,
  <-0.04,-0.28,1>,
  <-0.04,-0.3,1>,
  <-0.04,-0.32,1>,
  <-0.04,-0.34,1>,
  <-0.04,-0.36,1>,
  <-0.04,-0.38,1>,
  <-0.04,-0.4,1>,
  <-0.04,-0.42,1>,
  <-0.04,-0.44,1>,
  <-0.04,-0.46,1>,
  <-0.04,-0.48,1>,
  <-0.04,-0.5,1>,
  <-0.04,-0.52,1>,
  <-0.04,-0.54,1>,
  <-0.04,-0.56,1>,
  <-0.04,-0.58,1>,
  <-0.04,-0.6,1>,
  <-0.04,-0.62,1>,
  <-0.04,-0.64,1>,
  <-0.04,-0.66,1>,
  <-0.04,-0.68,1>,
  <-0.04,-0.7,1>,
  <-0.04,-0.72,1>,
  <-0.04,-0.74,1>,
  <-0.04,-0.76,1>,
  <-0.04,-0.78,1>,
  <-0.04,-0.8,1>,
  <-0.04,-0.82,1>,
  <-0.04,-0.84,1>,
  <-0.04,-0.86,1>,
  <-0.04,-0.88,1>,
  <-0.04,-0.9,1>,
  <-0.04,-0.92,1>,
  <-0.04,-0.94,1>,
  <-0.04,-0.96,1>,
  <-0.04,-0.98,1>,
  <-0.04,-1,1>,
  <-0.06,0,1>,
  <-0.06,-0.02,1>,
  <-0.06,-0.04,1>,
  <-0.06,-0.06,1>,
  <-0.06,-0.08,1>,
  <-0.06,-0.1,1>,
  <-0.06,-0.12,1>,
  <-0.06,-0.14,1>,
  <-0.06,-0.16,1>,
  <-0.06,-0.18,1>,
  <-0.06,-0.2,1>,
  <-0.06,-0.22,1>,
  <-0.06,-0.24,1>,
  <-0.06,-0.26,1>,
  <-0.06,-0.28,1>,
  <-0.06,-0.3,1>,
  <-0.06,-0.32,1>,
  <-0.06,-0.34,1>,
  <-0.06,-0.36,1>,
  <-0.06,-0.38,1>,
  <-0.06,-0.4,1>,
  <-0.06,-0.42,1>,
  <-0.06,-0.44,1>,
  <-0.06,-0.46,1>,
  <-0.06,-0.48,1>,
  <-0.06,-0.5,1>,
  <-0.06,-0.52,1>,
  <-0.06,-0.54,1>,
  <-0.06,-0.56,1>,
  <-0.06,-0.58,1>,
  <-0.06,-0.6,1>,
  <-0.06,-0.62,1>,
  <-0.06,-0.64,1>,
  <-0.06,-0.66,1>,
  <-0.06,-0.68,1>,
  <-0.06,-0.7,1>,
  <-0.06,-0.72,1>,
  <-0.06,-0.74,1>,
  <-0.06,-0.76,1>,
  <-0.06,-0.78,1>,
  <-0.06,-0.8,1>,
  <-0.06,-0.82,1>,
  <-0.06,-0.84,1>,
  <-0.06,-0.86,1>,
  <-0.06,-0.88,1>,
  <-0.06,-0.9,1>,
  <-0.06,-0.92,1>,
  <-0.06,-0.94,1>,
  <-0.06,-0.96,1>,
  <-0.06,-0.98,1>,
  <-0.06,-1,1>,
  <-0.08,0,1>,
  <-0.08,-0.02,1>,
  <-0.08,-0.04,1>,
  <-0.08,-0.06,1>,
  <-0.08,-0.08,1>,
  <-0.08,-0.1,1>,
  <-0.08,-0.12,1>,
  <-0.08,-0.14,1>,
  <-0.08,-0.16,1>,
  <-0.08,-0.18,1>,
  <-0.08,-0.2,1>,
  <-0.08,-0.22,1>,
  <-0.08,-0.24,1>,
  <-0.08,-0.26,1>,
  <-0.08,-0.28,1>,
  <-0.08,-0.3,1>,
  <-0.08,-0.32,1>,
  <-0.08,-0.34,1>,
  <-0.08,-0.36,1>,
  <-0.08,-0.38,1>,
  <-0.08,-0.4,1>,
  <-0.08,-0.42,1>,
  <-0.08,-0.44,1>,
  <-0.08,-0.46,1>,
  <-0.08,-0.48,1>,
  <-0.08,-0.5,1>,
  <-0.08,-0.52,1>,
  <-0.08,-0.54,1>,
  <-0.08,-0.56,1>,
  <-0.08,-0.58,1>,
  <-0.08,-0.6,1>,
  <-0.08,-0.62,1>,
  <-0.08,-0.64,1>,
  <-0.08,-0.66,1>,
  <-0.08,-0.68,1>,
  <-0.08,-0.7,1>,
  <-0.08,-0.72,1>,
  <-0.08,-0.74,1>,
  <-0.08,-0.76,1>,
  <-0.08,-0.78,1>,
  <-0.08,-0.8,1>,
  <-0.08,-0.82,1>,
  <-0.08,-0.84,1>,
  <-0.08,-0.86,1>,
  <-0.08,-0.88,1>,
  <-0.08,-0.9,1>,
  <-0.08,-0.92,1>,
  <-0.08,-0.94,1>,
  <-0.08,-0.96,1>,
  <-0.08,-0.98,1>,
  <-0.08,-1,1>,
  <-0.1,0,1>,
  <-0.1,-0.02,1>,
  <-0.1,-0.04,1>,
  <-0.1,-0.06,1>,
  <-0.1,-0.08,1>,
  <-0.1,-0.1,1>,
  <-0.1,-0.12,1>,
  <-0.1,-0.14,1>,
  <-0.1,-0.16,1>,
  <-0.1,-0.18,1>,
  <-0.1,-0.2,1>,
  <-0.1,-0.22,1>,
  <-0.1,-0.24,1>,
  <-0.1,-0.26,1>,
  <-0.1,-0.28,1>,
  <-0.1,-0.3,1>,
  <-0.1,-0.32,1>,
  <-0.1,-0.34,1>,
  <-0.1,-0.36,1>,
  <-0.1,-0.38,1>,
  <-0.1,-0.4,1>,
  <-0.1,-0.42,1>,
  <-0.1,-0.44,1>,
  <-0.1,-0.46,1>,
  <-0.1,-0.48,1>,
  <-0.1,-0.5,1>,
  <-0.1,-0.52,1>,
  <-0.1,-0.54,1>,
  <-0.1,-0.56,1>,
  <-0.1,-0.58,1>,
  <-0.1,-0.6,1>,
  <-0.1,-0.62,1>,
  <-0.1,-0.64,1>,
  <-0.1,-0.66,1>,
  <-0.1,-0.68,1>,
  <-0.1,-0.7,1>,
  <-0.1,-0.72,1>,
  <-0.1,-0.74,1>,
  <-0.1,-0.76,1>,
  <-0.1,-0.78,1>,
  <-0.1,-0.8,1>,
  <-0.1,-0.82,1>,
  <-0.1,-0.84,1>,
  <-0.1,-0.86,1>,
  <-0.1,-0.88,1>,
  <-0.1,-0.9,1>,
  <-0.1,-0.92,1>,
  <-0.1,-0.94,1>,
  <-0.1,-0.96,1>,
  <-0.1,-0.98,1>,
  <-0.1,-1,1>,
  <-0.12,0,1>,
  <-0.12,-0.02,1>,
  <-0.12,-0.04,1>,
  <-0.12,-0.06,1>,
  <-0.12,-0.08,1>,
  <-0.12,-0.1,1>,
  <-0.12,-0.12,1>,
  <-0.12,-0.14,1>,
  <-0.12,-0.16,1>,
  <-0.12,-0.18,1>,
  <-0.12,-0.2,1>,
  <-0.12,-0.22,1>,
  <-0.12,-0.24,1>,
  <-0.12,-0.26,1>,
  <-0.12,-0.28,1>,
  <-0.12,-0.3,1>,
  <-0.12,-0.32,1>,
  <-0.12,-0.34,1>,
  <-0.12,-0.36,1>,
  <-0.12,-0.38,1>,
  <-0.12,-0.4,1>,
  <-0.12,-0.42,1>,
  <-0.12,-0.44,1>,
  <-0.12,-0.46,1>,
  <-0.12,-0.48,1>,
  <-0.12,-0.5,1>,
  <-0.12,-0.52,1>,
  <-0.12,-0.54,1>,
  <-0.12,-0.56,1>,
  <-0.12,-0.58,1>,
  <-0.12,-0.6,1>,
  <-0.12,-0.62,1>,
  <-0.12,-0.64,1>,
  <-0.12,-0.66,1>,
  <-0.12,-0.68,1>,
  <-0.12,-0.7,1>,
  <-0.12,-0.72,1>,
  <-0.12,-0.74,1>,
  <-0.12,-0.76,1>,
  <-0.12,-0.78,1>,
  <-0.12,-0.8,1>,
  <-0.12,-0.82,1>,
  <-0.12,-0.84,1>,
  <-0.12,-0.86,1>,
  <-0.12,-0.88,1>,
  <-0.12,-0.9,1>,
  <-0.12,-0.92,1>,
  <-0.12,-0.94,1>,
  <-0.12,-0.96,1>,
  <-0.12,-0.98,1>,
  <-0.12,-1,1>,
  <-0.14,0,1>,
  <-0.14,-0.02,1>,
  <-0.14,-0.04,1>,
  <-0.14,-0.06,1>,
  <-0.14,-0.08,1>,
  <-0.14,-0.1,1>,
  <-0.14,-0.12,1>,
  <-0.14,-0.14,1>,
  <-0.14,-0.16,1>,
  <-0.14,-0.18,1>,
  <-0.14,-0.2,1>,
  <-0.14,-0.22,1>,
  <-0.14,-0.24,1>,
  <-0.14,-0.26,1>,
  <-0.14,-0.28,1>,
  <-0.14,-0.3,1>,
  <-0.14,-0.32,1>,
  <-0.14,-0.34,1>,
  <-0.14,-0.36,1>,
  <-0.14,-0.38,1>,
  <-0.14,-0.4,1>,
  <-0.14,-0.42,1>,
  <-0.14,-0.44,1>,
  <-0.14,-0.46,1>,
  <-0.14,-0.48,1>,
  <-0.14,-0.5,1>,
  <-0.14,-0.52,1>,
  <-0.14,-0.54,1>,
  <-0.14,-0.56,1>,
  <-0.14,-0.58,1>,
  <-0.14,-0.6,1>,
  <-0.14,-0.62,1>,
  <-0.14,-0.64,1>,
  <-0.14,-0.66,1>,
  <-0.14,-0.68,1>,
  <-0.14,-0.7,1>,
  <-0.14,-0.72,1>,
  <-0.14,-0.74,1>,
  <-0.14,-0.76,1>,
  <-0.14,-0.78,1>,
  <-0.14,-0.8,1>,
  <-0.14,-0.82,1>,
  <-0.14,-0.84,1>,
  <-0.14,-0.86,1>,
  <-0.14,-0.88,1>,
  <-0.14,-0.9,1>,
  <-0.14,-0.92,1>,
  <-0.14,-0.94,1>,
  <-0.14,-0.96,1>,
  <-0.14,-0.98,1>,
  <-0.14,-1,1>,
  <-0.16,0,1>,
  <-0.16,-0.02,1>,
  <-0.16,-0.04,1>,
  <-0.16,-0.06,1>,
  <-0.16,-0.08,1>,
  <-0.16,-0.1,1>,
  <-0.16,-0.12,1>,
  <-0.16,-0.14,1>,
  <-0.16,-0.16,1>,
  <-0.16,-0.18,1>,
  <-0.16,-0.2,1>,
  <-0.16,-0.22,1>,
  <-0.16,-0.24,1>,
  <-0.16,-0.26,1>,
  <-0.16,-0.28,1>,
  <-0.16,-0.3,1>,
  <-0.16,-0.32,1>,
  <-0.16,-0.34,1>,
  <-0.16,-0.36,1>,
  <-0.16,-0.38,1>,
  <-0.16,-0.4,1>,
  <-0.16,-0.42,1>,
  <-0.16,-0.44,1>,
  <-0.16,-0.46,1>,
  <-0.16,-0.48,1>,
  <-0.16,-0.5,1>,
  <-0.16,-0.52,1>,
  <-0.16,-0.54,1>,
  <-0.16,-0.56,1>,
  <-0.16,-0.58,1>,
  <-0.16,-0.6,1>,
  <-0.16,-0.62,1>,
  <-0.16,-0.64,1>,
  <-0.16,-0.66,1>,
  <-0.16,-0.68,1>,
  <-0.16,-0.7,1>,
  <-0.16,-0.72,1>,
  <-0.16,-0.74,1>,
  <-0.16,-0.76,1>,
  <-0.16,-0.78,1>,
  <-0.16,-0.8,1>,
  <-0.16,-0.82,1>,
  <-0.16,-0.84,1>,
  <-0.16,-0.86,1>,
  <-0.16,-0.88,1>,
  <-0.16,-0.9,1>,
  <-0.16,-0.92,1>,
  <-0.16,-0.94,1>,
  <-0.16,-0.96,1>,
  <-0.16,-0.98,1>,
  <-0.16,-1,1>,
  <-0.18,0,1>,
  <-0.18,-0.02,1>,
  <-0.18,-0.04,1>,
  <-0.18,-0.06,1>,
  <-0.18,-0.08,1>,
  <-0.18,-0.1,1>,
  <-0.18,-0.12,1>,
  <-0.18,-0.14,1>,
  <-0.18,-0.16,1>,
  <-0.18,-0.18,1>,
  <-0.18,-0.2,1>,
  <-0.18,-0.22,1>,
  <-0.18,-0.24,1>,
  <-0.18,-0.26,1>,
  <-0.18,-0.28,1>,
  <-0.18,-0.3,1>,
  <-0.18,-0.32,1>,
  <-0.18,-0.34,1>,
  <-0.18,-0.36,1>,
  <-0.18,-0.38,1>,
  <-0.18,-0.4,1>,
  <-0.18,-0.42,1>,
  <-0.18,-0.44,1>,
  <-0.18,-0.46,1>,
  <-0.18,-0.48,1>,
  <-0.18,-0.5,1>,
  <-0.18,-0.52,1>,
  <-0.18,-0.54,1>,
  <-0.18,-0.56,1>,
  <-0.18,-0.58,1>,
  <-0.18,-0.6,1>,
  <-0.18,-0.62,1>,
  <-0.18,-0.64,1>,
  <-0.18,-0.66,1>,
  <-0.18,-0.68,1>,
  <-0.18,-0.7,1>,
  <-0.18,-0.72,1>,
  <-0.18,-0.74,1>,
  <-0.18,-0.76,1>,
  <-0.18,-0.78,1>,
  <-0.18,-0.8,1>,
  <-0.18,-0.82,1>,
  <-0.18,-0.84,1>,
  <-0.18,-0.86,1>,
  <-0.18,-0.88,1>,
  <-0.18,-0.9,1>,
  <-0.18,-0.92,1>,
  <-0.18,-0.94,1>,
  <-0.18,-0.96,1>,
  <-0.18,-0.98,1>,
  <-0.18,-1,1>,
  <-0.2,0,1>,
  <-0.2,-0.02,1>,
  <-0.2,-0.04,1>,
  <-0.2,-0.06,1>,
  <-0.2,-0.08,1>,
  <-0.2,-0.1,1>,
  <-0.2,-0.12,1>,
  <-0.2,-0.14,1>,
  <-0.2,-0.16,1>,
  <-0.2,-0.18,1>,
  <-0.2,-0.2,1>,
  <-0.2,-0.22,1>,
  <-0.2,-0.24,1>,
  <-0.2,-0.26,1>,
  <-0.2,-0.28,1>,
  <-0.2,-0.3,1>,
  <-0.2,-0.32,1>,
  <-0.2,-0.34,1>,
  <-0.2,-0.36,1>,
  <-0.2,-0.38,1>,
  <-0.2,-0.4,1>,
  <-0.2,-0.42,1>,
  <-0.2,-0.44,1>,
  <-0.2,-0.46,1>,
  <-0.2,-0.48,1>,
  <-0.2,-0.5,1>,
  <-0.2,-0.52,1>,
  <-0.2,-0.54,1>,
  <-0.2,-0.56,1>,
  <-0.2,-0.58,1>,
  <-0.2,-0.6,1>,
  <-0.2,-0.62,1>,
  <-0.2,-0.64,1>,
  <-0.2,-0.66,1>,
  <-0.2,-0.68,1>,
  <-0.2,-0.7,1>,
  <-0.2,-0.72,1>,
  <-0.2,-0.74,1>,
  <-0.2,-0.76,1>,
  <-0.2,-0.78,1>,
  <-0.2,-0.8,1>,
  <-0.2,-0.82,1>,
  <-0.2,-0.84,1>,
  <-0.2,-0.86,1>,
  <-0.2,-0.88,1>,
  <-0.2,-0.9,1>,
  <-0.2,-0.92,1>,
  <-0.2,-0.94,1>,
  <-0.2,-0.96,1>,
  <-0.2,-0.98,1>,
  <-0.2,-1,1>,
  <-0.22,0,1>,
  <-0.22,-0.02,1>,
  <-0.22,-0.04,1>,
  <-0.22,-0.06,1>,
  <-0.22,-0.08,1>,
  <-0.22,-0.1,1>,
  <-0.22,-0.12,1>,
  <-0.22,-0.14,1>,
  <-0.22,-0.16,1>,
  <-0.22,-0.18,1>,
  <-0.22,-0.2,1>,
  <-0.22,-0.22,1>,
  <-0.22,-0.24,1>,
  <-0.22,-0.26,1>,
  <-0.22,-0.28,1>,
  <-0.22,-0.3,1>,
  <-0.22,-0.32,1>,
  <-0.22,-0.34,1>,
  <-0.22,-0.36,1>,
  <-0.22,-0.38,1>,
  <-0.22,-0.4,1>,
  <-0.22,-0.42,1>,
  <-0.22,-0.44,1>,
  <-0.22,-0.46,1>,
  <-0.22,-0.48,1>,
  <-0.22,-0.5,1>,
  <-0.22,-0.52,1>,
  <-0.22,-0.54,1>,
  <-0.22,-0.56,1>,
  <-0.22,-0.58,1>,
  <-0.22,-0.6,1>,
  <-0.22,-0.62,1>,
  <-0.22,-0.64,1>,
  <-0.22,-0.66,1>,
  <-0.22,-0.68,1>,
  <-0.22,-0.7,1>,
  <-0.22,-0.72,1>,
  <-0.22,-0.74,1>,
  <-0.22,-0.76,1>,
  <-0.22,-0.78,1>,
  <-0.22,-0.8,1>,
  <-0.22,-0.82,1>,
  <-0.22,-0.84,1>,
  <-0.22,-0.86,1>,
  <-0.22,-0.88,1>,
  <-0.22,-0.9,1>,
  <-0.22,-0.92,1>,
  <-0.22,-0.94,1>,
  <-0.22,-0.96,1>,
  <-0.22,-0.98,1>,
  <-0.22,-1,1>,
  <-0.24,0,1>,
  <-0.24,-0.02,1>,
  <-0.24,-0.04,1>,
  <-0.24,-0.06,1>,
  <-0.24,-0.08,1>,
  <-0.24,-0.1,1>,
  <-0.24,-0.12,1>,
  <-0.24,-0.14,1>,
  <-0.24,-0.16,1>,
  <-0.24,-0.18,1>,
  <-0.24,-0.2,1>,
  <-0.24,-0.22,1>,
  <-0.24,-0.24,1>,
  <-0.24,-0.26,1>,
  <-0.24,-0.28,1>,
  <-0.24,-0.3,1>,
  <-0.24,-0.32,1>,
  <-0.24,-0.34,1>,
  <-0.24,-0.36,1>,
  <-0.24,-0.38,1>,
  <-0.24,-0.4,1>,
  <-0.24,-0.42,1>,
  <-0.24,-0.44,1>,
  <-0.24,-0.46,1>,
  <-0.24,-0.48,1>,
  <-0.24,-0.5,1>,
  <-0.24,-0.52,1>,
  <-0.24,-0.54,1>,
  <-0.24,-0.56,1>,
  <-0.24,-0.58,1>,
  <-0.24,-0.6,1>,
  <-0.24,-0.62,1>,
  <-0.24,-0.64,1>,
  <-0.24,-0.66,1>,
  <-0.24,-0.68,1>,
  <-0.24,-0.7,1>,
  <-0.24,-0.72,1>,
  <-0.24,-0.74,1>,
  <-0.24,-0.76,1>,
  <-0.24,-0.78,1>,
  <-0.24,-0.8,1>,
  <-0.24,-0.82,1>,
  <-0.24,-0.84,1>,
  <-0.24,-0.86,1>,
  <-0.24,-0.88,1>,
  <-0.24,-0.9,1>,
  <-0.24,-0.92,1>,
  <-0.24,-0.94,1>,
  <-0.24,-0.96,1>,
  <-0.24,-0.98,1>,
  <-0.24,-1,1>,
  <-0.26,0,1>,
  <-0.26,-0.02,1>,
  <-0.26,-0.04,1>,
  <-0.26,-0.06,1>,
  <-0.26,-0.08,1>,
  <-0.26,-0.1,1>,
  <-0.26,-0.12,1>,
  <-0.26,-0.14,1>,
  <-0.26,-0.16,1>,
  <-0.26,-0.18,1>,
  <-0.26,-0.2,1>,
  <-0.26,-0.22,1>,
  <-0.26,-0.24,1>,
  <-0.26,-0.26,1>,
  <-0.26,-0.28,1>,
  <-0.26,-0.3,1>,
  <-0.26,-0.32,1>,
  <-0.26,-0.34,1>,
  <-0.26,-0.36,1>,
  <-0.26,-0.38,1>,
  <-0.26,-0.4,1>,
  <-0.26,-0.42,1>,
  <-0.26,-0.44,1>,
  <-0.26,-0.46,1>,
  <-0.26,-0.48,1>,
  <-0.26,-0.5,1>,
  <-0.26,-0.52,1>,
  <-0.26,-0.54,1>,
  <-0.26,-0.56,1>,
  <-0.26,-0.58,1>,
  <-0.26,-0.6,1>,
  <-0.26,-0.62,1>,
  <-0.26,-0.64,1>,
  <-0.26,-0.66,1>,
  <-0.26,-0.68,1>,
  <-0.26,-0.7,1>,
  <-0.26,-0.72,1>,
  <-0.26,-0.74,1>,
  <-0.26,-0.76,1>,
  <-0.26,-0.78,1>,
  <-0.26,-0.8,1>,
  <-0.26,-0.82,1>,
  <-0.26,-0.84,1>,
  <-0.26,-0.86,1>,
  <-0.26,-0.88,1>,
  <-0.26,-0.9,1>,
  <-0.26,-0.92,1>,
  <-0.26,-0.94,1>,
  <-0.26,-0.96,1>,
  <-0.26,-0.98,1>,
  <-0.26,-1,1>,
  <-0.28,0,1>,
  <-0.28,-0.02,1>,
  <-0.28,-0.04,1>,
  <-0.28,-0.06,1>,
  <-0.28,-0.08,1>,
  <-0.28,-0.1,1>,
  <-0.28,-0.12,1>,
  <-0.28,-0.14,1>,
  <-0.28,-0.16,1>,
  <-0.28,-0.18,1>,
  <-0.28,-0.2,1>,
  <-0.28,-0.22,1>,
  <-0.28,-0.24,1>,
  <-0.28,-0.26,1>,
  <-0.28,-0.28,1>,
  <-0.28,-0.3,1>,
  <-0.28,-0.32,1>,
  <-0.28,-0.34,1>,
  <-0.28,-0.36,1>,
  <-0.28,-0.38,1>,
  <-0.28,-0.4,1>,
  <-0.28,-0.42,1>,
  <-0.28,-0.44,1>,
  <-0.28,-0.46,1>,
  <-0.28,-0.48,1>,
  <-0.28,-0.5,1>,
  <-0.28,-0.52,1>,
  <-0.28,-0.54,1>,
  <-0.28,-0.56,1>,
  <-0.28,-0.58,1>,
  <-0.28,-0.6,1>,
  <-0.28,-0.62,1>,
  <-0.28,-0.64,1>,
  <-0.28,-0.66,1>,
  <-0.28,-0.68,1>,
  <-0.28,-0.7,1>,
  <-0.28,-0.72,1>,
  <-0.28,-0.74,1>,
  <-0.28,-0.76,1>,
  <-0.28,-0.78,1>,
  <-0.28,-0.8,1>,
  <-0.28,-0.82,1>,
  <-0.28,-0.84,1>,
  <-0.28,-0.86,1>,
  <-0.28,-0.88,1>,
  <-0.28,-0.9,1>,
  <-0.28,-0.92,1>,
  <-0.28,-0.94,1>,
  <-0.28,-0.96,1>,
  <-0.28,-0.98,1>,
  <-0.28,-1,1>,
  <-0.3,0,1>,
  <-0.3,-0.02,1>,
  <-0.3,-0.04,1>,
  <-0.3,-0.06,1>,
  <-0.3,-0.08,1>,
  <-0.3,-0.1,1>,
  <-0.3,-0.12,1>,
  <-0.3,-0.14,1>,
  <-0.3,-0.16,1>,
  <-0.3,-0.18,1>,
  <-0.3,-0.2,1>,
  <-0.3,-0.22,1>,
  <-0.3,-0.24,1>,
  <-0.3,-0.26,1>,
  <-0.3,-0.28,1>,
  <-0.3,-0.3,1>,
  <-0.3,-0.32,1>,
  <-0.3,-0.34,1>,
  <-0.3,-0.36,1>,
  <-0.3,-0.38,1>,
  <-0.3,-0.4,1>,
  <-0.3,-0.42,1>,
  <-0.3,-0.44,1>,
  <-0.3,-0.46,1>,
  <-0.3,-0.48,1>,
  <-0.3,-0.5,1>,
  <-0.3,-0.52,1>,
  <-0.3,-0.54,1>,
  <-0.3,-0.56,1>,
  <-0.3,-0.58,1>,
  <-0.3,-0.6,1>,
  <-0.3,-0.62,1>,
  <-0.3,-0.64,1>,
  <-0.3,-0.66,1>,
  <-0.3,-0.68,1>,
  <-0.3,-0.7,1>,
  <-0.3,-0.72,1>,
  <-0.3,-0.74,1>,
  <-0.3,-0.76,1>,
  <-0.3,-0.78,1>,
  <-0.3,-0.8,1>,
  <-0.3,-0.82,1>,
  <-0.3,-0.84,1>,
  <-0.3,-0.86,1>,
  <-0.3,-0.88,1>,
  <-0.3,-0.9,1>,
  <-0.3,-0.92,1>,
  <-0.3,-0.94,1>,
  <-0.3,-0.96,1>,
  <-0.3,-0.98,1>,
  <-0.3,-1,1>,
  <-0.32,0,1>,
  <-0.32,-0.02,1>,
  <-0.32,-0.04,1>,
  <-0.32,-0.06,1>,
  <-0.32,-0.08,1>,
  <-0.32,-0.1,1>,
  <-0.32,-0.12,1>,
  <-0.32,-0.14,1>,
  <-0.32,-0.16,1>,
  <-0.32,-0.18,1>,
  <-0.32,-0.2,1>,
  <-0.32,-0.22,1>,
  <-0.32,-0.24,1>,
  <-0.32,-0.26,1>,
  <-0.32,-0.28,1>,
  <-0.32,-0.3,1>,
  <-0.32,-0.32,1>,
  <-0.32,-0.34,1>,
  <-0.32,-0.36,1>,
  <-0.32,-0.38,1>,
  <-0.32,-0.4,1>,
  <-0.32,-0.42,1>,
  <-0.32,-0.44,1>,
  <-0.32,-0.46,1>,
  <-0.32,-0.48,1>,
  <-0.32,-0.5,1>,
  <-0.32,-0.52,1>,
  <-0.32,-0.54,1>,
  <-0.32,-0.56,1>,
  <-0.32,-0.58,1>,
  <-0.32,-0.6,1>,
  <-0.32,-0.62,1>,
  <-0.32,-0.64,1>,
  <-0.32,-0.66,1>,
  <-0.32,-0.68,1>,
  <-0.32,-0.7,1>,
  <-0.32,-0.72,1>,
  <-0.32,-0.74,1>,
  <-0.32,-0.76,1>,
  <-0.32,-0.78,1>,
  <-0.32,-0.8,1>,
  <-0.32,-0.82,1>,
  <-0.32,-0.84,1>,
  <-0.32,-0.86,1>,
  <-0.32,-0.88,1>,
  <-0.32,-0.9,1>,
  <-0.32,-0.92,1>,
  <-0.32,-0.94,1>,
  <-0.32,-0.96,1>,
  <-0.32,-0.98,1>,
  <-0.32,-1,1>,
  <-0.34,0,1>,
  <-0.34,-0.02,1>,
  <-0.34,-0.04,1>,
  <-0.34,-0.06,1>,
  <-0.34,-0.08,1>,
  <-0.34,-0.1,1>,
  <-0.34,-0.12,1>,
  <-0.34,-0.14,1>,
  <-0.34,-0.16,1>,
  <-0.34,-0.18,1>,
  <-0.34,-0.2,1>,
  <-0.34,-0.22,1>,
  <-0.34,-0.24,1>,
  <-0.34,-0.26,1>,
  <-0.34,-0.28,1>,
  <-0.34,-0.3,1>,
  <-0.34,-0.32,1>,
  <-0.34,-0.34,1>,
  <-0.34,-0.36,1>,
  <-0.34,-0.38,1>,
  <-0.34,-0.4,1>,
  <-0.34,-0.42,1>,
  <-0.34,-0.44,1>,
  <-0.34,-0.46,1>,
  <-0.34,-0.48,1>,
  <-0.34,-0.5,1>,
  <-0.34,-0.52,1>,
  <-0.34,-0.54,1>,
  <-0.34,-0.56,1>,
  <-0.34,-0.58,1>,
  <-0.34,-0.6,1>,
  <-0.34,-0.62,1>,
  <-0.34,-0.64,1>,
  <-0.34,-0.66,1>,
  <-0.34,-0.68,1>,
  <-0.34,-0.7,1>,
  <-0.34,-0.72,1>,
  <-0.34,-0.74,1>,
  <-0.34,-0.76,1>,
  <-0.34,-0.78,1>,
  <-0.34,-0.8,1>,
  <-0.34,-0.82,1>,
  <-0.34,-0.84,1>,
  <-0.34,-0.86,1>,
  <-0.34,-0.88,1>,
  <-0.34,-0.9,1>,
  <-0.34,-0.92,1>,
  <-0.34,-0.94,1>,
  <-0.34,-0.96,1>,
  <-0.34,-0.98,1>,
  <-0.34,-1,1>,
  <-0.36,0,1>,
  <-0.36,-0.02,1>,
  <-0.36,-0.04,1>,
  <-0.36,-0.06,1>,
  <-0.36,-0.08,1>,
  <-0.36,-0.1,1>,
  <-0.36,-0.12,1>,
  <-0.36,-0.14,1>,
  <-0.36,-0.16,1>,
  <-0.36,-0.18,1>,
  <-0.36,-0.2,1>,
  <-0.36,-0.22,1>,
  <-0.36,-0.24,1>,
  <-0.36,-0.26,1>,
  <-0.36,-0.28,1>,
  <-0.36,-0.3,1>,
  <-0.36,-0.32,1>,
  <-0.36,-0.34,1>,
  <-0.36,-0.36,1>,
  <-0.36,-0.38,1>,
  <-0.36,-0.4,1>,
  <-0.36,-0.42,1>,
  <-0.36,-0.44,1>,
  <-0.36,-0.46,1>,
  <-0.36,-0.48,1>,
  <-0.36,-0.5,1>,
  <-0.36,-0.52,1>,
  <-0.36,-0.54,1>,
  <-0.36,-0.56,1>,
  <-0.36,-0.58,1>,
  <-0.36,-0.6,1>,
  <-0.36,-0.62,1>,
  <-0.36,-0.64,1>,
  <-0.36,-0.66,1>,
  <-0.36,-0.68,1>,
  <-0.36,-0.7,1>,
  <-0.36,-0.72,1>,
  <-0.36,-0.74,1>,
  <-0.36,-0.76,1>,
  <-0.36,-0.78,1>,
  <-0.36,-0.8,1>,
  <-0.36,-0.82,1>,
  <-0.36,-0.84,1>,
  <-0.36,-0.86,1>,
  <-0.36,-0.88,1>,
  <-0.36,-0.9,1>,
  <-0.36,-0.92,1>,
  <-0.36,-0.94,1>,
  <-0.36,-0.96,1>,
  <-0.36,-0.98,1>,
  <-0.36,-1,1>,
  <-0.38,0,1>,
  <-0.38,-0.02,1>,
  <-0.38,-0.04,1>,
  <-0.38,-0.06,1>,
  <-0.38,-0.08,1>,
  <-0.38,-0.1,1>,
  <-0.38,-0.12,1>,
  <-0.38,-0.14,1>,
  <-0.38,-0.16,1>,
  <-0.38,-0.18,1>,
  <-0.38,-0.2,1>,
  <-0.38,-0.22,1>,
  <-0.38,-0.24,1>,
  <-0.38,-0.26,1>,
  <-0.38,-0.28,1>,
  <-0.38,-0.3,1>,
  <-0.38,-0.32,1>,
  <-0.38,-0.34,1>,
  <-0.38,-0.36,1>,
  <-0.38,-0.38,1>,
  <-0.38,-0.4,1>,
  <-0.38,-0.42,1>,
  <-0.38,-0.44,1>,
  <-0.38,-0.46,1>,
  <-0.38,-0.48,1>,
  <-0.38,-0.5,1>,
  <-0.38,-0.52,1>,
  <-0.38,-0.54,1>,
  <-0.38,-0.56,1>,
  <-0.38,-0.58,1>,
  <-0.38,-0.6,1>,
  <-0.38,-0.62,1>,
  <-0.38,-0.64,1>,
  <-0.38,-0.66,1>,
  <-0.38,-0.68,1>,
  <-0.38,-0.7,1>,
  <-0.38,-0.72,1>,
  <-0.38,-0.74,1>,
  <-0.38,-0.76,1>,
  <-0.38,-0.78,1>,
  <-0.38,-0.8,1>,
  <-0.38,-0.82,1>,
  <-0.38,-0.84,1>,
  <-0.38,-0.86,1>,
  <-0.38,-0.88,1>,
  <-0.38,-0.9,1>,
  <-0.38,-0.92,1>,
  <-0.38,-0.94,1>,
  <-0.38,-0.96,1>,
  <-0.38,-0.98,1>,
  <-0.38,-1,1>,
  <-0.4,0,1>,
  <-0.4,-0.02,1>,
  <-0.4,-0.04,1>,
  <-0.4,-0.06,1>,
  <-0.4,-0.08,1>,
  <-0.4,-0.1,1>,
  <-0.4,-0.12,1>,
  <-0.4,-0.14,1>,
  <-0.4,-0.16,1>,
  <-0.4,-0.18,1>,
  <-0.4,-0.2,1>,
  <-0.4,-0.22,1>,
  <-0.4,-0.24,1>,
  <-0.4,-0.26,1>,
  <-0.4,-0.28,1>,
  <-0.4,-0.3,1>,
  <-0.4,-0.32,1>,
  <-0.4,-0.34,1>,
  <-0.4,-0.36,1>,
  <-0.4,-0.38,1>,
  <-0.4,-0.4,1>,
  <-0.4,-0.42,1>,
  <-0.4,-0.44,1>,
  <-0.4,-0.46,1>,
  <-0.4,-0.48,1>,
  <-0.4,-0.5,1>,
  <-0.4,-0.52,1>,
  <-0.4,-0.54,1>,
  <-0.4,-0.56,1>,
  <-0.4,-0.58,1>,
  <-0.4,-0.6,1>,
  <-0.4,-0.62,1>,
  <-0.4,-0.64,1>,
  <-0.4,-0.66,1>,
  <-0.4,-0.68,1>,
  <-0.4,-0.7,1>,
  <-0.4,-0.72,1>,
  <-0.4,-0.74,1>,
  <-0.4,-0.76,1>,
  <-0.4,-0.78,1>,
  <-0.4,-0.8,1>,
  <-0.4,-0.82,1>,
  <-0.4,-0.84,1>,
  <-0.4,-0.86,1>,
  <-0.4,-0.88,1>,
  <-0.4,-0.9,1>,
  <-0.4,-0.92,1>,
  <-0.4,-0.94,1>,
  <-0.4,-0.96,1>,
  <-0.4,-0.98,1>,
  <-0.4,-1,1>,
  <-0.42,0,1>,
  <-0.42,-0.02,1>,
  <-0.42,-0.04,1>,
  <-0.42,-0.06,1>,
  <-0.42,-0.08,1>,
  <-0.42,-0.1,1>,
  <-0.42,-0.12,1>,
  <-0.42,-0.14,1>,
  <-0.42,-0.16,1>,
  <-0.42,-0.18,1>,
  <-0.42,-0.2,1>,
  <-0.42,-0.22,1>,
  <-0.42,-0.24,1>,
  <-0.42,-0.26,1>,
  <-0.42,-0.28,1>,
  <-0.42,-0.3,1>,
  <-0.42,-0.32,1>,
  <-0.42,-0.34,1>,
  <-0.42,-0.36,1>,
  <-0.42,-0.38,1>,
  <-0.42,-0.4,1>,
  <-0.42,-0.42,1>,
  <-0.42,-0.44,1>,
  <-0.42,-0.46,1>,
  <-0.42,-0.48,1>,
  <-0.42,-0.5,1>,
  <-0.42,-0.52,1>,
  <-0.42,-0.54,1>,
  <-0.42,-0.56,1>,
  <-0.42,-0.58,1>,
  <-0.42,-0.6,1>,
  <-0.42,-0.62,1>,
  <-0.42,-0.64,1>,
  <-0.42,-0.66,1>,
  <-0.42,-0.68,1>,
  <-0.42,-0.7,1>,
  <-0.42,-0.72,1>,
  <-0.42,-0.74,1>,
  <-0.42,-0.76,1>,
  <-0.42,-0.78,1>,
  <-0.42,-0.8,1>,
  <-0.42,-0.82,1>,
  <-0.42,-0.84,1>,
  <-0.42,-0.86,1>,
  <-0.42,-0.88,1>,
  <-0.42,-0.9,1>,
  <-0.42,-0.92,1>,
  <-0.42,-0.94,1>,
  <-0.42,-0.96,1>,
  <-0.42,-0.98,1>,
  <-0.42,-1,1>,
  <-0.44,0,1>,
  <-0.44,-0.02,1>,
  <-0.44,-0.04,1>,
  <-0.44,-0.06,1>,
  <-0.44,-0.08,1>,
  <-0.44,-0.1,1>,
  <-0.44,-0.12,1>,
  <-0.44,-0.14,1>,
  <-0.44,-0.16,1>,
  <-0.44,-0.18,1>,
  <-0.44,-0.2,1>,
  <-0.44,-0.22,1>,
  <-0.44,-0.24,1>,
  <-0.44,-0.26,1>,
  <-0.44,-0.28,1>,
  <-0.44,-0.3,1>,
  <-0.44,-0.32,1>,
  <-0.44,-0.34,1>,
  <-0.44,-0.36,1>,
  <-0.44,-0.38,1>,
  <-0.44,-0.4,1>,
  <-0.44,-0.42,1>,
  <-0.44,-0.44,1>,
  <-0.44,-0.46,1>,
  <-0.44,-0.48,1>,
  <-0.44,-0.5,1>,
  <-0.44,-0.52,1>,
  <-0.44,-0.54,1>,
  <-0.44,-0.56,1>,
  <-0.44,-0.58,1>,
  <-0.44,-0.6,1>,
  <-0.44,-0.62,1>,
  <-0.44,-0.64,1>,
  <-0.44,-0.66,1>,
  <-0.44,-0.68,1>,
  <-0.44,-0.7,1>,
  <-0.44,-0.72,1>,
  <-0.44,-0.74,1>,
  <-0.44,-0.76,1>,
  <-0.44,-0.78,1>,
  <-0.44,-0.8,1>,
  <-0.44,-0.82,1>,
  <-0.44,-0.84,1>,
  <-0.44,-0.86,1>,
  <-0.44,-0.88,1>,
  <-0.44,-0.9,1>,
  <-0.44,-0.92,1>,
  <-0.44,-0.94,1>,
  <-0.44,-0.96,1>,
  <-0.44,-0.98,1>,
  <-0.44,-1,1>,
  <-0.46,0,1>,
  <-0.46,-0.02,1>,
  <-0.46,-0.04,1>,
  <-0.46,-0.06,1>,
  <-0.46,-0.08,1>,
  <-0.46,-0.1,1>,
  <-0.46,-0.12,1>,
  <-0.46,-0.14,1>,
  <-0.46,-0.16,1>,
  <-0.46,-0.18,1>,
  <-0.46,-0.2,1>,
  <-0.46,-0.22,1>,
  <-0.46,-0.24,1>,
  <-0.46,-0.26,1>,
  <-0.46,-0.28,1>,
  <-0.46,-0.3,1>,
  <-0.46,-0.32,1>,
  <-0.46,-0.34,1>,
  <-0.46,-0.36,1>,
  <-0.46,-0.38,1>,
  <-0.46,-0.4,1>,
  <-0.46,-0.42,1>,
  <-0.46,-0.44,1>,
  <-0.46,-0.46,1>,
  <-0.46,-0.48,1>,
  <-0.46,-0.5,1>,
  <-0.46,-0.52,1>,
  <-0.46,-0.54,1>,
  <-0.46,-0.56,1>,
  <-0.46,-0.58,1>,
  <-0.46,-0.6,1>,
  <-0.46,-0.62,1>,
  <-0.46,-0.64,1>,
  <-0.46,-0.66,1>,
  <-0.46,-0.68,1>,
  <-0.46,-0.7,1>,
  <-0.46,-0.72,1>,
  <-0.46,-0.74,1>,
  <-0.46,-0.76,1>,
  <-0.46,-0.78,1>,
  <-0.46,-0.8,1>,
  <-0.46,-0.82,1>,
  <-0.46,-0.84,1>,
  <-0.46,-0.86,1>,
  <-0.46,-0.88,1>,
  <-0.46,-0.9,1>,
  <-0.46,-0.92,1>,
  <-0.46,-0.94,1>,
  <-0.46,-0.96,1>,
  <-0.46,-0.98,1>,
  <-0.46,-1,1>,
  <-0.48,0,1>,
  <-0.48,-0.02,1>,
  <-0.48,-0.04,1>,
  <-0.48,-0.06,1>,
  <-0.48,-0.08,1>,
  <-0.48,-0.1,1>,
  <-0.48,-0.12,1>,
  <-0.48,-0.14,1>,
  <-0.48,-0.16,1>,
  <-0.48,-0.18,1>,
  <-0.48,-0.2,1>,
  <-0.48,-0.22,1>,
  <-0.48,-0.24,1>,
  <-0.48,-0.26,1>,
  <-0.48,-0.28,1>,
  <-0.48,-0.3,1>,
  <-0.48,-0.32,1>,
  <-0.48,-0.34,1>,
  <-0.48,-0.36,1>,
  <-0.48,-0.38,1>,
  <-0.48,-0.4,1>,
  <-0.48,-0.42,1>,
  <-0.48,-0.44,1>,
  <-0.48,-0.46,1>,
  <-0.48,-0.48,1>,
  <-0.48,-0.5,1>,
  <-0.48,-0.52,1>,
  <-0.48,-0.54,1>,
  <-0.48,-0.56,1>,
  <-0.48,-0.58,1>,
  <-0.48,-0.6,1>,
  <-0.48,-0.62,1>,
  <-0.48,-0.64,1>,
  <-0.48,-0.66,1>,
  <-0.48,-0.68,1>,
  <-0.48,-0.7,1>,
  <-0.48,-0.72,1>,
  <-0.48,-0.74,1>,
  <-0.48,-0.76,1>,
  <-0.48,-0.78,1>,
  <-0.48,-0.8,1>,
  <-0.48,-0.82,1>,
  <-0.48,-0.84,1>,
  <-0.48,-0.86,1>,
  <-0.48,-0.88,1>,
  <-0.48,-0.9,1>,
  <-0.48,-0.92,1>,
  <-0.48,-0.94,1>,
  <-0.48,-0.96,1>,
  <-0.48,-0.98,1>,
  <-0.48,-1,1>,
  <-0.5,0,1>,
  <-0.5,-0.02,1>,
  <-0.5,-0.04,1>,
  <-0.5,-0.06,1>,
  <-0.5,-0.08,1>,
  <-0.5,-0.1,1>,
  <-0.5,-0.12,1>,
  <-0.5,-0.14,1>,
  <-0.5,-0.16,1>,
  <-0.5,-0.18,1>,
  <-0.5,-0.2,1>,
  <-0.5,-0.22,1>,
  <-0.5,-0.24,1>,
  <-0.5,-0.26,1>,
  <-0.5,-0.28,1>,
  <-0.5,-0.3,1>,
  <-0.5,-0.32,1>,
  <-0.5,-0.34,1>,
  <-0.5,-0.36,1>,
  <-0.5,-0.38,1>,
  <-0.5,-0.4,1>,
  <-0.5,-0.42,1>,
  <-0.5,-0.44,1>,
  <-0.5,-0.46,1>,
  <-0.5,-0.48,1>,
  <-0.5,-0.5,1>,
  <-0.5,-0.52,1>,
  <-0.5,-0.54,1>,
  <-0.5,-0.56,1>,
  <-0.5,-0.58,1>,
  <-0.5,-0.6,1>,
  <-0.5,-0.62,1>,
  <-0.5,-0.64,1>,
  <-0.5,-0.66,1>,
  <-0.5,-0.68,1>,
  <-0.5,-0.7,1>,
  <-0.5,-0.72,1>,
  <-0.5,-0.74,1>,
  <-0.5,-0.76,1>,
  <-0.5,-0.78,1>,
  <-0.5,-0.8,1>,
  <-0.5,-0.82,1>,
  <-0.5,-0.84,1>,
  <-0.5,-0.86,1>,
  <-0.5,-0.88,1>,
  <-0.5,-0.9,1>,
  <-0.5,-0.92,1>,
  <-0.5,-0.94,1>,
  <-0.5,-0.96,1>,
  <-0.5,-0.98,1>,
  <-0.5,-1,1>,
  <-0.52,0,1>,
  <-0.52,-0.02,1>,
  <-0.52,-0.04,1>,
  <-0.52,-0.06,1>,
  <-0.52,-0.08,1>,
  <-0.52,-0.1,1>,
  <-0.52,-0.12,1>,
  <-0.52,-0.14,1>,
  <-0.52,-0.16,1>,
  <-0.52,-0.18,1>,
  <-0.52,-0.2,1>,
  <-0.52,-0.22,1>,
  <-0.52,-0.24,1>,
  <-0.52,-0.26,1>,
  <-0.52,-0.28,1>,
  <-0.52,-0.3,1>,
  <-0.52,-0.32,1>,
  <-0.52,-0.34,1>,
  <-0.52,-0.36,1>,
  <-0.52,-0.38,1>,
  <-0.52,-0.4,1>,
  <-0.52,-0.42,1>,
  <-0.52,-0.44,1>,
  <-0.52,-0.46,1>,
  <-0.52,-0.48,1>,
  <-0.52,-0.5,1>,
  <-0.52,-0.52,1>,
  <-0.52,-0.54,1>,
  <-0.52,-0.56,1>,
  <-0.52,-0.58,1>,
  <-0.52,-0.6,1>,
  <-0.52,-0.62,1>,
  <-0.52,-0.64,1>,
  <-0.52,-0.66,1>,
  <-0.52,-0.68,1>,
  <-0.52,-0.7,1>,
  <-0.52,-0.72,1>,
  <-0.52,-0.74,1>,
  <-0.52,-0.76,1>,
  <-0.52,-0.78,1>,
  <-0.52,-0.8,1>,
  <-0.52,-0.82,1>,
  <-0.52,-0.84,1>,
  <-0.52,-0.86,1>,
  <-0.52,-0.88,1>,
  <-0.52,-0.9,1>,
  <-0.52,-0.92,1>,
  <-0.52,-0.94,1>,
  <-0.52,-0.96,1>,
  <-0.52,-0.98,1>,
  <-0.52,-1,1>,
  <-0.54,0,1>,
  <-0.54,-0.02,1>,
  <-0.54,-0.04,1>,
  <-0.54,-0.06,1>,
  <-0.54,-0.08,1>,
  <-0.54,-0.1,1>,
  <-0.54,-0.12,1>,
  <-0.54,-0.14,1>,
  <-0.54,-0.16,1>,
  <-0.54,-0.18,1>,
  <-0.54,-0.2,1>,
  <-0.54,-0.22,1>,
  <-0.54,-0.24,1>,
  <-0.54,-0.26,1>,
  <-0.54,-0.28,1>,
  <-0.54,-0.3,1>,
  <-0.54,-0.32,1>,
  <-0.54,-0.34,1>,
  <-0.54,-0.36,1>,
  <-0.54,-0.38,1>,
  <-0.54,-0.4,1>,
  <-0.54,-0.42,1>,
  <-0.54,-0.44,1>,
  <-0.54,-0.46,1>,
  <-0.54,-0.48,1>,
  <-0.54,-0.5,1>,
  <-0.54,-0.52,1>,
  <-0.54,-0.54,1>,
  <-0.54,-0.56,1>,
  <-0.54,-0.58,1>,
  <-0.54,-0.6,1>,
  <-0.54,-0.62,1>,
  <-0.54,-0.64,1>,
  <-0.54,-0.66,1>,
  <-0.54,-0.68,1>,
  <-0.54,-0.7,1>,
  <-0.54,-0.72,1>,
  <-0.54,-0.74,1>,
  <-0.54,-0.76,1>,
  <-0.54,-0.78,1>,
  <-0.54,-0.8,1>,
  <-0.54,-0.82,1>,
  <-0.54,-0.84,1>,
  <-0.54,-0.86,1>,
  <-0.54,-0.88,1>,
  <-0.54,-0.9,1>,
  <-0.54,-0.92,1>,
  <-0.54,-0.94,1>,
  <-0.54,-0.96,1>,
  <-0.54,-0.98,1>,
  <-0.54,-1,1>,
  <-0.56,0,1>,
  <-0.56,-0.02,1>,
  <-0.56,-0.04,1>,
  <-0.56,-0.06,1>,
  <-0.56,-0.08,1>,
  <-0.56,-0.1,1>,
  <-0.56,-0.12,1>,
  <-0.56,-0.14,1>,
  <-0.56,-0.16,1>,
  <-0.56,-0.18,1>,
  <-0.56,-0.2,1>,
  <-0.56,-0.22,1>,
  <-0.56,-0.24,1>,
  <-0.56,-0.26,1>,
  <-0.56,-0.28,1>,
  <-0.56,-0.3,1>,
  <-0.56,-0.32,1>,
  <-0.56,-0.34,1>,
  <-0.56,-0.36,1>,
  <-0.56,-0.38,1>,
  <-0.56,-0.4,1>,
  <-0.56,-0.42,1>,
  <-0.56,-0.44,1>,
  <-0.56,-0.46,1>,
  <-0.56,-0.48,1>,
  <-0.56,-0.5,1>,
  <-0.56,-0.52,1>,
  <-0.56,-0.54,1>,
  <-0.56,-0.56,1>,
  <-0.56,-0.58,1>,
  <-0.56,-0.6,1>,
  <-0.56,-0.62,1>,
  <-0.56,-0.64,1>,
  <-0.56,-0.66,1>,
  <-0.56,-0.68,1>,
  <-0.56,-0.7,1>,
  <-0.56,-0.72,1>,
  <-0.56,-0.74,1>,
  <-0.56,-0.76,1>,
  <-0.56,-0.78,1>,
  <-0.56,-0.8,1>,
  <-0.56,-0.82,1>,
  <-0.56,-0.84,1>,
  <-0.56,-0.86,1>,
  <-0.56,-0.88,1>,
  <-0.56,-0.9,1>,
  <-0.56,-0.92,1>,
  <-0.56,-0.94,1>,
  <-0.56,-0.96,1>,
  <-0.56,-0.98,1>,
  <-0.56,-1,1>,
  <-0.58,0,1>,
  <-0.58,-0.02,1>,
  <-0.58,-0.04,1>,
  <-0.58,-0.06,1>,
  <-0.58,-0.08,1>,
  <-0.58,-0.1,1>,
  <-0.58,-0.12,1>,
  <-0.58,-0.14,1>,
  <-0.58,-0.16,1>,
  <-0.58,-0.18,1>,
  <-0.58,-0.2,1>,
  <-0.58,-0.22,1>,
  <-0.58,-0.24,1>,
  <-0.58,-0.26,1>,
  <-0.58,-0.28,1>,
  <-0.58,-0.3,1>,
  <-0.58,-0.32,1>,
  <-0.58,-0.34,1>,
  <-0.58,-0.36,1>,
  <-0.58,-0.38,1>,
  <-0.58,-0.4,1>,
  <-0.58,-0.42,1>,
  <-0.58,-0.44,1>,
  <-0.58,-0.46,1>,
  <-0.58,-0.48,1>,
  <-0.58,-0.5,1>,
  <-0.58,-0.52,1>,
  <-0.58,-0.54,1>,
  <-0.58,-0.56,1>,
  <-0.58,-0.58,1>,
  <-0.58,-0.6,1>,
  <-0.58,-0.62,1>,
  <-0.58,-0.64,1>,
  <-0.58,-0.66,1>,
  <-0.58,-0.68,1>,
  <-0.58,-0.7,1>,
  <-0.58,-0.72,1>,
  <-0.58,-0.74,1>,
  <-0.58,-0.76,1>,
  <-0.58,-0.78,1>,
  <-0.58,-0.8,1>,
  <-0.58,-0.82,1>,
  <-0.58,-0.84,1>,
  <-0.58,-0.86,1>,
  <-0.58,-0.88,1>,
  <-0.58,-0.9,1>,
  <-0.58,-0.92,1>,
  <-0.58,-0.94,1>,
  <-0.58,-0.96,1>,
  <-0.58,-0.98,1>,
  <-0.58,-1,1>,
  <-0.6,0,1>,
  <-0.6,-0.02,1>,
  <-0.6,-0.04,1>,
  <-0.6,-0.06,1>,
  <-0.6,-0.08,1>,
  <-0.6,-0.1,1>,
  <-0.6,-0.12,1>,
  <-0.6,-0.14,1>,
  <-0.6,-0.16,1>,
  <-0.6,-0.18,1>,
  <-0.6,-0.2,1>,
  <-0.6,-0.22,1>,
  <-0.6,-0.24,1>,
  <-0.6,-0.26,1>,
  <-0.6,-0.28,1>,
  <-0.6,-0.3,1>,
  <-0.6,-0.32,1>,
  <-0.6,-0.34,1>,
  <-0.6,-0.36,1>,
  <-0.6,-0.38,1>,
  <-0.6,-0.4,1>,
  <-0.6,-0.42,1>,
  <-0.6,-0.44,1>,
  <-0.6,-0.46,1>,
  <-0.6,-0.48,1>,
  <-0.6,-0.5,1>,
  <-0.6,-0.52,1>,
  <-0.6,-0.54,1>,
  <-0.6,-0.56,1>,
  <-0.6,-0.58,1>,
  <-0.6,-0.6,1>,
  <-0.6,-0.62,1>,
  <-0.6,-0.64,1>,
  <-0.6,-0.66,1>,
  <-0.6,-0.68,1>,
  <-0.6,-0.7,1>,
  <-0.6,-0.72,1>,
  <-0.6,-0.74,1>,
  <-0.6,-0.76,1>,
  <-0.6,-0.78,1>,
  <-0.6,-0.8,1>,
  <-0.6,-0.82,1>,
  <-0.6,-0.84,1>,
  <-0.6,-0.86,1>,
  <-0.6,-0.88,1>,
  <-0.6,-0.9,1>,
  <-0.6,-0.92,1>,
  <-0.6,-0.94,1>,
  <-0.6,-0.96,1>,
  <-0.6,-0.98,1>,
  <-0.6,-1,1>,
  <-0.62,0,1>,
  <-0.62,-0.02,1>,
  <-0.62,-0.04,1>,
  <-0.62,-0.06,1>,
  <-0.62,-0.08,1>,
  <-0.62,-0.1,1>,
  <-0.62,-0.12,1>,
  <-0.62,-0.14,1>,
  <-0.62,-0.16,1>,
  <-0.62,-0.18,1>,
  <-0.62,-0.2,1>,
  <-0.62,-0.22,1>,
  <-0.62,-0.24,1>,
  <-0.62,-0.26,1>,
  <-0.62,-0.28,1>,
  <-0.62,-0.3,1>,
  <-0.62,-0.32,1>,
  <-0.62,-0.34,1>,
  <-0.62,-0.36,1>,
  <-0.62,-0.38,1>,
  <-0.62,-0.4,1>,
  <-0.62,-0.42,1>,
  <-0.62,-0.44,1>,
  <-0.62,-0.46,1>,
  <-0.62,-0.48,1>,
  <-0.62,-0.5,1>,
  <-0.62,-0.52,1>,
  <-0.62,-0.54,1>,
  <-0.62,-0.56,1>,
  <-0.62,-0.58,1>,
  <-0.62,-0.6,1>,
  <-0.62,-0.62,1>,
  <-0.62,-0.64,1>,
  <-0.62,-0.66,1>,
  <-0.62,-0.68,1>,
  <-0.62,-0.7,1>,
  <-0.62,-0.72,1>,
  <-0.62,-0.74,1>,
  <-0.62,-0.76,1>,
  <-0.62,-0.78,1>,
  <-0.62,-0.8,1>,
  <-0.62,-0.82,1>,
  <-0.62,-0.84,1>,
  <-0.62,-0.86,1>,
  <-0.62,-0.88,1>,
  <-0.62,-0.9,1>,
  <-0.62,-0.92,1>,
  <-0.62,-0.94,1>,
  <-0.62,-0.96,1>,
  <-0.62,-0.98,1>,
  <-0.62,-1,1>,
  <-0.64,0,1>,
  <-0.64,-0.02,1>,
  <-0.64,-0.04,1>,
  <-0.64,-0.06,1>,
  <-0.64,-0.08,1>,
  <-0.64,-0.1,1>,
  <-0.64,-0.12,1>,
  <-0.64,-0.14,1>,
  <-0.64,-0.16,1>,
  <-0.64,-0.18,1>,
  <-0.64,-0.2,1>,
  <-0.64,-0.22,1>,
  <-0.64,-0.24,1>,
  <-0.64,-0.26,1>,
  <-0.64,-0.28,1>,
  <-0.64,-0.3,1>,
  <-0.64,-0.32,1>,
  <-0.64,-0.34,1>,
  <-0.64,-0.36,1>,
  <-0.64,-0.38,1>,
  <-0.64,-0.4,1>,
  <-0.64,-0.42,1>,
  <-0.64,-0.44,1>,
  <-0.64,-0.46,1>,
  <-0.64,-0.48,1>,
  <-0.64,-0.5,1>,
  <-0.64,-0.52,1>,
  <-0.64,-0.54,1>,
  <-0.64,-0.56,1>,
  <-0.64,-0.58,1>,
  <-0.64,-0.6,1>,
  <-0.64,-0.62,1>,
  <-0.64,-0.64,1>,
  <-0.64,-0.66,1>,
  <-0.64,-0.68,1>,
  <-0.64,-0.7,1>,
  <-0.64,-0.72,1>,
  <-0.64,-0.74,1>,
  <-0.64,-0.76,1>,
  <-0.64,-0.78,1>,
  <-0.64,-0.8,1>,
  <-0.64,-0.82,1>,
  <-0.64,-0.84,1>,
  <-0.64,-0.86,1>,
  <-0.64,-0.88,1>,
  <-0.64,-0.9,1>,
  <-0.64,-0.92,1>,
  <-0.64,-0.94,1>,
  <-0.64,-0.96,1>,
  <-0.64,-0.98,1>,
  <-0.64,-1,1>,
  <-0.66,0,1>,
  <-0.66,-0.02,1>,
  <-0.66,-0.04,1>,
  <-0.66,-0.06,1>,
  <-0.66,-0.08,1>,
  <-0.66,-0.1,1>,
  <-0.66,-0.12,1>,
  <-0.66,-0.14,1>,
  <-0.66,-0.16,1>,
  <-0.66,-0.18,1>,
  <-0.66,-0.2,1>,
  <-0.66,-0.22,1>,
  <-0.66,-0.24,1>,
  <-0.66,-0.26,1>,
  <-0.66,-0.28,1>,
  <-0.66,-0.3,1>,
  <-0.66,-0.32,1>,
  <-0.66,-0.34,1>,
  <-0.66,-0.36,1>,
  <-0.66,-0.38,1>,
  <-0.66,-0.4,1>,
  <-0.66,-0.42,1>,
  <-0.66,-0.44,1>,
  <-0.66,-0.46,1>,
  <-0.66,-0.48,1>,
  <-0.66,-0.5,1>,
  <-0.66,-0.52,1>,
  <-0.66,-0.54,1>,
  <-0.66,-0.56,1>,
  <-0.66,-0.58,1>,
  <-0.66,-0.6,1>,
  <-0.66,-0.62,1>,
  <-0.66,-0.64,1>,
  <-0.66,-0.66,1>,
  <-0.66,-0.68,1>,
  <-0.66,-0.7,1>,
  <-0.66,-0.72,1>,
  <-0.66,-0.74,1>,
  <-0.66,-0.76,1>,
  <-0.66,-0.78,1>,
  <-0.66,-0.8,1>,
  <-0.66,-0.82,1>,
  <-0.66,-0.84,1>,
  <-0.66,-0.86,1>,
  <-0.66,-0.88,1>,
  <-0.66,-0.9,1>,
  <-0.66,-0.92,1>,
  <-0.66,-0.94,1>,
  <-0.66,-0.96,1>,
  <-0.66,-0.98,1>,
  <-0.66,-1,1>,
  <-0.68,0,1>,
  <-0.68,-0.02,1>,
  <-0.68,-0.04,1>,
  <-0.68,-0.06,1>,
  <-0.68,-0.08,1>,
  <-0.68,-0.1,1>,
  <-0.68,-0.12,1>,
  <-0.68,-0.14,1>,
  <-0.68,-0.16,1>,
  <-0.68,-0.18,1>,
  <-0.68,-0.2,1>,
  <-0.68,-0.22,1>,
  <-0.68,-0.24,1>,
  <-0.68,-0.26,1>,
  <-0.68,-0.28,1>,
  <-0.68,-0.3,1>,
  <-0.68,-0.32,1>,
  <-0.68,-0.34,1>,
  <-0.68,-0.36,1>,
  <-0.68,-0.38,1>,
  <-0.68,-0.4,1>,
  <-0.68,-0.42,1>,
  <-0.68,-0.44,1>,
  <-0.68,-0.46,1>,
  <-0.68,-0.48,1>,
  <-0.68,-0.5,1>,
  <-0.68,-0.52,1>,
  <-0.68,-0.54,1>,
  <-0.68,-0.56,1>,
  <-0.68,-0.58,1>,
  <-0.68,-0.6,1>,
  <-0.68,-0.62,1>,
  <-0.68,-0.64,1>,
  <-0.68,-0.66,1>,
  <-0.68,-0.68,1>,
  <-0.68,-0.7,1>,
  <-0.68,-0.72,1>,
  <-0.68,-0.74,1>,
  <-0.68,-0.76,1>,
  <-0.68,-0.78,1>,
  <-0.68,-0.8,1>,
  <-0.68,-0.82,1>,
  <-0.68,-0.84,1>,
  <-0.68,-0.86,1>,
  <-0.68,-0.88,1>,
  <-0.68,-0.9,1>,
  <-0.68,-0.92,1>,
  <-0.68,-0.94,1>,
  <-0.68,-0.96,1>,
  <-0.68,-0.98,1>,
  <-0.68,-1,1>,
  <-0.7,0,1>,
  <-0.7,-0.02,1>,
  <-0.7,-0.04,1>,
  <-0.7,-0.06,1>,
  <-0.7,-0.08,1>,
  <-0.7,-0.1,1>,
  <-0.7,-0.12,1>,
  <-0.7,-0.14,1>,
  <-0.7,-0.16,1>,
  <-0.7,-0.18,1>,
  <-0.7,-0.2,1>,
  <-0.7,-0.22,1>,
  <-0.7,-0.24,1>,
  <-0.7,-0.26,1>,
  <-0.7,-0.28,1>,
  <-0.7,-0.3,1>,
  <-0.7,-0.32,1>,
  <-0.7,-0.34,1>,
  <-0.7,-0.36,1>,
  <-0.7,-0.38,1>,
  <-0.7,-0.4,1>,
  <-0.7,-0.42,1>,
  <-0.7,-0.44,1>,
  <-0.7,-0.46,1>,
  <-0.7,-0.48,1>,
  <-0.7,-0.5,1>,
  <-0.7,-0.52,1>,
  <-0.7,-0.54,1>,
  <-0.7,-0.56,1>,
  <-0.7,-0.58,1>,
  <-0.7,-0.6,1>,
  <-0.7,-0.62,1>,
  <-0.7,-0.64,1>,
  <-0.7,-0.66,1>,
  <-0.7,-0.68,1>,
  <-0.7,-0.7,1>,
  <-0.7,-0.72,1>,
  <-0.7,-0.74,1>,
  <-0.7,-0.76,1>,
  <-0.7,-0.78,1>,
  <-0.7,-0.8,1>,
  <-0.7,-0.82,1>,
  <-0.7,-0.84,1>,
  <-0.7,-0.86,1>,
  <-0.7,-0.88,1>,
  <-0.7,-0.9,1>,
  <-0.7,-0.92,1>,
  <-0.7,-0.94,1>,
  <-0.7,-0.96,1>,
  <-0.7,-0.98,1>,
  <-0.7,-1,1>,
  <-0.72,0,1>,
  <-0.72,-0.02,1>,
  <-0.72,-0.04,1>,
  <-0.72,-0.06,1>,
  <-0.72,-0.08,1>,
  <-0.72,-0.1,1>,
  <-0.72,-0.12,1>,
  <-0.72,-0.14,1>,
  <-0.72,-0.16,1>,
  <-0.72,-0.18,1>,
  <-0.72,-0.2,1>,
  <-0.72,-0.22,1>,
  <-0.72,-0.24,1>,
  <-0.72,-0.26,1>,
  <-0.72,-0.28,1>,
  <-0.72,-0.3,1>,
  <-0.72,-0.32,1>,
  <-0.72,-0.34,1>,
  <-0.72,-0.36,1>,
  <-0.72,-0.38,1>,
  <-0.72,-0.4,1>,
  <-0.72,-0.42,1>,
  <-0.72,-0.44,1>,
  <-0.72,-0.46,1>,
  <-0.72,-0.48,1>,
  <-0.72,-0.5,1>,
  <-0.72,-0.52,1>,
  <-0.72,-0.54,1>,
  <-0.72,-0.56,1>,
  <-0.72,-0.58,1>,
  <-0.72,-0.6,1>,
  <-0.72,-0.62,1>,
  <-0.72,-0.64,1>,
  <-0.72,-0.66,1>,
  <-0.72,-0.68,1>,
  <-0.72,-0.7,1>,
  <-0.72,-0.72,1>,
  <-0.72,-0.74,1>,
  <-0.72,-0.76,1>,
  <-0.72,-0.78,1>,
  <-0.72,-0.8,1>,
  <-0.72,-0.82,1>,
  <-0.72,-0.84,1>,
  <-0.72,-0.86,1>,
  <-0.72,-0.88,1>,
  <-0.72,-0.9,1>,
  <-0.72,-0.92,1>,
  <-0.72,-0.94,1>,
  <-0.72,-0.96,1>,
  <-0.72,-0.98,1>,
  <-0.72,-1,1>,
  <-0.74,0,1>,
  <-0.74,-0.02,1>,
  <-0.74,-0.04,1>,
  <-0.74,-0.06,1>,
  <-0.74,-0.08,1>,
  <-0.74,-0.1,1>,
  <-0.74,-0.12,1>,
  <-0.74,-0.14,1>,
  <-0.74,-0.16,1>,
  <-0.74,-0.18,1>,
  <-0.74,-0.2,1>,
  <-0.74,-0.22,1>,
  <-0.74,-0.24,1>,
  <-0.74,-0.26,1>,
  <-0.74,-0.28,1>,
  <-0.74,-0.3,1>,
  <-0.74,-0.32,1>,
  <-0.74,-0.34,1>,
  <-0.74,-0.36,1>,
  <-0.74,-0.38,1>,
  <-0.74,-0.4,1>,
  <-0.74,-0.42,1>,
  <-0.74,-0.44,1>,
  <-0.74,-0.46,1>,
  <-0.74,-0.48,1>,
  <-0.74,-0.5,1>,
  <-0.74,-0.52,1>,
  <-0.74,-0.54,1>,
  <-0.74,-0.56,1>,
  <-0.74,-0.58,1>,
  <-0.74,-0.6,1>,
  <-0.74,-0.62,1>,
  <-0.74,-0.64,1>,
  <-0.74,-0.66,1>,
  <-0.74,-0.68,1>,
  <-0.74,-0.7,1>,
  <-0.74,-0.72,1>,
  <-0.74,-0.74,1>,
  <-0.74,-0.76,1>,
  <-0.74,-0.78,1>,
  <-0.74,-0.8,1>,
  <-0.74,-0.82,1>,
  <-0.74,-0.84,1>,
  <-0.74,-0.86,1>,
  <-0.74,-0.88,1>,
  <-0.74,-0.9,1>,
  <-0.74,-0.92,1>,
  <-0.74,-0.94,1>,
  <-0.74,-0.96,1>,
  <-0.74,-0.98,1>,
  <-0.74,-1,1>,
  <-0.76,0,1>,
  <-0.76,-0.02,1>,
  <-0.76,-0.04,1>,
  <-0.76,-0.06,1>,
  <-0.76,-0.08,1>,
  <-0.76,-0.1,1>,
  <-0.76,-0.12,1>,
  <-0.76,-0.14,1>,
  <-0.76,-0.16,1>,
  <-0.76,-0.18,1>,
  <-0.76,-0.2,1>,
  <-0.76,-0.22,1>,
  <-0.76,-0.24,1>,
  <-0.76,-0.26,1>,
  <-0.76,-0.28,1>,
  <-0.76,-0.3,1>,
  <-0.76,-0.32,1>,
  <-0.76,-0.34,1>,
  <-0.76,-0.36,1>,
  <-0.76,-0.38,1>,
  <-0.76,-0.4,1>,
  <-0.76,-0.42,1>,
  <-0.76,-0.44,1>,
  <-0.76,-0.46,1>,
  <-0.76,-0.48,1>,
  <-0.76,-0.5,1>,
  <-0.76,-0.52,1>,
  <-0.76,-0.54,1>,
  <-0.76,-0.56,1>,
  <-0.76,-0.58,1>,
  <-0.76,-0.6,1>,
  <-0.76,-0.62,1>,
  <-0.76,-0.64,1>,
  <-0.76,-0.66,1>,
  <-0.76,-0.68,1>,
  <-0.76,-0.7,1>,
  <-0.76,-0.72,1>,
  <-0.76,-0.74,1>,
  <-0.76,-0.76,1>,
  <-0.76,-0.78,1>,
  <-0.76,-0.8,1>,
  <-0.76,-0.82,1>,
  <-0.76,-0.84,1>,
  <-0.76,-0.86,1>,
  <-0.76,-0.88,1>,
  <-0.76,-0.9,1>,
  <-0.76,-0.92,1>,
  <-0.76,-0.94,1>,
  <-0.76,-0.96,1>,
  <-0.76,-0.98,1>,
  <-0.76,-1,1>,
  <-0.78,0,1>,
  <-0.78,-0.02,1>,
  <-0.78,-0.04,1>,
  <-0.78,-0.06,1>,
  <-0.78,-0.08,1>,
  <-0.78,-0.1,1>,
  <-0.78,-0.12,1>,
  <-0.78,-0.14,1>,
  <-0.78,-0.16,1>,
  <-0.78,-0.18,1>,
  <-0.78,-0.2,1>,
  <-0.78,-0.22,1>,
  <-0.78,-0.24,1>,
  <-0.78,-0.26,1>,
  <-0.78,-0.28,1>,
  <-0.78,-0.3,1>,
  <-0.78,-0.32,1>,
  <-0.78,-0.34,1>,
  <-0.78,-0.36,1>,
  <-0.78,-0.38,1>,
  <-0.78,-0.4,1>,
  <-0.78,-0.42,1>,
  <-0.78,-0.44,1>,
  <-0.78,-0.46,1>,
  <-0.78,-0.48,1>,
  <-0.78,-0.5,1>,
  <-0.78,-0.52,1>,
  <-0.78,-0.54,1>,
  <-0.78,-0.56,1>,
  <-0.78,-0.58,1>,
  <-0.78,-0.6,1>,
  <-0.78,-0.62,1>,
  <-0.78,-0.64,1>,
  <-0.78,-0.66,1>,
  <-0.78,-0.68,1>,
  <-0.78,-0.7,1>,
  <-0.78,-0.72,1>,
  <-0.78,-0.74,1>,
  <-0.78,-0.76,1>,
  <-0.78,-0.78,1>,
  <-0.78,-0.8,1>,
  <-0.78,-0.82,1>,
  <-0.78,-0.84,1>,
  <-0.78,-0.86,1>,
  <-0.78,-0.88,1>,
  <-0.78,-0.9,1>,
  <-0.78,-0.92,1>,
  <-0.78,-0.94,1>,
  <-0.78,-0.96,1>,
  <-0.78,-0.98,1>,
  <-0.78,-1,1>,
  <-0.8,0,1>,
  <-0.8,-0.02,1>,
  <-0.8,-0.04,1>,
  <-0.8,-0.06,1>,
  <-0.8,-0.08,1>,
  <-0.8,-0.1,1>,
  <-0.8,-0.12,1>,
  <-0.8,-0.14,1>,
  <-0.8,-0.16,1>,
  <-0.8,-0.18,1>,
  <-0.8,-0.2,1>,
  <-0.8,-0.22,1>,
  <-0.8,-0.24,1>,
  <-0.8,-0.26,1>,
  <-0.8,-0.28,1>,
  <-0.8,-0.3,1>,
  <-0.8,-0.32,1>,
  <-0.8,-0.34,1>,
  <-0.8,-0.36,1>,
  <-0.8,-0.38,1>,
  <-0.8,-0.4,1>,
  <-0.8,-0.42,1>,
  <-0.8,-0.44,1>,
  <-0.8,-0.46,1>,
  <-0.8,-0.48,1>,
  <-0.8,-0.5,1>,
  <-0.8,-0.52,1>,
  <-0.8,-0.54,1>,
  <-0.8,-0.56,1>,
  <-0.8,-0.58,1>,
  <-0.8,-0.6,1>,
  <-0.8,-0.62,1>,
  <-0.8,-0.64,1>,
  <-0.8,-0.66,1>,
  <-0.8,-0.68,1>,
  <-0.8,-0.7,1>,
  <-0.8,-0.72,1>,
  <-0.8,-0.74,1>,
  <-0.8,-0.76,1>,
  <-0.8,-0.78,1>,
  <-0.8,-0.8,1>,
  <-0.8,-0.82,1>,
  <-0.8,-0.84,1>,
  <-0.8,-0.86,1>,
  <-0.8,-0.88,1>,
  <-0.8,-0.9,1>,
  <-0.8,-0.92,1>,
  <-0.8,-0.94,1>,
  <-0.8,-0.96,1>,
  <-0.8,-0.98,1>,
  <-0.8,-1,1>,
  <-0.82,0,1>,
  <-0.82,-0.02,1>,
  <-0.82,-0.04,1>,
  <-0.82,-0.06,1>,
  <-0.82,-0.08,1>,
  <-0.82,-0.1,1>,
  <-0.82,-0.12,1>,
  <-0.82,-0.14,1>,
  <-0.82,-0.16,1>,
  <-0.82,-0.18,1>,
  <-0.82,-0.2,1>,
  <-0.82,-0.22,1>,
  <-0.82,-0.24,1>,
  <-0.82,-0.26,1>,
  <-0.82,-0.28,1>,
  <-0.82,-0.3,1>,
  <-0.82,-0.32,1>,
  <-0.82,-0.34,1>,
  <-0.82,-0.36,1>,
  <-0.82,-0.38,1>,
  <-0.82,-0.4,1>,
  <-0.82,-0.42,1>,
  <-0.82,-0.44,1>,
  <-0.82,-0.46,1>,
  <-0.82,-0.48,1>,
  <-0.82,-0.5,1>,
  <-0.82,-0.52,1>,
  <-0.82,-0.54,1>,
  <-0.82,-0.56,1>,
  <-0.82,-0.58,1>,
  <-0.82,-0.6,1>,
  <-0.82,-0.62,1>,
  <-0.82,-0.64,1>,
  <-0.82,-0.66,1>,
  <-0.82,-0.68,1>,
  <-0.82,-0.7,1>,
  <-0.82,-0.72,1>,
  <-0.82,-0.74,1>,
  <-0.82,-0.76,1>,
  <-0.82,-0.78,1>,
  <-0.82,-0.8,1>,
  <-0.82,-0.82,1>,
  <-0.82,-0.84,1>,
  <-0.82,-0.86,1>,
  <-0.82,-0.88,1>,
  <-0.82,-0.9,1>,
  <-0.82,-0.92,1>,
  <-0.82,-0.94,1>,
  <-0.82,-0.96,1>,
  <-0.82,-0.98,1>,
  <-0.82,-1,1>,
  <-0.84,0,1>,
  <-0.84,-0.02,1>,
  <-0.84,-0.04,1>,
  <-0.84,-0.06,1>,
  <-0.84,-0.08,1>,
  <-0.84,-0.1,1>,
  <-0.84,-0.12,1>,
  <-0.84,-0.14,1>,
  <-0.84,-0.16,1>,
  <-0.84,-0.18,1>,
  <-0.84,-0.2,1>,
  <-0.84,-0.22,1>,
  <-0.84,-0.24,1>,
  <-0.84,-0.26,1>,
  <-0.84,-0.28,1>,
  <-0.84,-0.3,1>,
  <-0.84,-0.32,1>,
  <-0.84,-0.34,1>,
  <-0.84,-0.36,1>,
  <-0.84,-0.38,1>,
  <-0.84,-0.4,1>,
  <-0.84,-0.42,1>,
  <-0.84,-0.44,1>,
  <-0.84,-0.46,1>,
  <-0.84,-0.48,1>,
  <-0.84,-0.5,1>,
  <-0.84,-0.52,1>,
  <-0.84,-0.54,1>,
  <-0.84,-0.56,1>,
  <-0.84,-0.58,1>,
  <-0.84,-0.6,1>,
  <-0.84,-0.62,1>,
  <-0.84,-0.64,1>,
  <-0.84,-0.66,1>,
  <-0.84,-0.68,1>,
  <-0.84,-0.7,1>,
  <-0.84,-0.72,1>,
  <-0.84,-0.74,1>,
  <-0.84,-0.76,1>,
  <-0.84,-0.78,1>,
  <-0.84,-0.8,1>,
  <-0.84,-0.82,1>,
  <-0.84,-0.84,1>,
  <-0.84,-0.86,1>,
  <-0.84,-0.88,1>,
  <-0.84,-0.9,1>,
  <-0.84,-0.92,1>,
  <-0.84,-0.94,1>,
  <-0.84,-0.96,1>,
  <-0.84,-0.98,1>,
  <-0.84,-1,1>,
  <-0.86,0,1>,
  <-0.86,-0.02,1>,
  <-0.86,-0.04,1>,
  <-0.86,-0.06,1>,
  <-0.86,-0.08,1>,
  <-0.86,-0.1,1>,
  <-0.86,-0.12,1>,
  <-0.86,-0.14,1>,
  <-0.86,-0.16,1>,
  <-0.86,-0.18,1>,
  <-0.86,-0.2,1>,
  <-0.86,-0.22,1>,
  <-0.86,-0.24,1>,
  <-0.86,-0.26,1>,
  <-0.86,-0.28,1>,
  <-0.86,-0.3,1>,
  <-0.86,-0.32,1>,
  <-0.86,-0.34,1>,
  <-0.86,-0.36,1>,
  <-0.86,-0.38,1>,
  <-0.86,-0.4,1>,
  <-0.86,-0.42,1>,
  <-0.86,-0.44,1>,
  <-0.86,-0.46,1>,
  <-0.86,-0.48,1>,
  <-0.86,-0.5,1>,
  <-0.86,-0.52,1>,
  <-0.86,-0.54,1>,
  <-0.86,-0.56,1>,
  <-0.86,-0.58,1>,
  <-0.86,-0.6,1>,
  <-0.86,-0.62,1>,
  <-0.86,-0.64,1>,
  <-0.86,-0.66,1>,
  <-0.86,-0.68,1>,
  <-0.86,-0.7,1>,
  <-0.86,-0.72,1>,
  <-0.86,-0.74,1>,
  <-0.86,-0.76,1>,
  <-0.86,-0.78,1>,
  <-0.86,-0.8,1>,
  <-0.86,-0.82,1>,
  <-0.86,-0.84,1>,
  <-0.86,-0.86,1>,
  <-0.86,-0.88,1>,
  <-0.86,-0.9,1>,
  <-0.86,-0.92,1>,
  <-0.86,-0.94,1>,
  <-0.86,-0.96,1>,
  <-0.86,-0.98,1>,
  <-0.86,-1,1>,
  <-0.88,0,1>,
  <-0.88,-0.02,1>,
  <-0.88,-0.04,1>,
  <-0.88,-0.06,1>,
  <-0.88,-0.08,1>,
  <-0.88,-0.1,1>,
  <-0.88,-0.12,1>,
  <-0.88,-0.14,1>,
  <-0.88,-0.16,1>,
  <-0.88,-0.18,1>,
  <-0.88,-0.2,1>,
  <-0.88,-0.22,1>,
  <-0.88,-0.24,1>,
  <-0.88,-0.26,1>,
  <-0.88,-0.28,1>,
  <-0.88,-0.3,1>,
  <-0.88,-0.32,1>,
  <-0.88,-0.34,1>,
  <-0.88,-0.36,1>,
  <-0.88,-0.38,1>,
  <-0.88,-0.4,1>,
  <-0.88,-0.42,1>,
  <-0.88,-0.44,1>,
  <-0.88,-0.46,1>,
  <-0.88,-0.48,1>,
  <-0.88,-0.5,1>,
  <-0.88,-0.52,1>,
  <-0.88,-0.54,1>,
  <-0.88,-0.56,1>,
  <-0.88,-0.58,1>,
  <-0.88,-0.6,1>,
  <-0.88,-0.62,1>,
  <-0.88,-0.64,1>,
  <-0.88,-0.66,1>,
  <-0.88,-0.68,1>,
  <-0.88,-0.7,1>,
  <-0.88,-0.72,1>,
  <-0.88,-0.74,1>,
  <-0.88,-0.76,1>,
  <-0.88,-0.78,1>,
  <-0.88,-0.8,1>,
  <-0.88,-0.82,1>,
  <-0.88,-0.84,1>,
  <-0.88,-0.86,1>,
  <-0.88,-0.88,1>,
  <-0.88,-0.9,1>,
  <-0.88,-0.92,1>,
  <-0.88,-0.94,1>,
  <-0.88,-0.96,1>,
  <-0.88,-0.98,1>,
  <-0.88,-1,1>,
  <-0.9,0,1>,
  <-0.9,-0.02,1>,
  <-0.9,-0.04,1>,
  <-0.9,-0.06,1>,
  <-0.9,-0.08,1>,
  <-0.9,-0.1,1>,
  <-0.9,-0.12,1>,
  <-0.9,-0.14,1>,
  <-0.9,-0.16,1>,
  <-0.9,-0.18,1>,
  <-0.9,-0.2,1>,
  <-0.9,-0.22,1>,
  <-0.9,-0.24,1>,
  <-0.9,-0.26,1>,
  <-0.9,-0.28,1>,
  <-0.9,-0.3,1>,
  <-0.9,-0.32,1>,
  <-0.9,-0.34,1>,
  <-0.9,-0.36,1>,
  <-0.9,-0.38,1>,
  <-0.9,-0.4,1>,
  <-0.9,-0.42,1>,
  <-0.9,-0.44,1>,
  <-0.9,-0.46,1>,
  <-0.9,-0.48,1>,
  <-0.9,-0.5,1>,
  <-0.9,-0.52,1>,
  <-0.9,-0.54,1>,
  <-0.9,-0.56,1>,
  <-0.9,-0.58,1>,
  <-0.9,-0.6,1>,
  <-0.9,-0.62,1>,
  <-0.9,-0.64,1>,
  <-0.9,-0.66,1>,
  <-0.9,-0.68,1>,
  <-0.9,-0.7,1>,
  <-0.9,-0.72,1>,
  <-0.9,-0.74,1>,
  <-0.9,-0.76,1>,
  <-0.9,-0.78,1>,
  <-0.9,-0.8,1>,
  <-0.9,-0.82,1>,
  <-0.9,-0.84,1>,
  <-0.9,-0.86,1>,
  <-0.9,-0.88,1>,
  <-0.9,-0.9,1>,
  <-0.9,-0.92,1>,
  <-0.9,-0.94,1>,
  <-0.9,-0.96,1>,
  <-0.9,-0.98,1>,
  <-0.9,-1,1>,
  <-0.92,0,1>,
  <-0.92,-0.02,1>,
  <-0.92,-0.04,1>,
  <-0.92,-0.06,1>,
  <-0.92,-0.08,1>,
  <-0.92,-0.1,1>,
  <-0.92,-0.12,1>,
  <-0.92,-0.14,1>,
  <-0.92,-0.16,1>,
  <-0.92,-0.18,1>,
  <-0.92,-0.2,1>,
  <-0.92,-0.22,1>,
  <-0.92,-0.24,1>,
  <-0.92,-0.26,1>,
  <-0.92,-0.28,1>,
  <-0.92,-0.3,1>,
  <-0.92,-0.32,1>,
  <-0.92,-0.34,1>,
  <-0.92,-0.36,1>,
  <-0.92,-0.38,1>,
  <-0.92,-0.4,1>,
  <-0.92,-0.42,1>,
  <-0.92,-0.44,1>,
  <-0.92,-0.46,1>,
  <-0.92,-0.48,1>,
  <-0.92,-0.5,1>,
  <-0.92,-0.52,1>,
  <-0.92,-0.54,1>,
  <-0.92,-0.56,1>,
  <-0.92,-0.58,1>,
  <-0.92,-0.6,1>,
  <-0.92,-0.62,1>,
  <-0.92,-0.64,1>,
  <-0.92,-0.66,1>,
  <-0.92,-0.68,1>,
  <-0.92,-0.7,1>,
  <-0.92,-0.72,1>,
  <-0.92,-0.74,1>,
  <-0.92,-0.76,1>,
  <-0.92,-0.78,1>,
  <-0.92,-0.8,1>,
  <-0.92,-0.82,1>,
  <-0.92,-0.84,1>,
  <-0.92,-0.86,1>,
  <-0.92,-0.88,1>,
  <-0.92,-0.9,1>,
  <-0.92,-0.92,1>,
  <-0.92,-0.94,1>,
  <-0.92,-0.96,1>,
  <-0.92,-0.98,1>,
  <-0.92,-1,1>,
  <-0.94,0,1>,
  <-0.94,-0.02,1>,
  <-0.94,-0.04,1>,
  <-0.94,-0.06,1>,
  <-0.94,-0.08,1>,
  <-0.94,-0.1,1>,
  <-0.94,-0.12,1>,
  <-0.94,-0.14,1>,
  <-0.94,-0.16,1>,
  <-0.94,-0.18,1>,
  <-0.94,-0.2,1>,
  <-0.94,-0.22,1>,
  <-0.94,-0.24,1>,
  <-0.94,-0.26,1>,
  <-0.94,-0.28,1>,
  <-0.94,-0.3,1>,
  <-0.94,-0.32,1>,
  <-0.94,-0.34,1>,
  <-0.94,-0.36,1>,
  <-0.94,-0.38,1>,
  <-0.94,-0.4,1>,
  <-0.94,-0.42,1>,
  <-0.94,-0.44,1>,
  <-0.94,-0.46,1>,
  <-0.94,-0.48,1>,
  <-0.94,-0.5,1>,
  <-0.94,-0.52,1>,
  <-0.94,-0.54,1>,
  <-0.94,-0.56,1>,
  <-0.94,-0.58,1>,
  <-0.94,-0.6,1>,
  <-0.94,-0.62,1>,
  <-0.94,-0.64,1>,
  <-0.94,-0.66,1>,
  <-0.94,-0.68,1>,
  <-0.94,-0.7,1>,
  <-0.94,-0.72,1>,
  <-0.94,-0.74,1>,
  <-0.94,-0.76,1>,
  <-0.94,-0.78,1>,
  <-0.94,-0.8,1>,
  <-0.94,-0.82,1>,
  <-0.94,-0.84,1>,
  <-0.94,-0.86,1>,
  <-0.94,-0.88,1>,
  <-0.94,-0.9,1>,
  <-0.94,-0.92,1>,
  <-0.94,-0.94,1>,
  <-0.94,-0.96,1>,
  <-0.94,-0.98,1>,
  <-0.94,-1,1>,
  <-0.96,0,1>,
  <-0.96,-0.02,1>,
  <-0.96,-0.04,1>,
  <-0.96,-0.06,1>,
  <-0.96,-0.08,1>,
  <-0.96,-0.1,1>,
  <-0.96,-0.12,1>,
  <-0.96,-0.14,1>,
  <-0.96,-0.16,1>,
  <-0.96,-0.18,1>,
  <-0.96,-0.2,1>,
  <-0.96,-0.22,1>,
  <-0.96,-0.24,1>,
  <-0.96,-0.26,1>,
  <-0.96,-0.28,1>,
  <-0.96,-0.3,1>,
  <-0.96,-0.32,1>,
  <-0.96,-0.34,1>,
  <-0.96,-0.36,1>,
  <-0.96,-0.38,1>,
  <-0.96,-0.4,1>,
  <-0.96,-0.42,1>,
  <-0.96,-0.44,1>,
  <-0.96,-0.46,1>,
  <-0.96,-0.48,1>,
  <-0.96,-0.5,1>,
  <-0.96,-0.52,1>,
  <-0.96,-0.54,1>,
  <-0.96,-0.56,1>,
  <-0.96,-0.58,1>,
  <-0.96,-0.6,1>,
  <-0.96,-0.62,1>,
  <-0.96,-0.64,1>,
  <-0.96,-0.66,1>,
  <-0.96,-0.68,1>,
  <-0.96,-0.7,1>,
  <-0.96,-0.72,1>,
  <-0.96,-0.74,1>,
  <-0.96,-0.76,1>,
  <-0.96,-0.78,1>,
  <-0.96,-0.8,1>,
  <-0.96,-0.82,1>,
  <-0.96,-0.84,1>,
  <-0.96,-0.86,1>,
  <-0.96,-0.88,1>,
  <-0.96,-0.9,1>,
  <-0.96,-0.92,1>,
  <-0.96,-0.94,1>,
  <-0.96,-0.96,1>,
  <-0.96,-0.98,1>,
  <-0.96,-1,1>,
  <-0.98,0,1>,
  <-0.98,-0.02,1>,
  <-0.98,-0.04,1>,
  <-0.98,-0.06,1>,
  <-0.98,-0.08,1>,
  <-0.98,-0.1,1>,
  <-0.98,-0.12,1>,
  <-0.98,-0.14,1>,
  <-0.98,-0.16,1>,
  <-0.98,-0.18,1>,
  <-0.98,-0.2,1>,
  <-0.98,-0.22,1>,
  <-0.98,-0.24,1>,
  <-0.98,-0.26,1>,
  <-0.98,-0.28,1>,
  <-0.98,-0.3,1>,
  <-0.98,-0.32,1>,
  <-0.98,-0.34,1>,
  <-0.98,-0.36,1>,
  <-0.98,-0.38,1>,
  <-0.98,-0.4,1>,
  <-0.98,-0.42,1>,
  <-0.98,-0.44,1>,
  <-0.98,-0.46,1>,
  <-0.98,-0.48,1>,
  <-0.98,-0.5,1>,
  <-0.98,-0.52,1>,
  <-0.98,-0.54,1>,
  <-0.98,-0.56,1>,
  <-0.98,-0.58,1>,
  <-0.98,-0.6,1>,
  <-0.98,-0.62,1>,
  <-0.98,-0.64,1>,
  <-0.98,-0.66,1>,
  <-0.98,-0.68,1>,
  <-0.98,-0.7,1>,
  <-0.98,-0.72,1>,
  <-0.98,-0.74,1>,
  <-0.98,-0.76,1>,
  <-0.98,-0.78,1>,
  <-0.98,-0.8,1>,
  <-0.98,-0.82,1>,
  <-0.98,-0.84,1>,
  <-0.98,-0.86,1>,
  <-0.98,-0.88,1>,
  <-0.98,-0.9,1>,
  <-0.98,-0.92,1>,
  <-0.98,-0.94,1>,
  <-0.98,-0.96,1>,
  <-0.98,-0.98,1>,
  <-0.98,-1,1>,
  <-1,0,1>,
  <-1,-0.02,1>,
  <-1,-0.04,1>,
  <-1,-0.06,1>,
  <-1,-0.08,1>,
  <-1,-0.1,1>,
  <-1,-0.12,1>,
  <-1,-0.14,1>,
  <-1,-0.16,1>,
  <-1,-0.18,1>,
  <-1,-0.2,1>,
  <-1,-0.22,1>,
  <-1,-0.24,1>,
  <-1,-0.26,1>,
  <-1,-0.28,1>,
  <-1,-0.3,1>,
  <-1,-0.32,1>,
  <-1,-0.34,1>,
  <-1,-0.36,1>,
  <-1,-0.38,1>,
  <-1,-0.4,1>,
  <-1,-0.42,1>,
  <-1,-0.44,1>,
  <-1,-0.46,1>,
  <-1,-0.48,1>,
  <-1,-0.5,1>,
  <-1,-0.52,1>,
  <-1,-0.54,1>,
  <-1,-0.56,1>,
  <-1,-0.58,1>,
  <-1,-0.6,1>,
  <-1,-0.62,1>,
  <-1,-0.64,1>,
  <-1,-0.66,1>,
  <-1,-0.68,1>,
  <-1,-0.7,1>,
  <-1,-0.72,1>,
  <-1,-0.74,1>,
  <-1,-0.76,1>,
  <-1,-0.78,1>,
  <-1,-0.8,1>,
  <-1,-0.82,1>,
  <-1,-0.84,1>,
  <-1,-0.86,1>,
  <-1,-0.88,1>,
  <-1,-0.9,1>,
  <-1,-0.92,1>,
  <-1,-0.94,1>,
  <-1,-0.96,1>,
  <-1,-0.98,1>,
  <-1,-1,1>,
  \}
  face_indices \{ 5000,
  <0,1,51>,
  <52,51,1>,
  <1,2,52>,
  <53,52,2>,
  <2,3,53>,
  <54,53,3>,
  <3,4,54>,
  <55,54,4>,
  <4,5,55>,
  <56,55,5>,
  <5,6,56>,
  <57,56,6>,
  <6,7,57>,
  <58,57,7>,
  <7,8,58>,
  <59,58,8>,
  <8,9,59>,
  <60,59,9>,
  <9,10,60>,
  <61,60,10>,
  <10,11,61>,
  <62,61,11>,
  <11,12,62>,
  <63,62,12>,
  <12,13,63>,
  <64,63,13>,
  <13,14,64>,
  <65,64,14>,
  <14,15,65>,
  <66,65,15>,
  <15,16,66>,
  <67,66,16>,
  <16,17,67>,
  <68,67,17>,
  <17,18,68>,
  <69,68,18>,
  <18,19,69>,
  <70,69,19>,
  <19,20,70>,
  <71,70,20>,
  <20,21,71>,
  <72,71,21>,
  <21,22,72>,
  <73,72,22>,
  <22,23,73>,
  <74,73,23>,
  <23,24,74>,
  <75,74,24>,
  <24,25,75>,
  <76,75,25>,
  <25,26,76>,
  <77,76,26>,
  <26,27,77>,
  <78,77,27>,
  <27,28,78>,
  <79,78,28>,
  <28,29,79>,
  <80,79,29>,
  <29,30,80>,
  <81,80,30>,
  <30,31,81>,
  <82,81,31>,
  <31,32,82>,
  <83,82,32>,
  <32,33,83>,
  <84,83,33>,
  <33,34,84>,
  <85,84,34>,
  <34,35,85>,
  <86,85,35>,
  <35,36,86>,
  <87,86,36>,
  <36,37,87>,
  <88,87,37>,
  <37,38,88>,
  <89,88,38>,
  <38,39,89>,
  <90,89,39>,
  <39,40,90>,
  <91,90,40>,
  <40,41,91>,
  <92,91,41>,
  <41,42,92>,
  <93,92,42>,
  <42,43,93>,
  <94,93,43>,
  <43,44,94>,
  <95,94,44>,
  <44,45,95>,
  <96,95,45>,
  <45,46,96>,
  <97,96,46>,
  <46,47,97>,
  <98,97,47>,
  <47,48,98>,
  <99,98,48>,
  <48,49,99>,
  <100,99,49>,
  <49,50,100>,
  <101,100,50>,
  <51,52,102>,
  <103,102,52>,
  <52,53,103>,
  <104,103,53>,
  <53,54,104>,
  <105,104,54>,
  <54,55,105>,
  <106,105,55>,
  <55,56,106>,
  <107,106,56>,
  <56,57,107>,
  <108,107,57>,
  <57,58,108>,
  <109,108,58>,
  <58,59,109>,
  <110,109,59>,
  <59,60,110>,
  <111,110,60>,
  <60,61,111>,
  <112,111,61>,
  <61,62,112>,
  <113,112,62>,
  <62,63,113>,
  <114,113,63>,
  <63,64,114>,
  <115,114,64>,
  <64,65,115>,
  <116,115,65>,
  <65,66,116>,
  <117,116,66>,
  <66,67,117>,
  <118,117,67>,
  <67,68,118>,
  <119,118,68>,
  <68,69,119>,
  <120,119,69>,
  <69,70,120>,
  <121,120,70>,
  <70,71,121>,
  <122,121,71>,
  <71,72,122>,
  <123,122,72>,
  <72,73,123>,
  <124,123,73>,
  <73,74,124>,
  <125,124,74>,
  <74,75,125>,
  <126,125,75>,
  <75,76,126>,
  <127,126,76>,
  <76,77,127>,
  <128,127,77>,
  <77,78,128>,
  <129,128,78>,
  <78,79,129>,
  <130,129,79>,
  <79,80,130>,
  <131,130,80>,
  <80,81,131>,
  <132,131,81>,
  <81,82,132>,
  <133,132,82>,
  <82,83,133>,
  <134,133,83>,
  <83,84,134>,
  <135,134,84>,
  <84,85,135>,
  <136,135,85>,
  <85,86,136>,
  <137,136,86>,
  <86,87,137>,
  <138,137,87>,
  <87,88,138>,
  <139,138,88>,
  <88,89,139>,
  <140,139,89>,
  <89,90,140>,
  <141,140,90>,
  <90,91,141>,
  <142,141,91>,
  <91,92,142>,
  <143,142,92>,
  <92,93,143>,
  <144,143,93>,
  <93,94,144>,
  <145,144,94>,
  <94,95,145>,
  <146,145,95>,
  <95,96,146>,
  <147,146,96>,
  <96,97,147>,
  <148,147,97>,
  <97,98,148>,
  <149,148,98>,
  <98,99,149>,
  <150,149,99>,
  <99,100,150>,
  <151,150,100>,
  <100,101,151>,
  <152,151,101>,
  <102,103,153>,
  <154,153,103>,
  <103,104,154>,
  <155,154,104>,
  <104,105,155>,
  <156,155,105>,
  <105,106,156>,
  <157,156,106>,
  <106,107,157>,
  <158,157,107>,
  <107,108,158>,
  <159,158,108>,
  <108,109,159>,
  <160,159,109>,
  <109,110,160>,
  <161,160,110>,
  <110,111,161>,
  <162,161,111>,
  <111,112,162>,
  <163,162,112>,
  <112,113,163>,
  <164,163,113>,
  <113,114,164>,
  <165,164,114>,
  <114,115,165>,
  <166,165,115>,
  <115,116,166>,
  <167,166,116>,
  <116,117,167>,
  <168,167,117>,
  <117,118,168>,
  <169,168,118>,
  <118,119,169>,
  <170,169,119>,
  <119,120,170>,
  <171,170,120>,
  <120,121,171>,
  <172,171,121>,
  <121,122,172>,
  <173,172,122>,
  <122,123,173>,
  <174,173,123>,
  <123,124,174>,
  <175,174,124>,
  <124,125,175>,
  <176,175,125>,
  <125,126,176>,
  <177,176,126>,
  <126,127,177>,
  <178,177,127>,
  <127,128,178>,
  <179,178,128>,
  <128,129,179>,
  <180,179,129>,
  <129,130,180>,
  <181,180,130>,
  <130,131,181>,
  <182,181,131>,
  <131,132,182>,
  <183,182,132>,
  <132,133,183>,
  <184,183,133>,
  <133,134,184>,
  <185,184,134>,
  <134,135,185>,
  <186,185,135>,
  <135,136,186>,
  <187,186,136>,
  <136,137,187>,
  <188,187,137>,
  <137,138,188>,
  <189,188,138>,
  <138,139,189>,
  <190,189,139>,
  <139,140,190>,
  <191,190,140>,
  <140,141,191>,
  <192,191,141>,
  <141,142,192>,
  <193,192,142>,
  <142,143,193>,
  <194,193,143>,
  <143,144,194>,
  <195,194,144>,
  <144,145,195>,
  <196,195,145>,
  <145,146,196>,
  <197,196,146>,
  <146,147,197>,
  <198,197,147>,
  <147,148,198>,
  <199,198,148>,
  <148,149,199>,
  <200,199,149>,
  <149,150,200>,
  <201,200,150>,
  <150,151,201>,
  <202,201,151>,
  <151,152,202>,
  <203,202,152>,
  <153,154,204>,
  <205,204,154>,
  <154,155,205>,
  <206,205,155>,
  <155,156,206>,
  <207,206,156>,
  <156,157,207>,
  <208,207,157>,
  <157,158,208>,
  <209,208,158>,
  <158,159,209>,
  <210,209,159>,
  <159,160,210>,
  <211,210,160>,
  <160,161,211>,
  <212,211,161>,
  <161,162,212>,
  <213,212,162>,
  <162,163,213>,
  <214,213,163>,
  <163,164,214>,
  <215,214,164>,
  <164,165,215>,
  <216,215,165>,
  <165,166,216>,
  <217,216,166>,
  <166,167,217>,
  <218,217,167>,
  <167,168,218>,
  <219,218,168>,
  <168,169,219>,
  <220,219,169>,
  <169,170,220>,
  <221,220,170>,
  <170,171,221>,
  <222,221,171>,
  <171,172,222>,
  <223,222,172>,
  <172,173,223>,
  <224,223,173>,
  <173,174,224>,
  <225,224,174>,
  <174,175,225>,
  <226,225,175>,
  <175,176,226>,
  <227,226,176>,
  <176,177,227>,
  <228,227,177>,
  <177,178,228>,
  <229,228,178>,
  <178,179,229>,
  <230,229,179>,
  <179,180,230>,
  <231,230,180>,
  <180,181,231>,
  <232,231,181>,
  <181,182,232>,
  <233,232,182>,
  <182,183,233>,
  <234,233,183>,
  <183,184,234>,
  <235,234,184>,
  <184,185,235>,
  <236,235,185>,
  <185,186,236>,
  <237,236,186>,
  <186,187,237>,
  <238,237,187>,
  <187,188,238>,
  <239,238,188>,
  <188,189,239>,
  <240,239,189>,
  <189,190,240>,
  <241,240,190>,
  <190,191,241>,
  <242,241,191>,
  <191,192,242>,
  <243,242,192>,
  <192,193,243>,
  <244,243,193>,
  <193,194,244>,
  <245,244,194>,
  <194,195,245>,
  <246,245,195>,
  <195,196,246>,
  <247,246,196>,
  <196,197,247>,
  <248,247,197>,
  <197,198,248>,
  <249,248,198>,
  <198,199,249>,
  <250,249,199>,
  <199,200,250>,
  <251,250,200>,
  <200,201,251>,
  <252,251,201>,
  <201,202,252>,
  <253,252,202>,
  <202,203,253>,
  <254,253,203>,
  <204,205,255>,
  <256,255,205>,
  <205,206,256>,
  <257,256,206>,
  <206,207,257>,
  <258,257,207>,
  <207,208,258>,
  <259,258,208>,
  <208,209,259>,
  <260,259,209>,
  <209,210,260>,
  <261,260,210>,
  <210,211,261>,
  <262,261,211>,
  <211,212,262>,
  <263,262,212>,
  <212,213,263>,
  <264,263,213>,
  <213,214,264>,
  <265,264,214>,
  <214,215,265>,
  <266,265,215>,
  <215,216,266>,
  <267,266,216>,
  <216,217,267>,
  <268,267,217>,
  <217,218,268>,
  <269,268,218>,
  <218,219,269>,
  <270,269,219>,
  <219,220,270>,
  <271,270,220>,
  <220,221,271>,
  <272,271,221>,
  <221,222,272>,
  <273,272,222>,
  <222,223,273>,
  <274,273,223>,
  <223,224,274>,
  <275,274,224>,
  <224,225,275>,
  <276,275,225>,
  <225,226,276>,
  <277,276,226>,
  <226,227,277>,
  <278,277,227>,
  <227,228,278>,
  <279,278,228>,
  <228,229,279>,
  <280,279,229>,
  <229,230,280>,
  <281,280,230>,
  <230,231,281>,
  <282,281,231>,
  <231,232,282>,
  <283,282,232>,
  <232,233,283>,
  <284,283,233>,
  <233,234,284>,
  <285,284,234>,
  <234,235,285>,
  <286,285,235>,
  <235,236,286>,
  <287,286,236>,
  <236,237,287>,
  <288,287,237>,
  <237,238,288>,
  <289,288,238>,
  <238,239,289>,
  <290,289,239>,
  <239,240,290>,
  <291,290,240>,
  <240,241,291>,
  <292,291,241>,
  <241,242,292>,
  <293,292,242>,
  <242,243,293>,
  <294,293,243>,
  <243,244,294>,
  <295,294,244>,
  <244,245,295>,
  <296,295,245>,
  <245,246,296>,
  <297,296,246>,
  <246,247,297>,
  <298,297,247>,
  <247,248,298>,
  <299,298,248>,
  <248,249,299>,
  <300,299,249>,
  <249,250,300>,
  <301,300,250>,
  <250,251,301>,
  <302,301,251>,
  <251,252,302>,
  <303,302,252>,
  <252,253,303>,
  <304,303,253>,
  <253,254,304>,
  <305,304,254>,
  <255,256,306>,
  <307,306,256>,
  <256,257,307>,
  <308,307,257>,
  <257,258,308>,
  <309,308,258>,
  <258,259,309>,
  <310,309,259>,
  <259,260,310>,
  <311,310,260>,
  <260,261,311>,
  <312,311,261>,
  <261,262,312>,
  <313,312,262>,
  <262,263,313>,
  <314,313,263>,
  <263,264,314>,
  <315,314,264>,
  <264,265,315>,
  <316,315,265>,
  <265,266,316>,
  <317,316,266>,
  <266,267,317>,
  <318,317,267>,
  <267,268,318>,
  <319,318,268>,
  <268,269,319>,
  <320,319,269>,
  <269,270,320>,
  <321,320,270>,
  <270,271,321>,
  <322,321,271>,
  <271,272,322>,
  <323,322,272>,
  <272,273,323>,
  <324,323,273>,
  <273,274,324>,
  <325,324,274>,
  <274,275,325>,
  <326,325,275>,
  <275,276,326>,
  <327,326,276>,
  <276,277,327>,
  <328,327,277>,
  <277,278,328>,
  <329,328,278>,
  <278,279,329>,
  <330,329,279>,
  <279,280,330>,
  <331,330,280>,
  <280,281,331>,
  <332,331,281>,
  <281,282,332>,
  <333,332,282>,
  <282,283,333>,
  <334,333,283>,
  <283,284,334>,
  <335,334,284>,
  <284,285,335>,
  <336,335,285>,
  <285,286,336>,
  <337,336,286>,
  <286,287,337>,
  <338,337,287>,
  <287,288,338>,
  <339,338,288>,
  <288,289,339>,
  <340,339,289>,
  <289,290,340>,
  <341,340,290>,
  <290,291,341>,
  <342,341,291>,
  <291,292,342>,
  <343,342,292>,
  <292,293,343>,
  <344,343,293>,
  <293,294,344>,
  <345,344,294>,
  <294,295,345>,
  <346,345,295>,
  <295,296,346>,
  <347,346,296>,
  <296,297,347>,
  <348,347,297>,
  <297,298,348>,
  <349,348,298>,
  <298,299,349>,
  <350,349,299>,
  <299,300,350>,
  <351,350,300>,
  <300,301,351>,
  <352,351,301>,
  <301,302,352>,
  <353,352,302>,
  <302,303,353>,
  <354,353,303>,
  <303,304,354>,
  <355,354,304>,
  <304,305,355>,
  <356,355,305>,
  <306,307,357>,
  <358,357,307>,
  <307,308,358>,
  <359,358,308>,
  <308,309,359>,
  <360,359,309>,
  <309,310,360>,
  <361,360,310>,
  <310,311,361>,
  <362,361,311>,
  <311,312,362>,
  <363,362,312>,
  <312,313,363>,
  <364,363,313>,
  <313,314,364>,
  <365,364,314>,
  <314,315,365>,
  <366,365,315>,
  <315,316,366>,
  <367,366,316>,
  <316,317,367>,
  <368,367,317>,
  <317,318,368>,
  <369,368,318>,
  <318,319,369>,
  <370,369,319>,
  <319,320,370>,
  <371,370,320>,
  <320,321,371>,
  <372,371,321>,
  <321,322,372>,
  <373,372,322>,
  <322,323,373>,
  <374,373,323>,
  <323,324,374>,
  <375,374,324>,
  <324,325,375>,
  <376,375,325>,
  <325,326,376>,
  <377,376,326>,
  <326,327,377>,
  <378,377,327>,
  <327,328,378>,
  <379,378,328>,
  <328,329,379>,
  <380,379,329>,
  <329,330,380>,
  <381,380,330>,
  <330,331,381>,
  <382,381,331>,
  <331,332,382>,
  <383,382,332>,
  <332,333,383>,
  <384,383,333>,
  <333,334,384>,
  <385,384,334>,
  <334,335,385>,
  <386,385,335>,
  <335,336,386>,
  <387,386,336>,
  <336,337,387>,
  <388,387,337>,
  <337,338,388>,
  <389,388,338>,
  <338,339,389>,
  <390,389,339>,
  <339,340,390>,
  <391,390,340>,
  <340,341,391>,
  <392,391,341>,
  <341,342,392>,
  <393,392,342>,
  <342,343,393>,
  <394,393,343>,
  <343,344,394>,
  <395,394,344>,
  <344,345,395>,
  <396,395,345>,
  <345,346,396>,
  <397,396,346>,
  <346,347,397>,
  <398,397,347>,
  <347,348,398>,
  <399,398,348>,
  <348,349,399>,
  <400,399,349>,
  <349,350,400>,
  <401,400,350>,
  <350,351,401>,
  <402,401,351>,
  <351,352,402>,
  <403,402,352>,
  <352,353,403>,
  <404,403,353>,
  <353,354,404>,
  <405,404,354>,
  <354,355,405>,
  <406,405,355>,
  <355,356,406>,
  <407,406,356>,
  <357,358,408>,
  <409,408,358>,
  <358,359,409>,
  <410,409,359>,
  <359,360,410>,
  <411,410,360>,
  <360,361,411>,
  <412,411,361>,
  <361,362,412>,
  <413,412,362>,
  <362,363,413>,
  <414,413,363>,
  <363,364,414>,
  <415,414,364>,
  <364,365,415>,
  <416,415,365>,
  <365,366,416>,
  <417,416,366>,
  <366,367,417>,
  <418,417,367>,
  <367,368,418>,
  <419,418,368>,
  <368,369,419>,
  <420,419,369>,
  <369,370,420>,
  <421,420,370>,
  <370,371,421>,
  <422,421,371>,
  <371,372,422>,
  <423,422,372>,
  <372,373,423>,
  <424,423,373>,
  <373,374,424>,
  <425,424,374>,
  <374,375,425>,
  <426,425,375>,
  <375,376,426>,
  <427,426,376>,
  <376,377,427>,
  <428,427,377>,
  <377,378,428>,
  <429,428,378>,
  <378,379,429>,
  <430,429,379>,
  <379,380,430>,
  <431,430,380>,
  <380,381,431>,
  <432,431,381>,
  <381,382,432>,
  <433,432,382>,
  <382,383,433>,
  <434,433,383>,
  <383,384,434>,
  <435,434,384>,
  <384,385,435>,
  <436,435,385>,
  <385,386,436>,
  <437,436,386>,
  <386,387,437>,
  <438,437,387>,
  <387,388,438>,
  <439,438,388>,
  <388,389,439>,
  <440,439,389>,
  <389,390,440>,
  <441,440,390>,
  <390,391,441>,
  <442,441,391>,
  <391,392,442>,
  <443,442,392>,
  <392,393,443>,
  <444,443,393>,
  <393,394,444>,
  <445,444,394>,
  <394,395,445>,
  <446,445,395>,
  <395,396,446>,
  <447,446,396>,
  <396,397,447>,
  <448,447,397>,
  <397,398,448>,
  <449,448,398>,
  <398,399,449>,
  <450,449,399>,
  <399,400,450>,
  <451,450,400>,
  <400,401,451>,
  <452,451,401>,
  <401,402,452>,
  <453,452,402>,
  <402,403,453>,
  <454,453,403>,
  <403,404,454>,
  <455,454,404>,
  <404,405,455>,
  <456,455,405>,
  <405,406,456>,
  <457,456,406>,
  <406,407,457>,
  <458,457,407>,
  <408,409,459>,
  <460,459,409>,
  <409,410,460>,
  <461,460,410>,
  <410,411,461>,
  <462,461,411>,
  <411,412,462>,
  <463,462,412>,
  <412,413,463>,
  <464,463,413>,
  <413,414,464>,
  <465,464,414>,
  <414,415,465>,
  <466,465,415>,
  <415,416,466>,
  <467,466,416>,
  <416,417,467>,
  <468,467,417>,
  <417,418,468>,
  <469,468,418>,
  <418,419,469>,
  <470,469,419>,
  <419,420,470>,
  <471,470,420>,
  <420,421,471>,
  <472,471,421>,
  <421,422,472>,
  <473,472,422>,
  <422,423,473>,
  <474,473,423>,
  <423,424,474>,
  <475,474,424>,
  <424,425,475>,
  <476,475,425>,
  <425,426,476>,
  <477,476,426>,
  <426,427,477>,
  <478,477,427>,
  <427,428,478>,
  <479,478,428>,
  <428,429,479>,
  <480,479,429>,
  <429,430,480>,
  <481,480,430>,
  <430,431,481>,
  <482,481,431>,
  <431,432,482>,
  <483,482,432>,
  <432,433,483>,
  <484,483,433>,
  <433,434,484>,
  <485,484,434>,
  <434,435,485>,
  <486,485,435>,
  <435,436,486>,
  <487,486,436>,
  <436,437,487>,
  <488,487,437>,
  <437,438,488>,
  <489,488,438>,
  <438,439,489>,
  <490,489,439>,
  <439,440,490>,
  <491,490,440>,
  <440,441,491>,
  <492,491,441>,
  <441,442,492>,
  <493,492,442>,
  <442,443,493>,
  <494,493,443>,
  <443,444,494>,
  <495,494,444>,
  <444,445,495>,
  <496,495,445>,
  <445,446,496>,
  <497,496,446>,
  <446,447,497>,
  <498,497,447>,
  <447,448,498>,
  <499,498,448>,
  <448,449,499>,
  <500,499,449>,
  <449,450,500>,
  <501,500,450>,
  <450,451,501>,
  <502,501,451>,
  <451,452,502>,
  <503,502,452>,
  <452,453,503>,
  <504,503,453>,
  <453,454,504>,
  <505,504,454>,
  <454,455,505>,
  <506,505,455>,
  <455,456,506>,
  <507,506,456>,
  <456,457,507>,
  <508,507,457>,
  <457,458,508>,
  <509,508,458>,
  <459,460,510>,
  <511,510,460>,
  <460,461,511>,
  <512,511,461>,
  <461,462,512>,
  <513,512,462>,
  <462,463,513>,
  <514,513,463>,
  <463,464,514>,
  <515,514,464>,
  <464,465,515>,
  <516,515,465>,
  <465,466,516>,
  <517,516,466>,
  <466,467,517>,
  <518,517,467>,
  <467,468,518>,
  <519,518,468>,
  <468,469,519>,
  <520,519,469>,
  <469,470,520>,
  <521,520,470>,
  <470,471,521>,
  <522,521,471>,
  <471,472,522>,
  <523,522,472>,
  <472,473,523>,
  <524,523,473>,
  <473,474,524>,
  <525,524,474>,
  <474,475,525>,
  <526,525,475>,
  <475,476,526>,
  <527,526,476>,
  <476,477,527>,
  <528,527,477>,
  <477,478,528>,
  <529,528,478>,
  <478,479,529>,
  <530,529,479>,
  <479,480,530>,
  <531,530,480>,
  <480,481,531>,
  <532,531,481>,
  <481,482,532>,
  <533,532,482>,
  <482,483,533>,
  <534,533,483>,
  <483,484,534>,
  <535,534,484>,
  <484,485,535>,
  <536,535,485>,
  <485,486,536>,
  <537,536,486>,
  <486,487,537>,
  <538,537,487>,
  <487,488,538>,
  <539,538,488>,
  <488,489,539>,
  <540,539,489>,
  <489,490,540>,
  <541,540,490>,
  <490,491,541>,
  <542,541,491>,
  <491,492,542>,
  <543,542,492>,
  <492,493,543>,
  <544,543,493>,
  <493,494,544>,
  <545,544,494>,
  <494,495,545>,
  <546,545,495>,
  <495,496,546>,
  <547,546,496>,
  <496,497,547>,
  <548,547,497>,
  <497,498,548>,
  <549,548,498>,
  <498,499,549>,
  <550,549,499>,
  <499,500,550>,
  <551,550,500>,
  <500,501,551>,
  <552,551,501>,
  <501,502,552>,
  <553,552,502>,
  <502,503,553>,
  <554,553,503>,
  <503,504,554>,
  <555,554,504>,
  <504,505,555>,
  <556,555,505>,
  <505,506,556>,
  <557,556,506>,
  <506,507,557>,
  <558,557,507>,
  <507,508,558>,
  <559,558,508>,
  <508,509,559>,
  <560,559,509>,
  <510,511,561>,
  <562,561,511>,
  <511,512,562>,
  <563,562,512>,
  <512,513,563>,
  <564,563,513>,
  <513,514,564>,
  <565,564,514>,
  <514,515,565>,
  <566,565,515>,
  <515,516,566>,
  <567,566,516>,
  <516,517,567>,
  <568,567,517>,
  <517,518,568>,
  <569,568,518>,
  <518,519,569>,
  <570,569,519>,
  <519,520,570>,
  <571,570,520>,
  <520,521,571>,
  <572,571,521>,
  <521,522,572>,
  <573,572,522>,
  <522,523,573>,
  <574,573,523>,
  <523,524,574>,
  <575,574,524>,
  <524,525,575>,
  <576,575,525>,
  <525,526,576>,
  <577,576,526>,
  <526,527,577>,
  <578,577,527>,
  <527,528,578>,
  <579,578,528>,
  <528,529,579>,
  <580,579,529>,
  <529,530,580>,
  <581,580,530>,
  <530,531,581>,
  <582,581,531>,
  <531,532,582>,
  <583,582,532>,
  <532,533,583>,
  <584,583,533>,
  <533,534,584>,
  <585,584,534>,
  <534,535,585>,
  <586,585,535>,
  <535,536,586>,
  <587,586,536>,
  <536,537,587>,
  <588,587,537>,
  <537,538,588>,
  <589,588,538>,
  <538,539,589>,
  <590,589,539>,
  <539,540,590>,
  <591,590,540>,
  <540,541,591>,
  <592,591,541>,
  <541,542,592>,
  <593,592,542>,
  <542,543,593>,
  <594,593,543>,
  <543,544,594>,
  <595,594,544>,
  <544,545,595>,
  <596,595,545>,
  <545,546,596>,
  <597,596,546>,
  <546,547,597>,
  <598,597,547>,
  <547,548,598>,
  <599,598,548>,
  <548,549,599>,
  <600,599,549>,
  <549,550,600>,
  <601,600,550>,
  <550,551,601>,
  <602,601,551>,
  <551,552,602>,
  <603,602,552>,
  <552,553,603>,
  <604,603,553>,
  <553,554,604>,
  <605,604,554>,
  <554,555,605>,
  <606,605,555>,
  <555,556,606>,
  <607,606,556>,
  <556,557,607>,
  <608,607,557>,
  <557,558,608>,
  <609,608,558>,
  <558,559,609>,
  <610,609,559>,
  <559,560,610>,
  <611,610,560>,
  <561,562,612>,
  <613,612,562>,
  <562,563,613>,
  <614,613,563>,
  <563,564,614>,
  <615,614,564>,
  <564,565,615>,
  <616,615,565>,
  <565,566,616>,
  <617,616,566>,
  <566,567,617>,
  <618,617,567>,
  <567,568,618>,
  <619,618,568>,
  <568,569,619>,
  <620,619,569>,
  <569,570,620>,
  <621,620,570>,
  <570,571,621>,
  <622,621,571>,
  <571,572,622>,
  <623,622,572>,
  <572,573,623>,
  <624,623,573>,
  <573,574,624>,
  <625,624,574>,
  <574,575,625>,
  <626,625,575>,
  <575,576,626>,
  <627,626,576>,
  <576,577,627>,
  <628,627,577>,
  <577,578,628>,
  <629,628,578>,
  <578,579,629>,
  <630,629,579>,
  <579,580,630>,
  <631,630,580>,
  <580,581,631>,
  <632,631,581>,
  <581,582,632>,
  <633,632,582>,
  <582,583,633>,
  <634,633,583>,
  <583,584,634>,
  <635,634,584>,
  <584,585,635>,
  <636,635,585>,
  <585,586,636>,
  <637,636,586>,
  <586,587,637>,
  <638,637,587>,
  <587,588,638>,
  <639,638,588>,
  <588,589,639>,
  <640,639,589>,
  <589,590,640>,
  <641,640,590>,
  <590,591,641>,
  <642,641,591>,
  <591,592,642>,
  <643,642,592>,
  <592,593,643>,
  <644,643,593>,
  <593,594,644>,
  <645,644,594>,
  <594,595,645>,
  <646,645,595>,
  <595,596,646>,
  <647,646,596>,
  <596,597,647>,
  <648,647,597>,
  <597,598,648>,
  <649,648,598>,
  <598,599,649>,
  <650,649,599>,
  <599,600,650>,
  <651,650,600>,
  <600,601,651>,
  <652,651,601>,
  <601,602,652>,
  <653,652,602>,
  <602,603,653>,
  <654,653,603>,
  <603,604,654>,
  <655,654,604>,
  <604,605,655>,
  <656,655,605>,
  <605,606,656>,
  <657,656,606>,
  <606,607,657>,
  <658,657,607>,
  <607,608,658>,
  <659,658,608>,
  <608,609,659>,
  <660,659,609>,
  <609,610,660>,
  <661,660,610>,
  <610,611,661>,
  <662,661,611>,
  <612,613,663>,
  <664,663,613>,
  <613,614,664>,
  <665,664,614>,
  <614,615,665>,
  <666,665,615>,
  <615,616,666>,
  <667,666,616>,
  <616,617,667>,
  <668,667,617>,
  <617,618,668>,
  <669,668,618>,
  <618,619,669>,
  <670,669,619>,
  <619,620,670>,
  <671,670,620>,
  <620,621,671>,
  <672,671,621>,
  <621,622,672>,
  <673,672,622>,
  <622,623,673>,
  <674,673,623>,
  <623,624,674>,
  <675,674,624>,
  <624,625,675>,
  <676,675,625>,
  <625,626,676>,
  <677,676,626>,
  <626,627,677>,
  <678,677,627>,
  <627,628,678>,
  <679,678,628>,
  <628,629,679>,
  <680,679,629>,
  <629,630,680>,
  <681,680,630>,
  <630,631,681>,
  <682,681,631>,
  <631,632,682>,
  <683,682,632>,
  <632,633,683>,
  <684,683,633>,
  <633,634,684>,
  <685,684,634>,
  <634,635,685>,
  <686,685,635>,
  <635,636,686>,
  <687,686,636>,
  <636,637,687>,
  <688,687,637>,
  <637,638,688>,
  <689,688,638>,
  <638,639,689>,
  <690,689,639>,
  <639,640,690>,
  <691,690,640>,
  <640,641,691>,
  <692,691,641>,
  <641,642,692>,
  <693,692,642>,
  <642,643,693>,
  <694,693,643>,
  <643,644,694>,
  <695,694,644>,
  <644,645,695>,
  <696,695,645>,
  <645,646,696>,
  <697,696,646>,
  <646,647,697>,
  <698,697,647>,
  <647,648,698>,
  <699,698,648>,
  <648,649,699>,
  <700,699,649>,
  <649,650,700>,
  <701,700,650>,
  <650,651,701>,
  <702,701,651>,
  <651,652,702>,
  <703,702,652>,
  <652,653,703>,
  <704,703,653>,
  <653,654,704>,
  <705,704,654>,
  <654,655,705>,
  <706,705,655>,
  <655,656,706>,
  <707,706,656>,
  <656,657,707>,
  <708,707,657>,
  <657,658,708>,
  <709,708,658>,
  <658,659,709>,
  <710,709,659>,
  <659,660,710>,
  <711,710,660>,
  <660,661,711>,
  <712,711,661>,
  <661,662,712>,
  <713,712,662>,
  <663,664,714>,
  <715,714,664>,
  <664,665,715>,
  <716,715,665>,
  <665,666,716>,
  <717,716,666>,
  <666,667,717>,
  <718,717,667>,
  <667,668,718>,
  <719,718,668>,
  <668,669,719>,
  <720,719,669>,
  <669,670,720>,
  <721,720,670>,
  <670,671,721>,
  <722,721,671>,
  <671,672,722>,
  <723,722,672>,
  <672,673,723>,
  <724,723,673>,
  <673,674,724>,
  <725,724,674>,
  <674,675,725>,
  <726,725,675>,
  <675,676,726>,
  <727,726,676>,
  <676,677,727>,
  <728,727,677>,
  <677,678,728>,
  <729,728,678>,
  <678,679,729>,
  <730,729,679>,
  <679,680,730>,
  <731,730,680>,
  <680,681,731>,
  <732,731,681>,
  <681,682,732>,
  <733,732,682>,
  <682,683,733>,
  <734,733,683>,
  <683,684,734>,
  <735,734,684>,
  <684,685,735>,
  <736,735,685>,
  <685,686,736>,
  <737,736,686>,
  <686,687,737>,
  <738,737,687>,
  <687,688,738>,
  <739,738,688>,
  <688,689,739>,
  <740,739,689>,
  <689,690,740>,
  <741,740,690>,
  <690,691,741>,
  <742,741,691>,
  <691,692,742>,
  <743,742,692>,
  <692,693,743>,
  <744,743,693>,
  <693,694,744>,
  <745,744,694>,
  <694,695,745>,
  <746,745,695>,
  <695,696,746>,
  <747,746,696>,
  <696,697,747>,
  <748,747,697>,
  <697,698,748>,
  <749,748,698>,
  <698,699,749>,
  <750,749,699>,
  <699,700,750>,
  <751,750,700>,
  <700,701,751>,
  <752,751,701>,
  <701,702,752>,
  <753,752,702>,
  <702,703,753>,
  <754,753,703>,
  <703,704,754>,
  <755,754,704>,
  <704,705,755>,
  <756,755,705>,
  <705,706,756>,
  <757,756,706>,
  <706,707,757>,
  <758,757,707>,
  <707,708,758>,
  <759,758,708>,
  <708,709,759>,
  <760,759,709>,
  <709,710,760>,
  <761,760,710>,
  <710,711,761>,
  <762,761,711>,
  <711,712,762>,
  <763,762,712>,
  <712,713,763>,
  <764,763,713>,
  <714,715,765>,
  <766,765,715>,
  <715,716,766>,
  <767,766,716>,
  <716,717,767>,
  <768,767,717>,
  <717,718,768>,
  <769,768,718>,
  <718,719,769>,
  <770,769,719>,
  <719,720,770>,
  <771,770,720>,
  <720,721,771>,
  <772,771,721>,
  <721,722,772>,
  <773,772,722>,
  <722,723,773>,
  <774,773,723>,
  <723,724,774>,
  <775,774,724>,
  <724,725,775>,
  <776,775,725>,
  <725,726,776>,
  <777,776,726>,
  <726,727,777>,
  <778,777,727>,
  <727,728,778>,
  <779,778,728>,
  <728,729,779>,
  <780,779,729>,
  <729,730,780>,
  <781,780,730>,
  <730,731,781>,
  <782,781,731>,
  <731,732,782>,
  <783,782,732>,
  <732,733,783>,
  <784,783,733>,
  <733,734,784>,
  <785,784,734>,
  <734,735,785>,
  <786,785,735>,
  <735,736,786>,
  <787,786,736>,
  <736,737,787>,
  <788,787,737>,
  <737,738,788>,
  <789,788,738>,
  <738,739,789>,
  <790,789,739>,
  <739,740,790>,
  <791,790,740>,
  <740,741,791>,
  <792,791,741>,
  <741,742,792>,
  <793,792,742>,
  <742,743,793>,
  <794,793,743>,
  <743,744,794>,
  <795,794,744>,
  <744,745,795>,
  <796,795,745>,
  <745,746,796>,
  <797,796,746>,
  <746,747,797>,
  <798,797,747>,
  <747,748,798>,
  <799,798,748>,
  <748,749,799>,
  <800,799,749>,
  <749,750,800>,
  <801,800,750>,
  <750,751,801>,
  <802,801,751>,
  <751,752,802>,
  <803,802,752>,
  <752,753,803>,
  <804,803,753>,
  <753,754,804>,
  <805,804,754>,
  <754,755,805>,
  <806,805,755>,
  <755,756,806>,
  <807,806,756>,
  <756,757,807>,
  <808,807,757>,
  <757,758,808>,
  <809,808,758>,
  <758,759,809>,
  <810,809,759>,
  <759,760,810>,
  <811,810,760>,
  <760,761,811>,
  <812,811,761>,
  <761,762,812>,
  <813,812,762>,
  <762,763,813>,
  <814,813,763>,
  <763,764,814>,
  <815,814,764>,
  <765,766,816>,
  <817,816,766>,
  <766,767,817>,
  <818,817,767>,
  <767,768,818>,
  <819,818,768>,
  <768,769,819>,
  <820,819,769>,
  <769,770,820>,
  <821,820,770>,
  <770,771,821>,
  <822,821,771>,
  <771,772,822>,
  <823,822,772>,
  <772,773,823>,
  <824,823,773>,
  <773,774,824>,
  <825,824,774>,
  <774,775,825>,
  <826,825,775>,
  <775,776,826>,
  <827,826,776>,
  <776,777,827>,
  <828,827,777>,
  <777,778,828>,
  <829,828,778>,
  <778,779,829>,
  <830,829,779>,
  <779,780,830>,
  <831,830,780>,
  <780,781,831>,
  <832,831,781>,
  <781,782,832>,
  <833,832,782>,
  <782,783,833>,
  <834,833,783>,
  <783,784,834>,
  <835,834,784>,
  <784,785,835>,
  <836,835,785>,
  <785,786,836>,
  <837,836,786>,
  <786,787,837>,
  <838,837,787>,
  <787,788,838>,
  <839,838,788>,
  <788,789,839>,
  <840,839,789>,
  <789,790,840>,
  <841,840,790>,
  <790,791,841>,
  <842,841,791>,
  <791,792,842>,
  <843,842,792>,
  <792,793,843>,
  <844,843,793>,
  <793,794,844>,
  <845,844,794>,
  <794,795,845>,
  <846,845,795>,
  <795,796,846>,
  <847,846,796>,
  <796,797,847>,
  <848,847,797>,
  <797,798,848>,
  <849,848,798>,
  <798,799,849>,
  <850,849,799>,
  <799,800,850>,
  <851,850,800>,
  <800,801,851>,
  <852,851,801>,
  <801,802,852>,
  <853,852,802>,
  <802,803,853>,
  <854,853,803>,
  <803,804,854>,
  <855,854,804>,
  <804,805,855>,
  <856,855,805>,
  <805,806,856>,
  <857,856,806>,
  <806,807,857>,
  <858,857,807>,
  <807,808,858>,
  <859,858,808>,
  <808,809,859>,
  <860,859,809>,
  <809,810,860>,
  <861,860,810>,
  <810,811,861>,
  <862,861,811>,
  <811,812,862>,
  <863,862,812>,
  <812,813,863>,
  <864,863,813>,
  <813,814,864>,
  <865,864,814>,
  <814,815,865>,
  <866,865,815>,
  <816,817,867>,
  <868,867,817>,
  <817,818,868>,
  <869,868,818>,
  <818,819,869>,
  <870,869,819>,
  <819,820,870>,
  <871,870,820>,
  <820,821,871>,
  <872,871,821>,
  <821,822,872>,
  <873,872,822>,
  <822,823,873>,
  <874,873,823>,
  <823,824,874>,
  <875,874,824>,
  <824,825,875>,
  <876,875,825>,
  <825,826,876>,
  <877,876,826>,
  <826,827,877>,
  <878,877,827>,
  <827,828,878>,
  <879,878,828>,
  <828,829,879>,
  <880,879,829>,
  <829,830,880>,
  <881,880,830>,
  <830,831,881>,
  <882,881,831>,
  <831,832,882>,
  <883,882,832>,
  <832,833,883>,
  <884,883,833>,
  <833,834,884>,
  <885,884,834>,
  <834,835,885>,
  <886,885,835>,
  <835,836,886>,
  <887,886,836>,
  <836,837,887>,
  <888,887,837>,
  <837,838,888>,
  <889,888,838>,
  <838,839,889>,
  <890,889,839>,
  <839,840,890>,
  <891,890,840>,
  <840,841,891>,
  <892,891,841>,
  <841,842,892>,
  <893,892,842>,
  <842,843,893>,
  <894,893,843>,
  <843,844,894>,
  <895,894,844>,
  <844,845,895>,
  <896,895,845>,
  <845,846,896>,
  <897,896,846>,
  <846,847,897>,
  <898,897,847>,
  <847,848,898>,
  <899,898,848>,
  <848,849,899>,
  <900,899,849>,
  <849,850,900>,
  <901,900,850>,
  <850,851,901>,
  <902,901,851>,
  <851,852,902>,
  <903,902,852>,
  <852,853,903>,
  <904,903,853>,
  <853,854,904>,
  <905,904,854>,
  <854,855,905>,
  <906,905,855>,
  <855,856,906>,
  <907,906,856>,
  <856,857,907>,
  <908,907,857>,
  <857,858,908>,
  <909,908,858>,
  <858,859,909>,
  <910,909,859>,
  <859,860,910>,
  <911,910,860>,
  <860,861,911>,
  <912,911,861>,
  <861,862,912>,
  <913,912,862>,
  <862,863,913>,
  <914,913,863>,
  <863,864,914>,
  <915,914,864>,
  <864,865,915>,
  <916,915,865>,
  <865,866,916>,
  <917,916,866>,
  <867,868,918>,
  <919,918,868>,
  <868,869,919>,
  <920,919,869>,
  <869,870,920>,
  <921,920,870>,
  <870,871,921>,
  <922,921,871>,
  <871,872,922>,
  <923,922,872>,
  <872,873,923>,
  <924,923,873>,
  <873,874,924>,
  <925,924,874>,
  <874,875,925>,
  <926,925,875>,
  <875,876,926>,
  <927,926,876>,
  <876,877,927>,
  <928,927,877>,
  <877,878,928>,
  <929,928,878>,
  <878,879,929>,
  <930,929,879>,
  <879,880,930>,
  <931,930,880>,
  <880,881,931>,
  <932,931,881>,
  <881,882,932>,
  <933,932,882>,
  <882,883,933>,
  <934,933,883>,
  <883,884,934>,
  <935,934,884>,
  <884,885,935>,
  <936,935,885>,
  <885,886,936>,
  <937,936,886>,
  <886,887,937>,
  <938,937,887>,
  <887,888,938>,
  <939,938,888>,
  <888,889,939>,
  <940,939,889>,
  <889,890,940>,
  <941,940,890>,
  <890,891,941>,
  <942,941,891>,
  <891,892,942>,
  <943,942,892>,
  <892,893,943>,
  <944,943,893>,
  <893,894,944>,
  <945,944,894>,
  <894,895,945>,
  <946,945,895>,
  <895,896,946>,
  <947,946,896>,
  <896,897,947>,
  <948,947,897>,
  <897,898,948>,
  <949,948,898>,
  <898,899,949>,
  <950,949,899>,
  <899,900,950>,
  <951,950,900>,
  <900,901,951>,
  <952,951,901>,
  <901,902,952>,
  <953,952,902>,
  <902,903,953>,
  <954,953,903>,
  <903,904,954>,
  <955,954,904>,
  <904,905,955>,
  <956,955,905>,
  <905,906,956>,
  <957,956,906>,
  <906,907,957>,
  <958,957,907>,
  <907,908,958>,
  <959,958,908>,
  <908,909,959>,
  <960,959,909>,
  <909,910,960>,
  <961,960,910>,
  <910,911,961>,
  <962,961,911>,
  <911,912,962>,
  <963,962,912>,
  <912,913,963>,
  <964,963,913>,
  <913,914,964>,
  <965,964,914>,
  <914,915,965>,
  <966,965,915>,
  <915,916,966>,
  <967,966,916>,
  <916,917,967>,
  <968,967,917>,
  <918,919,969>,
  <970,969,919>,
  <919,920,970>,
  <971,970,920>,
  <920,921,971>,
  <972,971,921>,
  <921,922,972>,
  <973,972,922>,
  <922,923,973>,
  <974,973,923>,
  <923,924,974>,
  <975,974,924>,
  <924,925,975>,
  <976,975,925>,
  <925,926,976>,
  <977,976,926>,
  <926,927,977>,
  <978,977,927>,
  <927,928,978>,
  <979,978,928>,
  <928,929,979>,
  <980,979,929>,
  <929,930,980>,
  <981,980,930>,
  <930,931,981>,
  <982,981,931>,
  <931,932,982>,
  <983,982,932>,
  <932,933,983>,
  <984,983,933>,
  <933,934,984>,
  <985,984,934>,
  <934,935,985>,
  <986,985,935>,
  <935,936,986>,
  <987,986,936>,
  <936,937,987>,
  <988,987,937>,
  <937,938,988>,
  <989,988,938>,
  <938,939,989>,
  <990,989,939>,
  <939,940,990>,
  <991,990,940>,
  <940,941,991>,
  <992,991,941>,
  <941,942,992>,
  <993,992,942>,
  <942,943,993>,
  <994,993,943>,
  <943,944,994>,
  <995,994,944>,
  <944,945,995>,
  <996,995,945>,
  <945,946,996>,
  <997,996,946>,
  <946,947,997>,
  <998,997,947>,
  <947,948,998>,
  <999,998,948>,
  <948,949,999>,
  <1000,999,949>,
  <949,950,1000>,
  <1001,1000,950>,
  <950,951,1001>,
  <1002,1001,951>,
  <951,952,1002>,
  <1003,1002,952>,
  <952,953,1003>,
  <1004,1003,953>,
  <953,954,1004>,
  <1005,1004,954>,
  <954,955,1005>,
  <1006,1005,955>,
  <955,956,1006>,
  <1007,1006,956>,
  <956,957,1007>,
  <1008,1007,957>,
  <957,958,1008>,
  <1009,1008,958>,
  <958,959,1009>,
  <1010,1009,959>,
  <959,960,1010>,
  <1011,1010,960>,
  <960,961,1011>,
  <1012,1011,961>,
  <961,962,1012>,
  <1013,1012,962>,
  <962,963,1013>,
  <1014,1013,963>,
  <963,964,1014>,
  <1015,1014,964>,
  <964,965,1015>,
  <1016,1015,965>,
  <965,966,1016>,
  <1017,1016,966>,
  <966,967,1017>,
  <1018,1017,967>,
  <967,968,1018>,
  <1019,1018,968>,
  <969,970,1020>,
  <1021,1020,970>,
  <970,971,1021>,
  <1022,1021,971>,
  <971,972,1022>,
  <1023,1022,972>,
  <972,973,1023>,
  <1024,1023,973>,
  <973,974,1024>,
  <1025,1024,974>,
  <974,975,1025>,
  <1026,1025,975>,
  <975,976,1026>,
  <1027,1026,976>,
  <976,977,1027>,
  <1028,1027,977>,
  <977,978,1028>,
  <1029,1028,978>,
  <978,979,1029>,
  <1030,1029,979>,
  <979,980,1030>,
  <1031,1030,980>,
  <980,981,1031>,
  <1032,1031,981>,
  <981,982,1032>,
  <1033,1032,982>,
  <982,983,1033>,
  <1034,1033,983>,
  <983,984,1034>,
  <1035,1034,984>,
  <984,985,1035>,
  <1036,1035,985>,
  <985,986,1036>,
  <1037,1036,986>,
  <986,987,1037>,
  <1038,1037,987>,
  <987,988,1038>,
  <1039,1038,988>,
  <988,989,1039>,
  <1040,1039,989>,
  <989,990,1040>,
  <1041,1040,990>,
  <990,991,1041>,
  <1042,1041,991>,
  <991,992,1042>,
  <1043,1042,992>,
  <992,993,1043>,
  <1044,1043,993>,
  <993,994,1044>,
  <1045,1044,994>,
  <994,995,1045>,
  <1046,1045,995>,
  <995,996,1046>,
  <1047,1046,996>,
  <996,997,1047>,
  <1048,1047,997>,
  <997,998,1048>,
  <1049,1048,998>,
  <998,999,1049>,
  <1050,1049,999>,
  <999,1000,1050>,
  <1051,1050,1000>,
  <1000,1001,1051>,
  <1052,1051,1001>,
  <1001,1002,1052>,
  <1053,1052,1002>,
  <1002,1003,1053>,
  <1054,1053,1003>,
  <1003,1004,1054>,
  <1055,1054,1004>,
  <1004,1005,1055>,
  <1056,1055,1005>,
  <1005,1006,1056>,
  <1057,1056,1006>,
  <1006,1007,1057>,
  <1058,1057,1007>,
  <1007,1008,1058>,
  <1059,1058,1008>,
  <1008,1009,1059>,
  <1060,1059,1009>,
  <1009,1010,1060>,
  <1061,1060,1010>,
  <1010,1011,1061>,
  <1062,1061,1011>,
  <1011,1012,1062>,
  <1063,1062,1012>,
  <1012,1013,1063>,
  <1064,1063,1013>,
  <1013,1014,1064>,
  <1065,1064,1014>,
  <1014,1015,1065>,
  <1066,1065,1015>,
  <1015,1016,1066>,
  <1067,1066,1016>,
  <1016,1017,1067>,
  <1068,1067,1017>,
  <1017,1018,1068>,
  <1069,1068,1018>,
  <1018,1019,1069>,
  <1070,1069,1019>,
  <1020,1021,1071>,
  <1072,1071,1021>,
  <1021,1022,1072>,
  <1073,1072,1022>,
  <1022,1023,1073>,
  <1074,1073,1023>,
  <1023,1024,1074>,
  <1075,1074,1024>,
  <1024,1025,1075>,
  <1076,1075,1025>,
  <1025,1026,1076>,
  <1077,1076,1026>,
  <1026,1027,1077>,
  <1078,1077,1027>,
  <1027,1028,1078>,
  <1079,1078,1028>,
  <1028,1029,1079>,
  <1080,1079,1029>,
  <1029,1030,1080>,
  <1081,1080,1030>,
  <1030,1031,1081>,
  <1082,1081,1031>,
  <1031,1032,1082>,
  <1083,1082,1032>,
  <1032,1033,1083>,
  <1084,1083,1033>,
  <1033,1034,1084>,
  <1085,1084,1034>,
  <1034,1035,1085>,
  <1086,1085,1035>,
  <1035,1036,1086>,
  <1087,1086,1036>,
  <1036,1037,1087>,
  <1088,1087,1037>,
  <1037,1038,1088>,
  <1089,1088,1038>,
  <1038,1039,1089>,
  <1090,1089,1039>,
  <1039,1040,1090>,
  <1091,1090,1040>,
  <1040,1041,1091>,
  <1092,1091,1041>,
  <1041,1042,1092>,
  <1093,1092,1042>,
  <1042,1043,1093>,
  <1094,1093,1043>,
  <1043,1044,1094>,
  <1095,1094,1044>,
  <1044,1045,1095>,
  <1096,1095,1045>,
  <1045,1046,1096>,
  <1097,1096,1046>,
  <1046,1047,1097>,
  <1098,1097,1047>,
  <1047,1048,1098>,
  <1099,1098,1048>,
  <1048,1049,1099>,
  <1100,1099,1049>,
  <1049,1050,1100>,
  <1101,1100,1050>,
  <1050,1051,1101>,
  <1102,1101,1051>,
  <1051,1052,1102>,
  <1103,1102,1052>,
  <1052,1053,1103>,
  <1104,1103,1053>,
  <1053,1054,1104>,
  <1105,1104,1054>,
  <1054,1055,1105>,
  <1106,1105,1055>,
  <1055,1056,1106>,
  <1107,1106,1056>,
  <1056,1057,1107>,
  <1108,1107,1057>,
  <1057,1058,1108>,
  <1109,1108,1058>,
  <1058,1059,1109>,
  <1110,1109,1059>,
  <1059,1060,1110>,
  <1111,1110,1060>,
  <1060,1061,1111>,
  <1112,1111,1061>,
  <1061,1062,1112>,
  <1113,1112,1062>,
  <1062,1063,1113>,
  <1114,1113,1063>,
  <1063,1064,1114>,
  <1115,1114,1064>,
  <1064,1065,1115>,
  <1116,1115,1065>,
  <1065,1066,1116>,
  <1117,1116,1066>,
  <1066,1067,1117>,
  <1118,1117,1067>,
  <1067,1068,1118>,
  <1119,1118,1068>,
  <1068,1069,1119>,
  <1120,1119,1069>,
  <1069,1070,1120>,
  <1121,1120,1070>,
  <1071,1072,1122>,
  <1123,1122,1072>,
  <1072,1073,1123>,
  <1124,1123,1073>,
  <1073,1074,1124>,
  <1125,1124,1074>,
  <1074,1075,1125>,
  <1126,1125,1075>,
  <1075,1076,1126>,
  <1127,1126,1076>,
  <1076,1077,1127>,
  <1128,1127,1077>,
  <1077,1078,1128>,
  <1129,1128,1078>,
  <1078,1079,1129>,
  <1130,1129,1079>,
  <1079,1080,1130>,
  <1131,1130,1080>,
  <1080,1081,1131>,
  <1132,1131,1081>,
  <1081,1082,1132>,
  <1133,1132,1082>,
  <1082,1083,1133>,
  <1134,1133,1083>,
  <1083,1084,1134>,
  <1135,1134,1084>,
  <1084,1085,1135>,
  <1136,1135,1085>,
  <1085,1086,1136>,
  <1137,1136,1086>,
  <1086,1087,1137>,
  <1138,1137,1087>,
  <1087,1088,1138>,
  <1139,1138,1088>,
  <1088,1089,1139>,
  <1140,1139,1089>,
  <1089,1090,1140>,
  <1141,1140,1090>,
  <1090,1091,1141>,
  <1142,1141,1091>,
  <1091,1092,1142>,
  <1143,1142,1092>,
  <1092,1093,1143>,
  <1144,1143,1093>,
  <1093,1094,1144>,
  <1145,1144,1094>,
  <1094,1095,1145>,
  <1146,1145,1095>,
  <1095,1096,1146>,
  <1147,1146,1096>,
  <1096,1097,1147>,
  <1148,1147,1097>,
  <1097,1098,1148>,
  <1149,1148,1098>,
  <1098,1099,1149>,
  <1150,1149,1099>,
  <1099,1100,1150>,
  <1151,1150,1100>,
  <1100,1101,1151>,
  <1152,1151,1101>,
  <1101,1102,1152>,
  <1153,1152,1102>,
  <1102,1103,1153>,
  <1154,1153,1103>,
  <1103,1104,1154>,
  <1155,1154,1104>,
  <1104,1105,1155>,
  <1156,1155,1105>,
  <1105,1106,1156>,
  <1157,1156,1106>,
  <1106,1107,1157>,
  <1158,1157,1107>,
  <1107,1108,1158>,
  <1159,1158,1108>,
  <1108,1109,1159>,
  <1160,1159,1109>,
  <1109,1110,1160>,
  <1161,1160,1110>,
  <1110,1111,1161>,
  <1162,1161,1111>,
  <1111,1112,1162>,
  <1163,1162,1112>,
  <1112,1113,1163>,
  <1164,1163,1113>,
  <1113,1114,1164>,
  <1165,1164,1114>,
  <1114,1115,1165>,
  <1166,1165,1115>,
  <1115,1116,1166>,
  <1167,1166,1116>,
  <1116,1117,1167>,
  <1168,1167,1117>,
  <1117,1118,1168>,
  <1169,1168,1118>,
  <1118,1119,1169>,
  <1170,1169,1119>,
  <1119,1120,1170>,
  <1171,1170,1120>,
  <1120,1121,1171>,
  <1172,1171,1121>,
  <1122,1123,1173>,
  <1174,1173,1123>,
  <1123,1124,1174>,
  <1175,1174,1124>,
  <1124,1125,1175>,
  <1176,1175,1125>,
  <1125,1126,1176>,
  <1177,1176,1126>,
  <1126,1127,1177>,
  <1178,1177,1127>,
  <1127,1128,1178>,
  <1179,1178,1128>,
  <1128,1129,1179>,
  <1180,1179,1129>,
  <1129,1130,1180>,
  <1181,1180,1130>,
  <1130,1131,1181>,
  <1182,1181,1131>,
  <1131,1132,1182>,
  <1183,1182,1132>,
  <1132,1133,1183>,
  <1184,1183,1133>,
  <1133,1134,1184>,
  <1185,1184,1134>,
  <1134,1135,1185>,
  <1186,1185,1135>,
  <1135,1136,1186>,
  <1187,1186,1136>,
  <1136,1137,1187>,
  <1188,1187,1137>,
  <1137,1138,1188>,
  <1189,1188,1138>,
  <1138,1139,1189>,
  <1190,1189,1139>,
  <1139,1140,1190>,
  <1191,1190,1140>,
  <1140,1141,1191>,
  <1192,1191,1141>,
  <1141,1142,1192>,
  <1193,1192,1142>,
  <1142,1143,1193>,
  <1194,1193,1143>,
  <1143,1144,1194>,
  <1195,1194,1144>,
  <1144,1145,1195>,
  <1196,1195,1145>,
  <1145,1146,1196>,
  <1197,1196,1146>,
  <1146,1147,1197>,
  <1198,1197,1147>,
  <1147,1148,1198>,
  <1199,1198,1148>,
  <1148,1149,1199>,
  <1200,1199,1149>,
  <1149,1150,1200>,
  <1201,1200,1150>,
  <1150,1151,1201>,
  <1202,1201,1151>,
  <1151,1152,1202>,
  <1203,1202,1152>,
  <1152,1153,1203>,
  <1204,1203,1153>,
  <1153,1154,1204>,
  <1205,1204,1154>,
  <1154,1155,1205>,
  <1206,1205,1155>,
  <1155,1156,1206>,
  <1207,1206,1156>,
  <1156,1157,1207>,
  <1208,1207,1157>,
  <1157,1158,1208>,
  <1209,1208,1158>,
  <1158,1159,1209>,
  <1210,1209,1159>,
  <1159,1160,1210>,
  <1211,1210,1160>,
  <1160,1161,1211>,
  <1212,1211,1161>,
  <1161,1162,1212>,
  <1213,1212,1162>,
  <1162,1163,1213>,
  <1214,1213,1163>,
  <1163,1164,1214>,
  <1215,1214,1164>,
  <1164,1165,1215>,
  <1216,1215,1165>,
  <1165,1166,1216>,
  <1217,1216,1166>,
  <1166,1167,1217>,
  <1218,1217,1167>,
  <1167,1168,1218>,
  <1219,1218,1168>,
  <1168,1169,1219>,
  <1220,1219,1169>,
  <1169,1170,1220>,
  <1221,1220,1170>,
  <1170,1171,1221>,
  <1222,1221,1171>,
  <1171,1172,1222>,
  <1223,1222,1172>,
  <1173,1174,1224>,
  <1225,1224,1174>,
  <1174,1175,1225>,
  <1226,1225,1175>,
  <1175,1176,1226>,
  <1227,1226,1176>,
  <1176,1177,1227>,
  <1228,1227,1177>,
  <1177,1178,1228>,
  <1229,1228,1178>,
  <1178,1179,1229>,
  <1230,1229,1179>,
  <1179,1180,1230>,
  <1231,1230,1180>,
  <1180,1181,1231>,
  <1232,1231,1181>,
  <1181,1182,1232>,
  <1233,1232,1182>,
  <1182,1183,1233>,
  <1234,1233,1183>,
  <1183,1184,1234>,
  <1235,1234,1184>,
  <1184,1185,1235>,
  <1236,1235,1185>,
  <1185,1186,1236>,
  <1237,1236,1186>,
  <1186,1187,1237>,
  <1238,1237,1187>,
  <1187,1188,1238>,
  <1239,1238,1188>,
  <1188,1189,1239>,
  <1240,1239,1189>,
  <1189,1190,1240>,
  <1241,1240,1190>,
  <1190,1191,1241>,
  <1242,1241,1191>,
  <1191,1192,1242>,
  <1243,1242,1192>,
  <1192,1193,1243>,
  <1244,1243,1193>,
  <1193,1194,1244>,
  <1245,1244,1194>,
  <1194,1195,1245>,
  <1246,1245,1195>,
  <1195,1196,1246>,
  <1247,1246,1196>,
  <1196,1197,1247>,
  <1248,1247,1197>,
  <1197,1198,1248>,
  <1249,1248,1198>,
  <1198,1199,1249>,
  <1250,1249,1199>,
  <1199,1200,1250>,
  <1251,1250,1200>,
  <1200,1201,1251>,
  <1252,1251,1201>,
  <1201,1202,1252>,
  <1253,1252,1202>,
  <1202,1203,1253>,
  <1254,1253,1203>,
  <1203,1204,1254>,
  <1255,1254,1204>,
  <1204,1205,1255>,
  <1256,1255,1205>,
  <1205,1206,1256>,
  <1257,1256,1206>,
  <1206,1207,1257>,
  <1258,1257,1207>,
  <1207,1208,1258>,
  <1259,1258,1208>,
  <1208,1209,1259>,
  <1260,1259,1209>,
  <1209,1210,1260>,
  <1261,1260,1210>,
  <1210,1211,1261>,
  <1262,1261,1211>,
  <1211,1212,1262>,
  <1263,1262,1212>,
  <1212,1213,1263>,
  <1264,1263,1213>,
  <1213,1214,1264>,
  <1265,1264,1214>,
  <1214,1215,1265>,
  <1266,1265,1215>,
  <1215,1216,1266>,
  <1267,1266,1216>,
  <1216,1217,1267>,
  <1268,1267,1217>,
  <1217,1218,1268>,
  <1269,1268,1218>,
  <1218,1219,1269>,
  <1270,1269,1219>,
  <1219,1220,1270>,
  <1271,1270,1220>,
  <1220,1221,1271>,
  <1272,1271,1221>,
  <1221,1222,1272>,
  <1273,1272,1222>,
  <1222,1223,1273>,
  <1274,1273,1223>,
  <1224,1225,1275>,
  <1276,1275,1225>,
  <1225,1226,1276>,
  <1277,1276,1226>,
  <1226,1227,1277>,
  <1278,1277,1227>,
  <1227,1228,1278>,
  <1279,1278,1228>,
  <1228,1229,1279>,
  <1280,1279,1229>,
  <1229,1230,1280>,
  <1281,1280,1230>,
  <1230,1231,1281>,
  <1282,1281,1231>,
  <1231,1232,1282>,
  <1283,1282,1232>,
  <1232,1233,1283>,
  <1284,1283,1233>,
  <1233,1234,1284>,
  <1285,1284,1234>,
  <1234,1235,1285>,
  <1286,1285,1235>,
  <1235,1236,1286>,
  <1287,1286,1236>,
  <1236,1237,1287>,
  <1288,1287,1237>,
  <1237,1238,1288>,
  <1289,1288,1238>,
  <1238,1239,1289>,
  <1290,1289,1239>,
  <1239,1240,1290>,
  <1291,1290,1240>,
  <1240,1241,1291>,
  <1292,1291,1241>,
  <1241,1242,1292>,
  <1293,1292,1242>,
  <1242,1243,1293>,
  <1294,1293,1243>,
  <1243,1244,1294>,
  <1295,1294,1244>,
  <1244,1245,1295>,
  <1296,1295,1245>,
  <1245,1246,1296>,
  <1297,1296,1246>,
  <1246,1247,1297>,
  <1298,1297,1247>,
  <1247,1248,1298>,
  <1299,1298,1248>,
  <1248,1249,1299>,
  <1300,1299,1249>,
  <1249,1250,1300>,
  <1301,1300,1250>,
  <1250,1251,1301>,
  <1302,1301,1251>,
  <1251,1252,1302>,
  <1303,1302,1252>,
  <1252,1253,1303>,
  <1304,1303,1253>,
  <1253,1254,1304>,
  <1305,1304,1254>,
  <1254,1255,1305>,
  <1306,1305,1255>,
  <1255,1256,1306>,
  <1307,1306,1256>,
  <1256,1257,1307>,
  <1308,1307,1257>,
  <1257,1258,1308>,
  <1309,1308,1258>,
  <1258,1259,1309>,
  <1310,1309,1259>,
  <1259,1260,1310>,
  <1311,1310,1260>,
  <1260,1261,1311>,
  <1312,1311,1261>,
  <1261,1262,1312>,
  <1313,1312,1262>,
  <1262,1263,1313>,
  <1314,1313,1263>,
  <1263,1264,1314>,
  <1315,1314,1264>,
  <1264,1265,1315>,
  <1316,1315,1265>,
  <1265,1266,1316>,
  <1317,1316,1266>,
  <1266,1267,1317>,
  <1318,1317,1267>,
  <1267,1268,1318>,
  <1319,1318,1268>,
  <1268,1269,1319>,
  <1320,1319,1269>,
  <1269,1270,1320>,
  <1321,1320,1270>,
  <1270,1271,1321>,
  <1322,1321,1271>,
  <1271,1272,1322>,
  <1323,1322,1272>,
  <1272,1273,1323>,
  <1324,1323,1273>,
  <1273,1274,1324>,
  <1325,1324,1274>,
  <1275,1276,1326>,
  <1327,1326,1276>,
  <1276,1277,1327>,
  <1328,1327,1277>,
  <1277,1278,1328>,
  <1329,1328,1278>,
  <1278,1279,1329>,
  <1330,1329,1279>,
  <1279,1280,1330>,
  <1331,1330,1280>,
  <1280,1281,1331>,
  <1332,1331,1281>,
  <1281,1282,1332>,
  <1333,1332,1282>,
  <1282,1283,1333>,
  <1334,1333,1283>,
  <1283,1284,1334>,
  <1335,1334,1284>,
  <1284,1285,1335>,
  <1336,1335,1285>,
  <1285,1286,1336>,
  <1337,1336,1286>,
  <1286,1287,1337>,
  <1338,1337,1287>,
  <1287,1288,1338>,
  <1339,1338,1288>,
  <1288,1289,1339>,
  <1340,1339,1289>,
  <1289,1290,1340>,
  <1341,1340,1290>,
  <1290,1291,1341>,
  <1342,1341,1291>,
  <1291,1292,1342>,
  <1343,1342,1292>,
  <1292,1293,1343>,
  <1344,1343,1293>,
  <1293,1294,1344>,
  <1345,1344,1294>,
  <1294,1295,1345>,
  <1346,1345,1295>,
  <1295,1296,1346>,
  <1347,1346,1296>,
  <1296,1297,1347>,
  <1348,1347,1297>,
  <1297,1298,1348>,
  <1349,1348,1298>,
  <1298,1299,1349>,
  <1350,1349,1299>,
  <1299,1300,1350>,
  <1351,1350,1300>,
  <1300,1301,1351>,
  <1352,1351,1301>,
  <1301,1302,1352>,
  <1353,1352,1302>,
  <1302,1303,1353>,
  <1354,1353,1303>,
  <1303,1304,1354>,
  <1355,1354,1304>,
  <1304,1305,1355>,
  <1356,1355,1305>,
  <1305,1306,1356>,
  <1357,1356,1306>,
  <1306,1307,1357>,
  <1358,1357,1307>,
  <1307,1308,1358>,
  <1359,1358,1308>,
  <1308,1309,1359>,
  <1360,1359,1309>,
  <1309,1310,1360>,
  <1361,1360,1310>,
  <1310,1311,1361>,
  <1362,1361,1311>,
  <1311,1312,1362>,
  <1363,1362,1312>,
  <1312,1313,1363>,
  <1364,1363,1313>,
  <1313,1314,1364>,
  <1365,1364,1314>,
  <1314,1315,1365>,
  <1366,1365,1315>,
  <1315,1316,1366>,
  <1367,1366,1316>,
  <1316,1317,1367>,
  <1368,1367,1317>,
  <1317,1318,1368>,
  <1369,1368,1318>,
  <1318,1319,1369>,
  <1370,1369,1319>,
  <1319,1320,1370>,
  <1371,1370,1320>,
  <1320,1321,1371>,
  <1372,1371,1321>,
  <1321,1322,1372>,
  <1373,1372,1322>,
  <1322,1323,1373>,
  <1374,1373,1323>,
  <1323,1324,1374>,
  <1375,1374,1324>,
  <1324,1325,1375>,
  <1376,1375,1325>,
  <1326,1327,1377>,
  <1378,1377,1327>,
  <1327,1328,1378>,
  <1379,1378,1328>,
  <1328,1329,1379>,
  <1380,1379,1329>,
  <1329,1330,1380>,
  <1381,1380,1330>,
  <1330,1331,1381>,
  <1382,1381,1331>,
  <1331,1332,1382>,
  <1383,1382,1332>,
  <1332,1333,1383>,
  <1384,1383,1333>,
  <1333,1334,1384>,
  <1385,1384,1334>,
  <1334,1335,1385>,
  <1386,1385,1335>,
  <1335,1336,1386>,
  <1387,1386,1336>,
  <1336,1337,1387>,
  <1388,1387,1337>,
  <1337,1338,1388>,
  <1389,1388,1338>,
  <1338,1339,1389>,
  <1390,1389,1339>,
  <1339,1340,1390>,
  <1391,1390,1340>,
  <1340,1341,1391>,
  <1392,1391,1341>,
  <1341,1342,1392>,
  <1393,1392,1342>,
  <1342,1343,1393>,
  <1394,1393,1343>,
  <1343,1344,1394>,
  <1395,1394,1344>,
  <1344,1345,1395>,
  <1396,1395,1345>,
  <1345,1346,1396>,
  <1397,1396,1346>,
  <1346,1347,1397>,
  <1398,1397,1347>,
  <1347,1348,1398>,
  <1399,1398,1348>,
  <1348,1349,1399>,
  <1400,1399,1349>,
  <1349,1350,1400>,
  <1401,1400,1350>,
  <1350,1351,1401>,
  <1402,1401,1351>,
  <1351,1352,1402>,
  <1403,1402,1352>,
  <1352,1353,1403>,
  <1404,1403,1353>,
  <1353,1354,1404>,
  <1405,1404,1354>,
  <1354,1355,1405>,
  <1406,1405,1355>,
  <1355,1356,1406>,
  <1407,1406,1356>,
  <1356,1357,1407>,
  <1408,1407,1357>,
  <1357,1358,1408>,
  <1409,1408,1358>,
  <1358,1359,1409>,
  <1410,1409,1359>,
  <1359,1360,1410>,
  <1411,1410,1360>,
  <1360,1361,1411>,
  <1412,1411,1361>,
  <1361,1362,1412>,
  <1413,1412,1362>,
  <1362,1363,1413>,
  <1414,1413,1363>,
  <1363,1364,1414>,
  <1415,1414,1364>,
  <1364,1365,1415>,
  <1416,1415,1365>,
  <1365,1366,1416>,
  <1417,1416,1366>,
  <1366,1367,1417>,
  <1418,1417,1367>,
  <1367,1368,1418>,
  <1419,1418,1368>,
  <1368,1369,1419>,
  <1420,1419,1369>,
  <1369,1370,1420>,
  <1421,1420,1370>,
  <1370,1371,1421>,
  <1422,1421,1371>,
  <1371,1372,1422>,
  <1423,1422,1372>,
  <1372,1373,1423>,
  <1424,1423,1373>,
  <1373,1374,1424>,
  <1425,1424,1374>,
  <1374,1375,1425>,
  <1426,1425,1375>,
  <1375,1376,1426>,
  <1427,1426,1376>,
  <1377,1378,1428>,
  <1429,1428,1378>,
  <1378,1379,1429>,
  <1430,1429,1379>,
  <1379,1380,1430>,
  <1431,1430,1380>,
  <1380,1381,1431>,
  <1432,1431,1381>,
  <1381,1382,1432>,
  <1433,1432,1382>,
  <1382,1383,1433>,
  <1434,1433,1383>,
  <1383,1384,1434>,
  <1435,1434,1384>,
  <1384,1385,1435>,
  <1436,1435,1385>,
  <1385,1386,1436>,
  <1437,1436,1386>,
  <1386,1387,1437>,
  <1438,1437,1387>,
  <1387,1388,1438>,
  <1439,1438,1388>,
  <1388,1389,1439>,
  <1440,1439,1389>,
  <1389,1390,1440>,
  <1441,1440,1390>,
  <1390,1391,1441>,
  <1442,1441,1391>,
  <1391,1392,1442>,
  <1443,1442,1392>,
  <1392,1393,1443>,
  <1444,1443,1393>,
  <1393,1394,1444>,
  <1445,1444,1394>,
  <1394,1395,1445>,
  <1446,1445,1395>,
  <1395,1396,1446>,
  <1447,1446,1396>,
  <1396,1397,1447>,
  <1448,1447,1397>,
  <1397,1398,1448>,
  <1449,1448,1398>,
  <1398,1399,1449>,
  <1450,1449,1399>,
  <1399,1400,1450>,
  <1451,1450,1400>,
  <1400,1401,1451>,
  <1452,1451,1401>,
  <1401,1402,1452>,
  <1453,1452,1402>,
  <1402,1403,1453>,
  <1454,1453,1403>,
  <1403,1404,1454>,
  <1455,1454,1404>,
  <1404,1405,1455>,
  <1456,1455,1405>,
  <1405,1406,1456>,
  <1457,1456,1406>,
  <1406,1407,1457>,
  <1458,1457,1407>,
  <1407,1408,1458>,
  <1459,1458,1408>,
  <1408,1409,1459>,
  <1460,1459,1409>,
  <1409,1410,1460>,
  <1461,1460,1410>,
  <1410,1411,1461>,
  <1462,1461,1411>,
  <1411,1412,1462>,
  <1463,1462,1412>,
  <1412,1413,1463>,
  <1464,1463,1413>,
  <1413,1414,1464>,
  <1465,1464,1414>,
  <1414,1415,1465>,
  <1466,1465,1415>,
  <1415,1416,1466>,
  <1467,1466,1416>,
  <1416,1417,1467>,
  <1468,1467,1417>,
  <1417,1418,1468>,
  <1469,1468,1418>,
  <1418,1419,1469>,
  <1470,1469,1419>,
  <1419,1420,1470>,
  <1471,1470,1420>,
  <1420,1421,1471>,
  <1472,1471,1421>,
  <1421,1422,1472>,
  <1473,1472,1422>,
  <1422,1423,1473>,
  <1474,1473,1423>,
  <1423,1424,1474>,
  <1475,1474,1424>,
  <1424,1425,1475>,
  <1476,1475,1425>,
  <1425,1426,1476>,
  <1477,1476,1426>,
  <1426,1427,1477>,
  <1478,1477,1427>,
  <1428,1429,1479>,
  <1480,1479,1429>,
  <1429,1430,1480>,
  <1481,1480,1430>,
  <1430,1431,1481>,
  <1482,1481,1431>,
  <1431,1432,1482>,
  <1483,1482,1432>,
  <1432,1433,1483>,
  <1484,1483,1433>,
  <1433,1434,1484>,
  <1485,1484,1434>,
  <1434,1435,1485>,
  <1486,1485,1435>,
  <1435,1436,1486>,
  <1487,1486,1436>,
  <1436,1437,1487>,
  <1488,1487,1437>,
  <1437,1438,1488>,
  <1489,1488,1438>,
  <1438,1439,1489>,
  <1490,1489,1439>,
  <1439,1440,1490>,
  <1491,1490,1440>,
  <1440,1441,1491>,
  <1492,1491,1441>,
  <1441,1442,1492>,
  <1493,1492,1442>,
  <1442,1443,1493>,
  <1494,1493,1443>,
  <1443,1444,1494>,
  <1495,1494,1444>,
  <1444,1445,1495>,
  <1496,1495,1445>,
  <1445,1446,1496>,
  <1497,1496,1446>,
  <1446,1447,1497>,
  <1498,1497,1447>,
  <1447,1448,1498>,
  <1499,1498,1448>,
  <1448,1449,1499>,
  <1500,1499,1449>,
  <1449,1450,1500>,
  <1501,1500,1450>,
  <1450,1451,1501>,
  <1502,1501,1451>,
  <1451,1452,1502>,
  <1503,1502,1452>,
  <1452,1453,1503>,
  <1504,1503,1453>,
  <1453,1454,1504>,
  <1505,1504,1454>,
  <1454,1455,1505>,
  <1506,1505,1455>,
  <1455,1456,1506>,
  <1507,1506,1456>,
  <1456,1457,1507>,
  <1508,1507,1457>,
  <1457,1458,1508>,
  <1509,1508,1458>,
  <1458,1459,1509>,
  <1510,1509,1459>,
  <1459,1460,1510>,
  <1511,1510,1460>,
  <1460,1461,1511>,
  <1512,1511,1461>,
  <1461,1462,1512>,
  <1513,1512,1462>,
  <1462,1463,1513>,
  <1514,1513,1463>,
  <1463,1464,1514>,
  <1515,1514,1464>,
  <1464,1465,1515>,
  <1516,1515,1465>,
  <1465,1466,1516>,
  <1517,1516,1466>,
  <1466,1467,1517>,
  <1518,1517,1467>,
  <1467,1468,1518>,
  <1519,1518,1468>,
  <1468,1469,1519>,
  <1520,1519,1469>,
  <1469,1470,1520>,
  <1521,1520,1470>,
  <1470,1471,1521>,
  <1522,1521,1471>,
  <1471,1472,1522>,
  <1523,1522,1472>,
  <1472,1473,1523>,
  <1524,1523,1473>,
  <1473,1474,1524>,
  <1525,1524,1474>,
  <1474,1475,1525>,
  <1526,1525,1475>,
  <1475,1476,1526>,
  <1527,1526,1476>,
  <1476,1477,1527>,
  <1528,1527,1477>,
  <1477,1478,1528>,
  <1529,1528,1478>,
  <1479,1480,1530>,
  <1531,1530,1480>,
  <1480,1481,1531>,
  <1532,1531,1481>,
  <1481,1482,1532>,
  <1533,1532,1482>,
  <1482,1483,1533>,
  <1534,1533,1483>,
  <1483,1484,1534>,
  <1535,1534,1484>,
  <1484,1485,1535>,
  <1536,1535,1485>,
  <1485,1486,1536>,
  <1537,1536,1486>,
  <1486,1487,1537>,
  <1538,1537,1487>,
  <1487,1488,1538>,
  <1539,1538,1488>,
  <1488,1489,1539>,
  <1540,1539,1489>,
  <1489,1490,1540>,
  <1541,1540,1490>,
  <1490,1491,1541>,
  <1542,1541,1491>,
  <1491,1492,1542>,
  <1543,1542,1492>,
  <1492,1493,1543>,
  <1544,1543,1493>,
  <1493,1494,1544>,
  <1545,1544,1494>,
  <1494,1495,1545>,
  <1546,1545,1495>,
  <1495,1496,1546>,
  <1547,1546,1496>,
  <1496,1497,1547>,
  <1548,1547,1497>,
  <1497,1498,1548>,
  <1549,1548,1498>,
  <1498,1499,1549>,
  <1550,1549,1499>,
  <1499,1500,1550>,
  <1551,1550,1500>,
  <1500,1501,1551>,
  <1552,1551,1501>,
  <1501,1502,1552>,
  <1553,1552,1502>,
  <1502,1503,1553>,
  <1554,1553,1503>,
  <1503,1504,1554>,
  <1555,1554,1504>,
  <1504,1505,1555>,
  <1556,1555,1505>,
  <1505,1506,1556>,
  <1557,1556,1506>,
  <1506,1507,1557>,
  <1558,1557,1507>,
  <1507,1508,1558>,
  <1559,1558,1508>,
  <1508,1509,1559>,
  <1560,1559,1509>,
  <1509,1510,1560>,
  <1561,1560,1510>,
  <1510,1511,1561>,
  <1562,1561,1511>,
  <1511,1512,1562>,
  <1563,1562,1512>,
  <1512,1513,1563>,
  <1564,1563,1513>,
  <1513,1514,1564>,
  <1565,1564,1514>,
  <1514,1515,1565>,
  <1566,1565,1515>,
  <1515,1516,1566>,
  <1567,1566,1516>,
  <1516,1517,1567>,
  <1568,1567,1517>,
  <1517,1518,1568>,
  <1569,1568,1518>,
  <1518,1519,1569>,
  <1570,1569,1519>,
  <1519,1520,1570>,
  <1571,1570,1520>,
  <1520,1521,1571>,
  <1572,1571,1521>,
  <1521,1522,1572>,
  <1573,1572,1522>,
  <1522,1523,1573>,
  <1574,1573,1523>,
  <1523,1524,1574>,
  <1575,1574,1524>,
  <1524,1525,1575>,
  <1576,1575,1525>,
  <1525,1526,1576>,
  <1577,1576,1526>,
  <1526,1527,1577>,
  <1578,1577,1527>,
  <1527,1528,1578>,
  <1579,1578,1528>,
  <1528,1529,1579>,
  <1580,1579,1529>,
  <1530,1531,1581>,
  <1582,1581,1531>,
  <1531,1532,1582>,
  <1583,1582,1532>,
  <1532,1533,1583>,
  <1584,1583,1533>,
  <1533,1534,1584>,
  <1585,1584,1534>,
  <1534,1535,1585>,
  <1586,1585,1535>,
  <1535,1536,1586>,
  <1587,1586,1536>,
  <1536,1537,1587>,
  <1588,1587,1537>,
  <1537,1538,1588>,
  <1589,1588,1538>,
  <1538,1539,1589>,
  <1590,1589,1539>,
  <1539,1540,1590>,
  <1591,1590,1540>,
  <1540,1541,1591>,
  <1592,1591,1541>,
  <1541,1542,1592>,
  <1593,1592,1542>,
  <1542,1543,1593>,
  <1594,1593,1543>,
  <1543,1544,1594>,
  <1595,1594,1544>,
  <1544,1545,1595>,
  <1596,1595,1545>,
  <1545,1546,1596>,
  <1597,1596,1546>,
  <1546,1547,1597>,
  <1598,1597,1547>,
  <1547,1548,1598>,
  <1599,1598,1548>,
  <1548,1549,1599>,
  <1600,1599,1549>,
  <1549,1550,1600>,
  <1601,1600,1550>,
  <1550,1551,1601>,
  <1602,1601,1551>,
  <1551,1552,1602>,
  <1603,1602,1552>,
  <1552,1553,1603>,
  <1604,1603,1553>,
  <1553,1554,1604>,
  <1605,1604,1554>,
  <1554,1555,1605>,
  <1606,1605,1555>,
  <1555,1556,1606>,
  <1607,1606,1556>,
  <1556,1557,1607>,
  <1608,1607,1557>,
  <1557,1558,1608>,
  <1609,1608,1558>,
  <1558,1559,1609>,
  <1610,1609,1559>,
  <1559,1560,1610>,
  <1611,1610,1560>,
  <1560,1561,1611>,
  <1612,1611,1561>,
  <1561,1562,1612>,
  <1613,1612,1562>,
  <1562,1563,1613>,
  <1614,1613,1563>,
  <1563,1564,1614>,
  <1615,1614,1564>,
  <1564,1565,1615>,
  <1616,1615,1565>,
  <1565,1566,1616>,
  <1617,1616,1566>,
  <1566,1567,1617>,
  <1618,1617,1567>,
  <1567,1568,1618>,
  <1619,1618,1568>,
  <1568,1569,1619>,
  <1620,1619,1569>,
  <1569,1570,1620>,
  <1621,1620,1570>,
  <1570,1571,1621>,
  <1622,1621,1571>,
  <1571,1572,1622>,
  <1623,1622,1572>,
  <1572,1573,1623>,
  <1624,1623,1573>,
  <1573,1574,1624>,
  <1625,1624,1574>,
  <1574,1575,1625>,
  <1626,1625,1575>,
  <1575,1576,1626>,
  <1627,1626,1576>,
  <1576,1577,1627>,
  <1628,1627,1577>,
  <1577,1578,1628>,
  <1629,1628,1578>,
  <1578,1579,1629>,
  <1630,1629,1579>,
  <1579,1580,1630>,
  <1631,1630,1580>,
  <1581,1582,1632>,
  <1633,1632,1582>,
  <1582,1583,1633>,
  <1634,1633,1583>,
  <1583,1584,1634>,
  <1635,1634,1584>,
  <1584,1585,1635>,
  <1636,1635,1585>,
  <1585,1586,1636>,
  <1637,1636,1586>,
  <1586,1587,1637>,
  <1638,1637,1587>,
  <1587,1588,1638>,
  <1639,1638,1588>,
  <1588,1589,1639>,
  <1640,1639,1589>,
  <1589,1590,1640>,
  <1641,1640,1590>,
  <1590,1591,1641>,
  <1642,1641,1591>,
  <1591,1592,1642>,
  <1643,1642,1592>,
  <1592,1593,1643>,
  <1644,1643,1593>,
  <1593,1594,1644>,
  <1645,1644,1594>,
  <1594,1595,1645>,
  <1646,1645,1595>,
  <1595,1596,1646>,
  <1647,1646,1596>,
  <1596,1597,1647>,
  <1648,1647,1597>,
  <1597,1598,1648>,
  <1649,1648,1598>,
  <1598,1599,1649>,
  <1650,1649,1599>,
  <1599,1600,1650>,
  <1651,1650,1600>,
  <1600,1601,1651>,
  <1652,1651,1601>,
  <1601,1602,1652>,
  <1653,1652,1602>,
  <1602,1603,1653>,
  <1654,1653,1603>,
  <1603,1604,1654>,
  <1655,1654,1604>,
  <1604,1605,1655>,
  <1656,1655,1605>,
  <1605,1606,1656>,
  <1657,1656,1606>,
  <1606,1607,1657>,
  <1658,1657,1607>,
  <1607,1608,1658>,
  <1659,1658,1608>,
  <1608,1609,1659>,
  <1660,1659,1609>,
  <1609,1610,1660>,
  <1661,1660,1610>,
  <1610,1611,1661>,
  <1662,1661,1611>,
  <1611,1612,1662>,
  <1663,1662,1612>,
  <1612,1613,1663>,
  <1664,1663,1613>,
  <1613,1614,1664>,
  <1665,1664,1614>,
  <1614,1615,1665>,
  <1666,1665,1615>,
  <1615,1616,1666>,
  <1667,1666,1616>,
  <1616,1617,1667>,
  <1668,1667,1617>,
  <1617,1618,1668>,
  <1669,1668,1618>,
  <1618,1619,1669>,
  <1670,1669,1619>,
  <1619,1620,1670>,
  <1671,1670,1620>,
  <1620,1621,1671>,
  <1672,1671,1621>,
  <1621,1622,1672>,
  <1673,1672,1622>,
  <1622,1623,1673>,
  <1674,1673,1623>,
  <1623,1624,1674>,
  <1675,1674,1624>,
  <1624,1625,1675>,
  <1676,1675,1625>,
  <1625,1626,1676>,
  <1677,1676,1626>,
  <1626,1627,1677>,
  <1678,1677,1627>,
  <1627,1628,1678>,
  <1679,1678,1628>,
  <1628,1629,1679>,
  <1680,1679,1629>,
  <1629,1630,1680>,
  <1681,1680,1630>,
  <1630,1631,1681>,
  <1682,1681,1631>,
  <1632,1633,1683>,
  <1684,1683,1633>,
  <1633,1634,1684>,
  <1685,1684,1634>,
  <1634,1635,1685>,
  <1686,1685,1635>,
  <1635,1636,1686>,
  <1687,1686,1636>,
  <1636,1637,1687>,
  <1688,1687,1637>,
  <1637,1638,1688>,
  <1689,1688,1638>,
  <1638,1639,1689>,
  <1690,1689,1639>,
  <1639,1640,1690>,
  <1691,1690,1640>,
  <1640,1641,1691>,
  <1692,1691,1641>,
  <1641,1642,1692>,
  <1693,1692,1642>,
  <1642,1643,1693>,
  <1694,1693,1643>,
  <1643,1644,1694>,
  <1695,1694,1644>,
  <1644,1645,1695>,
  <1696,1695,1645>,
  <1645,1646,1696>,
  <1697,1696,1646>,
  <1646,1647,1697>,
  <1698,1697,1647>,
  <1647,1648,1698>,
  <1699,1698,1648>,
  <1648,1649,1699>,
  <1700,1699,1649>,
  <1649,1650,1700>,
  <1701,1700,1650>,
  <1650,1651,1701>,
  <1702,1701,1651>,
  <1651,1652,1702>,
  <1703,1702,1652>,
  <1652,1653,1703>,
  <1704,1703,1653>,
  <1653,1654,1704>,
  <1705,1704,1654>,
  <1654,1655,1705>,
  <1706,1705,1655>,
  <1655,1656,1706>,
  <1707,1706,1656>,
  <1656,1657,1707>,
  <1708,1707,1657>,
  <1657,1658,1708>,
  <1709,1708,1658>,
  <1658,1659,1709>,
  <1710,1709,1659>,
  <1659,1660,1710>,
  <1711,1710,1660>,
  <1660,1661,1711>,
  <1712,1711,1661>,
  <1661,1662,1712>,
  <1713,1712,1662>,
  <1662,1663,1713>,
  <1714,1713,1663>,
  <1663,1664,1714>,
  <1715,1714,1664>,
  <1664,1665,1715>,
  <1716,1715,1665>,
  <1665,1666,1716>,
  <1717,1716,1666>,
  <1666,1667,1717>,
  <1718,1717,1667>,
  <1667,1668,1718>,
  <1719,1718,1668>,
  <1668,1669,1719>,
  <1720,1719,1669>,
  <1669,1670,1720>,
  <1721,1720,1670>,
  <1670,1671,1721>,
  <1722,1721,1671>,
  <1671,1672,1722>,
  <1723,1722,1672>,
  <1672,1673,1723>,
  <1724,1723,1673>,
  <1673,1674,1724>,
  <1725,1724,1674>,
  <1674,1675,1725>,
  <1726,1725,1675>,
  <1675,1676,1726>,
  <1727,1726,1676>,
  <1676,1677,1727>,
  <1728,1727,1677>,
  <1677,1678,1728>,
  <1729,1728,1678>,
  <1678,1679,1729>,
  <1730,1729,1679>,
  <1679,1680,1730>,
  <1731,1730,1680>,
  <1680,1681,1731>,
  <1732,1731,1681>,
  <1681,1682,1732>,
  <1733,1732,1682>,
  <1683,1684,1734>,
  <1735,1734,1684>,
  <1684,1685,1735>,
  <1736,1735,1685>,
  <1685,1686,1736>,
  <1737,1736,1686>,
  <1686,1687,1737>,
  <1738,1737,1687>,
  <1687,1688,1738>,
  <1739,1738,1688>,
  <1688,1689,1739>,
  <1740,1739,1689>,
  <1689,1690,1740>,
  <1741,1740,1690>,
  <1690,1691,1741>,
  <1742,1741,1691>,
  <1691,1692,1742>,
  <1743,1742,1692>,
  <1692,1693,1743>,
  <1744,1743,1693>,
  <1693,1694,1744>,
  <1745,1744,1694>,
  <1694,1695,1745>,
  <1746,1745,1695>,
  <1695,1696,1746>,
  <1747,1746,1696>,
  <1696,1697,1747>,
  <1748,1747,1697>,
  <1697,1698,1748>,
  <1749,1748,1698>,
  <1698,1699,1749>,
  <1750,1749,1699>,
  <1699,1700,1750>,
  <1751,1750,1700>,
  <1700,1701,1751>,
  <1752,1751,1701>,
  <1701,1702,1752>,
  <1753,1752,1702>,
  <1702,1703,1753>,
  <1754,1753,1703>,
  <1703,1704,1754>,
  <1755,1754,1704>,
  <1704,1705,1755>,
  <1756,1755,1705>,
  <1705,1706,1756>,
  <1757,1756,1706>,
  <1706,1707,1757>,
  <1758,1757,1707>,
  <1707,1708,1758>,
  <1759,1758,1708>,
  <1708,1709,1759>,
  <1760,1759,1709>,
  <1709,1710,1760>,
  <1761,1760,1710>,
  <1710,1711,1761>,
  <1762,1761,1711>,
  <1711,1712,1762>,
  <1763,1762,1712>,
  <1712,1713,1763>,
  <1764,1763,1713>,
  <1713,1714,1764>,
  <1765,1764,1714>,
  <1714,1715,1765>,
  <1766,1765,1715>,
  <1715,1716,1766>,
  <1767,1766,1716>,
  <1716,1717,1767>,
  <1768,1767,1717>,
  <1717,1718,1768>,
  <1769,1768,1718>,
  <1718,1719,1769>,
  <1770,1769,1719>,
  <1719,1720,1770>,
  <1771,1770,1720>,
  <1720,1721,1771>,
  <1772,1771,1721>,
  <1721,1722,1772>,
  <1773,1772,1722>,
  <1722,1723,1773>,
  <1774,1773,1723>,
  <1723,1724,1774>,
  <1775,1774,1724>,
  <1724,1725,1775>,
  <1776,1775,1725>,
  <1725,1726,1776>,
  <1777,1776,1726>,
  <1726,1727,1777>,
  <1778,1777,1727>,
  <1727,1728,1778>,
  <1779,1778,1728>,
  <1728,1729,1779>,
  <1780,1779,1729>,
  <1729,1730,1780>,
  <1781,1780,1730>,
  <1730,1731,1781>,
  <1782,1781,1731>,
  <1731,1732,1782>,
  <1783,1782,1732>,
  <1732,1733,1783>,
  <1784,1783,1733>,
  <1734,1735,1785>,
  <1786,1785,1735>,
  <1735,1736,1786>,
  <1787,1786,1736>,
  <1736,1737,1787>,
  <1788,1787,1737>,
  <1737,1738,1788>,
  <1789,1788,1738>,
  <1738,1739,1789>,
  <1790,1789,1739>,
  <1739,1740,1790>,
  <1791,1790,1740>,
  <1740,1741,1791>,
  <1792,1791,1741>,
  <1741,1742,1792>,
  <1793,1792,1742>,
  <1742,1743,1793>,
  <1794,1793,1743>,
  <1743,1744,1794>,
  <1795,1794,1744>,
  <1744,1745,1795>,
  <1796,1795,1745>,
  <1745,1746,1796>,
  <1797,1796,1746>,
  <1746,1747,1797>,
  <1798,1797,1747>,
  <1747,1748,1798>,
  <1799,1798,1748>,
  <1748,1749,1799>,
  <1800,1799,1749>,
  <1749,1750,1800>,
  <1801,1800,1750>,
  <1750,1751,1801>,
  <1802,1801,1751>,
  <1751,1752,1802>,
  <1803,1802,1752>,
  <1752,1753,1803>,
  <1804,1803,1753>,
  <1753,1754,1804>,
  <1805,1804,1754>,
  <1754,1755,1805>,
  <1806,1805,1755>,
  <1755,1756,1806>,
  <1807,1806,1756>,
  <1756,1757,1807>,
  <1808,1807,1757>,
  <1757,1758,1808>,
  <1809,1808,1758>,
  <1758,1759,1809>,
  <1810,1809,1759>,
  <1759,1760,1810>,
  <1811,1810,1760>,
  <1760,1761,1811>,
  <1812,1811,1761>,
  <1761,1762,1812>,
  <1813,1812,1762>,
  <1762,1763,1813>,
  <1814,1813,1763>,
  <1763,1764,1814>,
  <1815,1814,1764>,
  <1764,1765,1815>,
  <1816,1815,1765>,
  <1765,1766,1816>,
  <1817,1816,1766>,
  <1766,1767,1817>,
  <1818,1817,1767>,
  <1767,1768,1818>,
  <1819,1818,1768>,
  <1768,1769,1819>,
  <1820,1819,1769>,
  <1769,1770,1820>,
  <1821,1820,1770>,
  <1770,1771,1821>,
  <1822,1821,1771>,
  <1771,1772,1822>,
  <1823,1822,1772>,
  <1772,1773,1823>,
  <1824,1823,1773>,
  <1773,1774,1824>,
  <1825,1824,1774>,
  <1774,1775,1825>,
  <1826,1825,1775>,
  <1775,1776,1826>,
  <1827,1826,1776>,
  <1776,1777,1827>,
  <1828,1827,1777>,
  <1777,1778,1828>,
  <1829,1828,1778>,
  <1778,1779,1829>,
  <1830,1829,1779>,
  <1779,1780,1830>,
  <1831,1830,1780>,
  <1780,1781,1831>,
  <1832,1831,1781>,
  <1781,1782,1832>,
  <1833,1832,1782>,
  <1782,1783,1833>,
  <1834,1833,1783>,
  <1783,1784,1834>,
  <1835,1834,1784>,
  <1785,1786,1836>,
  <1837,1836,1786>,
  <1786,1787,1837>,
  <1838,1837,1787>,
  <1787,1788,1838>,
  <1839,1838,1788>,
  <1788,1789,1839>,
  <1840,1839,1789>,
  <1789,1790,1840>,
  <1841,1840,1790>,
  <1790,1791,1841>,
  <1842,1841,1791>,
  <1791,1792,1842>,
  <1843,1842,1792>,
  <1792,1793,1843>,
  <1844,1843,1793>,
  <1793,1794,1844>,
  <1845,1844,1794>,
  <1794,1795,1845>,
  <1846,1845,1795>,
  <1795,1796,1846>,
  <1847,1846,1796>,
  <1796,1797,1847>,
  <1848,1847,1797>,
  <1797,1798,1848>,
  <1849,1848,1798>,
  <1798,1799,1849>,
  <1850,1849,1799>,
  <1799,1800,1850>,
  <1851,1850,1800>,
  <1800,1801,1851>,
  <1852,1851,1801>,
  <1801,1802,1852>,
  <1853,1852,1802>,
  <1802,1803,1853>,
  <1854,1853,1803>,
  <1803,1804,1854>,
  <1855,1854,1804>,
  <1804,1805,1855>,
  <1856,1855,1805>,
  <1805,1806,1856>,
  <1857,1856,1806>,
  <1806,1807,1857>,
  <1858,1857,1807>,
  <1807,1808,1858>,
  <1859,1858,1808>,
  <1808,1809,1859>,
  <1860,1859,1809>,
  <1809,1810,1860>,
  <1861,1860,1810>,
  <1810,1811,1861>,
  <1862,1861,1811>,
  <1811,1812,1862>,
  <1863,1862,1812>,
  <1812,1813,1863>,
  <1864,1863,1813>,
  <1813,1814,1864>,
  <1865,1864,1814>,
  <1814,1815,1865>,
  <1866,1865,1815>,
  <1815,1816,1866>,
  <1867,1866,1816>,
  <1816,1817,1867>,
  <1868,1867,1817>,
  <1817,1818,1868>,
  <1869,1868,1818>,
  <1818,1819,1869>,
  <1870,1869,1819>,
  <1819,1820,1870>,
  <1871,1870,1820>,
  <1820,1821,1871>,
  <1872,1871,1821>,
  <1821,1822,1872>,
  <1873,1872,1822>,
  <1822,1823,1873>,
  <1874,1873,1823>,
  <1823,1824,1874>,
  <1875,1874,1824>,
  <1824,1825,1875>,
  <1876,1875,1825>,
  <1825,1826,1876>,
  <1877,1876,1826>,
  <1826,1827,1877>,
  <1878,1877,1827>,
  <1827,1828,1878>,
  <1879,1878,1828>,
  <1828,1829,1879>,
  <1880,1879,1829>,
  <1829,1830,1880>,
  <1881,1880,1830>,
  <1830,1831,1881>,
  <1882,1881,1831>,
  <1831,1832,1882>,
  <1883,1882,1832>,
  <1832,1833,1883>,
  <1884,1883,1833>,
  <1833,1834,1884>,
  <1885,1884,1834>,
  <1834,1835,1885>,
  <1886,1885,1835>,
  <1836,1837,1887>,
  <1888,1887,1837>,
  <1837,1838,1888>,
  <1889,1888,1838>,
  <1838,1839,1889>,
  <1890,1889,1839>,
  <1839,1840,1890>,
  <1891,1890,1840>,
  <1840,1841,1891>,
  <1892,1891,1841>,
  <1841,1842,1892>,
  <1893,1892,1842>,
  <1842,1843,1893>,
  <1894,1893,1843>,
  <1843,1844,1894>,
  <1895,1894,1844>,
  <1844,1845,1895>,
  <1896,1895,1845>,
  <1845,1846,1896>,
  <1897,1896,1846>,
  <1846,1847,1897>,
  <1898,1897,1847>,
  <1847,1848,1898>,
  <1899,1898,1848>,
  <1848,1849,1899>,
  <1900,1899,1849>,
  <1849,1850,1900>,
  <1901,1900,1850>,
  <1850,1851,1901>,
  <1902,1901,1851>,
  <1851,1852,1902>,
  <1903,1902,1852>,
  <1852,1853,1903>,
  <1904,1903,1853>,
  <1853,1854,1904>,
  <1905,1904,1854>,
  <1854,1855,1905>,
  <1906,1905,1855>,
  <1855,1856,1906>,
  <1907,1906,1856>,
  <1856,1857,1907>,
  <1908,1907,1857>,
  <1857,1858,1908>,
  <1909,1908,1858>,
  <1858,1859,1909>,
  <1910,1909,1859>,
  <1859,1860,1910>,
  <1911,1910,1860>,
  <1860,1861,1911>,
  <1912,1911,1861>,
  <1861,1862,1912>,
  <1913,1912,1862>,
  <1862,1863,1913>,
  <1914,1913,1863>,
  <1863,1864,1914>,
  <1915,1914,1864>,
  <1864,1865,1915>,
  <1916,1915,1865>,
  <1865,1866,1916>,
  <1917,1916,1866>,
  <1866,1867,1917>,
  <1918,1917,1867>,
  <1867,1868,1918>,
  <1919,1918,1868>,
  <1868,1869,1919>,
  <1920,1919,1869>,
  <1869,1870,1920>,
  <1921,1920,1870>,
  <1870,1871,1921>,
  <1922,1921,1871>,
  <1871,1872,1922>,
  <1923,1922,1872>,
  <1872,1873,1923>,
  <1924,1923,1873>,
  <1873,1874,1924>,
  <1925,1924,1874>,
  <1874,1875,1925>,
  <1926,1925,1875>,
  <1875,1876,1926>,
  <1927,1926,1876>,
  <1876,1877,1927>,
  <1928,1927,1877>,
  <1877,1878,1928>,
  <1929,1928,1878>,
  <1878,1879,1929>,
  <1930,1929,1879>,
  <1879,1880,1930>,
  <1931,1930,1880>,
  <1880,1881,1931>,
  <1932,1931,1881>,
  <1881,1882,1932>,
  <1933,1932,1882>,
  <1882,1883,1933>,
  <1934,1933,1883>,
  <1883,1884,1934>,
  <1935,1934,1884>,
  <1884,1885,1935>,
  <1936,1935,1885>,
  <1885,1886,1936>,
  <1937,1936,1886>,
  <1887,1888,1938>,
  <1939,1938,1888>,
  <1888,1889,1939>,
  <1940,1939,1889>,
  <1889,1890,1940>,
  <1941,1940,1890>,
  <1890,1891,1941>,
  <1942,1941,1891>,
  <1891,1892,1942>,
  <1943,1942,1892>,
  <1892,1893,1943>,
  <1944,1943,1893>,
  <1893,1894,1944>,
  <1945,1944,1894>,
  <1894,1895,1945>,
  <1946,1945,1895>,
  <1895,1896,1946>,
  <1947,1946,1896>,
  <1896,1897,1947>,
  <1948,1947,1897>,
  <1897,1898,1948>,
  <1949,1948,1898>,
  <1898,1899,1949>,
  <1950,1949,1899>,
  <1899,1900,1950>,
  <1951,1950,1900>,
  <1900,1901,1951>,
  <1952,1951,1901>,
  <1901,1902,1952>,
  <1953,1952,1902>,
  <1902,1903,1953>,
  <1954,1953,1903>,
  <1903,1904,1954>,
  <1955,1954,1904>,
  <1904,1905,1955>,
  <1956,1955,1905>,
  <1905,1906,1956>,
  <1957,1956,1906>,
  <1906,1907,1957>,
  <1958,1957,1907>,
  <1907,1908,1958>,
  <1959,1958,1908>,
  <1908,1909,1959>,
  <1960,1959,1909>,
  <1909,1910,1960>,
  <1961,1960,1910>,
  <1910,1911,1961>,
  <1962,1961,1911>,
  <1911,1912,1962>,
  <1963,1962,1912>,
  <1912,1913,1963>,
  <1964,1963,1913>,
  <1913,1914,1964>,
  <1965,1964,1914>,
  <1914,1915,1965>,
  <1966,1965,1915>,
  <1915,1916,1966>,
  <1967,1966,1916>,
  <1916,1917,1967>,
  <1968,1967,1917>,
  <1917,1918,1968>,
  <1969,1968,1918>,
  <1918,1919,1969>,
  <1970,1969,1919>,
  <1919,1920,1970>,
  <1971,1970,1920>,
  <1920,1921,1971>,
  <1972,1971,1921>,
  <1921,1922,1972>,
  <1973,1972,1922>,
  <1922,1923,1973>,
  <1974,1973,1923>,
  <1923,1924,1974>,
  <1975,1974,1924>,
  <1924,1925,1975>,
  <1976,1975,1925>,
  <1925,1926,1976>,
  <1977,1976,1926>,
  <1926,1927,1977>,
  <1978,1977,1927>,
  <1927,1928,1978>,
  <1979,1978,1928>,
  <1928,1929,1979>,
  <1980,1979,1929>,
  <1929,1930,1980>,
  <1981,1980,1930>,
  <1930,1931,1981>,
  <1982,1981,1931>,
  <1931,1932,1982>,
  <1983,1982,1932>,
  <1932,1933,1983>,
  <1984,1983,1933>,
  <1933,1934,1984>,
  <1985,1984,1934>,
  <1934,1935,1985>,
  <1986,1985,1935>,
  <1935,1936,1986>,
  <1987,1986,1936>,
  <1936,1937,1987>,
  <1988,1987,1937>,
  <1938,1939,1989>,
  <1990,1989,1939>,
  <1939,1940,1990>,
  <1991,1990,1940>,
  <1940,1941,1991>,
  <1992,1991,1941>,
  <1941,1942,1992>,
  <1993,1992,1942>,
  <1942,1943,1993>,
  <1994,1993,1943>,
  <1943,1944,1994>,
  <1995,1994,1944>,
  <1944,1945,1995>,
  <1996,1995,1945>,
  <1945,1946,1996>,
  <1997,1996,1946>,
  <1946,1947,1997>,
  <1998,1997,1947>,
  <1947,1948,1998>,
  <1999,1998,1948>,
  <1948,1949,1999>,
  <2000,1999,1949>,
  <1949,1950,2000>,
  <2001,2000,1950>,
  <1950,1951,2001>,
  <2002,2001,1951>,
  <1951,1952,2002>,
  <2003,2002,1952>,
  <1952,1953,2003>,
  <2004,2003,1953>,
  <1953,1954,2004>,
  <2005,2004,1954>,
  <1954,1955,2005>,
  <2006,2005,1955>,
  <1955,1956,2006>,
  <2007,2006,1956>,
  <1956,1957,2007>,
  <2008,2007,1957>,
  <1957,1958,2008>,
  <2009,2008,1958>,
  <1958,1959,2009>,
  <2010,2009,1959>,
  <1959,1960,2010>,
  <2011,2010,1960>,
  <1960,1961,2011>,
  <2012,2011,1961>,
  <1961,1962,2012>,
  <2013,2012,1962>,
  <1962,1963,2013>,
  <2014,2013,1963>,
  <1963,1964,2014>,
  <2015,2014,1964>,
  <1964,1965,2015>,
  <2016,2015,1965>,
  <1965,1966,2016>,
  <2017,2016,1966>,
  <1966,1967,2017>,
  <2018,2017,1967>,
  <1967,1968,2018>,
  <2019,2018,1968>,
  <1968,1969,2019>,
  <2020,2019,1969>,
  <1969,1970,2020>,
  <2021,2020,1970>,
  <1970,1971,2021>,
  <2022,2021,1971>,
  <1971,1972,2022>,
  <2023,2022,1972>,
  <1972,1973,2023>,
  <2024,2023,1973>,
  <1973,1974,2024>,
  <2025,2024,1974>,
  <1974,1975,2025>,
  <2026,2025,1975>,
  <1975,1976,2026>,
  <2027,2026,1976>,
  <1976,1977,2027>,
  <2028,2027,1977>,
  <1977,1978,2028>,
  <2029,2028,1978>,
  <1978,1979,2029>,
  <2030,2029,1979>,
  <1979,1980,2030>,
  <2031,2030,1980>,
  <1980,1981,2031>,
  <2032,2031,1981>,
  <1981,1982,2032>,
  <2033,2032,1982>,
  <1982,1983,2033>,
  <2034,2033,1983>,
  <1983,1984,2034>,
  <2035,2034,1984>,
  <1984,1985,2035>,
  <2036,2035,1985>,
  <1985,1986,2036>,
  <2037,2036,1986>,
  <1986,1987,2037>,
  <2038,2037,1987>,
  <1987,1988,2038>,
  <2039,2038,1988>,
  <1989,1990,2040>,
  <2041,2040,1990>,
  <1990,1991,2041>,
  <2042,2041,1991>,
  <1991,1992,2042>,
  <2043,2042,1992>,
  <1992,1993,2043>,
  <2044,2043,1993>,
  <1993,1994,2044>,
  <2045,2044,1994>,
  <1994,1995,2045>,
  <2046,2045,1995>,
  <1995,1996,2046>,
  <2047,2046,1996>,
  <1996,1997,2047>,
  <2048,2047,1997>,
  <1997,1998,2048>,
  <2049,2048,1998>,
  <1998,1999,2049>,
  <2050,2049,1999>,
  <1999,2000,2050>,
  <2051,2050,2000>,
  <2000,2001,2051>,
  <2052,2051,2001>,
  <2001,2002,2052>,
  <2053,2052,2002>,
  <2002,2003,2053>,
  <2054,2053,2003>,
  <2003,2004,2054>,
  <2055,2054,2004>,
  <2004,2005,2055>,
  <2056,2055,2005>,
  <2005,2006,2056>,
  <2057,2056,2006>,
  <2006,2007,2057>,
  <2058,2057,2007>,
  <2007,2008,2058>,
  <2059,2058,2008>,
  <2008,2009,2059>,
  <2060,2059,2009>,
  <2009,2010,2060>,
  <2061,2060,2010>,
  <2010,2011,2061>,
  <2062,2061,2011>,
  <2011,2012,2062>,
  <2063,2062,2012>,
  <2012,2013,2063>,
  <2064,2063,2013>,
  <2013,2014,2064>,
  <2065,2064,2014>,
  <2014,2015,2065>,
  <2066,2065,2015>,
  <2015,2016,2066>,
  <2067,2066,2016>,
  <2016,2017,2067>,
  <2068,2067,2017>,
  <2017,2018,2068>,
  <2069,2068,2018>,
  <2018,2019,2069>,
  <2070,2069,2019>,
  <2019,2020,2070>,
  <2071,2070,2020>,
  <2020,2021,2071>,
  <2072,2071,2021>,
  <2021,2022,2072>,
  <2073,2072,2022>,
  <2022,2023,2073>,
  <2074,2073,2023>,
  <2023,2024,2074>,
  <2075,2074,2024>,
  <2024,2025,2075>,
  <2076,2075,2025>,
  <2025,2026,2076>,
  <2077,2076,2026>,
  <2026,2027,2077>,
  <2078,2077,2027>,
  <2027,2028,2078>,
  <2079,2078,2028>,
  <2028,2029,2079>,
  <2080,2079,2029>,
  <2029,2030,2080>,
  <2081,2080,2030>,
  <2030,2031,2081>,
  <2082,2081,2031>,
  <2031,2032,2082>,
  <2083,2082,2032>,
  <2032,2033,2083>,
  <2084,2083,2033>,
  <2033,2034,2084>,
  <2085,2084,2034>,
  <2034,2035,2085>,
  <2086,2085,2035>,
  <2035,2036,2086>,
  <2087,2086,2036>,
  <2036,2037,2087>,
  <2088,2087,2037>,
  <2037,2038,2088>,
  <2089,2088,2038>,
  <2038,2039,2089>,
  <2090,2089,2039>,
  <2040,2041,2091>,
  <2092,2091,2041>,
  <2041,2042,2092>,
  <2093,2092,2042>,
  <2042,2043,2093>,
  <2094,2093,2043>,
  <2043,2044,2094>,
  <2095,2094,2044>,
  <2044,2045,2095>,
  <2096,2095,2045>,
  <2045,2046,2096>,
  <2097,2096,2046>,
  <2046,2047,2097>,
  <2098,2097,2047>,
  <2047,2048,2098>,
  <2099,2098,2048>,
  <2048,2049,2099>,
  <2100,2099,2049>,
  <2049,2050,2100>,
  <2101,2100,2050>,
  <2050,2051,2101>,
  <2102,2101,2051>,
  <2051,2052,2102>,
  <2103,2102,2052>,
  <2052,2053,2103>,
  <2104,2103,2053>,
  <2053,2054,2104>,
  <2105,2104,2054>,
  <2054,2055,2105>,
  <2106,2105,2055>,
  <2055,2056,2106>,
  <2107,2106,2056>,
  <2056,2057,2107>,
  <2108,2107,2057>,
  <2057,2058,2108>,
  <2109,2108,2058>,
  <2058,2059,2109>,
  <2110,2109,2059>,
  <2059,2060,2110>,
  <2111,2110,2060>,
  <2060,2061,2111>,
  <2112,2111,2061>,
  <2061,2062,2112>,
  <2113,2112,2062>,
  <2062,2063,2113>,
  <2114,2113,2063>,
  <2063,2064,2114>,
  <2115,2114,2064>,
  <2064,2065,2115>,
  <2116,2115,2065>,
  <2065,2066,2116>,
  <2117,2116,2066>,
  <2066,2067,2117>,
  <2118,2117,2067>,
  <2067,2068,2118>,
  <2119,2118,2068>,
  <2068,2069,2119>,
  <2120,2119,2069>,
  <2069,2070,2120>,
  <2121,2120,2070>,
  <2070,2071,2121>,
  <2122,2121,2071>,
  <2071,2072,2122>,
  <2123,2122,2072>,
  <2072,2073,2123>,
  <2124,2123,2073>,
  <2073,2074,2124>,
  <2125,2124,2074>,
  <2074,2075,2125>,
  <2126,2125,2075>,
  <2075,2076,2126>,
  <2127,2126,2076>,
  <2076,2077,2127>,
  <2128,2127,2077>,
  <2077,2078,2128>,
  <2129,2128,2078>,
  <2078,2079,2129>,
  <2130,2129,2079>,
  <2079,2080,2130>,
  <2131,2130,2080>,
  <2080,2081,2131>,
  <2132,2131,2081>,
  <2081,2082,2132>,
  <2133,2132,2082>,
  <2082,2083,2133>,
  <2134,2133,2083>,
  <2083,2084,2134>,
  <2135,2134,2084>,
  <2084,2085,2135>,
  <2136,2135,2085>,
  <2085,2086,2136>,
  <2137,2136,2086>,
  <2086,2087,2137>,
  <2138,2137,2087>,
  <2087,2088,2138>,
  <2139,2138,2088>,
  <2088,2089,2139>,
  <2140,2139,2089>,
  <2089,2090,2140>,
  <2141,2140,2090>,
  <2091,2092,2142>,
  <2143,2142,2092>,
  <2092,2093,2143>,
  <2144,2143,2093>,
  <2093,2094,2144>,
  <2145,2144,2094>,
  <2094,2095,2145>,
  <2146,2145,2095>,
  <2095,2096,2146>,
  <2147,2146,2096>,
  <2096,2097,2147>,
  <2148,2147,2097>,
  <2097,2098,2148>,
  <2149,2148,2098>,
  <2098,2099,2149>,
  <2150,2149,2099>,
  <2099,2100,2150>,
  <2151,2150,2100>,
  <2100,2101,2151>,
  <2152,2151,2101>,
  <2101,2102,2152>,
  <2153,2152,2102>,
  <2102,2103,2153>,
  <2154,2153,2103>,
  <2103,2104,2154>,
  <2155,2154,2104>,
  <2104,2105,2155>,
  <2156,2155,2105>,
  <2105,2106,2156>,
  <2157,2156,2106>,
  <2106,2107,2157>,
  <2158,2157,2107>,
  <2107,2108,2158>,
  <2159,2158,2108>,
  <2108,2109,2159>,
  <2160,2159,2109>,
  <2109,2110,2160>,
  <2161,2160,2110>,
  <2110,2111,2161>,
  <2162,2161,2111>,
  <2111,2112,2162>,
  <2163,2162,2112>,
  <2112,2113,2163>,
  <2164,2163,2113>,
  <2113,2114,2164>,
  <2165,2164,2114>,
  <2114,2115,2165>,
  <2166,2165,2115>,
  <2115,2116,2166>,
  <2167,2166,2116>,
  <2116,2117,2167>,
  <2168,2167,2117>,
  <2117,2118,2168>,
  <2169,2168,2118>,
  <2118,2119,2169>,
  <2170,2169,2119>,
  <2119,2120,2170>,
  <2171,2170,2120>,
  <2120,2121,2171>,
  <2172,2171,2121>,
  <2121,2122,2172>,
  <2173,2172,2122>,
  <2122,2123,2173>,
  <2174,2173,2123>,
  <2123,2124,2174>,
  <2175,2174,2124>,
  <2124,2125,2175>,
  <2176,2175,2125>,
  <2125,2126,2176>,
  <2177,2176,2126>,
  <2126,2127,2177>,
  <2178,2177,2127>,
  <2127,2128,2178>,
  <2179,2178,2128>,
  <2128,2129,2179>,
  <2180,2179,2129>,
  <2129,2130,2180>,
  <2181,2180,2130>,
  <2130,2131,2181>,
  <2182,2181,2131>,
  <2131,2132,2182>,
  <2183,2182,2132>,
  <2132,2133,2183>,
  <2184,2183,2133>,
  <2133,2134,2184>,
  <2185,2184,2134>,
  <2134,2135,2185>,
  <2186,2185,2135>,
  <2135,2136,2186>,
  <2187,2186,2136>,
  <2136,2137,2187>,
  <2188,2187,2137>,
  <2137,2138,2188>,
  <2189,2188,2138>,
  <2138,2139,2189>,
  <2190,2189,2139>,
  <2139,2140,2190>,
  <2191,2190,2140>,
  <2140,2141,2191>,
  <2192,2191,2141>,
  <2142,2143,2193>,
  <2194,2193,2143>,
  <2143,2144,2194>,
  <2195,2194,2144>,
  <2144,2145,2195>,
  <2196,2195,2145>,
  <2145,2146,2196>,
  <2197,2196,2146>,
  <2146,2147,2197>,
  <2198,2197,2147>,
  <2147,2148,2198>,
  <2199,2198,2148>,
  <2148,2149,2199>,
  <2200,2199,2149>,
  <2149,2150,2200>,
  <2201,2200,2150>,
  <2150,2151,2201>,
  <2202,2201,2151>,
  <2151,2152,2202>,
  <2203,2202,2152>,
  <2152,2153,2203>,
  <2204,2203,2153>,
  <2153,2154,2204>,
  <2205,2204,2154>,
  <2154,2155,2205>,
  <2206,2205,2155>,
  <2155,2156,2206>,
  <2207,2206,2156>,
  <2156,2157,2207>,
  <2208,2207,2157>,
  <2157,2158,2208>,
  <2209,2208,2158>,
  <2158,2159,2209>,
  <2210,2209,2159>,
  <2159,2160,2210>,
  <2211,2210,2160>,
  <2160,2161,2211>,
  <2212,2211,2161>,
  <2161,2162,2212>,
  <2213,2212,2162>,
  <2162,2163,2213>,
  <2214,2213,2163>,
  <2163,2164,2214>,
  <2215,2214,2164>,
  <2164,2165,2215>,
  <2216,2215,2165>,
  <2165,2166,2216>,
  <2217,2216,2166>,
  <2166,2167,2217>,
  <2218,2217,2167>,
  <2167,2168,2218>,
  <2219,2218,2168>,
  <2168,2169,2219>,
  <2220,2219,2169>,
  <2169,2170,2220>,
  <2221,2220,2170>,
  <2170,2171,2221>,
  <2222,2221,2171>,
  <2171,2172,2222>,
  <2223,2222,2172>,
  <2172,2173,2223>,
  <2224,2223,2173>,
  <2173,2174,2224>,
  <2225,2224,2174>,
  <2174,2175,2225>,
  <2226,2225,2175>,
  <2175,2176,2226>,
  <2227,2226,2176>,
  <2176,2177,2227>,
  <2228,2227,2177>,
  <2177,2178,2228>,
  <2229,2228,2178>,
  <2178,2179,2229>,
  <2230,2229,2179>,
  <2179,2180,2230>,
  <2231,2230,2180>,
  <2180,2181,2231>,
  <2232,2231,2181>,
  <2181,2182,2232>,
  <2233,2232,2182>,
  <2182,2183,2233>,
  <2234,2233,2183>,
  <2183,2184,2234>,
  <2235,2234,2184>,
  <2184,2185,2235>,
  <2236,2235,2185>,
  <2185,2186,2236>,
  <2237,2236,2186>,
  <2186,2187,2237>,
  <2238,2237,2187>,
  <2187,2188,2238>,
  <2239,2238,2188>,
  <2188,2189,2239>,
  <2240,2239,2189>,
  <2189,2190,2240>,
  <2241,2240,2190>,
  <2190,2191,2241>,
  <2242,2241,2191>,
  <2191,2192,2242>,
  <2243,2242,2192>,
  <2193,2194,2244>,
  <2245,2244,2194>,
  <2194,2195,2245>,
  <2246,2245,2195>,
  <2195,2196,2246>,
  <2247,2246,2196>,
  <2196,2197,2247>,
  <2248,2247,2197>,
  <2197,2198,2248>,
  <2249,2248,2198>,
  <2198,2199,2249>,
  <2250,2249,2199>,
  <2199,2200,2250>,
  <2251,2250,2200>,
  <2200,2201,2251>,
  <2252,2251,2201>,
  <2201,2202,2252>,
  <2253,2252,2202>,
  <2202,2203,2253>,
  <2254,2253,2203>,
  <2203,2204,2254>,
  <2255,2254,2204>,
  <2204,2205,2255>,
  <2256,2255,2205>,
  <2205,2206,2256>,
  <2257,2256,2206>,
  <2206,2207,2257>,
  <2258,2257,2207>,
  <2207,2208,2258>,
  <2259,2258,2208>,
  <2208,2209,2259>,
  <2260,2259,2209>,
  <2209,2210,2260>,
  <2261,2260,2210>,
  <2210,2211,2261>,
  <2262,2261,2211>,
  <2211,2212,2262>,
  <2263,2262,2212>,
  <2212,2213,2263>,
  <2264,2263,2213>,
  <2213,2214,2264>,
  <2265,2264,2214>,
  <2214,2215,2265>,
  <2266,2265,2215>,
  <2215,2216,2266>,
  <2267,2266,2216>,
  <2216,2217,2267>,
  <2268,2267,2217>,
  <2217,2218,2268>,
  <2269,2268,2218>,
  <2218,2219,2269>,
  <2270,2269,2219>,
  <2219,2220,2270>,
  <2271,2270,2220>,
  <2220,2221,2271>,
  <2272,2271,2221>,
  <2221,2222,2272>,
  <2273,2272,2222>,
  <2222,2223,2273>,
  <2274,2273,2223>,
  <2223,2224,2274>,
  <2275,2274,2224>,
  <2224,2225,2275>,
  <2276,2275,2225>,
  <2225,2226,2276>,
  <2277,2276,2226>,
  <2226,2227,2277>,
  <2278,2277,2227>,
  <2227,2228,2278>,
  <2279,2278,2228>,
  <2228,2229,2279>,
  <2280,2279,2229>,
  <2229,2230,2280>,
  <2281,2280,2230>,
  <2230,2231,2281>,
  <2282,2281,2231>,
  <2231,2232,2282>,
  <2283,2282,2232>,
  <2232,2233,2283>,
  <2284,2283,2233>,
  <2233,2234,2284>,
  <2285,2284,2234>,
  <2234,2235,2285>,
  <2286,2285,2235>,
  <2235,2236,2286>,
  <2287,2286,2236>,
  <2236,2237,2287>,
  <2288,2287,2237>,
  <2237,2238,2288>,
  <2289,2288,2238>,
  <2238,2239,2289>,
  <2290,2289,2239>,
  <2239,2240,2290>,
  <2291,2290,2240>,
  <2240,2241,2291>,
  <2292,2291,2241>,
  <2241,2242,2292>,
  <2293,2292,2242>,
  <2242,2243,2293>,
  <2294,2293,2243>,
  <2244,2245,2295>,
  <2296,2295,2245>,
  <2245,2246,2296>,
  <2297,2296,2246>,
  <2246,2247,2297>,
  <2298,2297,2247>,
  <2247,2248,2298>,
  <2299,2298,2248>,
  <2248,2249,2299>,
  <2300,2299,2249>,
  <2249,2250,2300>,
  <2301,2300,2250>,
  <2250,2251,2301>,
  <2302,2301,2251>,
  <2251,2252,2302>,
  <2303,2302,2252>,
  <2252,2253,2303>,
  <2304,2303,2253>,
  <2253,2254,2304>,
  <2305,2304,2254>,
  <2254,2255,2305>,
  <2306,2305,2255>,
  <2255,2256,2306>,
  <2307,2306,2256>,
  <2256,2257,2307>,
  <2308,2307,2257>,
  <2257,2258,2308>,
  <2309,2308,2258>,
  <2258,2259,2309>,
  <2310,2309,2259>,
  <2259,2260,2310>,
  <2311,2310,2260>,
  <2260,2261,2311>,
  <2312,2311,2261>,
  <2261,2262,2312>,
  <2313,2312,2262>,
  <2262,2263,2313>,
  <2314,2313,2263>,
  <2263,2264,2314>,
  <2315,2314,2264>,
  <2264,2265,2315>,
  <2316,2315,2265>,
  <2265,2266,2316>,
  <2317,2316,2266>,
  <2266,2267,2317>,
  <2318,2317,2267>,
  <2267,2268,2318>,
  <2319,2318,2268>,
  <2268,2269,2319>,
  <2320,2319,2269>,
  <2269,2270,2320>,
  <2321,2320,2270>,
  <2270,2271,2321>,
  <2322,2321,2271>,
  <2271,2272,2322>,
  <2323,2322,2272>,
  <2272,2273,2323>,
  <2324,2323,2273>,
  <2273,2274,2324>,
  <2325,2324,2274>,
  <2274,2275,2325>,
  <2326,2325,2275>,
  <2275,2276,2326>,
  <2327,2326,2276>,
  <2276,2277,2327>,
  <2328,2327,2277>,
  <2277,2278,2328>,
  <2329,2328,2278>,
  <2278,2279,2329>,
  <2330,2329,2279>,
  <2279,2280,2330>,
  <2331,2330,2280>,
  <2280,2281,2331>,
  <2332,2331,2281>,
  <2281,2282,2332>,
  <2333,2332,2282>,
  <2282,2283,2333>,
  <2334,2333,2283>,
  <2283,2284,2334>,
  <2335,2334,2284>,
  <2284,2285,2335>,
  <2336,2335,2285>,
  <2285,2286,2336>,
  <2337,2336,2286>,
  <2286,2287,2337>,
  <2338,2337,2287>,
  <2287,2288,2338>,
  <2339,2338,2288>,
  <2288,2289,2339>,
  <2340,2339,2289>,
  <2289,2290,2340>,
  <2341,2340,2290>,
  <2290,2291,2341>,
  <2342,2341,2291>,
  <2291,2292,2342>,
  <2343,2342,2292>,
  <2292,2293,2343>,
  <2344,2343,2293>,
  <2293,2294,2344>,
  <2345,2344,2294>,
  <2295,2296,2346>,
  <2347,2346,2296>,
  <2296,2297,2347>,
  <2348,2347,2297>,
  <2297,2298,2348>,
  <2349,2348,2298>,
  <2298,2299,2349>,
  <2350,2349,2299>,
  <2299,2300,2350>,
  <2351,2350,2300>,
  <2300,2301,2351>,
  <2352,2351,2301>,
  <2301,2302,2352>,
  <2353,2352,2302>,
  <2302,2303,2353>,
  <2354,2353,2303>,
  <2303,2304,2354>,
  <2355,2354,2304>,
  <2304,2305,2355>,
  <2356,2355,2305>,
  <2305,2306,2356>,
  <2357,2356,2306>,
  <2306,2307,2357>,
  <2358,2357,2307>,
  <2307,2308,2358>,
  <2359,2358,2308>,
  <2308,2309,2359>,
  <2360,2359,2309>,
  <2309,2310,2360>,
  <2361,2360,2310>,
  <2310,2311,2361>,
  <2362,2361,2311>,
  <2311,2312,2362>,
  <2363,2362,2312>,
  <2312,2313,2363>,
  <2364,2363,2313>,
  <2313,2314,2364>,
  <2365,2364,2314>,
  <2314,2315,2365>,
  <2366,2365,2315>,
  <2315,2316,2366>,
  <2367,2366,2316>,
  <2316,2317,2367>,
  <2368,2367,2317>,
  <2317,2318,2368>,
  <2369,2368,2318>,
  <2318,2319,2369>,
  <2370,2369,2319>,
  <2319,2320,2370>,
  <2371,2370,2320>,
  <2320,2321,2371>,
  <2372,2371,2321>,
  <2321,2322,2372>,
  <2373,2372,2322>,
  <2322,2323,2373>,
  <2374,2373,2323>,
  <2323,2324,2374>,
  <2375,2374,2324>,
  <2324,2325,2375>,
  <2376,2375,2325>,
  <2325,2326,2376>,
  <2377,2376,2326>,
  <2326,2327,2377>,
  <2378,2377,2327>,
  <2327,2328,2378>,
  <2379,2378,2328>,
  <2328,2329,2379>,
  <2380,2379,2329>,
  <2329,2330,2380>,
  <2381,2380,2330>,
  <2330,2331,2381>,
  <2382,2381,2331>,
  <2331,2332,2382>,
  <2383,2382,2332>,
  <2332,2333,2383>,
  <2384,2383,2333>,
  <2333,2334,2384>,
  <2385,2384,2334>,
  <2334,2335,2385>,
  <2386,2385,2335>,
  <2335,2336,2386>,
  <2387,2386,2336>,
  <2336,2337,2387>,
  <2388,2387,2337>,
  <2337,2338,2388>,
  <2389,2388,2338>,
  <2338,2339,2389>,
  <2390,2389,2339>,
  <2339,2340,2390>,
  <2391,2390,2340>,
  <2340,2341,2391>,
  <2392,2391,2341>,
  <2341,2342,2392>,
  <2393,2392,2342>,
  <2342,2343,2393>,
  <2394,2393,2343>,
  <2343,2344,2394>,
  <2395,2394,2344>,
  <2344,2345,2395>,
  <2396,2395,2345>,
  <2346,2347,2397>,
  <2398,2397,2347>,
  <2347,2348,2398>,
  <2399,2398,2348>,
  <2348,2349,2399>,
  <2400,2399,2349>,
  <2349,2350,2400>,
  <2401,2400,2350>,
  <2350,2351,2401>,
  <2402,2401,2351>,
  <2351,2352,2402>,
  <2403,2402,2352>,
  <2352,2353,2403>,
  <2404,2403,2353>,
  <2353,2354,2404>,
  <2405,2404,2354>,
  <2354,2355,2405>,
  <2406,2405,2355>,
  <2355,2356,2406>,
  <2407,2406,2356>,
  <2356,2357,2407>,
  <2408,2407,2357>,
  <2357,2358,2408>,
  <2409,2408,2358>,
  <2358,2359,2409>,
  <2410,2409,2359>,
  <2359,2360,2410>,
  <2411,2410,2360>,
  <2360,2361,2411>,
  <2412,2411,2361>,
  <2361,2362,2412>,
  <2413,2412,2362>,
  <2362,2363,2413>,
  <2414,2413,2363>,
  <2363,2364,2414>,
  <2415,2414,2364>,
  <2364,2365,2415>,
  <2416,2415,2365>,
  <2365,2366,2416>,
  <2417,2416,2366>,
  <2366,2367,2417>,
  <2418,2417,2367>,
  <2367,2368,2418>,
  <2419,2418,2368>,
  <2368,2369,2419>,
  <2420,2419,2369>,
  <2369,2370,2420>,
  <2421,2420,2370>,
  <2370,2371,2421>,
  <2422,2421,2371>,
  <2371,2372,2422>,
  <2423,2422,2372>,
  <2372,2373,2423>,
  <2424,2423,2373>,
  <2373,2374,2424>,
  <2425,2424,2374>,
  <2374,2375,2425>,
  <2426,2425,2375>,
  <2375,2376,2426>,
  <2427,2426,2376>,
  <2376,2377,2427>,
  <2428,2427,2377>,
  <2377,2378,2428>,
  <2429,2428,2378>,
  <2378,2379,2429>,
  <2430,2429,2379>,
  <2379,2380,2430>,
  <2431,2430,2380>,
  <2380,2381,2431>,
  <2432,2431,2381>,
  <2381,2382,2432>,
  <2433,2432,2382>,
  <2382,2383,2433>,
  <2434,2433,2383>,
  <2383,2384,2434>,
  <2435,2434,2384>,
  <2384,2385,2435>,
  <2436,2435,2385>,
  <2385,2386,2436>,
  <2437,2436,2386>,
  <2386,2387,2437>,
  <2438,2437,2387>,
  <2387,2388,2438>,
  <2439,2438,2388>,
  <2388,2389,2439>,
  <2440,2439,2389>,
  <2389,2390,2440>,
  <2441,2440,2390>,
  <2390,2391,2441>,
  <2442,2441,2391>,
  <2391,2392,2442>,
  <2443,2442,2392>,
  <2392,2393,2443>,
  <2444,2443,2393>,
  <2393,2394,2444>,
  <2445,2444,2394>,
  <2394,2395,2445>,
  <2446,2445,2395>,
  <2395,2396,2446>,
  <2447,2446,2396>,
  <2397,2398,2448>,
  <2449,2448,2398>,
  <2398,2399,2449>,
  <2450,2449,2399>,
  <2399,2400,2450>,
  <2451,2450,2400>,
  <2400,2401,2451>,
  <2452,2451,2401>,
  <2401,2402,2452>,
  <2453,2452,2402>,
  <2402,2403,2453>,
  <2454,2453,2403>,
  <2403,2404,2454>,
  <2455,2454,2404>,
  <2404,2405,2455>,
  <2456,2455,2405>,
  <2405,2406,2456>,
  <2457,2456,2406>,
  <2406,2407,2457>,
  <2458,2457,2407>,
  <2407,2408,2458>,
  <2459,2458,2408>,
  <2408,2409,2459>,
  <2460,2459,2409>,
  <2409,2410,2460>,
  <2461,2460,2410>,
  <2410,2411,2461>,
  <2462,2461,2411>,
  <2411,2412,2462>,
  <2463,2462,2412>,
  <2412,2413,2463>,
  <2464,2463,2413>,
  <2413,2414,2464>,
  <2465,2464,2414>,
  <2414,2415,2465>,
  <2466,2465,2415>,
  <2415,2416,2466>,
  <2467,2466,2416>,
  <2416,2417,2467>,
  <2468,2467,2417>,
  <2417,2418,2468>,
  <2469,2468,2418>,
  <2418,2419,2469>,
  <2470,2469,2419>,
  <2419,2420,2470>,
  <2471,2470,2420>,
  <2420,2421,2471>,
  <2472,2471,2421>,
  <2421,2422,2472>,
  <2473,2472,2422>,
  <2422,2423,2473>,
  <2474,2473,2423>,
  <2423,2424,2474>,
  <2475,2474,2424>,
  <2424,2425,2475>,
  <2476,2475,2425>,
  <2425,2426,2476>,
  <2477,2476,2426>,
  <2426,2427,2477>,
  <2478,2477,2427>,
  <2427,2428,2478>,
  <2479,2478,2428>,
  <2428,2429,2479>,
  <2480,2479,2429>,
  <2429,2430,2480>,
  <2481,2480,2430>,
  <2430,2431,2481>,
  <2482,2481,2431>,
  <2431,2432,2482>,
  <2483,2482,2432>,
  <2432,2433,2483>,
  <2484,2483,2433>,
  <2433,2434,2484>,
  <2485,2484,2434>,
  <2434,2435,2485>,
  <2486,2485,2435>,
  <2435,2436,2486>,
  <2487,2486,2436>,
  <2436,2437,2487>,
  <2488,2487,2437>,
  <2437,2438,2488>,
  <2489,2488,2438>,
  <2438,2439,2489>,
  <2490,2489,2439>,
  <2439,2440,2490>,
  <2491,2490,2440>,
  <2440,2441,2491>,
  <2492,2491,2441>,
  <2441,2442,2492>,
  <2493,2492,2442>,
  <2442,2443,2493>,
  <2494,2493,2443>,
  <2443,2444,2494>,
  <2495,2494,2444>,
  <2444,2445,2495>,
  <2496,2495,2445>,
  <2445,2446,2496>,
  <2497,2496,2446>,
  <2446,2447,2497>,
  <2498,2497,2447>,
  <2448,2449,2499>,
  <2500,2499,2449>,
  <2449,2450,2500>,
  <2501,2500,2450>,
  <2450,2451,2501>,
  <2502,2501,2451>,
  <2451,2452,2502>,
  <2503,2502,2452>,
  <2452,2453,2503>,
  <2504,2503,2453>,
  <2453,2454,2504>,
  <2505,2504,2454>,
  <2454,2455,2505>,
  <2506,2505,2455>,
  <2455,2456,2506>,
  <2507,2506,2456>,
  <2456,2457,2507>,
  <2508,2507,2457>,
  <2457,2458,2508>,
  <2509,2508,2458>,
  <2458,2459,2509>,
  <2510,2509,2459>,
  <2459,2460,2510>,
  <2511,2510,2460>,
  <2460,2461,2511>,
  <2512,2511,2461>,
  <2461,2462,2512>,
  <2513,2512,2462>,
  <2462,2463,2513>,
  <2514,2513,2463>,
  <2463,2464,2514>,
  <2515,2514,2464>,
  <2464,2465,2515>,
  <2516,2515,2465>,
  <2465,2466,2516>,
  <2517,2516,2466>,
  <2466,2467,2517>,
  <2518,2517,2467>,
  <2467,2468,2518>,
  <2519,2518,2468>,
  <2468,2469,2519>,
  <2520,2519,2469>,
  <2469,2470,2520>,
  <2521,2520,2470>,
  <2470,2471,2521>,
  <2522,2521,2471>,
  <2471,2472,2522>,
  <2523,2522,2472>,
  <2472,2473,2523>,
  <2524,2523,2473>,
  <2473,2474,2524>,
  <2525,2524,2474>,
  <2474,2475,2525>,
  <2526,2525,2475>,
  <2475,2476,2526>,
  <2527,2526,2476>,
  <2476,2477,2527>,
  <2528,2527,2477>,
  <2477,2478,2528>,
  <2529,2528,2478>,
  <2478,2479,2529>,
  <2530,2529,2479>,
  <2479,2480,2530>,
  <2531,2530,2480>,
  <2480,2481,2531>,
  <2532,2531,2481>,
  <2481,2482,2532>,
  <2533,2532,2482>,
  <2482,2483,2533>,
  <2534,2533,2483>,
  <2483,2484,2534>,
  <2535,2534,2484>,
  <2484,2485,2535>,
  <2536,2535,2485>,
  <2485,2486,2536>,
  <2537,2536,2486>,
  <2486,2487,2537>,
  <2538,2537,2487>,
  <2487,2488,2538>,
  <2539,2538,2488>,
  <2488,2489,2539>,
  <2540,2539,2489>,
  <2489,2490,2540>,
  <2541,2540,2490>,
  <2490,2491,2541>,
  <2542,2541,2491>,
  <2491,2492,2542>,
  <2543,2542,2492>,
  <2492,2493,2543>,
  <2544,2543,2493>,
  <2493,2494,2544>,
  <2545,2544,2494>,
  <2494,2495,2545>,
  <2546,2545,2495>,
  <2495,2496,2546>,
  <2547,2546,2496>,
  <2496,2497,2547>,
  <2548,2547,2497>,
  <2497,2498,2548>,
  <2549,2548,2498>,
  <2499,2500,2550>,
  <2551,2550,2500>,
  <2500,2501,2551>,
  <2552,2551,2501>,
  <2501,2502,2552>,
  <2553,2552,2502>,
  <2502,2503,2553>,
  <2554,2553,2503>,
  <2503,2504,2554>,
  <2555,2554,2504>,
  <2504,2505,2555>,
  <2556,2555,2505>,
  <2505,2506,2556>,
  <2557,2556,2506>,
  <2506,2507,2557>,
  <2558,2557,2507>,
  <2507,2508,2558>,
  <2559,2558,2508>,
  <2508,2509,2559>,
  <2560,2559,2509>,
  <2509,2510,2560>,
  <2561,2560,2510>,
  <2510,2511,2561>,
  <2562,2561,2511>,
  <2511,2512,2562>,
  <2563,2562,2512>,
  <2512,2513,2563>,
  <2564,2563,2513>,
  <2513,2514,2564>,
  <2565,2564,2514>,
  <2514,2515,2565>,
  <2566,2565,2515>,
  <2515,2516,2566>,
  <2567,2566,2516>,
  <2516,2517,2567>,
  <2568,2567,2517>,
  <2517,2518,2568>,
  <2569,2568,2518>,
  <2518,2519,2569>,
  <2570,2569,2519>,
  <2519,2520,2570>,
  <2571,2570,2520>,
  <2520,2521,2571>,
  <2572,2571,2521>,
  <2521,2522,2572>,
  <2573,2572,2522>,
  <2522,2523,2573>,
  <2574,2573,2523>,
  <2523,2524,2574>,
  <2575,2574,2524>,
  <2524,2525,2575>,
  <2576,2575,2525>,
  <2525,2526,2576>,
  <2577,2576,2526>,
  <2526,2527,2577>,
  <2578,2577,2527>,
  <2527,2528,2578>,
  <2579,2578,2528>,
  <2528,2529,2579>,
  <2580,2579,2529>,
  <2529,2530,2580>,
  <2581,2580,2530>,
  <2530,2531,2581>,
  <2582,2581,2531>,
  <2531,2532,2582>,
  <2583,2582,2532>,
  <2532,2533,2583>,
  <2584,2583,2533>,
  <2533,2534,2584>,
  <2585,2584,2534>,
  <2534,2535,2585>,
  <2586,2585,2535>,
  <2535,2536,2586>,
  <2587,2586,2536>,
  <2536,2537,2587>,
  <2588,2587,2537>,
  <2537,2538,2588>,
  <2589,2588,2538>,
  <2538,2539,2589>,
  <2590,2589,2539>,
  <2539,2540,2590>,
  <2591,2590,2540>,
  <2540,2541,2591>,
  <2592,2591,2541>,
  <2541,2542,2592>,
  <2593,2592,2542>,
  <2542,2543,2593>,
  <2594,2593,2543>,
  <2543,2544,2594>,
  <2595,2594,2544>,
  <2544,2545,2595>,
  <2596,2595,2545>,
  <2545,2546,2596>,
  <2597,2596,2546>,
  <2546,2547,2597>,
  <2598,2597,2547>,
  <2547,2548,2598>,
  <2599,2598,2548>,
  <2548,2549,2599>,
  <2600,2599,2549>,
  \}
  
  \}
\end{euleroutput}
\begin{eulercomment}
Sekarang kita mendefinisikan dua cakram, yang akan berpotongan dengan
permukaan.
\end{eulercomment}
\begin{eulerprompt}
>cl=povdisc([0.5,0.5,0],[1,1,0],2); ...
>ll=povdisc([0,0,1/4],[0,0,1],2);
\end{eulerprompt}
\begin{eulercomment}
Tulis permukaan dikurangi dua cakram.
\end{eulercomment}
\begin{eulerprompt}
>writeln(povdifference(mesh,povunion([cl,ll]),povlook(green)));
\end{eulerprompt}
\begin{euleroutput}
  difference \{ 
  object \{ mesh1 \}
  union \{ cylinder \{ <0.492929,0.492929,0>, <0.507071,0.507071,0>, 2
   texture \{ pigment \{ color rgb <0.470588,0.470588,0.470588> \}  \} 
   finish \{ ambient 0.2 \} 
   \}
  cylinder \{ <0,0,0.24>, <0,0,0.26>, 2
   texture \{ pigment \{ color rgb <0.470588,0.470588,0.470588> \}  \} 
   finish \{ ambient 0.2 \} 
   \}
   texture \{ pigment \{ color rgb <0.470588,0.470588,0.470588> \}  \} 
   finish \{ ambient 0.2 \} 
   \}
   texture \{ pigment \{ color rgb <0.0627451,0.564706,0.0627451> \}  \} 
   finish \{ ambient 0.2 \} 
   \}
\end{euleroutput}
\begin{eulercomment}
Tulis dua persimpangan.
\end{eulercomment}
\begin{eulerprompt}
>writeln(povintersection([mesh,cl],povlook(red))); ...
>writeln(povintersection([mesh,ll],povlook(gray)));
\end{eulerprompt}
\begin{euleroutput}
  intersection \{ object \{ mesh1 \}
  cylinder \{ <0.492929,0.492929,0>, <0.507071,0.507071,0>, 2
   texture \{ pigment \{ color rgb <0.470588,0.470588,0.470588> \}  \} 
   finish \{ ambient 0.2 \} 
   \}
   texture \{ pigment \{ color rgb <0.564706,0.0627451,0.0627451> \}  \} 
   finish \{ ambient 0.2 \} 
   \}
  intersection \{ object \{ mesh1 \}
  cylinder \{ <0,0,0.24>, <0,0,0.26>, 2
   texture \{ pigment \{ color rgb <0.470588,0.470588,0.470588> \}  \} 
   finish \{ ambient 0.2 \} 
   \}
   texture \{ pigment \{ color rgb <0.470588,0.470588,0.470588> \}  \} 
   finish \{ ambient 0.2 \} 
   \}
\end{euleroutput}
\begin{eulercomment}
Write a point at the maximum.
\end{eulercomment}
\begin{eulerprompt}
>writeln(povpoint([1/2,1/2,1/4],povlook(gray),size=2*defaultpointsize));
\end{eulerprompt}
\begin{euleroutput}
  sphere \{ <0.5,0.5,0.25>, 0.04 texture \{ pigment \{ color rgb <0.470588,0.470588,0.470588> \}  \} 
   finish \{ ambient 0.2 \} 
  
  \}
\end{euleroutput}
\begin{eulercomment}
Tambahkan sumbu dan selesaikan.
\end{eulercomment}
\begin{eulerprompt}
>writeAxes(0,1,0,1,0,1,d=0.015); ...
>povend();
\end{eulerprompt}
\begin{euleroutput}
  union \{
    cylinder \{ <-0.1,0,0>,<1.1,0,0>,0.015 \}
    cone \{ 
      <1.19,0,0>,0
      <1.1,0,0>,0.06
    \}
    texture \{ pigment \{ color rgb <0.470588,0.470588,0.470588> \} \}
  \}
  union \{
    cylinder \{ <-0.1,0,0>,<1.1,0,0>,0.015 \}
    cone \{ 
      <1.19,0,0>,0
      <1.1,0,0>,0.06
    \}
    rotate 90*z
    texture \{ pigment \{ color rgb <0.470588,0.470588,0.470588> \} \}
  \}
  union \{
    cylinder \{ <-0.1,0,0>,<1.1,0,0>,0.015 \}
    cone \{ 
      <1.19,0,0>,0
      <1.1,0,0>,0.06
    \}
    rotate -90*y
    texture \{ pigment \{ color rgb <0.470588,0.470588,0.470588> \} \}
  \}
\end{euleroutput}
\eulerimg{27}{images/Shalih Abdillah_22305144009_ApKom4-089.png}
\eulerheading{Anaglyph di Povray}
\begin{eulercomment}
Untuk menghasilkan anaglyph untuk kacamata merah/sian, Povray harus
berjalan dua kali dari posisi kamera yang berbeda. Ini menghasilkan
dua file Povray dan dua file PNG, yang dimuat dengan fungsi
loadanaglyph().

Tentu saja, Anda memerlukan kacamata merah/sian untuk melihat contoh
berikut dengan benar.

Fungsi pov3d() memiliki sakelar sederhana untuk menghasilkan
anaglyphs.
\end{eulercomment}
\begin{eulerprompt}
>pov3d("-exp(-x^2-y^2)/2",r=2,height=45°,>anaglyph, ...
>  center=[0,0,0.5],zoom=3.5);
\end{eulerprompt}
\eulerimg{27}{images/Shalih Abdillah_22305144009_ApKom4-090.png}
\begin{eulercomment}
Jika Anda membuat adegan dengan objek, Anda perlu menempatkan generasi
adegan ke dalam fungsi, dan menjalankannya dua kali dengan nilai yang
berbeda untuk parameter anaglyph.
\end{eulercomment}
\begin{eulerprompt}
>function myscene ...
\end{eulerprompt}
\begin{eulerudf}
    s=povsphere(povc,1);
    cl=povcylinder(-povz,povz,0.5);
    clx=povobject(cl,rotate=xrotate(90°));
    cly=povobject(cl,rotate=yrotate(90°));
    c=povbox([-1,-1,0],1);
    un=povunion([cl,clx,cly,c]);
    obj=povdifference(s,un,povlook(red));
    writeln(obj);
    writeAxes();
  endfunction
\end{eulerudf}
\begin{eulercomment}
Fungsi povanaglyph() melakukan semua ini. Parameternya seperti di
povstart() dan povend() digabungkan.
\end{eulercomment}
\begin{eulerprompt}
>povanaglyph("myscene",zoom=4.5);
\end{eulerprompt}
\eulerheading{Mendefinisikan Objek sendiri}
\begin{eulercomment}
Antarmuka povray Euler berisi banyak objek. Tapi Anda tidak terbatas
pada ini. Anda dapat membuat objek sendiri, yang menggabungkan objek
lain, atau objek yang sama sekali baru.

Kami mendemonstrasikan sebuah torus. Perintah Povray untuk ini adalah
"torus". Jadi kami mengembalikan string dengan perintah ini dan
parameternya. Perhatikan bahwa torus selalu berpusat di titik asal.
\end{eulercomment}
\begin{eulerprompt}
>function povdonat (r1,r2,look="") ...
\end{eulerprompt}
\begin{eulerudf}
    return "torus \{"+r1+","+r2+look+"\}";
  endfunction
\end{eulerudf}
\begin{eulercomment}
Inilah torus pertama kita
\end{eulercomment}
\begin{eulerprompt}
>t1=povdonat(0.8,0.2)
\end{eulerprompt}
\begin{euleroutput}
  torus \{0.8,0.2\}
\end{euleroutput}
\begin{eulercomment}
Mari kita gunakan objek ini untuk membuat torus kedua, diterjemahkan
dan diputar.
\end{eulercomment}
\begin{eulerprompt}
>t2=povobject(t1,rotate=xrotate(90°),translate=[0.8,0,0])
\end{eulerprompt}
\begin{euleroutput}
  object \{ torus \{0.8,0.2\}
   rotate 90 *x 
   translate <0.8,0,0>
   \}
\end{euleroutput}
\begin{eulercomment}
Sekarang kita menempatkan objek-objek ini ke dalam sebuah adegan.
Untuk tampilan, kami menggunakan Phong Shading.
\end{eulercomment}
\begin{eulerprompt}
>povstart(center=[0.4,0,0],angle=0°,zoom=3.8,aspect=1.5); ...
>writeln(povobject(t1,povlook(green,phong=1))); ...
>writeln(povobject(t2,povlook(green,phong=1)));
>povstart(center=[0.4,0,0],angle=0°,zoom=3.8,aspect=1.5); ...
>writeln(povobject(t1,povlook(green,phong=1))); ...
>writeln(povobject(t2,povlook(red,phong=1))); ...
> ...
\end{eulerprompt}
\begin{eulercomment}
\textgreater{}povend();

memanggil program Povray. Namun, jika terjadi kesalahan, itu tidak
menampilkan kesalahan. Karena itu Anda harus menggunakan

\end{eulercomment}
\begin{eulerttcomment}
 >povend(<keluar);
\end{eulerttcomment}
\begin{eulercomment}

jika ada yang tidak berhasil. Ini akan membiarkan jendela Povray
terbuka.
\end{eulercomment}
\begin{eulerprompt}
>povend(h=320,w=480);
\end{eulerprompt}
\eulerimg{18}{images/Shalih Abdillah_22305144009_ApKom4-091.png}
\begin{eulerprompt}
>povend(h=320,w=320);
\end{eulerprompt}
\eulerimg{19}{images/Shalih Abdillah_22305144009_ApKom4-092.png}
\begin{eulerprompt}
>povend(h=480,w=480);
\end{eulerprompt}
\eulerimg{27}{images/Shalih Abdillah_22305144009_ApKom4-093.png}
\begin{eulerprompt}
>povend(h=100,w=120);
\end{eulerprompt}
\eulerimg{5}{images/Shalih Abdillah_22305144009_ApKom4-094.png}
\begin{eulerprompt}
>povend(h=120,w=80);
\end{eulerprompt}
\begin{eulercomment}
Berikut adalah contoh yang lebih rumit. Kami memecahkan

\end{eulercomment}
\begin{eulerformula}
\[
Ax \le b, \quad x \ge 0, \quad c.x \to \text{Max.}
\]
\end{eulerformula}
\begin{eulercomment}
dan menunjukkan titik layak dan optimal dalam plot 3D.
\end{eulercomment}
\begin{eulerprompt}
>A=[10,8,4;5,6,8;6,3,2;9,5,6];
>b=[10,10,10,10]';
>c=[1,1,1];
\end{eulerprompt}
\begin{eulercomment}
Pertama, mari kita periksa, apakah contoh ini memiliki solusi sama
sekali.
\end{eulercomment}
\begin{eulerprompt}
>x=simplex(A,b,c,>max,>check)'
\end{eulerprompt}
\begin{euleroutput}
  [0,  1,  0.5]
\end{euleroutput}
\begin{eulercomment}
Ya, sudah.

Selanjutnya kita mendefinisikan dua objek. Yang pertama adalah pesawat

\end{eulercomment}
\begin{eulerformula}
\[
a \cdot x \le b
\]
\end{eulerformula}
\begin{eulerprompt}
>function oneplane (a,b,look="") ...
\end{eulerprompt}
\begin{eulerudf}
    return povplane(a,b,look)
  endfunction
\end{eulerudf}
\begin{eulercomment}
Kemudian kita mendefinisikan persimpangan dari semua setengah ruang
dan sebuah kubus.
\end{eulercomment}
\begin{eulerprompt}
>function adm (A, b, r, look="") ...
\end{eulerprompt}
\begin{eulerudf}
    ol=[];
    loop 1 to rows(A); ol=ol|oneplane(A[#],b[#]); end;
    ol=ol|povbox([0,0,0],[r,r,r]);
    return povintersection(ol,look);
  endfunction
\end{eulerudf}
\begin{eulercomment}
Kita sekarang dapat merencanakan adegannya.
\end{eulercomment}
\begin{eulerprompt}
>povstart(angle=120°,center=[0.5,0.5,0.5],zoom=3.5); ...
>writeln(adm(A,b,2,povlook(green,0.4))); ...
>writeAxes(0,1.3,0,1.6,0,1.5); ...
\end{eulerprompt}
\begin{eulercomment}
Berikut ini adalah lingkaran di sekitar optimal.
\end{eulercomment}
\begin{eulerprompt}
>writeln(povintersection([povsphere(x,0.5),povplane(c,c.x')], ...
>  povlook(red,0.9)));
\end{eulerprompt}
\begin{eulercomment}
Dan kesalahan ke arah yang optimal.
\end{eulercomment}
\begin{eulerprompt}
>writeln(povarrow(x,c*0.5,povlook(red)));
\end{eulerprompt}
\begin{eulercomment}
Kami menambahkan teks ke layar. Teks hanyalah objek 3D. Kita perlu
menempatkan dan memutarnya menurut pandangan kita.
\end{eulercomment}
\begin{eulerprompt}
>writeln(povtext("LINEAR PROBLEM",[0,0.2,1.3],size=0.05,rotate=125°)); ...
>povend();
\end{eulerprompt}
\eulerimg{27}{images/Shalih Abdillah_22305144009_ApKom4-097.png}
\begin{eulerprompt}
>writeln(povtext("LINEAR PROBLEM",[0,0.5,1.3],size=0.01,rotate=360°)); ...
>povend();
\end{eulerprompt}
\begin{euleroutput}
  object \{ text \{ ttf "timrom.ttf" "LINEAR PROBLEM" 0.04 ,0  texture \{ pigment \{ color rgb <0.470588,0.470588,0.470588> \}  \} 
   finish \{ ambient 0.2 \} 
   \}
   scale 0.01
   rotate 90 *x  rotate 360 *z 
   translate <0,0.5,1.3>
   \}
\end{euleroutput}
\eulerimg{27}{images/Shalih Abdillah_22305144009_ApKom4-098.png}
\eulerheading{Lebih Banyak Contoh}
\begin{eulercomment}
Anda dapat menemukan beberapa contoh lagi untuk Povray di Euler di
file berikut.

Lihat: Examples/Dandelin Spheres\\
Lihat: Examples/Donat Math\\
Lihat: Examples/Trefoil Knot\\
Lihat: Examples/Optimization by Affine Scaling
\end{eulercomment}
\end{eulernotebook}


\newpage
\chapter{KB Pekan 6-7: Menggunakan EMT untuk kalkulus}
\documentclass{article}

\usepackage{eumat}

\begin{eulernotebook}
\eulerheading{Kalkulus dengan EMT}
\begin{eulercomment}
Shalih Abdillah\\
22305144009\\
Matematika B 2022

\end{eulercomment}
\eulersubheading{}
\begin{eulercomment}
Materi Kalkulus mencakup di antaranya:

- Fungsi (fungsi aljabar, trigonometri, eksponensial, logaritma,
komposisi fungsi)\\
- Limit Fungsi,\\
- Turunan Fungsi,\\
- Integral Tak Tentu,\\
- Integral Tentu dan Aplikasinya,\\
- Barisan dan Deret (kekonvergenan barisan dan deret).

EMT (bersama Maxima) dapat digunakan untuk melakukan semua perhitungan
di dalam kalkulus, baik secara numerik maupun analitik (eksak).

\end{eulercomment}
\eulersubheading{Mendefinisikan Fungsi}
\begin{eulercomment}
Terdapat beberapa cara mendefinisikan fungsi pada EMT, yakni:

- Menggunakan format nama\_fungsi := rumus fungsi (untuk fungsi
numerik),\\
- Menggunakan format nama\_fungsi \&= rumus fungsi (untuk fungsi
simbolik, namun dapat dihitung secara numerik),\\
- Menggunakan format nama\_fungsi \&\&= rumus fungsi (untuk fungsi
simbolik murni, tidak dapat dihitung langsung),\\
- Fungsi sebagai program EMT.

Setiap format harus diawali dengan perintah function (bukan sebagai
ekspresi).

Berikut adalah adalah beberapa contoh cara mendefinisikan fungsi:

\end{eulercomment}
\begin{eulerformula}
\[
f(x)=2x^2+e^{\sin(x)}.
\]
\end{eulerformula}
\begin{eulerprompt}
>function f(x) := 2*x^2+exp(sin(x)) // fungsi numerik
>f(0), f(1), f(pi)
\end{eulerprompt}
\begin{euleroutput}
  1
  4.31977682472
  20.7392088022
\end{euleroutput}
\begin{eulerprompt}
>f(a) // tidak dapat dihitung nilainya
\end{eulerprompt}
\begin{euleroutput}
  Variable or function a not found.
  Error in:
  f(a) // tidak dapat dihitung nilainya ...
     ^
\end{euleroutput}
\begin{eulercomment}
Silakan Anda plot kurva fungsi di atas!

Berikutnya kita definisikan fungsi:

\end{eulercomment}
\begin{eulerformula}
\[
g(x)=\frac{\sqrt{x^2-3x}}{x+1}.
\]
\end{eulerformula}
\begin{eulerprompt}
>function g(x) := sqrt(x^2-3*x)/(x+1)
>g(3)
\end{eulerprompt}
\begin{euleroutput}
  0
\end{euleroutput}
\begin{eulerprompt}
>g(0)
\end{eulerprompt}
\begin{euleroutput}
  0
\end{euleroutput}
\begin{eulerprompt}
>g(1) // kompleks, tidak dapat dihitung oleh fungsi numerik
\end{eulerprompt}
\begin{euleroutput}
  Floating point error!
  Error in sqrt
  Try "trace errors" to inspect local variables after errors.
  g:
      useglobal; return sqrt(x^2-3*x)/(x+1) 
  Error in:
  g(1) // kompleks, tidak dapat dihitung oleh fungsi numerik ...
      ^
\end{euleroutput}
\begin{eulercomment}
Silakan Anda plot kurva fungsi di atas!
\end{eulercomment}
\begin{eulerprompt}
>f(g(5)) // komposisi fungsi
\end{eulerprompt}
\begin{euleroutput}
  2.20920171961
\end{euleroutput}
\begin{eulerprompt}
>g(f(5))
\end{eulerprompt}
\begin{euleroutput}
  0.950898070639
\end{euleroutput}
\begin{eulerprompt}
>function h(x) := f(g(x)) // definisi komposisi fungsi 
>h(5) // sama dengan f(g(5))
\end{eulerprompt}
\begin{euleroutput}
  2.20920171961
\end{euleroutput}
\begin{eulercomment}
Silakan Anda plot kurva fungsi komposisi fungsi f dan g:

\end{eulercomment}
\begin{eulerformula}
\[
h(x)=f(g(x))
\]
\end{eulerformula}
\begin{eulercomment}
dan

\end{eulercomment}
\begin{eulerformula}
\[
u(x)=g(f(x))
\]
\end{eulerformula}
\begin{eulercomment}
bersama-sama kurva fungsi f dan g dalam satu bidang koordinat.
\end{eulercomment}
\begin{eulerprompt}
>f(0:10) // nilai-nilai f(0), f(1), f(2), ..., f(10)
\end{eulerprompt}
\begin{euleroutput}
  [1,  4.31978,  10.4826,  19.1516,  32.4692,  50.3833,  72.7562,
  99.929,  130.69,  163.51,  200.58]
\end{euleroutput}
\begin{eulerprompt}
>fmap(0:10) // sama dengan f(0:10), berlaku untuk semua fungsi
\end{eulerprompt}
\begin{euleroutput}
  [1,  4.31978,  10.4826,  19.1516,  32.4692,  50.3833,  72.7562,
  99.929,  130.69,  163.51,  200.58]
\end{euleroutput}
\begin{eulerprompt}
>gmap(200:210)
\end{eulerprompt}
\begin{euleroutput}
  [0.987534,  0.987596,  0.987657,  0.987718,  0.987778,  0.987837,
  0.987896,  0.987954,  0.988012,  0.988069,  0.988126]
\end{euleroutput}
\begin{eulercomment}
Misalkan kita akan mendefinisikan fungsi

\end{eulercomment}
\begin{eulerformula}
\[
f(x) = \begin{cases} x^3 & x>0 \\ x^2 & x\le 0. \end{cases}
\]
\end{eulerformula}
\begin{eulercomment}
Fungsi tersebut tidak dapat didefinisikan sebagai fungsi numerik
secara "inline" menggunakan format :=, melainkan didefinisikan sebagai
program. Perhatikan, kata "map" digunakan agar fungsi dapat menerima
vektor sebagai input, dan hasilnya berupa vektor. Jika tanpa kata
"map" fungsinya hanya dapat menerima input satu nilai.
\end{eulercomment}
\begin{eulerprompt}
>function map f(x) ...
\end{eulerprompt}
\begin{eulerudf}
    if x>0 then return x^3
    else return x^2
    endif;
  endfunction
\end{eulerudf}
\begin{eulerprompt}
>f(1)
\end{eulerprompt}
\begin{euleroutput}
  1
\end{euleroutput}
\begin{eulerprompt}
>f(-2)
\end{eulerprompt}
\begin{euleroutput}
  4
\end{euleroutput}
\begin{eulerprompt}
>f(-5:5)
\end{eulerprompt}
\begin{euleroutput}
  [25,  16,  9,  4,  1,  0,  1,  8,  27,  64,  125]
\end{euleroutput}
\begin{eulerprompt}
>aspect(1.5); plot2d("f(x)",-5,5):
\end{eulerprompt}
\eulerimg{17}{images/Shalih Abdillah_22305144009_ApKom5-006.png}
\begin{eulerprompt}
>function f(x) &= 2*E^x // fungsi simbolik
\end{eulerprompt}
\begin{euleroutput}
  
                                      x
                                   2 E
  
\end{euleroutput}
\begin{eulerprompt}
>$f(a) // nilai fungsi secara simbolik
\end{eulerprompt}
\begin{eulerformula}
\[
2\,e^{a}
\]
\end{eulerformula}
\begin{eulerprompt}
>f(E) // nilai fungsi berupa bilangan desimal
\end{eulerprompt}
\begin{euleroutput}
  30.308524483
\end{euleroutput}
\begin{eulerprompt}
>$f(E), $float(%)
\end{eulerprompt}
\begin{eulerformula}
\[
30.30852448295852
\]
\end{eulerformula}
\eulerimg{0}{images/Shalih Abdillah_22305144009_ApKom5-009-large.png}
\begin{eulerprompt}
>function g(x) &= 3*x+1
\end{eulerprompt}
\begin{euleroutput}
  
                                 3 x + 1
  
\end{euleroutput}
\begin{eulerprompt}
>function h(x) &= f(g(x)) // komposisi fungsi
\end{eulerprompt}
\begin{euleroutput}
  
                                   3 x + 1
                                2 E
  
\end{euleroutput}
\begin{eulerprompt}
>plot2d("h(x)",-1,1):
\end{eulerprompt}
\eulerimg{17}{images/Shalih Abdillah_22305144009_ApKom5-010.png}
\eulerheading{Latihan}
\begin{eulercomment}
Bukalah buku Kalkulus. Cari dan pilih beberapa (paling sedikit 5
fungsi berbeda tipe/bentuk/jenis) fungsi dari buku tersebut, kemudian
definisikan fungsi-fungsi tersebut dan komposisinya di EMT pada
baris-baris perintah berikut (jika perlu tambahkan lagi). Untuk setiap
fungsi, hitung beberapa nilainya, baik untuk satu nilai maupun vektor.
Gambar grafik fungsi-fungsi tersebut dan komposisi-komposisi 2 fungsi.

Juga, carilah fungsi beberapa (dua) variabel. Lakukan hal sama seperti
di atas.

Soal 1
\end{eulercomment}
\begin{eulerprompt}
>function f(x) := 3*x + 10*x - 4
>f(3)
\end{eulerprompt}
\begin{euleroutput}
  35
\end{euleroutput}
\begin{eulerprompt}
>f(0:8)
\end{eulerprompt}
\begin{euleroutput}
  [-4,  9,  22,  35,  48,  61,  74,  87,  100]
\end{euleroutput}
\begin{eulerprompt}
>aspect(2); plot2d("f(x)",0,8):
\end{eulerprompt}
\eulerimg{13}{images/Shalih Abdillah_22305144009_ApKom5-011.png}
\begin{eulercomment}
Soal 2
\end{eulercomment}
\begin{eulerprompt}
>function f(y) := 6*y^2 + 3*y^3 - 2
>f(3)
\end{eulerprompt}
\begin{euleroutput}
  133
\end{euleroutput}
\begin{eulerprompt}
>f(1:5) 
\end{eulerprompt}
\begin{euleroutput}
  [7,  46,  133,  286,  523]
\end{euleroutput}
\begin{eulerprompt}
> aspect(2); plot2d("f(x)",-20,20):
\end{eulerprompt}
\eulerimg{13}{images/Shalih Abdillah_22305144009_ApKom5-012.png}
\begin{eulercomment}
Soal 3
\end{eulercomment}
\begin{eulerprompt}
> function f(x) := (3*x+4) / (x^2)
>f(4) 
\end{eulerprompt}
\begin{euleroutput}
  1
\end{euleroutput}
\begin{eulerprompt}
>f(1:10) 
\end{eulerprompt}
\begin{euleroutput}
  [7,  2.5,  1.44444,  1,  0.76,  0.611111,  0.510204,  0.4375,
  0.382716,  0.34]
\end{euleroutput}
\begin{eulerprompt}
> aspect(2); plot2d("f(x)",0,5):
\end{eulerprompt}
\eulerimg{13}{images/Shalih Abdillah_22305144009_ApKom5-013.png}
\begin{eulercomment}
Soal 4
\end{eulercomment}
\begin{eulerprompt}
> function f(x) := sin(x)^2 + 13
> f(30)
\end{eulerprompt}
\begin{euleroutput}
  13.9762064902
\end{euleroutput}
\begin{eulerprompt}
> f(0)
\end{eulerprompt}
\begin{euleroutput}
  13
\end{euleroutput}
\begin{eulerprompt}
> aspect(2); plot2d("f(x)",0,10):
\end{eulerprompt}
\eulerimg{13}{images/Shalih Abdillah_22305144009_ApKom5-014.png}
\begin{eulercomment}
Soal 5
\end{eulercomment}
\begin{eulerprompt}
> function f(x,y) :=  -x^2 + 3*y^2 -10
> f(2,5)
\end{eulerprompt}
\begin{euleroutput}
  61
\end{euleroutput}
\begin{eulerprompt}
> f(0:10,0:10)
\end{eulerprompt}
\begin{euleroutput}
  [-10,  -8,  -2,  8,  22,  40,  62,  88,  118,  152,  190]
\end{euleroutput}
\begin{eulerprompt}
> aspect(2); plot2d("f(x,x)",-10,10):
\end{eulerprompt}
\eulerimg{13}{images/Shalih Abdillah_22305144009_ApKom5-015.png}
\begin{eulercomment}
Soal 6
\end{eulercomment}
\begin{eulerprompt}
> function f(x,y,z) := x^3 + 2*y^2 + 3*z^3 + 4
> f(1,2,3)
\end{eulerprompt}
\begin{euleroutput}
  94
\end{euleroutput}
\begin{eulerprompt}
> f(3,4,5)
\end{eulerprompt}
\begin{euleroutput}
  438
\end{euleroutput}
\begin{eulerprompt}
> aspect(2); plot3d("f(x,x,x)",angle=40, height=45°):
\end{eulerprompt}
\eulerimg{13}{images/Shalih Abdillah_22305144009_ApKom5-016.png}
\begin{eulercomment}
\begin{eulercomment}
\eulerheading{Menghitung Limit}
\begin{eulercomment}
Perhitungan limit pada EMT dapat dilakukan dengan menggunakan fungsi Maxima, yakni "limit".
Fungsi "limit" dapat digunakan untuk menghitung limit fungsi dalam bentuk ekspresi maupun fungsi
yang sudah didefinisikan sebelumnya. Nilai limit dapat dihitung pada sebarang nilai atau pada tak
hingga (-inf, minf, dan inf). Limit kiri dan limit kanan juga dapat dihitung, dengan cara memberi
opsi "plus" atau "minus". Hasil limit dapat berupa nilai, "und" (tak definisi), "ind" (tak tentu
namun terbatas), "infinity" (kompleks tak hingga).

Perhatikan beberapa contoh berikut. Perhatikan cara menampilkan perhitungan secara lengkap, tidak
hanya menampilkan hasilnya saja.
\end{eulercomment}
\begin{eulerprompt}
>$showev('limit(sqrt(x^2-3*x)/(x+1),x,inf))
\end{eulerprompt}
\begin{euleroutput}
  Maxima said:
  limit: variable must be a symbol or subscripted symbol; found: 
   errexp1
  #0: showev(f='limit([0,sqrt(-4.999975000072321e-7*r+2.7777500001498e-14*r^2)/(1+1.66665833335744e-7*r),sqrt(-3.99...)
   -- an error. To debug this try: debugmode(true);
  
  Error in:
   $showev('limit(sqrt(x^2-3*x)/(x+1),x,inf)) ...
                                            ^
\end{euleroutput}
\begin{eulerprompt}
>$limit((x^3-13*x^2+51*x-63)/(x^3-4*x^2-3*x+18),x,3)
\end{eulerprompt}
\begin{eulerformula}
\[
-\frac{4}{5}
\]
\end{eulerformula}
\begin{eulercomment}
maxima: limit((x\textasciicircum{}3-13*x\textasciicircum{}2+51*x-63)/(x\textasciicircum{}3-4*x\textasciicircum{}2-3*x+18),x,3)=limit((x\textasciicircum{}3-13*x\textasciicircum{}2+51*x-63)/(x\textasciicircum{}3-4*x\textasciicircum{}2-3*x+18),x,3)

Fungsi tersebut diskontinu di titik x=3. Berikut adalah grafik
fungsinya.
\end{eulercomment}
\begin{eulerprompt}
>aspect(1.5); plot2d("(x^3-13*x^2+51*x-63)/(x^3-4*x^2-3*x+18)",0,4); plot2d(3,-4/5,>points,style="ow",>add):
\end{eulerprompt}
\eulerimg{17}{images/Shalih Abdillah_22305144009_ApKom5-018.png}
\begin{eulerprompt}
>$limit(2*x*sin(x)/(1-cos(x)),x,0)
\end{eulerprompt}
\begin{eulerformula}
\[
4
\]
\end{eulerformula}
\begin{eulercomment}
maxima: 'limit(2*x*sin(x)/(1-cos(x)),x,0)

Fungsi tersebut diskontinu di titik x=0. Berikut adalah grafik
fungsinya.
\end{eulercomment}
\begin{eulerprompt}
>plot2d("2*x*sin(x)/(1-cos(x))",-pi,pi); plot2d(0,4,>points,style="ow",>add):
\end{eulerprompt}
\eulerimg{17}{images/Shalih Abdillah_22305144009_ApKom5-020.png}
\begin{eulerprompt}
>$limit(cot(7*h)/cot(5*h),h,0)
\end{eulerprompt}
\begin{eulerformula}
\[
\frac{5}{7}
\]
\end{eulerformula}
\begin{eulerformula}
\[
\lim_{h\rightarrow 0}{\frac{\cot \left(7\,h\right)}{\cot \left(5\,h  \right)}}=\frac{5}{7}
\]
\end{eulerformula}
\begin{eulercomment}
Fungsi tersebut juga diskontinu (karena tidak terdefinisi) di x=0.
Berikut adalah grafiknya.
\end{eulercomment}
\begin{eulerprompt}
>plot2d("cot(7*x)/cot(5*x)",-0.001,0.001); plot2d(0,5/7,>points,style="ow",>add):
\end{eulerprompt}
\eulerimg{17}{images/Shalih Abdillah_22305144009_ApKom5-023.png}
\begin{eulerprompt}
>$showev('limit(((x/8)^(1/3)-1)/(x-8),x,8))
\end{eulerprompt}
\begin{eulerformula}
\[
\lim_{x\rightarrow 8}{\frac{\frac{x^{\frac{1}{3}}}{2}-1}{x-8}}=  \frac{1}{24}
\]
\end{eulerformula}
\begin{eulercomment}
Tunjukkan limit tersebut dengan grafik, seperti contoh-contoh sebelumnya.
\end{eulercomment}
\begin{eulerprompt}
>plot2d("(((x/8)^(1/3)-1)/(x-8))",-3,3); plot2d(0,5/7,>points,style="ow",>add):
\end{eulerprompt}
\eulerimg{13}{images/Shalih Abdillah_22305144009_ApKom5-025.png}
\begin{eulerprompt}
>$showev('limit(1/(2*x-1),x,0))
\end{eulerprompt}
\begin{eulerformula}
\[
\lim_{x\rightarrow 0}{\frac{1}{2\,x-1}}=-1
\]
\end{eulerformula}
\begin{eulercomment}
Tunjukkan limit tersebut dengan grafik, seperti contoh-contoh sebelumnya.
\end{eulercomment}
\begin{eulerprompt}
>plot2d("1/(2*x-1)",-3,3); plot2d(0,5/7,>points,style="ow",>add):
\end{eulerprompt}
\eulerimg{13}{images/Shalih Abdillah_22305144009_ApKom5-027.png}
\begin{eulerprompt}
>$showev('limit((x^2-3*x-10)/(x-5),x,5))
\end{eulerprompt}
\begin{eulerformula}
\[
\lim_{x\rightarrow 5}{\frac{x^2-3\,x-10}{x-5}}=7
\]
\end{eulerformula}
\begin{eulercomment}
Tunjukkan limit tersebut dengan grafik, seperti contoh-contoh sebelumnya.
\end{eulercomment}
\begin{eulerprompt}
>plot2d("(x^2-3*x-10)/(x-5)",-3,3); plot2d(0,5/7,>points,style="ow",>add):
\end{eulerprompt}
\eulerimg{13}{images/Shalih Abdillah_22305144009_ApKom5-029.png}
\begin{eulerprompt}
>$showev('limit(sqrt(x^2+x)-x,x,inf))
\end{eulerprompt}
\begin{eulerformula}
\[
\lim_{x\rightarrow \infty }{\sqrt{x^2+x}-x}=\frac{1}{2}
\]
\end{eulerformula}
\begin{eulercomment}
Tunjukkan limit tersebut dengan grafik, seperti contoh-contoh sebelumnya.
\end{eulercomment}
\begin{eulerprompt}
>plot2d("sqrt(x^2+x)-x",-3,3); plot2d(0,5/7,>points,style="ow",>add):
\end{eulerprompt}
\eulerimg{13}{images/Shalih Abdillah_22305144009_ApKom5-031.png}
\begin{eulerprompt}
>$showev('limit(abs(x-1)/(x-1),x,1,minus))
\end{eulerprompt}
\begin{eulerformula}
\[
\lim_{x\uparrow 1}{\frac{\left| x-1\right| }{x-1}}=-1
\]
\end{eulerformula}
\begin{eulercomment}
Hitung limit di atas untuk x menuju 1 dari kanan.\\
Tunjukkan limit tersebut dengan grafik, seperti contoh-contoh sebelumnya.
\end{eulercomment}
\begin{eulerprompt}
>plot2d("abs(x-1)/(x-1)",-pi,pi); plot2d(0,1,>points,style="ow",>add):
\end{eulerprompt}
\eulerimg{13}{images/Shalih Abdillah_22305144009_ApKom5-033.png}
\begin{eulerprompt}
>$showev('limit(sin(x)/x,x,0))
\end{eulerprompt}
\begin{eulerformula}
\[
\lim_{x\rightarrow 0}{\frac{\sin x}{x}}=1
\]
\end{eulerformula}
\begin{eulerprompt}
>plot2d("sin(x)/x",-pi,pi); plot2d(0,1,>points,style="ow",>add):
\end{eulerprompt}
\eulerimg{17}{images/Shalih Abdillah_22305144009_ApKom5-035.png}
\begin{eulerprompt}
>$showev('limit(sin(x^3)/x,x,0))
\end{eulerprompt}
\begin{eulerformula}
\[
\lim_{x\rightarrow 0}{\frac{\sin x^3}{x}}=0
\]
\end{eulerformula}
\begin{eulerprompt}
>plot2d("sin(x^3)/x",-3,3); plot2d(0,1,>points,style="ow",>add):
\end{eulerprompt}
\eulerimg{13}{images/Shalih Abdillah_22305144009_ApKom5-037.png}
\begin{eulercomment}
Tunjukkan limit tersebut dengan grafik, seperti contoh-contoh sebelumnya.
\end{eulercomment}
\begin{eulerprompt}
>$showev('limit(log(x), x, minf))
\end{eulerprompt}
\begin{eulerformula}
\[
\lim_{x\rightarrow  -\infty }{\log x}={\it infinity}
\]
\end{eulerformula}
\begin{eulerprompt}
>$showev('limit((-2)^x,x, inf))
\end{eulerprompt}
\begin{eulerformula}
\[
\lim_{x\rightarrow \infty }{\left(-2\right)^{x}}={\it infinity}
\]
\end{eulerformula}
\begin{eulerprompt}
>$showev('limit(t-sqrt(2-t),t,2,minus))
\end{eulerprompt}
\begin{eulerformula}
\[
\lim_{t\uparrow 2}{t-\sqrt{2-t}}=2
\]
\end{eulerformula}
\begin{eulerprompt}
>$showev('limit(t-sqrt(2-t),t,2,plus))
\end{eulerprompt}
\begin{eulerformula}
\[
\lim_{t\downarrow 2}{t-\sqrt{2-t}}=2
\]
\end{eulerformula}
\begin{eulerprompt}
>$showev('limit(t-sqrt(2-t),t,5,plus)) // Perhatikan hasilnya
\end{eulerprompt}
\begin{eulerformula}
\[
\lim_{t\downarrow 5}{t-\sqrt{2-t}}=5-\sqrt{3}\,i
\]
\end{eulerformula}
\begin{eulerprompt}
>plot2d("x-sqrt(2-x)",0,2):
\end{eulerprompt}
\eulerimg{17}{images/Shalih Abdillah_22305144009_ApKom5-043.png}
\begin{eulerprompt}
>$showev('limit((x^2-9)/(2*x^2-5*x-3),x,3))
\end{eulerprompt}
\begin{eulerformula}
\[
\lim_{x\rightarrow 3}{\frac{x^2-9}{2\,x^2-5\,x-3}}=\frac{6}{7}
\]
\end{eulerformula}
\begin{eulercomment}
Tunjukkan limit tersebut dengan grafik, seperti contoh-contoh sebelumnya.
\end{eulercomment}
\begin{eulerprompt}
>plot2d("(x^2-9)/(2*x^2-5*x-3)",0,2):
\end{eulerprompt}
\eulerimg{13}{images/Shalih Abdillah_22305144009_ApKom5-045.png}
\begin{eulerprompt}
>$showev('limit((1-cos(x))/x,x,0))
\end{eulerprompt}
\begin{eulerformula}
\[
\lim_{x\rightarrow 0}{\frac{1-\cos x}{x}}=0
\]
\end{eulerformula}
\begin{eulerprompt}
>plot2d("(1-cos(x))/x",0,2):
\end{eulerprompt}
\eulerimg{13}{images/Shalih Abdillah_22305144009_ApKom5-047.png}
\begin{eulerprompt}
>$showev('limit((x^2+abs(x))/(x^2-abs(x)),x,0))
\end{eulerprompt}
\begin{eulerformula}
\[
\lim_{x\rightarrow 0}{\frac{\left| x\right| +x^2}{x^2-\left| x  \right| }}=-1
\]
\end{eulerformula}
\begin{eulercomment}
Tunjukkan limit tersebut dengan grafik, seperti contoh-contoh sebelumnya.
\end{eulercomment}
\begin{eulerprompt}
>plot2d("(x^2+abs(x))/(x^2-abs(x))",0,2):
\end{eulerprompt}
\eulerimg{13}{images/Shalih Abdillah_22305144009_ApKom5-049.png}
\begin{eulerprompt}
>$showev('limit((1+1/x)^x,x,inf))
\end{eulerprompt}
\begin{eulerformula}
\[
\lim_{x\rightarrow \infty }{\left(\frac{1}{x}+1\right)^{x}}=e
\]
\end{eulerformula}
\begin{eulerprompt}
>plot2d("(1+1/x)^x",0,1000):
\end{eulerprompt}
\eulerimg{17}{images/Shalih Abdillah_22305144009_ApKom5-051.png}
\begin{eulerprompt}
>$showev('limit((1+k/x)^x,x,inf))
\end{eulerprompt}
\begin{eulerformula}
\[
\lim_{x\rightarrow \infty }{\left(\frac{k}{x}+1\right)^{x}}=e^{k}
\]
\end{eulerformula}
\begin{eulerprompt}
>$showev('limit((1+x)^(1/x),x,0))
\end{eulerprompt}
\begin{eulerformula}
\[
\lim_{x\rightarrow 0}{\left(x+1\right)^{\frac{1}{x}}}=e
\]
\end{eulerformula}
\begin{eulerprompt}
>plot2d("(1+x)^(1/x)",0,1000):
\end{eulerprompt}
\eulerimg{13}{images/Shalih Abdillah_22305144009_ApKom5-054.png}
\begin{eulerprompt}
>$showev('limit((x/(x+k))^x,x,inf))
\end{eulerprompt}
\begin{eulerformula}
\[
\lim_{x\rightarrow \infty }{\left(\frac{x}{x+k}\right)^{x}}=e^ {- k   }
\]
\end{eulerformula}
\begin{eulerprompt}
>$showev('limit((E^x-E^2)/(x-2),x,2))
\end{eulerprompt}
\begin{eulerformula}
\[
\lim_{x\rightarrow 2}{\frac{e^{x}-e^2}{x-2}}=e^2
\]
\end{eulerformula}
\begin{eulerprompt}
>plot2d("(E^x-E^2)/(x-2)",0,1000):
\end{eulerprompt}
\eulerimg{13}{images/Shalih Abdillah_22305144009_ApKom5-057.png}
\begin{eulerprompt}
>$showev('limit(sin(1/x),x,0))
\end{eulerprompt}
\begin{eulerformula}
\[
\lim_{x\rightarrow 0}{\sin \left(\frac{1}{x}\right)}={\it ind}
\]
\end{eulerformula}
\begin{eulerprompt}
>$showev('limit(sin(1/x),x,inf))
\end{eulerprompt}
\begin{eulerformula}
\[
\lim_{x\rightarrow \infty }{\sin \left(\frac{1}{x}\right)}=0
\]
\end{eulerformula}
\begin{eulerprompt}
>plot2d("sin(1/x)",-0.001,0.001):
\end{eulerprompt}
\eulerimg{17}{images/Shalih Abdillah_22305144009_ApKom5-060.png}
\eulerheading{Latihan}
\begin{eulercomment}
Bukalah buku Kalkulus. Cari dan pilih beberapa (paling sedikit 5
fungsi berbeda tipe/bentuk/jenis) fungsi dari buku tersebut, kemudian
definisikan di EMT pada baris-baris perintah berikut (jika perlu
tambahkan lagi). Untuk setiap fungsi, hitung nilai limit fungsi
tersebut di beberapa nilai dan di tak hingga. Gambar grafik fungsi
tersebut untuk mengkonfirmasi nilai-nilai limit tersebut.

Fungsi 1\\
\end{eulercomment}
\begin{eulerformula}
\[
\text{$f(x)=\frac{8x-7}{x-2}$}
\]
\end{eulerformula}
\begin{eulerprompt}
>$showev('limit((8*x-7)/(x-2),x,1))
\end{eulerprompt}
\begin{eulerformula}
\[
\lim_{x\rightarrow 1}{\frac{8\,x-7}{x-2}}=-1
\]
\end{eulerformula}
\begin{eulerprompt}
>plot2d("(8*x-7)/(x-2)",-5.5,5.5):
\end{eulerprompt}
\eulerimg{27}{images/Shalih Abdillah_22305144009_ApKom5-063.png}
\begin{eulercomment}
Fungsi 2\\
\end{eulercomment}
\begin{eulerformula}
\[
\text{$f(x)=\frac{4x^5+3x^3+2x^2+x+5}{x-2}$}
\]
\end{eulerformula}
\begin{eulerprompt}
>$showev('limit((4*x^5+3*x^3+2*x^2+x+5)/(x-2),x,1))
\end{eulerprompt}
\begin{eulerformula}
\[
\lim_{x\rightarrow 1}{\frac{4\,x^5+3\,x^3+2\,x^2+x+5}{x-2}}=-15
\]
\end{eulerformula}
\begin{eulerprompt}
>plot2d("(4*x^5+3*x^3+2*x^2+x+5)/(x-2)",-3,2,-5,-1):
\end{eulerprompt}
\eulerimg{27}{images/Shalih Abdillah_22305144009_ApKom5-066.png}
\begin{eulercomment}
Fungsi 3\\
\end{eulercomment}
\begin{eulerformula}
\[
\text{$f(x)=\frac{sin x+cos x}{sin(2x)}$}
\]
\end{eulerformula}
\begin{eulerprompt}
>$showev('limit(((sin(x)+cos(x))/sin(2*x)),x,2))
\end{eulerprompt}
\begin{eulerformula}
\[
\lim_{x\rightarrow 2}{\frac{\sin x+\cos x}{\sin \left(2\,x\right)}}=  \frac{\sin 2}{\sin 4}+\frac{\cos 2}{\sin 4}
\]
\end{eulerformula}
\begin{eulerprompt}
>plot2d("(sin(x)+cos(x))/(sin(2*x))",-5,5,-5,5):
\end{eulerprompt}
\eulerimg{27}{images/Shalih Abdillah_22305144009_ApKom5-069.png}
\begin{eulercomment}
Fungsi 4\\
\end{eulercomment}
\begin{eulerformula}
\[
\text{$f(x)=\frac{|x+x^2|}{|x-x^2|}$}
\]
\end{eulerformula}
\begin{eulerprompt}
>$showev('limit((abs(x+x^2))/(abs(x-x^2)),x,0))
\end{eulerprompt}
\begin{eulerformula}
\[
\lim_{x\rightarrow 0}{\frac{\left| x^2+x\right| }{\left| x^2-x  \right| }}=1
\]
\end{eulerformula}
\begin{eulerprompt}
>plot2d("((x+x^2)/(x-x^2))",-10,20)
\end{eulerprompt}
\begin{eulercomment}
Fungsi 5\\
\end{eulercomment}
\begin{eulerformula}
\[
\text{$f(x)=\frac{1}{sin^x2}$}
\]
\end{eulerformula}
\begin{eulerprompt}
>$showev('limit(sin(2)^(-x),x,0))
\end{eulerprompt}
\begin{eulerformula}
\[
\lim_{x\rightarrow 0}{\frac{1}{\sin ^{x}2}}=1
\]
\end{eulerformula}
\begin{eulerprompt}
>plot2d("(sin(2)^(-x))",-100,100,-100,100):
\end{eulerprompt}
\eulerimg{27}{images/Shalih Abdillah_22305144009_ApKom5-074.png}
\begin{eulercomment}
\begin{eulercomment}
\eulerheading{Turunan Fungsi}
\begin{eulercomment}
Definisi turunan:

\end{eulercomment}
\begin{eulerformula}
\[
f'(x) = \lim_{h\to 0} \frac{f(x+h)-f(x)}{h}
\]
\end{eulerformula}
\begin{eulercomment}
Berikut adalah contoh-contoh menentukan turunan fungsi dengan
menggunakan definisi turunan (limit).
\end{eulercomment}
\begin{eulerprompt}
>$showev('limit(((x+h)^2-x^2)/h,h,0)) // turunan x^2
\end{eulerprompt}
\begin{eulerformula}
\[
\lim_{h\rightarrow 0}{\frac{\left(x+h\right)^2-x^2}{h}}=2\,x
\]
\end{eulerformula}
\begin{eulerprompt}
>p &= expand((x+h)^2-x^2)|simplify; $p //pembilang dijabarkan dan disederhanakan
\end{eulerprompt}
\begin{eulerformula}
\[
2\,h\,x+h^2
\]
\end{eulerformula}
\begin{eulerprompt}
>q &=ratsimp(p/h); $q // ekspresi yang akan dihitung limitnya disederhanakan
\end{eulerprompt}
\begin{eulerformula}
\[
2\,x+h
\]
\end{eulerformula}
\begin{eulerprompt}
>$limit(q,h,0) // nilai limit sebagai turunan
\end{eulerprompt}
\begin{eulerformula}
\[
2\,x
\]
\end{eulerformula}
\begin{eulerprompt}
>$showev('limit(((x+h)^n-x^n)/h,h,0)) // turunan x^n
\end{eulerprompt}
\begin{eulerformula}
\[
\lim_{h\rightarrow 0}{\frac{\left(x+h\right)^{n}-x^{n}}{h}}=n\,x^{n  -1}
\]
\end{eulerformula}
\begin{eulercomment}
Mengapa hasilnya seperti itu? Tuliskan atau tunjukkan bahwa hasil limit tersebut
benar, sehingga benar turunan fungsinya benar.  Tulis penjelasan Anda di komentar
ini.

Sebagai petunjuk, ekspansikan (x+h)\textasciicircum{}n dengan menggunakan teorema binomial.
\end{eulercomment}
\begin{eulerprompt}
>$showev('limit((sin(x+h)-sin(x))/h,h,0)) // turunan sin(x)
\end{eulerprompt}
\begin{eulerformula}
\[
\lim_{h\rightarrow 0}{\frac{\sin \left(x+h\right)-\sin x}{h}}=\cos   x
\]
\end{eulerformula}
\begin{eulercomment}
Mengapa hasilnya seperti itu? Tuliskan atau tunjukkan bahwa hasil limit tersebut\\
benar, sehingga benar turunan fungsinya benar.  Tulis penjelasan Anda di komentar
ini.

Sebagai petunjuk, ekspansikan sin(x+h) dengan menggunakan rumus jumlah dua sudut.
\end{eulercomment}
\begin{eulerprompt}
>$showev('limit((log(x+h)-log(x))/h,h,0)) // turunan log(x)
\end{eulerprompt}
\begin{eulerformula}
\[
\lim_{h\rightarrow 0}{\frac{\log \left(x+h\right)-\log x}{h}}=  \frac{1}{x}
\]
\end{eulerformula}
\begin{eulercomment}
Mengapa hasilnya seperti itu? Tuliskan atau tunjukkan bahwa hasil limit tersebut\\
benar, sehingga benar turunan fungsinya benar.  Tulis penjelasan Anda di komentar
ini.

Sebagai petunjuk, gunakan sifat-sifat logaritma dan hasil limit pada bagian
sebelumnya di atas.
\end{eulercomment}
\begin{eulerprompt}
>$showev('limit((1/(x+h)-1/x)/h,h,0)) // turunan 1/x
\end{eulerprompt}
\begin{eulerformula}
\[
\lim_{h\rightarrow 0}{\frac{\frac{1}{x+h}-\frac{1}{x}}{h}}=-\frac{1  }{x^2}
\]
\end{eulerformula}
\begin{eulerprompt}
>$showev('limit((E^(x+h)-E^x)/h,h,0)) // turunan f(x)=e^x
\end{eulerprompt}
\begin{euleroutput}
  Answering "Is x an integer?" with "integer"
  Answering "Is x an integer?" with "integer"
  Answering "Is x an integer?" with "integer"
  Answering "Is x an integer?" with "integer"
  Answering "Is x an integer?" with "integer"
  Maxima is asking
  Acceptable answers are: yes, y, Y, no, n, N, unknown, uk
  Is x an integer?
  
  Use assume!
  Error in:
   $showev('limit((E^(x+h)-E^x)/h,h,0)) // turunan f(x)=e^x ...
                                       ^
\end{euleroutput}
\begin{eulercomment}
Maxima bermasalah dengan limit:

\end{eulercomment}
\begin{eulerformula}
\[
\lim_{h\to 0}\frac{e^{x+h}-e^x}{h}.
\]
\end{eulerformula}
\begin{eulercomment}
Oleh karena itu diperlukan trik khusus agar hasilnya benar.
\end{eulercomment}
\begin{eulerprompt}
>$showev('limit((E^h-1)/h,h,0))
\end{eulerprompt}
\begin{eulerformula}
\[
\lim_{h\rightarrow 0}{\frac{e^{h}-1}{h}}=1
\]
\end{eulerformula}
\begin{eulerprompt}
>$showev('factor(E^(x+h)-E^x))
\end{eulerprompt}
\begin{eulerformula}
\[
{\it factor}\left(e^{x+h}-e^{x}\right)=\left(e^{h}-1\right)\,e^{x}
\]
\end{eulerformula}
\begin{eulerprompt}
>$showev('limit(factor((E^(x+h)-E^x)/h),h,0)) // turunan f(x)=e^x
\end{eulerprompt}
\begin{eulerformula}
\[
\left(\lim_{h\rightarrow 0}{\frac{e^{h}-1}{h}}\right)\,e^{x}=e^{x}
\]
\end{eulerformula}
\begin{eulerprompt}
>function f(x) &= x^x
\end{eulerprompt}
\begin{euleroutput}
  
                                     x
                                    x
  
\end{euleroutput}
\begin{eulerprompt}
>$showev('limit(f(x),x,0))
\end{eulerprompt}
\begin{eulerformula}
\[
\lim_{x\rightarrow 0}{x^{x}}=1
\]
\end{eulerformula}
\begin{eulercomment}
Silakan Anda gambar kurva

\end{eulercomment}
\begin{eulerformula}
\[
y=x^x.
\]
\end{eulerformula}
\begin{eulerprompt}
>$showev('limit((f(x+h)-f(x))/h,h,0)) // turunan f(x)=x^x
\end{eulerprompt}
\begin{eulerformula}
\[
\lim_{h\rightarrow 0}{\frac{\left(x+h\right)^{x+h}-x^{x}}{h}}=  {\it infinity}
\]
\end{eulerformula}
\begin{eulercomment}
Di sini Maxima juga bermasalah terkait limit:

\end{eulercomment}
\begin{eulerformula}
\[
\lim_{h\to 0} \frac{(x+h)^{x+h}-x^x}{h}.
\]
\end{eulerformula}
\begin{eulercomment}
Dalam hal ini diperlukan asumsi nilai x.
\end{eulercomment}
\begin{eulerprompt}
>&assume(x>0); $showev('limit((f(x+h)-f(x))/h,h,0)) // turunan f(x)=x^x
\end{eulerprompt}
\begin{eulerformula}
\[
\lim_{h\rightarrow 0}{\frac{\left(x+h\right)^{x+h}-x^{x}}{h}}=x^{x}  \,\left(\log x+1\right)
\]
\end{eulerformula}
\begin{eulercomment}
Mengapa hasilnya seperti itu? Tuliskan atau tunjukkan bahwa hasil limit tersebut benar, sehingga benar turunan fungsinya benar.
Tulis penjelasan Anda di komentar ini.
\end{eulercomment}
\begin{eulerprompt}
>&forget(x>0) // jangan lupa, lupakan asumsi untuk kembali ke semula
\end{eulerprompt}
\begin{euleroutput}
  
                                 [x > 0]
  
\end{euleroutput}
\begin{eulerprompt}
>&forget(x<0)
\end{eulerprompt}
\begin{euleroutput}
  
                                 [x < 0]
  
\end{euleroutput}
\begin{eulerprompt}
>&facts()
\end{eulerprompt}
\begin{euleroutput}
  
                                    []
  
\end{euleroutput}
\begin{eulerprompt}
>$showev('limit((asin(x+h)-asin(x))/h,h,0)) // turunan arcsin(x)
\end{eulerprompt}
\begin{eulerformula}
\[
\lim_{h\rightarrow 0}{\frac{\arcsin \left(x+h\right)-\arcsin x}{h}}=  \frac{1}{\sqrt{1-x^2}}
\]
\end{eulerformula}
\begin{eulercomment}
Mengapa hasilnya seperti itu? Tuliskan atau tunjukkan bahwa hasil limit tersebut benar, sehingga
benar turunan fungsinya benar. Tulis penjelasan Anda di komentar ini.
\end{eulercomment}
\begin{eulerprompt}
>$showev('limit((tan(x+h)-tan(x))/h,h,0)) // turunan tan(x)
\end{eulerprompt}
\begin{eulerformula}
\[
\lim_{h\rightarrow 0}{\frac{\tan \left(x+h\right)-\tan x}{h}}=  \frac{1}{\cos ^2x}
\]
\end{eulerformula}
\begin{eulercomment}
Mengapa hasilnya seperti itu? Tuliskan atau tunjukkan bahwa hasil limit tersebut benar, sehingga
benar turunan fungsinya benar. Tulis penjelasan Anda di komentar ini.
\end{eulercomment}
\begin{eulerprompt}
>function f(x) &= sinh(x) // definisikan f(x)=sinh(x)
\end{eulerprompt}
\begin{euleroutput}
  
                                 sinh(x)
  
\end{euleroutput}
\begin{eulerprompt}
>function df(x) &= limit((f(x+h)-f(x))/h,h,0); $df(x) // df(x) = f'(x)
\end{eulerprompt}
\begin{eulerformula}
\[
\frac{e^ {- x }\,\left(e^{2\,x}+1\right)}{2}
\]
\end{eulerformula}
\begin{eulercomment}
Hasilnya adalah cosh(x), karena

\end{eulercomment}
\begin{eulerformula}
\[
\frac{e^x+e^{-x}}{2}=\cosh(x).
\]
\end{eulerformula}
\begin{eulerprompt}
>plot2d(["f(x)","df(x)"],-pi,pi,color=[blue,red]):
\end{eulerprompt}
\eulerimg{17}{images/Shalih Abdillah_22305144009_ApKom5-095.png}
\begin{eulerprompt}
>function f(x) &= sin(3*x^5+7)^2
\end{eulerprompt}
\begin{euleroutput}
  
                                 2    5
                              sin (3 x  + 7)
  
\end{euleroutput}
\begin{eulerprompt}
>diff(f,3), diffc(f,3)
\end{eulerprompt}
\begin{euleroutput}
  1198.32948904
  1198.72863721
\end{euleroutput}
\begin{eulercomment}
Apakah perbedaan diff dan diffc?
\end{eulercomment}
\begin{eulerprompt}
>$showev('diff(f(x),x))
\end{eulerprompt}
\begin{eulerformula}
\[
\frac{d}{d\,x}\,\sin ^2\left(3\,x^5+7\right)=30\,x^4\,\cos \left(3  \,x^5+7\right)\,\sin \left(3\,x^5+7\right)
\]
\end{eulerformula}
\begin{eulerprompt}
>$% with x=3
\end{eulerprompt}
\begin{eulerformula}
\[
{\it \%at}\left(\frac{d}{d\,x}\,\sin ^2\left(3\,x^5+7\right) , x=3  \right)=2430\,\cos 736\,\sin 736
\]
\end{eulerformula}
\begin{eulerprompt}
>$float(%)
\end{eulerprompt}
\begin{eulerformula}
\[
{\it \%at}\left(\frac{d^{1.0}}{d\,x^{1.0}}\,\sin ^2\left(3.0\,x^5+  7.0\right) , x=3.0\right)=1198.728637211748
\]
\end{eulerformula}
\begin{eulerprompt}
>plot2d(f,0,3.1):
\end{eulerprompt}
\eulerimg{17}{images/Shalih Abdillah_22305144009_ApKom5-099.png}
\begin{eulerprompt}
>function f(x) &=5*cos(2*x)-2*x*sin(2*x) // mendifinisikan fungsi f
\end{eulerprompt}
\begin{euleroutput}
  
                        5 cos(2 x) - 2 x sin(2 x)
  
\end{euleroutput}
\begin{eulerprompt}
>function df(x) &=diff(f(x),x) // fd(x) = f'(x)
\end{eulerprompt}
\begin{euleroutput}
  
                       - 12 sin(2 x) - 4 x cos(2 x)
  
\end{euleroutput}
\begin{eulerprompt}
>$'f(1)=f(1), $float(f(1)), $'f(2)=f(2), $float(f(2)) // nilai f(1) dan f(2)
\end{eulerprompt}
\begin{eulerformula}
\[
-0.2410081230863468
\]
\end{eulerformula}
\eulerimg{0}{images/Shalih Abdillah_22305144009_ApKom5-101-large.png}
\eulerimg{0}{images/Shalih Abdillah_22305144009_ApKom5-102-large.png}
\eulerimg{0}{images/Shalih Abdillah_22305144009_ApKom5-103-large.png}
\begin{eulerprompt}
>xp=solve("df(x)",1,2,0) // solusi f'(x)=0 pada interval [1, 2]
\end{eulerprompt}
\begin{euleroutput}
  1.35822987384
\end{euleroutput}
\begin{eulerprompt}
>df(xp), f(xp) // cek bahwa f'(xp)=0 dan nilai ekstrim di titik tersebut
\end{eulerprompt}
\begin{euleroutput}
  0
  -5.67530133759
\end{euleroutput}
\begin{eulerprompt}
>plot2d(["f(x)","df(x)"],0,2*pi,color=[blue,red]): //grafik fungsi dan turunannya
\end{eulerprompt}
\eulerimg{17}{images/Shalih Abdillah_22305144009_ApKom5-104.png}
\begin{eulercomment}
Perhatikan titik-titik "puncak" grafik y=f(x) dan nilai turunan pada
saat grafik fungsinya mencapai titik "puncak" tersebut.
\end{eulercomment}
\begin{eulerprompt}
>  
\end{eulerprompt}
\eulerheading{Latihan}
\begin{eulercomment}
Bukalah buku Kalkulus. Cari dan pilih beberapa (paling sedikit 5
fungsi berbeda tipe/bentuk/jenis) fungsi dari buku tersebut, kemudian
definisikan di EMT pada baris-baris perintah berikut (jika perlu
tambahkan lagi). Untuk setiap fungsi, tentukan turunannya dengan
menggunakan definisi turunan (limit), menggunakan perintah diff, dan
secara manual (langkah demi langkah yang dihitung dengan Maxima)
seperti contoh-contoh di atas. Gambar grafik fungsi asli dan fungsi
turunannya pada sumbu koordinat yang sama.

Fungsi 1
\end{eulercomment}
\begin{eulerprompt}
>function f(x) &= (x^3+4)*(x-6); $f(x)
\end{eulerprompt}
\begin{eulerformula}
\[
\left(x-6\right)\,\left(x^3+4\right)
\]
\end{eulerformula}
\begin{eulerprompt}
>function df(x) &= limit((f(x+h)-f(x))/h,h,0); $df(x)
\end{eulerprompt}
\begin{eulerformula}
\[
4\,x^3-18\,x^2+4
\]
\end{eulerformula}
\begin{eulerprompt}
>aspect(1.5); plot2d(["f(x)","df(x)"],-10,10,color=[blue,red]):
\end{eulerprompt}
\eulerimg{17}{images/Shalih Abdillah_22305144009_ApKom5-107.png}
\begin{eulercomment}
Fungsi 2
\end{eulercomment}
\begin{eulerprompt}
>function f(x) &= (3*x^3+2*x^2-x)*(x^4-3*x^3-4*x^2+1); $f(x)
\end{eulerprompt}
\begin{eulerformula}
\[
\left(3\,x^3+2\,x^2-x\right)\,\left(x^4-3\,x^3-4\,x^2+1\right)
\]
\end{eulerformula}
\begin{eulerprompt}
>function df(x) &= limit((f(x+h)-f(x))/h,h,0); $df(x)
\end{eulerprompt}
\begin{eulerformula}
\[
21\,x^6-42\,x^5-95\,x^4-20\,x^3+21\,x^2+4\,x-1
\]
\end{eulerformula}
\begin{eulerprompt}
>aspect(1.5); plot2d(["f(x)","df(x)"],-10,10,color=[blue,red]):
\end{eulerprompt}
\eulerimg{17}{images/Shalih Abdillah_22305144009_ApKom5-110.png}
\begin{eulercomment}
Fungsi 3
\end{eulercomment}
\begin{eulerprompt}
>function f(x) &= (16)/(4*x^2-3*x+9); $f(x)
\end{eulerprompt}
\begin{eulerformula}
\[
\frac{16}{4\,x^2-3\,x+9}
\]
\end{eulerformula}
\begin{eulerprompt}
>function df(x) &= limit((f(x+h)-f(x))/h,h,0); $df(x)
\end{eulerprompt}
\begin{eulerformula}
\[
\frac{48-128\,x}{16\,x^4-24\,x^3+81\,x^2-54\,x+81}
\]
\end{eulerformula}
\begin{eulerprompt}
>aspect(1.5); plot2d(["f(x)","df(x)"],-10,10,color=[green,red]):
\end{eulerprompt}
\eulerimg{17}{images/Shalih Abdillah_22305144009_ApKom5-113.png}
\begin{eulercomment}
Fungsi 4
\end{eulercomment}
\begin{eulerprompt}
>function f(x) &= (2*x^2-1)/(x-1); $f(x)
\end{eulerprompt}
\begin{eulerformula}
\[
\frac{2\,x^2-1}{x-1}
\]
\end{eulerformula}
\begin{eulerprompt}
>function df(x) &= limit((f(x+h)-f(x))/h,h,0); $df(x)
\end{eulerprompt}
\begin{eulerformula}
\[
\frac{2\,x^2-4\,x+1}{x^2-2\,x+1}
\]
\end{eulerformula}
\begin{eulerprompt}
>aspect(1.5); plot2d(["f(x)","df(x)"],-0.5,0.5,color=[green,red]):
\end{eulerprompt}
\eulerimg{17}{images/Shalih Abdillah_22305144009_ApKom5-116.png}
\begin{eulercomment}
Fungsi 5
\end{eulercomment}
\begin{eulerprompt}
>function f(x) &= (sin(x)+cos(x))/(sin(x)); $f(x)
\end{eulerprompt}
\begin{eulerformula}
\[
\frac{\sin x+\cos x}{\sin x}
\]
\end{eulerformula}
\begin{eulerprompt}
>function df(x) &= limit((f(x+h)-f(x))/h,h,0); $df(x)
\end{eulerprompt}
\begin{eulerformula}
\[
\frac{-\sin ^2x-\cos ^2x}{\sin ^2x}
\]
\end{eulerformula}
\begin{eulerprompt}
>aspect(1.5); plot2d(["f(x)","df(x)"],-pi/3,pi/3,color=[green,red]):
\end{eulerprompt}
\eulerimg{17}{images/Shalih Abdillah_22305144009_ApKom5-119.png}
\begin{eulerprompt}
>   
\end{eulerprompt}
\begin{euleroutput}
  
\end{euleroutput}
\eulerheading{Integral}
\begin{eulercomment}
EMT dapat digunakan untuk menghitung integral, baik integral tak tentu
maupun integral tentu. Untuk integral tak tentu (simbolik) sudah tentu
EMT menggunakan Maxima, sedangkan untuk perhitungan integral tentu EMT
sudah menyediakan beberapa fungsi yang mengimplementasikan algoritma
kuadratur (perhitungan integral tentu menggunakan metode numerik).

Pada notebook ini akan ditunjukkan perhitungan integral tentu dengan
menggunakan Teorema Dasar Kalkulus:

\end{eulercomment}
\begin{eulerformula}
\[
\int_a^b f(x)\ dx = F(b)-F(a), \quad \text{ dengan  } F'(x) = f(x).
\]
\end{eulerformula}
\begin{eulercomment}
Fungsi untuk menentukan integral adalah integrate. Fungsi ini dapat
digunakan untuk menentukan, baik integral tentu maupun tak tentu (jika
fungsinya memiliki antiderivatif). Untuk perhitungan integral tentu
fungsi integrate menggunakan metode numerik (kecuali fungsinya tidak
integrabel, kita tidak akan menggunakan metode ini).
\end{eulercomment}
\begin{eulerprompt}
>$showev('integrate(x^n,x))
\end{eulerprompt}
\begin{euleroutput}
  Answering "Is n equal to -1?" with "no"
\end{euleroutput}
\begin{eulerformula}
\[
\int {x^{n}}{\;dx}=\frac{x^{n+1}}{n+1}
\]
\end{eulerformula}
\begin{eulerprompt}
>$showev('integrate(1/(1+x),x))
\end{eulerprompt}
\begin{eulerformula}
\[
\int {\frac{1}{x+1}}{\;dx}=\log \left(x+1\right)
\]
\end{eulerformula}
\begin{eulerprompt}
>$showev('integrate(1/(1+x^2),x))
\end{eulerprompt}
\begin{eulerformula}
\[
\int {\frac{1}{x^2+1}}{\;dx}=\arctan x
\]
\end{eulerformula}
\begin{eulerprompt}
>$showev('integrate(1/sqrt(1-x^2),x))
\end{eulerprompt}
\begin{eulerformula}
\[
\int {\frac{1}{\sqrt{1-x^2}}}{\;dx}=\arcsin x
\]
\end{eulerformula}
\begin{eulerprompt}
>$showev('integrate(sin(x),x,0,pi))
\end{eulerprompt}
\begin{eulerformula}
\[
\int_{0}^{\pi}{\sin x\;dx}=2
\]
\end{eulerformula}
\begin{eulerprompt}
>plot2d("sin(x)",0,2*pi):
\end{eulerprompt}
\eulerimg{17}{images/Shalih Abdillah_22305144009_ApKom5-126.png}
\begin{eulerprompt}
>$showev('integrate(sin(x),x,a,b))
\end{eulerprompt}
\begin{eulerformula}
\[
\int_{a}^{b}{\sin x\;dx}=\cos a-\cos b
\]
\end{eulerformula}
\begin{eulerprompt}
>$showev('integrate(x^n,x,a,b))
\end{eulerprompt}
\begin{euleroutput}
  Answering "Is n positive, negative or zero?" with "positive"
\end{euleroutput}
\begin{eulerformula}
\[
\int_{a}^{b}{x^{n}\;dx}=\frac{b^{n+1}}{n+1}-\frac{a^{n+1}}{n+1}
\]
\end{eulerformula}
\begin{eulerprompt}
>$showev('integrate(x^2*sqrt(2*x+1),x))
\end{eulerprompt}
\begin{eulerformula}
\[
\int {x^2\,\sqrt{2\,x+1}}{\;dx}=\frac{\left(2\,x+1\right)^{\frac{7  }{2}}}{28}-\frac{\left(2\,x+1\right)^{\frac{5}{2}}}{10}+\frac{\left(  2\,x+1\right)^{\frac{3}{2}}}{12}
\]
\end{eulerformula}
\begin{eulerprompt}
>$showev('integrate(x^2*sqrt(2*x+1),x,0,2))
\end{eulerprompt}
\begin{eulerformula}
\[
\int_{0}^{2}{x^2\,\sqrt{2\,x+1}\;dx}=\frac{2\,5^{\frac{5}{2}}}{21}-  \frac{2}{105}
\]
\end{eulerformula}
\begin{eulerprompt}
>$ratsimp(%)
\end{eulerprompt}
\begin{eulerformula}
\[
\int_{0}^{2}{x^2\,\sqrt{2\,x+1}\;dx}=\frac{2\,5^{\frac{7}{2}}-2}{  105}
\]
\end{eulerformula}
\begin{eulerprompt}
>$showev('integrate((sin(sqrt(x)+a)*E^sqrt(x))/sqrt(x),x,0,pi^2))
\end{eulerprompt}
\begin{eulerformula}
\[
\int_{0}^{\pi^2}{\frac{\sin \left(\sqrt{x}+a\right)\,e^{\sqrt{x}}}{  \sqrt{x}}\;dx}=\left(-e^{\pi}-1\right)\,\sin a+\left(e^{\pi}+1  \right)\,\cos a
\]
\end{eulerformula}
\begin{eulerprompt}
>$factor(%)
\end{eulerprompt}
\begin{eulerformula}
\[
\int_{0}^{\pi^2}{\frac{\sin \left(\sqrt{x}+a\right)\,e^{\sqrt{x}}}{  \sqrt{x}}\;dx}=\left(-e^{\pi}-1\right)\,\left(\sin a-\cos a\right)
\]
\end{eulerformula}
\begin{eulerprompt}
>function map f(x) &= E^(-x^2)
\end{eulerprompt}
\begin{euleroutput}
  
                                      2
                                   - x
                                  E
  
\end{euleroutput}
\begin{eulerprompt}
>$showev('integrate(f(x),x))
\end{eulerprompt}
\begin{eulerformula}
\[
\int {e^ {- x^2 }}{\;dx}=\frac{\sqrt{\pi}\,\mathrm{erf}\left(x  \right)}{2}
\]
\end{eulerformula}
\begin{eulercomment}
Fungsi f tidak memiliki antiturunan, integralnya masih memuat integral
lain.

\end{eulercomment}
\begin{eulerformula}
\[
erf(x) = \int \frac{e^{-x^2}}{\sqrt{\pi}} \ dx.
\]
\end{eulerformula}
\begin{eulercomment}
Kita tidak dapat menggunakan teorema Dasar kalkulus untuk menghitung
integral tentu fungsi tersebut jika semua batasnya berhingga. Dalam
hal ini dapat digunakan metode numerik (rumus kuadratur).

Misalkan kita akan menghitung:

\end{eulercomment}
\begin{eulerformula}
\[
\int_{0}^{\pi}{\frac{\sin x+\cos x}{\sin x}\;dx}
\]
\end{eulerformula}
\begin{eulerprompt}
>x=0:0.1:pi-0.1; plot2d(x,f(x+0.1),>bar); plot2d("f(x)",0,pi,>add):
\end{eulerprompt}
\eulerimg{17}{images/Shalih Abdillah_22305144009_ApKom5-137.png}
\begin{eulercomment}
Integral tentu

\end{eulercomment}
\begin{eulerformula}
\[
\int_{0}^{\pi}{\frac{\sin x+\cos x}{\sin x}\;dx}
\]
\end{eulerformula}
\begin{eulercomment}
dapat dihampiri dengan jumlah luas persegi-persegi panjang di bawah
kurva y=f(x) tersebut. Langkah-langkahnya adalah sebagai berikut.
\end{eulercomment}
\begin{eulerprompt}
>t &= makelist(a,a,0,pi-0.1,0.1); // t sebagai list untuk menyimpan nilai-nilai x
>fx &= makelist(f(t[i]+0.1),i,1,length(t)); // simpan nilai-nilai f(x)
>// jangan menggunakan x sebagai list, kecuali Anda pakar Maxima!
\end{eulerprompt}
\begin{eulercomment}
Hasilnya adalah:

\end{eulercomment}
\begin{eulerformula}
\[
\int_{0}^{\pi}{\frac{\sin x+\cos x}{\sin x}\;dx}=1.22817137704692
\]
\end{eulerformula}
\begin{eulercomment}
Jumlah tersebut diperoleh dari hasil kali lebar sub-subinterval (=0.1)
dan jumlah nilai-nilai f(x) untuk x = 0.1, 0.2, 0.3, ..., 3.2.
\end{eulercomment}
\begin{eulerprompt}
>0.1*sum(f(x+0.1)) // cek langsung dengan perhitungan numerik EMT
\end{eulerprompt}
\begin{euleroutput}
  0.836219610253
\end{euleroutput}
\begin{eulercomment}
Untuk mendapatkan nilai integral tentu yang mendekati nilai
sebenarnya, lebar sub-intervalnya dapat diperkecil lagi, sehingga
daerah di bawah kurva tertutup semuanya, misalnya dapat digunakan
lebar subinterval 0.001. (Silakan dicoba!)

Meskipun Maxima tidak dapat menghitung integral tentu fungsi tersebut
untuk batas-batas yang berhingga, namun integral tersebut dapat
dihitung secara eksak jika batas-batasnya tak hingga. Ini adalah salah
satu keajaiban di dalam matematika, yang terbatas tidak dapat dihitung
secara eksak, namun yang tak hingga malah dapat dihitung secara eksak.
\end{eulercomment}
\begin{eulerprompt}
>$showev('integrate(f(x),x,0,inf))
\end{eulerprompt}
\begin{eulerformula}
\[
\int_{0}^{\infty }{e^ {- x^2 }\;dx}=\frac{\sqrt{\pi}}{2}
\]
\end{eulerformula}
\begin{eulercomment}
Tunjukkan kebenaran hasil di atas!

Berikut adalah contoh lain fungsi yang tidak memiliki antiderivatif,
sehingga integral tentunya hanya dapat dihitung dengan metode numerik.
\end{eulercomment}
\begin{eulerprompt}
>function f(x) &= x^x
\end{eulerprompt}
\begin{euleroutput}
  
                                     x
                                    x
  
\end{euleroutput}
\begin{eulerprompt}
>$showev('integrate(f(x),x,0,1))
\end{eulerprompt}
\begin{eulerformula}
\[
\int_{0}^{1}{x^{x}\;dx}=\int_{0}^{1}{x^{x}\;dx}
\]
\end{eulerformula}
\begin{eulerprompt}
>x=0:0.1:1-0.01; plot2d(x,f(x+0.01),>bar); plot2d("f(x)",0,1,>add):
\end{eulerprompt}
\eulerimg{17}{images/Shalih Abdillah_22305144009_ApKom5-142.png}
\begin{eulercomment}
Maxima gagal menghitung integral tentu tersebut secara langsung menggunakan perintah
integrate. Berikut kita lakukan seperti contoh sebelumnya untuk mendapat hasil atau
pendekatan nilai integral tentu tersebut.
\end{eulercomment}
\begin{eulerprompt}
>t &= makelist(a,a,0,1-0.01,0.01);
>fx &= makelist(f(t[i]+0.01),i,1,length(t));
\end{eulerprompt}
\begin{eulerformula}
\[
\int_{0}^{1}{\frac{\sin x+\cos x}{\sin x}\;dx}=0.122817137704692
\]
\end{eulerformula}
\begin{eulercomment}
Apakah hasil tersebut cukup baik? perhatikan gambarnya.
\end{eulercomment}
\begin{eulerprompt}
>function f(x) &= sin(3*x^5+7)^2
\end{eulerprompt}
\begin{euleroutput}
  
                                 2    5
                              sin (3 x  + 7)
  
\end{euleroutput}
\begin{eulerprompt}
>integrate(f,0,1)
\end{eulerprompt}
\begin{euleroutput}
  0.542581176074
\end{euleroutput}
\begin{eulerprompt}
>&showev('integrate(f(x),x,0,1))
\end{eulerprompt}
\begin{euleroutput}
  
           1                           1              pi
          /                      gamma(-) sin(14) sin(--)
          [     2    5                 5              10
          I  sin (3 x  + 7) dx = ------------------------
          ]                                  1/5
          /                              10 6
           0
         4/5                  1          4/5                  1
   - (((6    gamma_incomplete(-, 6 I) + 6    gamma_incomplete(-, - 6 I))
                              5                               5
               4/5                    1
   sin(14) + (6    I gamma_incomplete(-, 6 I)
                                      5
      4/5                    1                       pi
   - 6    I gamma_incomplete(-, - 6 I)) cos(14)) sin(--) - 60)/120
                             5                       10
  
\end{euleroutput}
\begin{eulerprompt}
>&float(%)
\end{eulerprompt}
\begin{euleroutput}
  
           1.0
          /
          [       2      5
          I    sin (3.0 x  + 7.0) dx = 
          ]
          /
           0.0
  0.09820784258795788 - 0.008333333333333333
   (0.3090169943749474 (0.1367372182078336
   (4.192962712629476 I gamma__incomplete(0.2, 6.0 I)
   - 4.192962712629476 I gamma__incomplete(0.2, - 6.0 I))
   + 0.9906073556948704 (4.192962712629476 gamma__incomplete(0.2, 6.0 I)
   + 4.192962712629476 gamma__incomplete(0.2, - 6.0 I))) - 60.0)
  
\end{euleroutput}
\begin{eulerprompt}
>$showev('integrate(x*exp(-x),x,0,1)) // Integral tentu (eksak)
\end{eulerprompt}
\begin{eulerformula}
\[
\int_{0}^{1}{x\,e^ {- x }\;dx}=1-2\,e^ {- 1 }
\]
\end{eulerformula}
\eulerheading{Aplikasi Integral Tentu}
\begin{eulerprompt}
>plot2d("x^3-x",-0.1,1.1); plot2d("-x^2",>add);  ...
>b=solve("x^3-x+x^2",0.5); x=linspace(0,b,200); xi=flipx(x); ...
>plot2d(x|xi,x^3-x|-xi^2,>filled,style="|",fillcolor=1,>add): // Plot daerah antara 2 kurva
\end{eulerprompt}
\eulerimg{17}{images/Shalih Abdillah_22305144009_ApKom5-145.png}
\begin{eulerprompt}
>a=solve("x^3-x+x^2",0), b=solve("x^3-x+x^2",1) // absis titik-titik potong kedua kurva
\end{eulerprompt}
\begin{euleroutput}
  0
  0.61803398875
\end{euleroutput}
\begin{eulerprompt}
>integrate("(-x^2)-(x^3-x)",a,b) // luas daerah yang diarsir
\end{eulerprompt}
\begin{euleroutput}
  0.0758191713542
\end{euleroutput}
\begin{eulercomment}
Hasil tersebut akan kita bandingkan dengan perhitungan secara analitik.
\end{eulercomment}
\begin{eulerprompt}
>a &= solve((-x^2)-(x^3-x),x); $a // menentukan absis titik potong kedua kurva secara eksak
\end{eulerprompt}
\begin{eulerformula}
\[
\left[ x=\frac{-\sqrt{5}-1}{2} , x=\frac{\sqrt{5}-1}{2} , x=0   \right] 
\]
\end{eulerformula}
\begin{eulerprompt}
>$showev('integrate(-x^2-x^3+x,x,0,(sqrt(5)-1)/2)) // Nilai integral secara eksak
\end{eulerprompt}
\begin{eulerformula}
\[
\int_{0}^{\frac{\sqrt{5}-1}{2}}{-x^3-x^2+x\;dx}=\frac{13-5^{\frac{3  }{2}}}{24}
\]
\end{eulerformula}
\begin{eulerprompt}
>$float(%)
\end{eulerprompt}
\begin{eulerformula}
\[
\int_{0.0}^{0.6180339887498949}{-1.0\,x^3-1.0\,x^2+x\;dx}=  0.07581917135421037
\]
\end{eulerformula}
\eulersubheading{Panjang Kurva}
\begin{eulercomment}
Hitunglah panjang kurva berikut ini dan luas daerah di dalam kurva
tersebut.

\end{eulercomment}
\begin{eulerformula}
\[
\gamma(t) = (r(t) \cos(t), r(t) \sin(t))
\]
\end{eulerformula}
\begin{eulercomment}
dengan

\end{eulercomment}
\begin{eulerformula}
\[
r(t) = 1 + \dfrac{\sin(3t)}{2},\quad 0\le t\le 2\pi.
\]
\end{eulerformula}
\begin{eulerprompt}
>t=linspace(0,2pi,1000); r=1+sin(3*t)/2; x=r*cos(t); y=r*sin(t); ...
>plot2d(x,y,>filled,fillcolor=red,style="/",r=1.5): // Kita gambar kurvanya terlebih dahulu
\end{eulerprompt}
\eulerimg{27}{images/Shalih Abdillah_22305144009_ApKom5-151.png}
\begin{eulerprompt}
>function r(t) &= 1+sin(3*t)/2; $'r(t)=r(t)
\end{eulerprompt}
\begin{eulerformula}
\[
r\left(\left[ 0 , 0.1 , 0.2 , 0.3 , 0.4 , 0.5 , 0.6 , 0.7 ,   0.7999999999999999 , 0.8999999999999999 , 0.9999999999999999 , 1.1   , 1.2 , 1.3 , 1.4 , 1.5 , 1.6 , 1.7 , 1.8 , 1.900000000000001 , 2.0   , 2.100000000000001 , 2.200000000000001 , 2.300000000000001 ,   2.400000000000001 , 2.500000000000001 , 2.600000000000001 ,   2.700000000000001 , 2.800000000000001 , 2.900000000000001 ,   3.000000000000001 \right] \right)=\left[ 1 , 1.14776010333067 ,   1.282321236697518 , 1.391663454813742 , 1.466019542983613 ,   1.498747493302027 , 1.486923815439098 , 1.431604683324437 ,   1.337731590275576 , 1.213689940116915 , 1.070560004029934 ,   0.9211271529283759 , 0.778739778352574 , 0.656116920408013 ,   0.5642121137932059 , 0.5112349411674514 , 0.5019176955820797 ,   0.5370926588361341 , 0.6136177562220066 , 0.7246572287011819 ,   0.860292250900538 , 1.008406950242176 , 1.15577068175669 ,   1.289219882194101 , 1.396833931924577 , 1.46899998838737 ,   1.499271672687303 , 1.484944905422543 , 1.427299454044139 ,   1.33148461504109 , 1.206059242620877 \right] 
\]
\end{eulerformula}
\begin{eulerprompt}
>function fx(t) &= r(t)*cos(t); $'fx(t)=fx(t)
\end{eulerprompt}
\begin{eulerformula}
\[
{\it fx}\left(\left[ 0 , 0.1 , 0.2 , 0.3 , 0.4 , 0.5 , 0.6 , 0.7 ,   0.7999999999999999 , 0.8999999999999999 , 0.9999999999999999 , 1.1   , 1.2 , 1.3 , 1.4 , 1.5 , 1.6 , 1.7 , 1.8 , 1.900000000000001 , 2.0   , 2.100000000000001 , 2.200000000000001 , 2.300000000000001 ,   2.400000000000001 , 2.500000000000001 , 2.600000000000001 ,   2.700000000000001 , 2.800000000000001 , 2.900000000000001 ,   3.000000000000001 \right] \right)=\left[ 1 , 1.142026083553954 ,   1.256760186143285 , 1.329506878966172 , 1.350293417488142 ,   1.315274664798767 , 1.227211181539273 , 1.09495165732058 ,   0.9320065742506471 , 0.7544417651665002 , 0.5784260387475783 ,   0.4178197039080961 , 0.2821823974055013 , 0.1755105076499151 ,   0.09589752097113675 , 0.03616332913293813 , -0.01465575694556054 , -  0.06920143201758068 , -0.1394152395545122 , -0.2342741215913122 , -  0.3580078987182797 , -0.5090903206811853 , -0.6801723375047852 , -  0.85897629366313 , -1.030016563055952 , -1.176879961935062 , -  1.2847090345704 , -1.342497321422681 , -1.344833432424567 , -  1.292815858745206 , -1.193989600650285 \right] 
\]
\end{eulerformula}
\begin{eulerprompt}
>function fy(t) &= r(t)*sin(t); $'fy(t)=fy(t)
\end{eulerprompt}
\begin{eulerformula}
\[
{\it fy}\left(\left[ 0 , 0.1 , 0.2 , 0.3 , 0.4 , 0.5 , 0.6 , 0.7 ,   0.7999999999999999 , 0.8999999999999999 , 0.9999999999999999 , 1.1   , 1.2 , 1.3 , 1.4 , 1.5 , 1.6 , 1.7 , 1.8 , 1.900000000000001 , 2.0   , 2.100000000000001 , 2.200000000000001 , 2.300000000000001 ,   2.400000000000001 , 2.500000000000001 , 2.600000000000001 ,   2.700000000000001 , 2.800000000000001 , 2.900000000000001 ,   3.000000000000001 \right] \right)=\left[ 0 , 0.1145848126064173 ,   0.2547579019589912 , 0.4112646717695909 , 0.570894900220764 ,   0.7185378242080235 , 0.8395803408995154 , 0.9222650581299158 ,   0.9596299042728891 , 0.9507159900377483 , 0.900845180887014 ,   0.8209152982422039 , 0.7258159112220535 , 0.6322068292498618 ,   0.5560026751937331 , 0.5099542907913515 , 0.5017036794032687 ,   0.5326158897201473 , 0.597570198161595 , 0.685743199063247 ,   0.7822615300619328 , 0.8704663248428708 , 0.9344364398403995 ,   0.9613779857939977 , 0.9435098903595438 , 0.8791555727388753 ,   0.7728766040033648 , 0.6346355758333209 , 0.4781284038287776 ,   0.3185568010073166 , 0.1701990900393339 \right] 
\]
\end{eulerformula}
\begin{eulerprompt}
>function ds(t) &= trigreduce(radcan(sqrt(diff(fx(t),t)^2+diff(fy(t),t)^2))); $'ds(t)=ds(t)
\end{eulerprompt}
\begin{euleroutput}
  Maxima said:
  diff: second argument must be a variable; found 
      [0,0.1,0.2,0.3,0.4,0.5,0.6,0.7,0.7999999999999999,
       0.8999999999999999,0.9999999999999999,1.1,1.2,1.3,1.4,1.5,1.6,
       1.7,1.8,1.900000000000001,2.0,2.100000000000001,
       2.200000000000001,2.300000000000001,2.400000000000001,
       2.500000000000001,2.600000000000001,2.700000000000001,
       2.800000000000001,2.900000000000001,3.000000000000001]
   -- an error. To debug this try: debugmode(true);
  
  Error in:
  ... e(radcan(sqrt(diff(fx(t),t)^2+diff(fy(t),t)^2))); $'ds(t)=ds(t ...
                                                       ^
\end{euleroutput}
\begin{eulerprompt}
>$integrate(ds(x),x,0,2*pi) //panjang (keliling) kurva
\end{eulerprompt}
\begin{euleroutput}
  Maxima said:
  defint: variable of integration must be a simple or subscripted variable.
  defint: found errexp1
   -- an error. To debug this try: debugmode(true);
  
  Error in:
   $integrate(ds(x),x,0,2*pi) //panjang (keliling) kurva ...
                             ^
\end{euleroutput}
\begin{eulercomment}
Maxima gagal melakukan perhitungan eksak integral tersebut.

Berikut kita hitung integralnya secara umerik dengan perintah EMT.
\end{eulercomment}
\begin{eulerprompt}
>integrate("ds(x)",0,2*pi)
\end{eulerprompt}
\begin{euleroutput}
  Function ds not found.
  Try list ... to find functions!
  Error in expression: ds(x)
   %mapexpression1:
      return expr(x,args());
  Error in map.
   %evalexpression:
      if maps then return %mapexpression1(x,f$;args());
  gauss:
      if maps then y=%evalexpression(f$,a+h-(h*xn)',maps;args());
  adaptivegauss:
      t1=gauss(f$,c,c+h;args(),=maps);
  Try "trace errors" to inspect local variables after errors.
  integrate:
      return adaptivegauss(f$,a,b,eps*1000;args(),=maps);
\end{euleroutput}
\begin{eulercomment}
Spiral Logaritmik

\end{eulercomment}
\begin{eulerformula}
\[
x=e^{ax}\cos x,\ y=e^{ax}\sin x.
\]
\end{eulerformula}
\begin{eulerprompt}
>a=0.1; plot2d("exp(a*x)*cos(x)","exp(a*x)*sin(x)",r=2,xmin=0,xmax=2*pi):
\end{eulerprompt}
\eulerimg{17}{images/Shalih Abdillah_22305144009_ApKom5-155.png}
\begin{eulerprompt}
>&kill(a) // hapus expresi a
\end{eulerprompt}
\begin{euleroutput}
  
                                   done
  
\end{euleroutput}
\begin{eulerprompt}
>function fx(t) &= exp(a*t)*cos(t); $'fx(t)=fx(t)
\end{eulerprompt}
\begin{eulerformula}
\[
{\it fx}\left(\left[ 0 , 0.1 , 0.2 , 0.3 , 0.4 , 0.5 , 0.6 , 0.7 ,   0.7999999999999999 , 0.8999999999999999 , 0.9999999999999999 , 1.1   , 1.2 , 1.3 , 1.4 , 1.5 , 1.6 , 1.7 , 1.8 , 1.900000000000001 , 2.0   , 2.100000000000001 , 2.200000000000001 , 2.300000000000001 ,   2.400000000000001 , 2.500000000000001 , 2.600000000000001 ,   2.700000000000001 , 2.800000000000001 , 2.900000000000001 ,   3.000000000000001 \right] \right)=\left[ 1 , 0.9950041652780258\,e^{  0.1\,a} , 0.9800665778412416\,e^{0.2\,a} , 0.955336489125606\,e^{0.3  \,a} , 0.9210609940028851\,e^{0.4\,a} , 0.8775825618903728\,e^{0.5\,  a} , 0.8253356149096783\,e^{0.6\,a} , 0.7648421872844885\,e^{0.7\,a}   , 0.6967067093471655\,e^{0.7999999999999999\,a} ,   0.6216099682706645\,e^{0.8999999999999999\,a} , 0.5403023058681398\,  e^{0.9999999999999999\,a} , 0.4535961214255775\,e^{1.1\,a} ,   0.3623577544766736\,e^{1.2\,a} , 0.2674988286245873\,e^{1.3\,a} ,   0.1699671429002408\,e^{1.4\,a} , 0.07073720166770268\,e^{1.5\,a} , -  0.02919952230128904\,e^{1.6\,a} , -0.1288444942955251\,e^{1.7\,a} ,   -0.2272020946930875\,e^{1.8\,a} , -0.323289566863504\,e^{  1.900000000000001\,a} , -0.4161468365471428\,e^{2.0\,a} , -  0.5048461045998579\,e^{2.100000000000001\,a} , -0.5885011172553463\,  e^{2.200000000000001\,a} , -0.6662760212798248\,e^{2.300000000000001  \,a} , -0.737393715541246\,e^{2.400000000000001\,a} , -  0.8011436155469343\,e^{2.500000000000001\,a} , -0.8568887533689478\,  e^{2.600000000000001\,a} , -0.9040721420170617\,e^{2.700000000000001  \,a} , -0.9422223406686585\,e^{2.800000000000001\,a} , -  0.9709581651495908\,e^{2.900000000000001\,a} , -0.9899924966004456\,  e^{3.000000000000001\,a} \right] 
\]
\end{eulerformula}
\begin{eulerprompt}
>function fy(t) &= exp(a*t)*sin(t); $'fy(t)=fy(t)
\end{eulerprompt}
\begin{eulerformula}
\[
{\it fy}\left(\left[ 0 , 0.1 , 0.2 , 0.3 , 0.4 , 0.5 , 0.6 , 0.7 ,   0.7999999999999999 , 0.8999999999999999 , 0.9999999999999999 , 1.1   , 1.2 , 1.3 , 1.4 , 1.5 , 1.6 , 1.7 , 1.8 , 1.900000000000001 , 2.0   , 2.100000000000001 , 2.200000000000001 , 2.300000000000001 ,   2.400000000000001 , 2.500000000000001 , 2.600000000000001 ,   2.700000000000001 , 2.800000000000001 , 2.900000000000001 ,   3.000000000000001 \right] \right)=\left[ 0 , 0.09983341664682815\,e  ^{0.1\,a} , 0.1986693307950612\,e^{0.2\,a} , 0.2955202066613396\,e^{  0.3\,a} , 0.3894183423086505\,e^{0.4\,a} , 0.479425538604203\,e^{0.5  \,a} , 0.5646424733950354\,e^{0.6\,a} , 0.644217687237691\,e^{0.7\,a  } , 0.7173560908995227\,e^{0.7999999999999999\,a} ,   0.7833269096274833\,e^{0.8999999999999999\,a} , 0.8414709848078964\,  e^{0.9999999999999999\,a} , 0.8912073600614353\,e^{1.1\,a} ,   0.9320390859672263\,e^{1.2\,a} , 0.963558185417193\,e^{1.3\,a} ,   0.9854497299884603\,e^{1.4\,a} , 0.9974949866040544\,e^{1.5\,a} ,   0.9995736030415051\,e^{1.6\,a} , 0.9916648104524686\,e^{1.7\,a} ,   0.973847630878195\,e^{1.8\,a} , 0.9463000876874142\,e^{  1.900000000000001\,a} , 0.9092974268256815\,e^{2.0\,a} ,   0.8632093666488735\,e^{2.100000000000001\,a} , 0.8084964038195899\,e  ^{2.200000000000001\,a} , 0.7457052121767197\,e^{2.300000000000001\,  a} , 0.6754631805511503\,e^{2.400000000000001\,a} ,   0.5984721441039558\,e^{2.500000000000001\,a} , 0.5155013718214634\,e  ^{2.600000000000001\,a} , 0.4273798802338289\,e^{2.700000000000001\,  a} , 0.3349881501559038\,e^{2.800000000000001\,a} ,   0.2392493292139811\,e^{2.900000000000001\,a} , 0.1411200080598659\,e  ^{3.000000000000001\,a} \right] 
\]
\end{eulerformula}
\begin{eulerprompt}
>function df(t) &= trigreduce(radcan(sqrt(diff(fx(t),t)^2+diff(fy(t),t)^2))); $'df(t)=df(t)
\end{eulerprompt}
\begin{euleroutput}
  Maxima said:
  diff: second argument must be a variable; found 
      [0,0.1,0.2,0.3,0.4,0.5,0.6,0.7,0.7999999999999999,
       0.8999999999999999,0.9999999999999999,1.1,1.2,1.3,1.4,1.5,1.6,
       1.7,1.8,1.900000000000001,2.0,2.100000000000001,
       2.200000000000001,2.300000000000001,2.400000000000001,
       2.500000000000001,2.600000000000001,2.700000000000001,
       2.800000000000001,2.900000000000001,3.000000000000001]
   -- an error. To debug this try: debugmode(true);
  
  Error in:
  ... e(radcan(sqrt(diff(fx(t),t)^2+diff(fy(t),t)^2))); $'df(t)=df(t ...
                                                       ^
\end{euleroutput}
\begin{eulerprompt}
>S &=integrate(df(t),t,0,2*%pi); $S // panjang kurva (spiral)
\end{eulerprompt}
\begin{euleroutput}
  Maxima said:
  expt: undefined: 0 to a negative exponent.
  #0: df(x=[0,0.1,0.2,0.3,0.4,0.5,0.6,0.7,0.7999999999999999,0.8999999999999999,0.9999999999999999,1.1,1.2,1.3,...)
   -- an error. To debug this try: debugmode(true);
  
  Error in:
  S &=integrate(df(t),t,0,2*%pi); $S // panjang kurva (spiral) ...
                                ^
\end{euleroutput}
\begin{eulerprompt}
>S(a=0.1) // Panjang kurva untuk a=0.1
\end{eulerprompt}
\begin{euleroutput}
  Function S not found.
  Try list ... to find functions!
  Error in:
  S(a=0.1) // Panjang kurva untuk a=0.1 ...
          ^
\end{euleroutput}
\begin{eulercomment}
Soal:

Tunjukkan bahwa keliling lingkaran dengan jari-jari r adalah K=2.pi.r.

Berikut adalah contoh menghitung panjang parabola.
\end{eulercomment}
\begin{eulerprompt}
>plot2d("x^2",xmin=-1,xmax=1):
\end{eulerprompt}
\eulerimg{17}{images/Shalih Abdillah_22305144009_ApKom5-158.png}
\begin{eulerprompt}
>$showev('integrate(sqrt(1+diff(x^2,x)^2),x,-1,1))
\end{eulerprompt}
\begin{eulerformula}
\[
\int_{-1}^{1}{\sqrt{4\,x^2+1}\;dx}=\frac{{\rm asinh}\; 2+2\,\sqrt{5  }}{2}
\]
\end{eulerformula}
\begin{eulerprompt}
>$float(%)
\end{eulerprompt}
\begin{eulerformula}
\[
\int_{-1.0}^{1.0}{\sqrt{4.0\,x^2+1.0}\;dx}=2.957885715089195
\]
\end{eulerformula}
\begin{eulerprompt}
>x=-1:0.2:1; y=x^2; plot2d(x,y);  ...
>  plot2d(x,y,points=1,style="o#",add=1):
\end{eulerprompt}
\eulerimg{17}{images/Shalih Abdillah_22305144009_ApKom5-161.png}
\begin{eulercomment}
Panjang tersebut dapat dihampiri dengan menggunakan jumlah panjang ruas-ruas garis yang menghubungkan titik-titik pada parabola
tersebut.
\end{eulercomment}
\begin{eulerprompt}
>i=1:cols(x)-1; sum(sqrt((x[i+1]-x[i])^2+(y[i+1]-y[i])^2))
\end{eulerprompt}
\begin{euleroutput}
  2.95191957027
\end{euleroutput}
\begin{eulercomment}
Hasilnya mendekati panjang yang dihitung secara eksak. Untuk mendapatkan hampiran yang cukup akurat, jarak antar titik dapat
diperkecil, misalnya 0.1, 0.05, 0.01, dan seterusnya. Cobalah Anda ulangi perhitungannya dengan nilai-nilai tersebut.

\end{eulercomment}
\eulersubheading{Koordinat Kartesius}
\begin{eulercomment}
Berikut diberikan contoh perhitungan panjang kurva menggunakan koordinat Kartesius. Kita akan hitung panjang kurva dengan
persamaan implisit:

\end{eulercomment}
\begin{eulerformula}
\[
x^3+y^3-3xy=0.
\]
\end{eulerformula}
\begin{eulerprompt}
>z &= x^3+y^3-3*x*y; $z
\end{eulerprompt}
\begin{eulerformula}
\[
y^3-3\,x\,y+x^3
\]
\end{eulerformula}
\begin{eulerprompt}
>plot2d(z,r=2,level=0,n=100):
\end{eulerprompt}
\eulerimg{17}{images/Shalih Abdillah_22305144009_ApKom5-163.png}
\begin{eulercomment}
Kita tertarik pada kurva di kuadran pertama.
\end{eulercomment}
\begin{eulerprompt}
>plot2d(z,a=0,b=2,c=0,d=2,level=[-10;0],n=100,contourwidth=3,style="/"):
\end{eulerprompt}
\eulerimg{17}{images/Shalih Abdillah_22305144009_ApKom5-164.png}
\begin{eulercomment}
Kita selesaikan persamaannya untuk x.
\end{eulercomment}
\begin{eulerprompt}
>$z with y=l*x, sol &= solve(%,x); $sol
\end{eulerprompt}
\begin{eulerformula}
\[
\left[ x=\frac{3\,l}{l^3+1} , x=0 \right] 
\]
\end{eulerformula}
\eulerimg{1}{images/Shalih Abdillah_22305144009_ApKom5-166-large.png}
\begin{eulercomment}
Kita gunakan solusi tersebut untuk mendefinisikan fungsi dengan Maxima.
\end{eulercomment}
\begin{eulerprompt}
>function f(l) &= rhs(sol[1]); $'f(l)=f(l)
\end{eulerprompt}
\begin{eulerformula}
\[
f\left(l\right)=\frac{3\,l}{l^3+1}
\]
\end{eulerformula}
\begin{eulercomment}
Fungsi tersebut juga dapat digunaka untuk menggambar kurvanya. Ingat, bahwa fungsi tersebut adalah nilai x dan nilai y=l*x, yakni
x=f(l) dan y=l*f(l).
\end{eulercomment}
\begin{eulerprompt}
>plot2d(&f(x),&x*f(x),xmin=-0.5,xmax=2,a=0,b=2,c=0,d=2,r=1.5):
\end{eulerprompt}
\eulerimg{17}{images/Shalih Abdillah_22305144009_ApKom5-168.png}
\begin{eulercomment}
Elemen panjang kurva adalah:

\end{eulercomment}
\begin{eulerformula}
\[
ds=\sqrt{f'(l)^2+(lf'(l)+f(l))^2}.
\]
\end{eulerformula}
\begin{eulerprompt}
>function ds(l) &= ratsimp(sqrt(diff(f(l),l)^2+diff(l*f(l),l)^2)); $'ds(l)=ds(l)
\end{eulerprompt}
\begin{eulerformula}
\[
{\it ds}\left(l\right)=\frac{\sqrt{9\,l^8+36\,l^6-36\,l^5-36\,l^3+  36\,l^2+9}}{\sqrt{l^{12}+4\,l^9+6\,l^6+4\,l^3+1}}
\]
\end{eulerformula}
\begin{eulerprompt}
>$integrate(ds(l),l,0,1)
\end{eulerprompt}
\begin{eulerformula}
\[
\int_{0}^{1}{\frac{\sqrt{9\,l^8+36\,l^6-36\,l^5-36\,l^3+36\,l^2+9}  }{\sqrt{l^{12}+4\,l^9+6\,l^6+4\,l^3+1}}\;dl}
\]
\end{eulerformula}
\begin{eulercomment}
Integral tersebut tidak dapat dihitung secara eksak menggunakan Maxima. Kita hitung integral etrsebut secara numerik dengan Euler.
Karena kurva simetris, kita hitung untuk nilai variabel integrasi dari 0 sampai 1, kemudian hasilnya dikalikan 2. 
\end{eulercomment}
\begin{eulerprompt}
>2*integrate("ds(x)",0,1)
\end{eulerprompt}
\begin{euleroutput}
  4.91748872168
\end{euleroutput}
\begin{eulerprompt}
>2*romberg(&ds(x),0,1)// perintah Euler lain untuk menghitung nilai hampiran integral yang sama
\end{eulerprompt}
\begin{euleroutput}
  4.91748872168
\end{euleroutput}
\begin{eulercomment}
Perhitungan di datas dapat dilakukan untuk sebarang fungsi x dan y dengan mendefinisikan fungsi EMT, misalnya kita beri nama
panjangkurva. Fungsi ini selalu memanggil Maxima untuk menurunkan fungsi yang diberikan.
\end{eulercomment}
\begin{eulerprompt}
>function panjangkurva(fx,fy,a,b) ...
\end{eulerprompt}
\begin{eulerudf}
  ds=mxm("sqrt(diff(@fx,x)^2+diff(@fy,x)^2)");
  return romberg(ds,a,b);
  endfunction
\end{eulerudf}
\begin{eulerprompt}
>panjangkurva("x","x^2",-1,1) // cek untuk menghitung panjang kurva parabola sebelumnya
\end{eulerprompt}
\begin{euleroutput}
  2.95788571509
\end{euleroutput}
\begin{eulercomment}
Bandingkan dengan nilai eksak di atas.
\end{eulercomment}
\begin{eulerprompt}
>2*panjangkurva(mxm("f(x)"),mxm("x*f(x)"),0,1) // cek contoh terakhir, bandingkan hasilnya!
\end{eulerprompt}
\begin{euleroutput}
  4.91748872168
\end{euleroutput}
\begin{eulercomment}
Kita hitung panjang spiral Archimides berikut ini dengan fungsi tersebut.
\end{eulercomment}
\begin{eulerprompt}
>plot2d("x*cos(x)","x*sin(x)",xmin=0,xmax=2*pi,square=1):
\end{eulerprompt}
\eulerimg{17}{images/Shalih Abdillah_22305144009_ApKom5-171.png}
\begin{eulerprompt}
>panjangkurva("x*cos(x)","x*sin(x)",0,2*pi)
\end{eulerprompt}
\begin{euleroutput}
  21.2562941482
\end{euleroutput}
\begin{eulercomment}
Berikut kita definisikan fungsi yang sama namun dengan Maxima, untuk perhitungan eksak.
\end{eulercomment}
\begin{eulerprompt}
>&kill(ds,x,fx,fy)
\end{eulerprompt}
\begin{euleroutput}
  
                                   done
  
\end{euleroutput}
\begin{eulerprompt}
>function ds(fx,fy) &&= sqrt(diff(fx,x)^2+diff(fy,x)^2)
\end{eulerprompt}
\begin{euleroutput}
  
                             2              2
                    sqrt(diff (fy, x) + diff (fx, x))
  
\end{euleroutput}
\begin{eulerprompt}
>sol &= ds(x*cos(x),x*sin(x)); $sol // Kita gunakan untuk menghitung panjang kurva terakhir di atas
\end{eulerprompt}
\begin{eulerformula}
\[
\sqrt{\left(\cos x-x\,\sin x\right)^2+\left(\sin x+x\,\cos x\right)  ^2}
\]
\end{eulerformula}
\begin{eulerprompt}
>$sol | trigreduce | expand, $integrate(%,x,0,2*pi), %()
\end{eulerprompt}
\begin{eulerformula}
\[
\frac{{\rm asinh}\; \left(2\,\pi\right)+2\,\pi\,\sqrt{4\,\pi^2+1}}{  2}
\]
\end{eulerformula}
\eulerimg{1}{images/Shalih Abdillah_22305144009_ApKom5-174-large.png}
\begin{euleroutput}
  21.2562941482
\end{euleroutput}
\begin{eulercomment}
Hasilnya sama dengan perhitungan menggunakan fungsi EMT.

Berikut adalah contoh lain penggunaan fungsi Maxima tersebut.
\end{eulercomment}
\begin{eulerprompt}
>plot2d("3*x^2-1","3*x^3-1",xmin=-1/sqrt(3),xmax=1/sqrt(3),square=1):
\end{eulerprompt}
\eulerimg{17}{images/Shalih Abdillah_22305144009_ApKom5-175.png}
\begin{eulerprompt}
>sol &= radcan(ds(3*x^2-1,3*x^3-1)); $sol
\end{eulerprompt}
\begin{eulerformula}
\[
3\,x\,\sqrt{9\,x^2+4}
\]
\end{eulerformula}
\begin{eulerprompt}
>$showev('integrate(sol,x,0,1/sqrt(3))), $2*float(%) // panjang kurva di atas
\end{eulerprompt}
\begin{eulerformula}
\[
6.0\,\int_{0.0}^{0.5773502691896258}{x\,\sqrt{9.0\,x^2+4.0}\;dx}=  2.337835372767141
\]
\end{eulerformula}
\eulerimg{1}{images/Shalih Abdillah_22305144009_ApKom5-178-large.png}
\eulersubheading{Sikloid}
\begin{eulercomment}
Berikut kita akan menghitung panjang kurva lintasan (sikloid) suatu
titik pada lingkaran yang berputar ke kanan pada permukaan datar.
Misalkan jari-jari lingkaran tersebut adalah r. Posisi titik pusat
lingkaran pada saat t adalah:

\end{eulercomment}
\begin{eulerformula}
\[
(rt,r).
\]
\end{eulerformula}
\begin{eulercomment}
Misalkan posisi titik pada lingkaran tersebut mula-mula (0,0) dan
posisinya pada saat t adalah:

\end{eulercomment}
\begin{eulerformula}
\[
(r(t-\sin(t)),r(1-\cos(t))).
\]
\end{eulerformula}
\begin{eulercomment}
Berikut kita plot lintasan tersebut dan beberapa posisi lingkaran
ketika t=0, t=pi/2, t=r*pi.
\end{eulercomment}
\begin{eulerprompt}
>x &= r*(t-sin(t))
\end{eulerprompt}
\begin{euleroutput}
  
          [0, 1.665833531718508e-4 r, 0.001330669204938795 r, 
  0.004479793338660443 r, 0.0105816576913495 r, 0.02057446139579699 r, 
  0.03535752660496461 r, 0.05578231276230894 r, 0.08264390910047725 r, 
  0.1166730903725166 r, 0.1585290151921035 r, 0.2087926399385646 r, 
  0.2679609140327737 r, 0.3364418145828071 r, 0.4145502700115399 r, 
  0.5025050133959458 r, 0.6004263969584952 r, 0.7083351895475318 r, 
  0.8261523691218055 r, 0.9536999123125863 r, 1.090702573174319 r, 
  1.236790633351127 r, 1.391503596180411 r, 1.554294787823281 r, 
  1.724536819448851 r, 1.901527855896045 r, 2.084498628178538 r, 
  2.272620119766172 r, 2.465011849844097 r, 2.66075067078602 r, 
  2.858879991940135 r]
  
\end{euleroutput}
\begin{eulerprompt}
>y &= r*(1-cos(t))
\end{eulerprompt}
\begin{euleroutput}
  
          [0, 0.004995834721974179 r, 0.01993342215875837 r, 
  0.04466351087439402 r, 0.0789390059971149 r, 0.1224174381096272 r, 
  0.1746643850903217 r, 0.2351578127155115 r, 0.3032932906528345 r, 
  0.3783900317293355 r, 0.4596976941318602 r, 0.5464038785744225 r, 
  0.6376422455233264 r, 0.7325011713754126 r, 0.8300328570997592 r, 
  0.9292627983322973 r, 1.029199522301289 r, 1.128844494295525 r, 
  1.227202094693087 r, 1.323289566863504 r, 1.416146836547143 r, 
  1.504846104599858 r, 1.588501117255346 r, 1.666276021279825 r, 
  1.737393715541246 r, 1.801143615546934 r, 1.856888753368948 r, 
  1.904072142017062 r, 1.942222340668659 r, 1.970958165149591 r, 
  1.989992496600446 r]
  
\end{euleroutput}
\begin{eulerprompt}
>1
\end{eulerprompt}
\begin{euleroutput}
  1
\end{euleroutput}
\begin{eulercomment}
Berikut kita gambar sikloid untuk r=1.
\end{eulercomment}
\begin{eulerprompt}
>ex &= x-sin(x); ey &= 1-cos(x); aspect(1);
>plot2d(ex,ey,xmin=0,xmax=4pi,square=1); ...
>  plot2d("2+cos(x)","1+sin(x)",xmin=0,xmax=2pi,>add,color=blue); ...
>  plot2d([2,ex(2)],[1,ey(2)],color=red,>add); ...
>  plot2d(ex(2),ey(2),>points,>add,color=red); ...
>  plot2d("2pi+cos(x)","1+sin(x)",xmin=0,xmax=2pi,>add,color=blue); ...
>  plot2d([2pi,ex(2pi)],[1,ey(2pi)],color=red,>add);  ...
>  plot2d(ex(2pi),ey(2pi),>points,>add,color=red):
\end{eulerprompt}
\eulerimg{27}{images/Shalih Abdillah_22305144009_ApKom5-181.png}
\begin{eulercomment}
Berikut dihitung panjang lintasan untuk 1 putaran penuh. (Jangan salah menduga bahwa panjang lintasan 1 putaran penuh sama dengan
keliling lingkaran!)
\end{eulercomment}
\begin{eulerprompt}
>ds &= radcan(sqrt(diff(ex,x)^2+diff(ey,x)^2)); $ds=trigsimp(ds) // elemen panjang kurva sikloid
\end{eulerprompt}
\begin{eulerformula}
\[
\sqrt{\sin ^2x+\cos ^2x-2\,\cos x+1}=\sqrt{2-2\,\cos x}
\]
\end{eulerformula}
\begin{eulerprompt}
>ds &= trigsimp(ds); $ds
\end{eulerprompt}
\begin{eulerformula}
\[
\sqrt{2-2\,\cos x}
\]
\end{eulerformula}
\begin{eulerprompt}
>$showev('integrate(ds,x,0,2*pi)) // hitung panjang sikloid satu putaran penuh
\end{eulerprompt}
\begin{eulerformula}
\[
\int_{0}^{2\,\pi}{\sqrt{2-2\,\cos x}\;dx}=8
\]
\end{eulerformula}
\begin{eulerprompt}
>integrate(mxm("ds"),0,2*pi) // hitung secara numerik
\end{eulerprompt}
\begin{euleroutput}
  8
\end{euleroutput}
\begin{eulerprompt}
>romberg(mxm("ds"),0,2*pi) // cara lain hitung secara numerik
\end{eulerprompt}
\begin{euleroutput}
  8
\end{euleroutput}
\begin{eulercomment}
Perhatikan, seperti terlihat pada gambar, panjang sikloid lebih besar
daripada keliling lingkarannya, yakni:

\end{eulercomment}
\begin{eulerformula}
\[
2\pi.
\]
\end{eulerformula}
\eulersubheading{Kurvatur (Kelengkungan) Kurva}
\begin{eulercomment}
image: Osculating.png

Aslinya, kelengkungan kurva diferensiabel (yakni, kurva mulus yang
tidak lancip) di titik P didefinisikan melalui lingkaran oskulasi
(yaitu, lingkaran yang melalui titik P dan terbaik memperkirakan,
paling banyak menyinggung kurva di sekitar P). Pusat dan radius
kelengkungan kurva di P adalah pusat dan radius lingkaran oskulasi.
Kelengkungan adalah kebalikan dari radius kelengkungan:

\end{eulercomment}
\begin{eulerformula}
\[
\kappa =\frac {1}{R}
\]
\end{eulerformula}
\begin{eulercomment}
dengan R adalah radius kelengkungan. (Setiap lingkaran memiliki
kelengkungan ini pada setiap titiknya, dapat diartikan, setiap
lingkaran berputar 2pi sejauh 2piR.)\\
Definisi ini sulit dimanipulasi dan dinyatakan ke dalam rumus untuk
kurva umum. Oleh karena itu digunakan definisi lain yang ekivalen.

\end{eulercomment}
\eulersubheading{Definisi Kurvatur dengan Fungsi Parametrik Panjang Kurva}
\begin{eulercomment}
Setiap kurva diferensiabel dapat dinyatakan dengan persamaan
parametrik terhadap panjang kurva s:

\end{eulercomment}
\begin{eulerformula}
\[
\gamma(s) = (x(s),\ y(s)),
\]
\end{eulerformula}
\begin{eulercomment}
dengan x dan y adalah fungsi riil yang diferensiabel, yang memenuhi:

\end{eulercomment}
\begin{eulerformula}
\[
\|\gamma'(s)\|=\sqrt{x'(s)^2+y'(s)^2}=1.
\]
\end{eulerformula}
\begin{eulercomment}
Ini berarti bahwa vektor singgung


\end{eulercomment}
\begin{eulerformula}
\[
\mathbf{T}(s)=(x'(s),\ y'(s))
\]
\end{eulerformula}
\begin{eulercomment}
memiliki norm 1 dan merupakan vektor singgung satuan.

Apabila kurvanya memiliki turunan kedua, artinya turunan kedua x dan y
ada, maka T'(s) ada. Vektor ini merupakan normal kurva yang arahnya
menuju pusat kurvatur, norm-nya merupakan nilai kurvatur
(kelengkungan):

\end{eulercomment}
\begin{eulerformula}
\[
 \begin{aligned}\mathbf{T}(s) &= \mathbf{\gamma}'(s),\\ \mathbf{T}^{2}(s) &=1\ \text{(konstanta)}\Rightarrow \mathbf{T}'(s)\cdot \mathbf{T}(s)=0\\ \kappa(s) &=\|\mathbf {T}'(s)\|= \|\mathbf{\gamma}''(s)\|=\sqrt{x''(s)^{2}+y''(s)^{2}}.\end{aligned}
\]
\end{eulerformula}
\begin{eulercomment}
Nilai

\end{eulercomment}
\begin{eulerformula}
\[
R(s)=\frac{1}{\kappa(s)}
\]
\end{eulerformula}
\begin{eulercomment}
disebut jari-jari (radius) kelengkungan kurva.

Bilangan riil

\end{eulercomment}
\begin{eulerformula}
\[
 k(s) = \pm\kappa(s)
\]
\end{eulerformula}
\begin{eulercomment}
disebut nilai kelengkungan bertanda.

Contoh:\\
Akan ditentukan kurvatur lingkaran

\end{eulercomment}
\begin{eulerformula}
\[
x=r\cos t,\ y= r\sin t.
\]
\end{eulerformula}
\begin{eulerprompt}
>fx &= r*cos(t); fy &=r*sin(t);
>&assume(t>0,r>0); s &=integrate(sqrt(diff(fx,t)^2+diff(fy,t)^2),t,0,t); s // elemen panjang kurva, panjang busur lingkaran (s)
\end{eulerprompt}
\begin{euleroutput}
  Maxima said:
  diff: second argument must be a variable; found 
      [0,0.1,0.2,0.3,0.4,0.5,0.6,0.7,0.7999999999999999,
       0.8999999999999999,0.9999999999999999,1.1,1.2,1.3,1.4,1.5,1.6,
       1.7,1.8,1.900000000000001,2.0,2.100000000000001,
       2.200000000000001,2.300000000000001,2.400000000000001,
       2.500000000000001,2.600000000000001,2.700000000000001,
       2.800000000000001,2.900000000000001,3.000000000000001]
   -- an error. To debug this try: debugmode(true);
  
  Error in:
  ... =integrate(sqrt(diff(fx,t)^2+diff(fy,t)^2),t,0,t); s // elemen ...
                                                       ^
\end{euleroutput}
\begin{eulerprompt}
>&kill(s); fx &= r*cos(s/r); fy &=r*sin(s/r); // definisi ulang persamaan parametrik terhadap s dengan substitusi t=s/r
>k &= trigsimp(sqrt(diff(fx,s,2)^2+diff(fy,s,2)^2)); $k // nilai kurvatur lingkaran dengan menggunakan definisi di atas
\end{eulerprompt}
\begin{eulerformula}
\[
\frac{1}{r}
\]
\end{eulerformula}
\begin{eulercomment}
Untuk representasi parametrik umum, misalkan

\end{eulercomment}
\begin{eulerformula}
\[
x = x(t),\ y= y(t)
\]
\end{eulerformula}
\begin{eulercomment}
merupakan persamaan parametrik untuk kurva bidang yang
terdiferensialkan dua kali. Kurvatur untuk kurva tersebut
didefinisikan sebagai

\end{eulercomment}
\begin{eulerformula}
\[
\begin{aligned}\kappa &= \frac{d\phi}{ds}=\frac{\frac{d\phi}{dt}}{\frac{ds}{dt}}\quad (\phi \text{ adalah sudut kemiringan garis singgung dan }s \text{ adalah panjang kurva})\\ &=\frac{\frac{d\phi}{dt}}{\sqrt{(\frac{dx}{dt})^2+(\frac{dy}{dt})^2}}= \frac{\frac{d\phi}{dt}}{\sqrt{x'(t)^2+y'(t)^2}}.\end{aligned}.
\]
\end{eulerformula}
\begin{eulercomment}
Selanjutnya, pembilang pada persamaan di atas dapat dicari sebagai
berikut.

\end{eulercomment}
\begin{eulerformula}
\[
\begin{aligned}\sec^2\phi\frac{d\phi}{dt} &= \frac{d}{dt}\left(\tan\phi\right)= \frac{d}{dt}\left(\frac{dy}{dx}\right)= \frac{d}{dt}\left(\frac{dy/dt}{dx/dt}\right)= \frac{d}{dt}\left(\frac{y'(t)}{x'(t)}\right)=\frac{x'(t)y''(t)-x''(t)y'(t)}{x'(t)^2}.\\ & \\ \frac{d\phi}{dt} &= \frac{1}{\sec^2\phi}\frac{x'(t)y''(t)-x''(t)y'(t)}{x'(t)^2}\\ &= \frac{1}{1+\tan^2\phi}\frac{x'(t)y''(t)-x''(t)y'(t)}{x'(t)^2}\\ &= \frac{1}{1+\left(\frac{y'(t)}{x'(t)}\right)^2}\frac{x'(t)y''(t)-x''(t)y'(t)}{x'(t)^2}\\ &= \frac{x'(t)y''(t)-x''(t)y'(t)}{x'(t)^2+y'(t)^2}.\end{aligned}
\]
\end{eulerformula}
\begin{eulercomment}
Jadi, rumus kurvatur untuk kurva parametrik

\end{eulercomment}
\begin{eulerformula}
\[
x=x(t),\ y=y(t)
\]
\end{eulerformula}
\begin{eulercomment}
adalah

\end{eulercomment}
\begin{eulerformula}
\[
\kappa(t) = \frac{x'(t)y''(t)-x''(t)y'(t)}{\left(x'(t)^2+y'(t)^2\right)^{3/2}}.
\]
\end{eulerformula}
\begin{eulercomment}
Jika kurvanya dinyatakan dengan persamaan parametrik pada koordinat
kutub

\end{eulercomment}
\begin{eulerformula}
\[
x=r(\theta)\cos\theta,\ y=r(\theta)\sin\theta,
\]
\end{eulerformula}
\begin{eulercomment}
maka rumus kurvaturnya adalah

\end{eulercomment}
\begin{eulerformula}
\[
\kappa(\theta) = \frac{r(\theta)^2+2r'(\theta)^2-r(\theta)r''(\theta)}{\left(r'(\theta)^2+r'(\theta)^2\right)^{3/2}}.
\]
\end{eulerformula}
\begin{eulercomment}
(Silakan Anda turunkan rumus tersebut!)

Contoh:\\
Lingkaran dengan pusat (0,0) dan jari-jari r dapat dinyatakan dengan
persamaan parametrik

\end{eulercomment}
\begin{eulerformula}
\[
x=r\cos t,\ y=r\sin t.
\]
\end{eulerformula}
\begin{eulercomment}
Nilai kelengkungan lingkaran tersebut adalah

\end{eulercomment}
\begin{eulerformula}
\[
\kappa(t)=\frac{x'(t)y''(t)-x''(t)y'(t)}{\left(x'(t)^2+y'(t)^2\right)^{3/2}}=\frac{r^2}{r^3}=\frac 1 r.
\]
\end{eulerformula}
\begin{eulercomment}
Hasil cocok dengan definisi kurvatur suatu kelengkungan.
\end{eulercomment}
\begin{eulercomment}
Kurva

\end{eulercomment}
\begin{eulerformula}
\[
y=f(x)
\]
\end{eulerformula}
\begin{eulercomment}
dapat dinyatakan ke dalam persamaan parametrik

\end{eulercomment}
\begin{eulerformula}
\[
x=t,\ y=f(t),\ \text{ dengan } x'(t)=1,\ x''(t)=0,
\]
\end{eulerformula}
\begin{eulercomment}
sehingga kurvaturnya adalah

\end{eulercomment}
\begin{eulerformula}
\[
\kappa(t) = \frac{y''(t)}{\left(1+y'(t)^2\right)^{3/2}}.
\]
\end{eulerformula}
\begin{eulercomment}
Contoh:\\
Akan ditentukan kurvatur parabola

\end{eulercomment}
\begin{eulerformula}
\[
y=ax^2+bx+c.
\]
\end{eulerformula}
\begin{eulerprompt}
>function f(x) &= a*x^2+b*x+c; $y=f(x)
\end{eulerprompt}
\begin{eulerformula}
\[
y=a\,x^2+b\,x+c
\]
\end{eulerformula}
\begin{eulerprompt}
>function k(x) &= (diff(f(x),x,2))/(1+diff(f(x),x)^2)^(3/2); $'k(x)=k(x) // kelengkungan parabola 
\end{eulerprompt}
\begin{eulerformula}
\[
k\left(x\right)=\frac{2\,a}{\left(\left(2\,a\,x+b\right)^2+1\right)  ^{\frac{3}{2}}}
\]
\end{eulerformula}
\begin{eulerprompt}
>function f(x) &= x^2+x+1; $y=f(x) // akan kita plot kelengkungan parabola untuk a=b=c=1
\end{eulerprompt}
\begin{eulerformula}
\[
y=x^2+x+1
\]
\end{eulerformula}
\begin{eulerprompt}
>function k(x) &= (diff(f(x),x,2))/(1+diff(f(x),x)^2)^(3/2); $'k(x)=k(x) // kelengkungan parabola 
\end{eulerprompt}
\begin{eulerformula}
\[
k\left(x\right)=\frac{2}{\left(\left(2\,x+1\right)^2+1\right)^{  \frac{3}{2}}}
\]
\end{eulerformula}
\begin{eulercomment}
Berikut kita gambar parabola tersebut beserta kurva kelengkungan,
kurva jari-jari kelengkungan dan salah satu lingkaran oskulasi di
titik puncak parabola. Perhatikan, puncak parabola dan jari-jari
lingkaran oskulasi di puncak parabola adalah

\end{eulercomment}
\begin{eulerformula}
\[
(-1/2,3/4),\ 1/k(2)=1/2,
\]
\end{eulerformula}
\begin{eulercomment}
sehingga pusat lingkaran oskulasi adalah (-1/2, 5/4).
\end{eulercomment}
\begin{eulerprompt}
>plot2d(["f(x)", "k(x)"],-2,1, color=[blue,red]); plot2d("1/k(x)",-1.5,1,color=green,>add); ...
>plot2d("-1/2+1/k(-1/2)*cos(x)","5/4+1/k(-1/2)*sin(x)",xmin=0,xmax=2pi,>add,color=blue):
\end{eulerprompt}
\eulerimg{17}{images/Shalih Abdillah_22305144009_ApKom5-213.png}
\begin{eulercomment}
Untuk kurva yang dinyatakan dengan fungsi implisit

\end{eulercomment}
\begin{eulerformula}
\[
F(x,y)=0
\]
\end{eulerformula}
\begin{eulercomment}
dengan turunan-turunan parsial

\end{eulercomment}
\begin{eulerformula}
\[
F_x=\frac{\partial F}{\partial x},\ F_y=\frac{\partial F}{\partial y},\ F_{xy}=\frac{\partial}{\partial y}\left(\frac{\partial F}{\partial x}\right),\ F_{xx}=\frac{\partial}{\partial x}\left(\frac{\partial F}{\partial x}\right),\ F_{yy}=\frac{\partial}{\partial y}\left(\frac{\partial F}{\partial y}\right),
\]
\end{eulerformula}
\begin{eulercomment}
berlaku

\end{eulercomment}
\begin{eulerformula}
\[
F_x dx+ F_y dy = 0\text{ atau } \frac{dy}{dx}=-\frac{F_x}{F_y},
\]
\end{eulerformula}
\begin{eulercomment}
sehingga kurvaturnya adalah

\end{eulercomment}
\begin{eulerformula}
\[
\kappa =\frac {F_y^2F_{xx}-2F_xF_yF_{xy}+F_x^2F_{yy}}{\left(F_x^2+F_y^2\right)^{3/2}}.
\]
\end{eulerformula}
\begin{eulercomment}
(Silakan Anda turunkan sendiri!)

Contoh 1:\\
Parabola

\end{eulercomment}
\begin{eulerformula}
\[
y=ax^2+bx+c
\]
\end{eulerformula}
\begin{eulercomment}
dapat dinyatakan ke dalam persamaan implisit

\end{eulercomment}
\begin{eulerformula}
\[
ax^2+bx+c-y=0.
\]
\end{eulerformula}
\begin{eulerprompt}
>function F(x,y) &=a*x^2+b*x+c-y; $F(x,y)
\end{eulerprompt}
\begin{eulerformula}
\[
-y+a\,x^2+b\,x+c
\]
\end{eulerformula}
\begin{eulerprompt}
>Fx &= diff(F(x,y),x), Fxx &=diff(F(x,y),x,2), Fy &=diff(F(x,y),y), Fxy &=diff(diff(F(x,y),x),y), Fyy &=diff(F(x,y),y,2)  
\end{eulerprompt}
\begin{euleroutput}
  
                                2 a x + b
  
  
                                   2 a
  
  
                                   - 1
  
  
                                    0
  
  
                                    0
  
\end{euleroutput}
\begin{eulerprompt}
>function k(x) &= (Fy^2*Fxx-2*Fx*Fy*Fxy+Fx^2*Fyy)/(Fx^2+Fy^2)^(3/2); $'k(x)=k(x) // kurvatur parabola tersebut
\end{eulerprompt}
\begin{eulerformula}
\[
k\left(x\right)=\frac{2\,a}{\left(\left(2\,a\,x+b\right)^2+1\right)  ^{\frac{3}{2}}}
\]
\end{eulerformula}
\begin{eulercomment}
Hasilnya sama dengan sebelumnya yang menggunakan persamaan parabola biasa.
\end{eulercomment}
\eulerheading{Latihan}
\begin{eulercomment}
- Bukalah buku Kalkulus.\\
- Cari dan pilih beberapa (paling sedikit 5 fungsi berbeda
tipe/bentuk/jenis) fungsi dari buku tersebut, kemudian definisikan di
EMT pada baris-baris perintah berikut (jika perlu tambahkan lagi).\\
- Untuk setiap fungsi, tentukan anti turunannya (jika ada), hitunglah
integral tentu dengan batas-batas yang menarik (Anda tentukan
sendiri), seperti contoh-contoh tersebut.\\
- Lakukan hal yang sama untuk fungsi-fungsi yang tidak dapat
diintegralkan (cari sedikitnya 3 fungsi).\\
- Gambar grafik fungsi dan daerah integrasinya pada sumbu koordinat
yang sama.\\
- Gunakan integral tentu untuk mencari luas daerah yang dibatasi oleh
dua kurva yang berpotongan di dua titik. (Cari dan gambar kedua kurva
dan arsir (warnai) daerah yang dibatasi oleh keduanya.)\\
- Gunakan integral tentu untuk menghitung volume benda putar kurva y=
f(x) yang diputar mengelilingi sumbu x dari x=a sampai x=b, yakni

\end{eulercomment}
\begin{eulerformula}
\[
V = \int_a^b \pi (f(x)^2\ dx.
\]
\end{eulerformula}
\begin{eulercomment}
(Pilih fungsinya dan gambar kurva dan benda putar yang dihasilkan.
Anda dapat mencari contoh-contoh bagaimana cara menggambar benda hasil
perputaran suatu kurva.)\\
- Gunakan integral tentu untuk menghitung panjang kurva y=f(x) dari
x=a sampai x=b dengan menggunakan rumus:

\end{eulercomment}
\begin{eulerformula}
\[
S = \int_a^b \sqrt{1+(f'(x))^2} \ dx.
\]
\end{eulerformula}
\begin{eulercomment}
(Pilih fungsi dan gambar kurvanya.)\\
- Apabila fungsi dinyatakan dalam koordinat kutub x=f(r,t), y=g(r,t),
r=h(t), x=a bersesuaian dengan t=t0 dan x=b bersesuian dengan t=t1,
maka rumus di atas akan menjadi:

\end{eulercomment}
\begin{eulerformula}
\[
S=\int_{t_0}^{t_1} \sqrt{x'(t)^2+y'(t)^2}\ dt.
\]
\end{eulerformula}
\begin{eulercomment}
- Pilih beberapa kurva menarik (selain lingkaran dan parabola) dari
buku  kalkulus. Nyatakan setiap kurva tersebut dalam bentuk:\\
\end{eulercomment}
\begin{eulerttcomment}
  a. koordinat Kartesius (persamaan y=f(x))
  b. koordinat kutub ( r=r(theta))
  c. persamaan parametrik x=x(t), y=y(t)
  d. persamaan implit F(x,y)=0
\end{eulerttcomment}
\begin{eulercomment}
- Tentukan kurvatur masing-masing kurva dengan menggunakan keempat
representasi tersebut (hasilnya harus sama).\\
- Gambarlah kurva asli, kurva kurvatur, kurva jari-jari lingkaran
oskulasi, dan salah satu lingkaran oskulasinya.

Fungsi 1
\end{eulercomment}
\begin{eulerprompt}
>function f(x) &= x^(2*x); $f(x)
\end{eulerprompt}
\begin{eulerformula}
\[
\left[ 0 , 1.665833531718508 \times 10^{-4}\,r ,   0.001330669204938795\,r , 0.004479793338660443\,r ,   0.0105816576913495\,r , 0.02057446139579699\,r , 0.03535752660496461  \,r , 0.05578231276230894\,r , 0.08264390910047725\,r ,   0.1166730903725166\,r , 0.1585290151921035\,r , 0.2087926399385646\,  r , 0.2679609140327737\,r , 0.3364418145828071\,r ,   0.4145502700115399\,r , 0.5025050133959458\,r , 0.6004263969584952\,  r , 0.7083351895475318\,r , 0.8261523691218055\,r ,   0.9536999123125863\,r , 1.090702573174319\,r , 1.236790633351127\,r   , 1.391503596180411\,r , 1.554294787823281\,r , 1.724536819448851\,  r , 1.901527855896045\,r , 2.084498628178538\,r , 2.272620119766172  \,r , 2.465011849844097\,r , 2.66075067078602\,r , 2.858879991940135  \,r \right] ^{\left[ 0 , 3.331667063437016 \times 10^{-4}\,r ,   0.002661338409877589\,r , 0.008959586677320885\,r ,   0.021163315382699\,r , 0.04114892279159399\,r , 0.07071505320992921  \,r , 0.1115646255246179\,r , 0.1652878182009545\,r ,   0.2333461807450332\,r , 0.317058030384207\,r , 0.4175852798771291\,r   , 0.5359218280655473\,r , 0.6728836291656142\,r ,   0.8291005400230798\,r , 1.005010026791892\,r , 1.20085279391699\,r   , 1.416670379095064\,r , 1.652304738243611\,r , 1.907399824625173\,  r , 2.181405146348638\,r , 2.473581266702254\,r , 2.783007192360822  \,r , 3.108589575646562\,r , 3.449073638897701\,r , 3.80305571179209  \,r , 4.168997256357075\,r , 4.545240239532344\,r ,   4.930023699688195\,r , 5.32150134157204\,r , 5.717759983880271\,r   \right] }
\]
\end{eulerformula}
\begin{eulerprompt}
>function intf(x) &= integrate(f(x),x); $intf(x)
\end{eulerprompt}
\begin{eulerformula}
\[
\left[ 0 , 1.665833531718508 \times 10^{-4}\,r ,   0.001330669204938795\,r , 0.004479793338660443\,r ,   0.0105816576913495\,r , 0.02057446139579699\,r , 0.03535752660496461  \,r , 0.05578231276230894\,r , 0.08264390910047725\,r ,   0.1166730903725166\,r , 0.1585290151921035\,r , 0.2087926399385646\,  r , 0.2679609140327737\,r , 0.3364418145828071\,r ,   0.4145502700115399\,r , 0.5025050133959458\,r , 0.6004263969584952\,  r , 0.7083351895475318\,r , 0.8261523691218055\,r ,   0.9536999123125863\,r , 1.090702573174319\,r , 1.236790633351127\,r   , 1.391503596180411\,r , 1.554294787823281\,r , 1.724536819448851\,  r , 1.901527855896045\,r , 2.084498628178538\,r , 2.272620119766172  \,r , 2.465011849844097\,r , 2.66075067078602\,r , 2.858879991940135  \,r \right] ^{\left[ 0 , 3.331667063437016 \times 10^{-4}\,r ,   0.002661338409877589\,r , 0.008959586677320885\,r ,   0.021163315382699\,r , 0.04114892279159399\,r , 0.07071505320992921  \,r , 0.1115646255246179\,r , 0.1652878182009545\,r ,   0.2333461807450332\,r , 0.317058030384207\,r , 0.4175852798771291\,r   , 0.5359218280655473\,r , 0.6728836291656142\,r ,   0.8291005400230798\,r , 1.005010026791892\,r , 1.20085279391699\,r   , 1.416670379095064\,r , 1.652304738243611\,r , 1.907399824625173\,  r , 2.181405146348638\,r , 2.473581266702254\,r , 2.783007192360822  \,r , 3.108589575646562\,r , 3.449073638897701\,r , 3.80305571179209  \,r , 4.168997256357075\,r , 4.545240239532344\,r ,   4.930023699688195\,r , 5.32150134157204\,r , 5.717759983880271\,r   \right] }\,\left[ 0 , 1.665833531718508 \times 10^{-4}\,r ,   0.001330669204938795\,r , 0.004479793338660443\,r ,   0.0105816576913495\,r , 0.02057446139579699\,r , 0.03535752660496461  \,r , 0.05578231276230894\,r , 0.08264390910047725\,r ,   0.1166730903725166\,r , 0.1585290151921035\,r , 0.2087926399385646\,  r , 0.2679609140327737\,r , 0.3364418145828071\,r ,   0.4145502700115399\,r , 0.5025050133959458\,r , 0.6004263969584952\,  r , 0.7083351895475318\,r , 0.8261523691218055\,r ,   0.9536999123125863\,r , 1.090702573174319\,r , 1.236790633351127\,r   , 1.391503596180411\,r , 1.554294787823281\,r , 1.724536819448851\,  r , 1.901527855896045\,r , 2.084498628178538\,r , 2.272620119766172  \,r , 2.465011849844097\,r , 2.66075067078602\,r , 2.858879991940135  \,r \right] 
\]
\end{eulerformula}
\begin{eulerprompt}
>x=0:0.01:2.5-0.01; plot2d(x,f(x+0.01),>bar); plot2d("f(x)",0,2.5,>add):
\end{eulerprompt}
\begin{euleroutput}
  Error : f(x) does not produce a real or column vector
  
  Error generated by error() command
  
   %ploteval:
      error(f$|" does not produce a real or column vector"); 
  adaptiveevalone:
      s=%ploteval(g$,t;args());
  Try "trace errors" to inspect local variables after errors.
  plot2d:
      dw/n,dw/n^2,dw/n,auto;args());
\end{euleroutput}
\begin{eulerprompt}
>t &= makelist(a,a,0,2.5-0.01,0.01);
>fx &= makelist(f(t[i]+0.01),i,1,length(t));
\end{eulerprompt}
\begin{euleroutput}
  
  Maxima output too long!
  Error in:
  fx &= makelist(f(t[i]+0.01),i,1,length(t)); ...
                                            ^
\end{euleroutput}
\begin{eulerprompt}
>>0.01*sum(f(x+0.01))
\end{eulerprompt}
\begin{euleroutput}
  Syntax error in expression, or unfinished expression!
  Error in:
   >0.01*sum(f(x+0.01)) ...
  ^
\end{euleroutput}
\begin{eulerprompt}
>  
\end{eulerprompt}
\eulerheading{Barisan dan Deret}
\begin{eulercomment}
(Catatan: bagian ini belum lengkap. Anda dapat membaca contoh-contoh pengguanaan EMT dan
Maxima untuk menghitung limit barisan, rumus jumlah parsial suatu deret, jumlah tak hingga
suatu deret konvergen, dan sebagainya. Anda dapat mengeksplor contoh-contoh di EMT atau
perbagai panduan penggunaan Maxima di software Maxima atau dari Internet.)

Barisan dapat didefinisikan dengan beberapa cara di dalam EMT, di antaranya:

- dengan cara yang sama seperti mendefinisikan vektor dengan elemen-elemen beraturan
(menggunakan titik dua ":");\\
- menggunakan perintah "sequence" dan rumus barisan (suku ke -n);\\
- menggunakan perintah "iterate" atau "niterate";\\
- menggunakan fungsi Maxima "create\_list" atau "makelist" untuk menghasilkan barisan
simbolik;\\
- menggunakan fungsi biasa yang inputnya vektor atau barisan;\\
- menggunakan fungsi rekursif.

EMT menyediakan beberapa perintah (fungsi) terkait barisan, yakni:

- sum: menghitung jumlah semua elemen suatu barisan\\
- cumsum: jumlah kumulatif suatu barisan\\
- differences: selisih antar elemen-elemen berturutan

EMT juga dapat digunakan untuk menghitung jumlah deret berhingga maupun deret tak hingga,
dengan menggunakan perintah (fungsi) "sum". Perhitungan dapat dilakukan secara numerik
maupun simbolik dan eksak.

Berikut adalah beberapa contoh perhitungan barisan dan deret menggunakan EMT.
\end{eulercomment}
\begin{eulerprompt}
>1:10 // barisan sederhana
\end{eulerprompt}
\begin{euleroutput}
  [1,  2,  3,  4,  5,  6,  7,  8,  9,  10]
\end{euleroutput}
\begin{eulerprompt}
>1:2:30
\end{eulerprompt}
\begin{euleroutput}
  [1,  3,  5,  7,  9,  11,  13,  15,  17,  19,  21,  23,  25,  27,  29]
\end{euleroutput}
\eulerheading{Iterasi dan Barisan}
\begin{eulercomment}
EMT menyediakan fungsi iterate("g(x)", x0, n) untuk melakukan iterasi

\end{eulercomment}
\begin{eulerformula}
\[
x_{k+1}=g(x_k), \ x_0=x_0, k= 1, 2, 3, ..., n.
\]
\end{eulerformula}
\begin{eulercomment}
Berikut ini disajikan contoh-contoh penggunaan iterasi dan rekursi
dengan EMT. Contoh pertama menunjukkan pertumbuhan dari nilai awal
1000 dengan laju pertambahan 5\%, selama 10 periode.
\end{eulercomment}
\begin{eulerprompt}
>q=1.05; iterate("x*q",1000,n=10)'
\end{eulerprompt}
\begin{euleroutput}
           1000 
           1050 
         1102.5 
        1157.63 
        1215.51 
        1276.28 
         1340.1 
         1407.1 
        1477.46 
        1551.33 
        1628.89 
\end{euleroutput}
\begin{eulercomment}
Contoh berikutnya memperlihatkan bahaya menabung di bank pada masa sekarang! Dengan bunga
tabungan sebesar 6\% per tahun atau 0.5\% per bulan dipotong pajak 20\%, dan biaya administrasi
10000 per bulan, tabungan sebesar 1 juta tanpa diambil selama sekitar 10 tahunan akan habis
diambil oleh bank!
\end{eulercomment}
\begin{eulerprompt}
>r=0.005; plot2d(iterate("(1+0.8*r)*x-10000",1000000,n=130)):
\end{eulerprompt}
\eulerimg{27}{images/Shalih Abdillah_22305144009_ApKom5-228.png}
\begin{eulercomment}
Silakan Anda coba-coba, dengan tabungan minimal berapa agar tidak akan habis diambil oleh
bank dengan ketentuan bunga dan biaya administrasi seperti di atas.

Berikut adalah perhitungan minimal tabungan agar aman di bank dengan bunga sebesar r dan
biaya administrasi a, pajak bunga 20\%.
\end{eulercomment}
\begin{eulerprompt}
>$solve(0.8*r*A-a,A), $% with [r=0.005, a=10] 
\end{eulerprompt}
\begin{eulerformula}
\[
\left[ A=2500.0 \right] 
\]
\end{eulerformula}
\eulerimg{0}{images/Shalih Abdillah_22305144009_ApKom5-230-large.png}
\begin{eulercomment}
Berikut didefinisikan fungsi untuk menghitung saldo tabungan, kemudian dilakukan iterasi.
\end{eulercomment}
\begin{eulerprompt}
>function saldo(x,r,a) := round((1+0.8*r)*x-a,2);
>iterate(\{\{"saldo",0.005,10\}\},1000,n=6)
\end{eulerprompt}
\begin{euleroutput}
  [1000,  994,  987.98,  981.93,  975.86,  969.76,  963.64]
\end{euleroutput}
\begin{eulerprompt}
>iterate(\{\{"saldo",0.005,10\}\},2000,n=6)
\end{eulerprompt}
\begin{euleroutput}
  [2000,  1998,  1995.99,  1993.97,  1991.95,  1989.92,  1987.88]
\end{euleroutput}
\begin{eulerprompt}
>iterate(\{\{"saldo",0.005,10\}\},2500,n=6)
\end{eulerprompt}
\begin{euleroutput}
  [2500,  2500,  2500,  2500,  2500,  2500,  2500]
\end{euleroutput}
\begin{eulercomment}
Tabungan senilai 2,5 juta akan aman dan tidak akan berubah nilai (jika tidak ada penarikan),
sedangkan jika tabungan awal kurang dari 2,5 juta, lama kelamaan akan berkurang meskipun
tidak pernah dilakukan penarikan uang tabungan.
\end{eulercomment}
\begin{eulerprompt}
>iterate(\{\{"saldo",0.005,10\}\},3000,n=6)
\end{eulerprompt}
\begin{euleroutput}
  [3000,  3002,  3004.01,  3006.03,  3008.05,  3010.08,  3012.12]
\end{euleroutput}
\begin{eulercomment}
Tabungan yang lebih dari 2,5 juta baru akan bertambah jika tidak ada
penarikan.

Untuk barisan yang lebih kompleks dapat digunakan fungsi "sequence()".
Fungsi ini menghitung nilai-nilai x[n] dari semua nilai sebelumnya,
x[1],...,x[n-1] yang diketahui.\\
Berikut adalah contoh barisan Fibonacci.

\end{eulercomment}
\begin{eulerformula}
\[
x_n = x_{n-1}+x_{n-2}, \quad x_1=1, \quad x_2 =1
\]
\end{eulerformula}
\begin{eulerprompt}
>sequence("x[n-1]+x[n-2]",[1,1],15)
\end{eulerprompt}
\begin{euleroutput}
  [1,  1,  2,  3,  5,  8,  13,  21,  34,  55,  89,  144,  233,  377,  610]
\end{euleroutput}
\begin{eulercomment}
Barisan Fibonacci memiliki banyak sifat menarik, salah satunya adalah akar pangkat ke-n suku
ke-n akan konvergen ke pecahan emas:
\end{eulercomment}
\begin{eulerprompt}
>$'(1+sqrt(5))/2=float((1+sqrt(5))/2)
\end{eulerprompt}
\begin{eulerformula}
\[
\frac{\sqrt{5}+1}{2}=1.618033988749895
\]
\end{eulerformula}
\begin{eulerprompt}
>plot2d(sequence("x[n-1]+x[n-2]",[1,1],250)^(1/(1:250))):
\end{eulerprompt}
\eulerimg{27}{images/Shalih Abdillah_22305144009_ApKom5-233.png}
\begin{eulercomment}
Barisan yang sama juga dapat dihasilkan dengan menggunakan loop.
\end{eulercomment}
\begin{eulerprompt}
>x=ones(500); for k=3 to 500; x[k]=x[k-1]+x[k-2]; end;
\end{eulerprompt}
\begin{eulercomment}
Rekursi dapat dilakukan dengan menggunakan rumus yang tergantung pada semua elemen
sebelumnya. Pada contoh berikut, elemen ke-n merupakan jumlah (n-1) elemen sebelumnya,
dimulai dengan 1 (elemen ke-1). Jelas, nilai elemen ke-n adalah 2\textasciicircum{}(n-2), untuk n=2, 4, 5,
....
\end{eulercomment}
\begin{eulerprompt}
>sequence("sum(x)",1,10)
\end{eulerprompt}
\begin{euleroutput}
  [1,  1,  2,  4,  8,  16,  32,  64,  128,  256]
\end{euleroutput}
\begin{eulercomment}
Selain menggunakan ekspresi dalam x dan n, kita juga dapat menggunakan
fungsi.

Pada contoh berikut, digunakan iterasi

\end{eulercomment}
\begin{eulerformula}
\[
x_n =A \cdot x_{n-1},
\]
\end{eulerformula}
\begin{eulercomment}
dengan A suatu matriks 2x2, dan setiap x[n] merupakan matriks/vektor
2x1.
\end{eulercomment}
\begin{eulerprompt}
>A=[1,1;1,2]; function suku(x,n) := A.x[,n-1]
>sequence("suku",[1;1],6)
\end{eulerprompt}
\begin{euleroutput}
  Real 2 x 6 matrix
  
              1             2             5            13     ...
              1             3             8            21     ...
\end{euleroutput}
\begin{eulercomment}
Hasil yang sama juga dapat diperoleh dengan menggunakan fungsi
perpangkatan matriks "matrixpower()". Cara ini lebih cepat, karena
hanya menggunakan perkalian matriks sebanyak log\_2(n).

\end{eulercomment}
\begin{eulerformula}
\[
x_n=A.x_{n-1}=A^2.x_{n-2}=A^3.x_{n-3}= ... = A^{n-1}.x_1.
\]
\end{eulerformula}
\begin{eulerprompt}
>sequence("matrixpower(A,n).[1;1]",1,6)
\end{eulerprompt}
\begin{euleroutput}
  Real 2 x 6 matrix
  
              1             5            13            34     ...
              1             8            21            55     ...
\end{euleroutput}
\eulerheading{Spiral Theodorus}
\begin{eulercomment}
\end{eulercomment}
\eulerimg{15}{images/Shalih Abdillah_22305144009_ApKom5-236.png}
\begin{eulercomment}
Spiral Theodorus (spiral segitiga siku-siku) dapat digambar secara
rekursif. Rumus rekursifnya adalah:

\end{eulercomment}
\begin{eulerformula}
\[
x_n = \left( 1 + \frac{i}{\sqrt{n-1}} \right) \, x_{n-1}, \quad x_1=1,
\]
\end{eulerformula}
\begin{eulercomment}
yang menghasilkan barisan bilangan kompleks.
\end{eulercomment}
\begin{eulerprompt}
>function g(n) := 1+I/sqrt(n)
\end{eulerprompt}
\begin{eulercomment}
Rekursinya dapat dijalankan sebanyak 17 untuk menghasilkan barisan 17 bilangan kompleks,
kemudian digambar bilangan-bilangan kompleksnya.
\end{eulercomment}
\begin{eulerprompt}
>x=sequence("g(n-1)*x[n-1]",1,17); plot2d(x,r=3.5); textbox(latex("Spiral\(\backslash\) Theodorus"),0.4):
\end{eulerprompt}
\eulerimg{27}{images/Shalih Abdillah_22305144009_ApKom5-238.png}
\begin{eulercomment}
Selanjutnya dihubungan titik 0 dengan titik-titik kompleks tersebut menggunakan loop.
\end{eulercomment}
\begin{eulerprompt}
>for i=1:cols(x); plot2d([0,x[i]],>add); end:
\end{eulerprompt}
\eulerimg{27}{images/Shalih Abdillah_22305144009_ApKom5-239.png}
\begin{eulerprompt}
> 
\end{eulerprompt}
\begin{eulercomment}
Spiral tersebut juga dapat didefinisikan menggunakan fungsi rekursif, yang tidak memmerlukan
indeks dan bilangan kompleks. Dalam hal ini diigunakan vektor kolom pada bidang.
\end{eulercomment}
\begin{eulerprompt}
>function gstep (v) ...
\end{eulerprompt}
\begin{eulerudf}
  w=[-v[2];v[1]];
  return v+w/norm(w);
  endfunction
\end{eulerudf}
\begin{eulercomment}
Jika dilakukan iterasi 16 kali dimulai dari [1;0] akan didapatkan matriks yang memuat
vektor-vektor dari setiap iterasi.
\end{eulercomment}
\begin{eulerprompt}
>x=iterate("gstep",[1;0],16); plot2d(x[1],x[2],r=3.5,>points):
\end{eulerprompt}
\eulerimg{27}{images/Shalih Abdillah_22305144009_ApKom5-240.png}
\begin{eulercomment}
\begin{eulercomment}
\eulerheading{Kekonvergenan}
\begin{eulercomment}
Terkadang kita ingin melakukan iterasi sampai konvergen. Apabila iterasinya tidak konvergen
setelah ditunggu lama, Anda dapat menghentikannya dengan menekan tombol [ESC].
\end{eulercomment}
\begin{eulerprompt}
>iterate("cos(x)",1) // iterasi x(n+1)=cos(x(n)), dengan x(0)=1.
\end{eulerprompt}
\begin{euleroutput}
  0.739085133216
\end{euleroutput}
\begin{eulercomment}
Iterasi tersebut konvergen ke penyelesaian persamaan

\end{eulercomment}
\begin{eulerformula}
\[
x = \cos(x).
\]
\end{eulerformula}
\begin{eulercomment}
Iterasi ini juga dapat dilakukan pada interval, hasilnya adalah
barisan interval yang memuat akar tersebut.
\end{eulercomment}
\begin{eulerprompt}
>hasil := iterate("cos(x)",~1,2~) //iterasi x(n+1)=cos(x(n)), dengan interval awal (1, 2)
\end{eulerprompt}
\begin{euleroutput}
  ~0.739085133211,0.7390851332133~
\end{euleroutput}
\begin{eulercomment}
Jika interval hasil tersebut sedikit diperlebar, akan terlihat bahwa interval tersebut
memuat akar persamaan x=cos(x).
\end{eulercomment}
\begin{eulerprompt}
>h=expand(hasil,100), cos(h) << h
\end{eulerprompt}
\begin{euleroutput}
  ~0.73908513309,0.73908513333~
  1
\end{euleroutput}
\begin{eulercomment}
Iterasi juga dapat digunakan pada fungsi yang didefinisikan.
\end{eulercomment}
\begin{eulerprompt}
>function f(x) := (x+2/x)/2
\end{eulerprompt}
\begin{eulercomment}
Iterasi x(n+1)=f(x(n)) akan konvergen ke akar kuadrat 2.
\end{eulercomment}
\begin{eulerprompt}
>iterate("f",2), sqrt(2)
\end{eulerprompt}
\begin{euleroutput}
  1.41421356237
  1.41421356237
\end{euleroutput}
\begin{eulercomment}
Jika pada perintah iterate diberikan tambahan parameter n, maka hasil iterasinya akan
ditampilkan mulai dari iterasi pertama sampai ke-n.
\end{eulercomment}
\begin{eulerprompt}
>iterate("f",2,5)
\end{eulerprompt}
\begin{euleroutput}
  [2,  1.5,  1.41667,  1.41422,  1.41421,  1.41421]
\end{euleroutput}
\begin{eulercomment}
Untuk iterasi ini tidak dapat dilakukan terhadap interval.
\end{eulercomment}
\begin{eulerprompt}
>niterate("f",~1,2~,5)
\end{eulerprompt}
\begin{euleroutput}
  [ ~1,2~,  ~1,2~,  ~1,2~,  ~1,2~,  ~1,2~,  ~1,2~ ]
\end{euleroutput}
\begin{eulercomment}
Perhatikan, hasil iterasinya sama dengan interval awal. Alasannya adalah perhitungan dengan
interval bersifat terlalu longgar. Untuk meingkatkan perhitungan pada ekspresi dapat
digunakan pembagian intervalnya, menggunakan fungsi ieval().
\end{eulercomment}
\begin{eulerprompt}
>function s(x) := ieval("(x+2/x)/2",x,10)
\end{eulerprompt}
\begin{eulercomment}
Selanjutnya dapat dilakukan iterasi hingga diperoleh hasil optimal,
dan intervalnya tidak semakin mengecil. Hasilnya berupa interval yang
memuat akar persamaan:

\end{eulercomment}
\begin{eulerformula}
\[
x = \frac{1}{2} \left( x + \frac{2}{x} \right).
\]
\end{eulerformula}
\begin{eulercomment}
Satu-satunya solusi adalah\\
\end{eulercomment}
\begin{eulerformula}
\[
x = \sqrt2.
\]
\end{eulerformula}
\begin{eulerprompt}
>iterate("s",~1,2~)
\end{eulerprompt}
\begin{euleroutput}
  ~1.41421356236,1.41421356239~
\end{euleroutput}
\begin{eulercomment}
Fungsi "iterate()" juga dapat bekerja pada vektor. Berikut adalah
contoh fungsi vektor, yang menghasilkan rata-rata aritmetika dan
rata-rata geometri.

\end{eulercomment}
\begin{eulerformula}
\[
(a_{n+1},b_{n+1}) = \left( \frac{a_n+b_n}{2}, \sqrt{a_nb_n} \right)
\]
\end{eulerformula}
\begin{eulercomment}
Iterasi ke-n disimpan pada vektor kolom x[n].
\end{eulercomment}
\begin{eulerprompt}
>function g(x) := [(x[1]+x[2])/2;sqrt(x[1]*x[2])]
\end{eulerprompt}
\begin{eulercomment}
Iterasi dengan menggunakan fungsi tersebut akan konvergen ke rata-rata aritmetika dan
geometri dari nilai-nilai awal. 
\end{eulercomment}
\begin{eulerprompt}
>iterate("g",[1;5])
\end{eulerprompt}
\begin{euleroutput}
        2.60401 
        2.60401 
\end{euleroutput}
\begin{eulercomment}
Hasil tersebut konvergen agak cepat, seperti kita cek sebagai berikut.
\end{eulercomment}
\begin{eulerprompt}
>iterate("g",[1;5],4)
\end{eulerprompt}
\begin{euleroutput}
              1             3       2.61803       2.60403       2.60401 
              5       2.23607       2.59002       2.60399       2.60401 
\end{euleroutput}
\begin{eulercomment}
Iterasi pada interval dapat dilakukan dan stabil, namun tidak menunjukkan bahwa limitnya
pada batas-batas yang dihitung.
\end{eulercomment}
\begin{eulerprompt}
>iterate("g",[~1~;~5~],4)
\end{eulerprompt}
\begin{euleroutput}
  Interval 2 x 5 matrix
  
  ~0.999999999999999778,1.00000000000000022~     ...
  ~4.99999999999999911,5.00000000000000089~     ...
\end{euleroutput}
\begin{eulercomment}
Iterasi berikut konvergen sangat lambat.

\end{eulercomment}
\begin{eulerformula}
\[
x_{n+1} = \sqrt{x_n}.
\]
\end{eulerformula}
\begin{eulerprompt}
>iterate("sqrt(x)",2,10)
\end{eulerprompt}
\begin{euleroutput}
  [2,  1.41421,  1.18921,  1.09051,  1.04427,  1.0219,  1.01089,
  1.00543,  1.00271,  1.00135,  1.00068]
\end{euleroutput}
\begin{eulercomment}
Kekonvergenan iterasi tersebut dapat dipercepatdengan percepatan Steffenson:
\end{eulercomment}
\begin{eulerprompt}
>steffenson("sqrt(x)",2,10)
\end{eulerprompt}
\begin{euleroutput}
  [1.04888,  1.00028,  1,  1]
\end{euleroutput}
\eulerheading{Iterasi menggunakan Loop yang ditulis Langsung}
\begin{eulercomment}
Berikut adalah beberapa contoh penggunaan loop untuk melakukan iterasi yang ditulis langsung
pada baris perintah.
\end{eulercomment}
\begin{eulerprompt}
>x=2; repeat x=(x+2/x)/2; until x^2~=2; end; x,
\end{eulerprompt}
\begin{euleroutput}
  1.41421356237
\end{euleroutput}
\begin{eulercomment}
Penggabungan matriks menggunakan tanda "\textbar{}" dapat digunakan untuk menyimpan semua hasil
iterasi.
\end{eulercomment}
\begin{eulerprompt}
>v=[1]; for i=2 to 8; v=v|(v[i-1]*i); end; v,
\end{eulerprompt}
\begin{euleroutput}
  [1,  2,  6,  24,  120,  720,  5040,  40320]
\end{euleroutput}
\begin{eulercomment}
hasil iterasi juga dapat disimpan pada vektor yang sudah ada.
\end{eulercomment}
\begin{eulerprompt}
>v=ones(1,100); for i=2 to cols(v); v[i]=v[i-1]*i; end; ...
>plot2d(v,logplot=1); textbox(latex(&log(n)),x=0.5):
\end{eulerprompt}
\eulerimg{27}{images/Shalih Abdillah_22305144009_ApKom5-246.png}
\begin{eulerprompt}
>A =[0.5,0.2;0.7,0.1]; b=[2;2]; ...
>x=[1;1]; repeat xnew=A.x-b; until all(xnew~=x); x=xnew; end; ...
>x,
\end{eulerprompt}
\begin{euleroutput}
       -7.09677 
       -7.74194 
\end{euleroutput}
\eulerheading{Iterasi di dalam Fungsi}
\begin{eulercomment}
Fungsi atau program juga dapat menggunakan iterasi dan dapat digunakan untuk melakukan iterasi. Berikut adalah beberapa contoh
iterasi di dalam fungsi.

Contoh berikut adalah suatu fungsi untuk menghitung berapa lama suatu iterasi konvergen. Nilai fungsi tersebut adalah hasil akhir
iterasi dan banyak iterasi sampai konvergen.
\end{eulercomment}
\begin{eulerprompt}
>function map hiter(f$,x0) ...
\end{eulerprompt}
\begin{eulerudf}
  x=x0;
  maxiter=0;
  repeat
    xnew=f$(x);
    maxiter=maxiter+1;
    until xnew~=x;
    x=xnew;
  end;
  return maxiter;
  endfunction
\end{eulerudf}
\begin{eulercomment}
Misalnya, berikut adalah iterasi untuk mendapatkan hampiran akar kuadrat 2, cukup cepat,
konvergen pada iterasi ke-5, jika dimulai dari hampiran awal 2.
\end{eulercomment}
\begin{eulerprompt}
>hiter("(x+2/x)/2",2)
\end{eulerprompt}
\begin{euleroutput}
  5
\end{euleroutput}
\begin{eulercomment}
Karena fungsinya didefinisikan menggunakan "map". maka nilai awalnya dapat berupa vektor.
\end{eulercomment}
\begin{eulerprompt}
>x=1.5:0.1:10; hasil=hiter("(x+2/x)/2",x); ...
>  plot2d(x,hasil):
\end{eulerprompt}
\eulerimg{27}{images/Shalih Abdillah_22305144009_ApKom5-247.png}
\begin{eulercomment}
Dari gambar di atas terlihat bahwa kekonvergenan iterasinya semakin lambat, untuk nilai awal
semakin besar, namun penambahnnya tidak kontinu. Kita dapat menemukan kapan maksimum
iterasinya bertambah.
\end{eulercomment}
\begin{eulerprompt}
>hasil[1:10]
\end{eulerprompt}
\begin{euleroutput}
  [4,  5,  5,  5,  5,  5,  6,  6,  6,  6]
\end{euleroutput}
\begin{eulerprompt}
>x[nonzeros(differences(hasil))]
\end{eulerprompt}
\begin{euleroutput}
  [1.5,  2,  3.4,  6.6]
\end{euleroutput}
\begin{eulercomment}
maksimum iterasi sampai konvergen meningkat pada saat nilai awalnya 1.5, 2, 3.4, dan 6.6.

Contoh berikutnya adalah metode Newton pada polinomial kompleks berderajat 3.
\end{eulercomment}
\begin{eulerprompt}
>p &= x^3-1; newton &= x-p/diff(p,x); $newton
\end{eulerprompt}
\begin{euleroutput}
  Maxima said:
  diff: second argument must be a variable; found errexp1
   -- an error. To debug this try: debugmode(true);
  
  Error in:
  p &= x^3-1; newton &= x-p/diff(p,x); $newton ...
                                     ^
\end{euleroutput}
\begin{eulercomment}
Selanjutnya didefinisikan fungsi untuk melakukan iterasi (aslinya 10 kali).
\end{eulercomment}
\begin{eulerprompt}
>function iterasi(f$,x,n=10) ...
\end{eulerprompt}
\begin{eulerudf}
  loop 1 to n; x=f$(x); end;
  return x;
  endfunction
\end{eulerudf}
\begin{eulercomment}
Kita mulai dengan menentukan titik-titik grid pada bidang kompleksnya.
\end{eulercomment}
\begin{eulerprompt}
>r=1.5; x=linspace(-r,r,501); Z=x+I*x'; W=iterasi(newton,Z);
\end{eulerprompt}
\begin{euleroutput}
  Function newton needs at least 3 arguments!
  Use: newton (f$: call, df$: call, x: scalar complex \{, y: number, eps: none\}) 
  Error in:
  ...  x=linspace(-r,r,501); Z=x+I*x'; W=iterasi(newton,Z); ...
                                                       ^
\end{euleroutput}
\begin{eulercomment}
Berikut adalah akar-akar polinomial di atas.
\end{eulercomment}
\begin{eulerprompt}
>z=&solve(p)()
\end{eulerprompt}
\begin{euleroutput}
  Maxima said:
  solve: more equations than unknowns.
  Unknowns given :  
  [r]
  Equations given:  
  errexp1
   -- an error. To debug this try: debugmode(true);
  
  Error in:
  z=&solve(p)() ...
             ^
\end{euleroutput}
\begin{eulercomment}
Untuk menggambar hasil iterasinya, dihitung jarak dari hasil iterasi ke-10 ke masing-masing
akar, kemudian digunakan untuk menghitung warna yang akan digambar, yang menunjukkan limit
untuk masing-masing nilai awal. 

Fungsi plotrgb() menggunakan jendela gambar terkini untuk menggambar warna RGB sebagai
matriks.
\end{eulercomment}
\begin{eulerprompt}
>C=rgb(max(abs(W-z[1]),1),max(abs(W-z[2]),1),max(abs(W-z[3]),1)); ...
>  plot2d(none,-r,r,-r,r); plotrgb(C):
\end{eulerprompt}
\begin{euleroutput}
  Variable W not found!
  Error in:
  C=rgb(max(abs(W-z[1]),1),max(abs(W-z[2]),1),max(abs(W-z[3]),1) ...
                      ^
\end{euleroutput}
\eulerheading{Iterasi Simbolik}
\begin{eulercomment}
Seperti sudah dibahas sebelumnya, untuk menghasilkan barisan ekspresi simbolik dengan Maxima
dapat digunakan fungsi makelist().
\end{eulercomment}
\begin{eulerprompt}
>&powerdisp:true // untuk menampilkan deret pangkat mulai dari suku berpangkat terkecil
\end{eulerprompt}
\begin{euleroutput}
  
                                   true
  
\end{euleroutput}
\begin{eulerprompt}
>deret &= makelist(taylor(exp(x),x,0,k),k,1,3); $deret // barisan deret Taylor untuk e^x
\end{eulerprompt}
\begin{euleroutput}
  Maxima said:
  taylor: 0.1539740213994798*r cannot be a variable.
   -- an error. To debug this try: debugmode(true);
  
  Error in:
  deret &= makelist(taylor(exp(x),x,0,k),k,1,3); $deret // baris ...
                                               ^
\end{euleroutput}
\begin{eulercomment}
Untuk mengubah barisan deret tersebut menjadi vektor string di EMT digunakan fungsi
mxm2str(). Selanjutnya, vektor string/ekspresi hasilnya dapat digambar seperti menggambar
vektor eskpresi pada EMT.
\end{eulercomment}
\begin{eulerprompt}
>plot2d("exp(x)",0,3); // plot fungsi aslinya, e^x
>plot2d(mxm2str("deret"),>add,color=4:6): // plot ketiga deret taylor hampiran fungsi tersebut
\end{eulerprompt}
\begin{euleroutput}
  Maxima said:
  length: argument cannot be a symbol; found deret
   -- an error. To debug this try: debugmode(true);
  
  mxmeval:
      return evaluate(mxm(s));
  Try "trace errors" to inspect local variables after errors.
  mxm2str:
      n=mxmeval("length(VVV)");
\end{euleroutput}
\begin{eulercomment}
Selain cara di atas dapat juga dengan cara menggunakan indeks pada vektor/list yang
dihasilkan.
\end{eulercomment}
\begin{eulerprompt}
>$deret[3]
\end{eulerprompt}
\begin{eulerformula}
\[
{\it deret}_{3}
\]
\end{eulerformula}
\begin{eulerprompt}
>plot2d(["exp(x)",&deret[1],&deret[2],&deret[3]],0,3,color=1:4):
\end{eulerprompt}
\begin{euleroutput}
  deret is not a variable!
  Error in expression: deret[1]
   %ploteval:
      y0=f$(x[1],args());
  Try "trace errors" to inspect local variables after errors.
  plot2d:
      u=u_(%ploteval(xx[#],t;args()));
\end{euleroutput}
\begin{eulerprompt}
>$sum(sin(k*x)/k,k,1,5)
\end{eulerprompt}
\begin{eulerformula}
\[
\left[ 0 , \sin \left(1.66665833335744 \times 10^{-7}\,r\right)+  \frac{\sin \left(3.333316666714881 \times 10^{-7}\,r\right)}{2}+  \frac{\sin \left(4.999975000072321 \times 10^{-7}\,r\right)}{3}+  \frac{\sin \left(6.666633333429761 \times 10^{-7}\,r\right)}{4}+  \frac{\sin \left(8.333291666787201 \times 10^{-7}\,r\right)}{5} ,   \sin \left(1.33330666692022 \times 10^{-6}\,r\right)+\frac{\sin   \left(2.66661333384044 \times 10^{-6}\,r\right)}{2}+\frac{\sin   \left(3.999920000760659 \times 10^{-6}\,r\right)}{3}+\frac{\sin   \left(5.333226667680879 \times 10^{-6}\,r\right)}{4}+\frac{\sin   \left(6.666533334601099 \times 10^{-6}\,r\right)}{5} , \sin \left(  4.499797504338432 \times 10^{-6}\,r\right)+\frac{\sin \left(  8.999595008676864 \times 10^{-6}\,r\right)}{2}+\frac{\sin \left(  1.34993925130153 \times 10^{-5}\,r\right)}{3}+\frac{\sin \left(  1.799919001735373 \times 10^{-5}\,r\right)}{4}+\frac{\sin \left(  2.249898752169216 \times 10^{-5}\,r\right)}{5} , \sin \left(  1.066581336583994 \times 10^{-5}\,r\right)+\frac{\sin \left(  2.133162673167988 \times 10^{-5}\,r\right)}{2}+\frac{\sin \left(  3.199744009751981 \times 10^{-5}\,r\right)}{3}+\frac{\sin \left(  4.266325346335975 \times 10^{-5}\,r\right)}{4}+\frac{\sin \left(  5.332906682919969 \times 10^{-5}\,r\right)}{5} , \sin \left(  2.083072932167196 \times 10^{-5}\,r\right)+\frac{\sin \left(  4.166145864334392 \times 10^{-5}\,r\right)}{2}+\frac{\sin \left(  6.249218796501588 \times 10^{-5}\,r\right)}{3}+\frac{\sin \left(  8.332291728668784 \times 10^{-5}\,r\right)}{4}+\frac{\sin \left(  1.041536466083598 \times 10^{-4}\,r\right)}{5} , \sin \left(  3.599352055540239 \times 10^{-5}\,r\right)+\frac{\sin \left(  7.198704111080478 \times 10^{-5}\,r\right)}{2}+\frac{\sin \left(  1.079805616662072 \times 10^{-4}\,r\right)}{3}+\frac{\sin \left(  1.439740822216096 \times 10^{-4}\,r\right)}{4}+\frac{\sin \left(  1.79967602777012 \times 10^{-4}\,r\right)}{5} , \sin \left(  5.71526624672386 \times 10^{-5}\,r\right)+\frac{\sin \left(  1.143053249344772 \times 10^{-4}\,r\right)}{2}+\frac{\sin \left(  1.714579874017158 \times 10^{-4}\,r\right)}{3}+\frac{\sin \left(  2.286106498689544 \times 10^{-4}\,r\right)}{4}+\frac{\sin \left(  2.85763312336193 \times 10^{-4}\,r\right)}{5} , \sin \left(  8.530603082730626 \times 10^{-5}\,r\right)+\frac{\sin \left(  1.706120616546125 \times 10^{-4}\,r\right)}{2}+\frac{\sin \left(  2.559180924819188 \times 10^{-4}\,r\right)}{3}+\frac{\sin \left(  3.41224123309225 \times 10^{-4}\,r\right)}{4}+\frac{\sin \left(  4.265301541365313 \times 10^{-4}\,r\right)}{5} , \sin \left(  1.214508019889565 \times 10^{-4}\,r\right)+\frac{\sin \left(  2.42901603977913 \times 10^{-4}\,r\right)}{2}+\frac{\sin \left(  3.643524059668696 \times 10^{-4}\,r\right)}{3}+\frac{\sin \left(  4.858032079558261 \times 10^{-4}\,r\right)}{4}+\frac{\sin \left(  6.072540099447826 \times 10^{-4}\,r\right)}{5} , \sin \left(  1.665833531718508 \times 10^{-4}\,r\right)+\frac{\sin \left(  3.331667063437016 \times 10^{-4}\,r\right)}{2}+\frac{\sin \left(  4.997500595155524 \times 10^{-4}\,r\right)}{3}+\frac{\sin \left(  6.663334126874032 \times 10^{-4}\,r\right)}{4}+\frac{\sin \left(  8.32916765859254 \times 10^{-4}\,r\right)}{5} , \sin \left(  2.216991628251896 \times 10^{-4}\,r\right)+\frac{\sin \left(  4.433983256503793 \times 10^{-4}\,r\right)}{2}+\frac{\sin \left(  6.650974884755689 \times 10^{-4}\,r\right)}{3}+\frac{\sin \left(  8.867966513007586 \times 10^{-4}\,r\right)}{4}+\frac{\sin \left(  0.001108495814125948\,r\right)}{5} , \sin \left(  2.877927110806339 \times 10^{-4}\,r\right)+\frac{\sin \left(  5.755854221612677 \times 10^{-4}\,r\right)}{2}+\frac{\sin \left(  8.633781332419016 \times 10^{-4}\,r\right)}{3}+\frac{\sin \left(  0.001151170844322535\,r\right)}{4}+\frac{\sin \left(  0.001438963555403169\,r\right)}{5} , \sin \left(  3.658573803051457 \times 10^{-4}\,r\right)+\frac{\sin \left(  7.317147606102914 \times 10^{-4}\,r\right)}{2}+\frac{\sin \left(  0.001097572140915437\,r\right)}{3}+\frac{\sin \left(  0.001463429521220583\,r\right)}{4}+\frac{\sin \left(  0.001829286901525728\,r\right)}{5} , \sin \left(  4.568853557635201 \times 10^{-4}\,r\right)+\frac{\sin \left(  9.137707115270399 \times 10^{-4}\,r\right)}{2}+\frac{\sin \left(  0.00137065606729056\,r\right)}{3}+\frac{\sin \left(  0.00182754142305408\,r\right)}{4}+\frac{\sin \left(  0.0022844267788176\,r\right)}{5} , \sin \left(  5.618675264007778 \times 10^{-4}\,r\right)+\frac{\sin \left(  0.001123735052801556\,r\right)}{2}+\frac{\sin \left(  0.001685602579202333\,r\right)}{3}+\frac{\sin \left(  0.002247470105603111\,r\right)}{4}+\frac{\sin \left(  0.002809337632003889\,r\right)}{5} , \sin \left(  6.817933857540259 \times 10^{-4}\,r\right)+\frac{\sin \left(  0.001363586771508052\,r\right)}{2}+\frac{\sin \left(  0.002045380157262078\,r\right)}{3}+\frac{\sin \left(  0.002727173543016104\,r\right)}{4}+\frac{\sin \left(  0.00340896692877013\,r\right)}{5} , \sin \left(  8.176509330039827 \times 10^{-4}\,r\right)+\frac{\sin \left(  0.001635301866007965\,r\right)}{2}+\frac{\sin \left(  0.002452952799011948\,r\right)}{3}+\frac{\sin \left(  0.003270603732015931\,r\right)}{4}+\frac{\sin \left(  0.004088254665019914\,r\right)}{5} , \sin \left(  9.704265741758145 \times 10^{-4}\,r\right)+\frac{\sin \left(  0.001940853148351629\,r\right)}{2}+\frac{\sin \left(  0.002911279722527443\,r\right)}{3}+\frac{\sin \left(  0.003881706296703258\,r\right)}{4}+\frac{\sin \left(  0.004852132870879072\,r\right)}{5} , \sin \left(0.001141105023499428  \,r\right)+\frac{\sin \left(0.002282210046998856\,r\right)}{2}+  \frac{\sin \left(0.003423315070498284\,r\right)}{3}+\frac{\sin   \left(0.004564420093997712\,r\right)}{4}+\frac{\sin \left(  0.00570552511749714\,r\right)}{5} , \sin \left(0.001330669204938795  \,r\right)+\frac{\sin \left(0.002661338409877589\,r\right)}{2}+  \frac{\sin \left(0.003992007614816384\,r\right)}{3}+\frac{\sin   \left(0.005322676819755179\,r\right)}{4}+\frac{\sin \left(  0.006653346024693974\,r\right)}{5} , \sin \left(0.001540100153900437  \,r\right)+\frac{\sin \left(0.003080200307800873\,r\right)}{2}+  \frac{\sin \left(0.00462030046170131\,r\right)}{3}+\frac{\sin \left(  0.006160400615601747\,r\right)}{4}+\frac{\sin \left(  0.007700500769502183\,r\right)}{5} , \sin \left(0.001770376919130678  \,r\right)+\frac{\sin \left(0.003540753838261357\,r\right)}{2}+  \frac{\sin \left(0.005311130757392035\,r\right)}{3}+\frac{\sin   \left(0.007081507676522714\,r\right)}{4}+\frac{\sin \left(  0.008851884595653392\,r\right)}{5} , \sin \left(0.002022476464811601  \,r\right)+\frac{\sin \left(0.004044952929623202\,r\right)}{2}+  \frac{\sin \left(0.006067429394434803\,r\right)}{3}+\frac{\sin   \left(0.008089905859246405\,r\right)}{4}+\frac{\sin \left(  0.01011238232405801\,r\right)}{5} , \sin \left(0.002297373572865413  \,r\right)+\frac{\sin \left(0.004594747145730826\,r\right)}{2}+  \frac{\sin \left(0.00689212071859624\,r\right)}{3}+\frac{\sin \left(  0.009189494291461653\,r\right)}{4}+\frac{\sin \left(  0.01148686786432707\,r\right)}{5} , \sin \left(0.002596040745477063  \,r\right)+\frac{\sin \left(0.005192081490954126\,r\right)}{2}+  \frac{\sin \left(0.007788122236431189\,r\right)}{3}+\frac{\sin   \left(0.01038416298190825\,r\right)}{4}+\frac{\sin \left(  0.01298020372738531\,r\right)}{5} , \sin \left(0.002919448107844891  \,r\right)+\frac{\sin \left(0.005838896215689782\,r\right)}{2}+  \frac{\sin \left(0.008758344323534673\,r\right)}{3}+\frac{\sin   \left(0.01167779243137956\,r\right)}{4}+\frac{\sin \left(  0.01459724053922445\,r\right)}{5} , \sin \left(0.003268563311168871  \,r\right)+\frac{\sin \left(0.006537126622337741\,r\right)}{2}+  \frac{\sin \left(0.009805689933506612\,r\right)}{3}+\frac{\sin   \left(0.01307425324467548\,r\right)}{4}+\frac{\sin \left(  0.01634281655584435\,r\right)}{5} , \sin \left(0.003644351435886262  \,r\right)+\frac{\sin \left(0.007288702871772523\,r\right)}{2}+  \frac{\sin \left(0.01093305430765878\,r\right)}{3}+\frac{\sin \left(  0.01457740574354505\,r\right)}{4}+\frac{\sin \left(  0.01822175717943131\,r\right)}{5} , \sin \left(0.004047774895164447  \,r\right)+\frac{\sin \left(0.008095549790328893\,r\right)}{2}+  \frac{\sin \left(0.01214332468549334\,r\right)}{3}+\frac{\sin \left(  0.01619109958065779\,r\right)}{4}+\frac{\sin \left(  0.02023887447582223\,r\right)}{5} , \sin \left(0.004479793338660443  \,r\right)+\frac{\sin \left(0.008959586677320885\,r\right)}{2}+  \frac{\sin \left(0.01343938001598133\,r\right)}{3}+\frac{\sin \left(  0.01791917335464177\,r\right)}{4}+\frac{\sin \left(  0.02239896669330221\,r\right)}{5} , \sin \left(0.0049413635565565\,r  \right)+\frac{\sin \left(0.009882727113112999\,r\right)}{2}+\frac{  \sin \left(0.0148240906696695\,r\right)}{3}+\frac{\sin \left(  0.019765454226226\,r\right)}{4}+\frac{\sin \left(0.0247068177827825  \,r\right)}{5} , \sin \left(0.005433439383882244\,r\right)+\frac{  \sin \left(0.01086687876776449\,r\right)}{2}+\frac{\sin \left(  0.01630031815164673\,r\right)}{3}+\frac{\sin \left(  0.02173375753552897\,r\right)}{4}+\frac{\sin \left(  0.02716719691941122\,r\right)}{5} , \sin \left(0.005956971605131645  \,r\right)+\frac{\sin \left(0.01191394321026329\,r\right)}{2}+\frac{  \sin \left(0.01787091481539493\,r\right)}{3}+\frac{\sin \left(  0.02382788642052658\,r\right)}{4}+\frac{\sin \left(  0.02978485802565822\,r\right)}{5} , \sin \left(0.006512907859185624  \,r\right)+\frac{\sin \left(0.01302581571837125\,r\right)}{2}+\frac{  \sin \left(0.01953872357755687\,r\right)}{3}+\frac{\sin \left(  0.0260516314367425\,r\right)}{4}+\frac{\sin \left(  0.03256453929592812\,r\right)}{5} , \sin \left(0.007102192544548636  \,r\right)+\frac{\sin \left(0.01420438508909727\,r\right)}{2}+\frac{  \sin \left(0.02130657763364591\,r\right)}{3}+\frac{\sin \left(  0.02840877017819454\,r\right)}{4}+\frac{\sin \left(  0.03551096272274318\,r\right)}{5} , \sin \left(0.007725766724910044  \,r\right)+\frac{\sin \left(0.01545153344982009\,r\right)}{2}+\frac{  \sin \left(0.02317730017473013\,r\right)}{3}+\frac{\sin \left(  0.03090306689964017\,r\right)}{4}+\frac{\sin \left(  0.03862883362455022\,r\right)}{5} , \sin \left(0.00838456803503801\,  r\right)+\frac{\sin \left(0.01676913607007602\,r\right)}{2}+\frac{  \sin \left(0.02515370410511403\,r\right)}{3}+\frac{\sin \left(  0.03353827214015204\,r\right)}{4}+\frac{\sin \left(  0.04192284017519005\,r\right)}{5} , \sin \left(0.009079530587017326  \,r\right)+\frac{\sin \left(0.01815906117403465\,r\right)}{2}+\frac{  \sin \left(0.02723859176105198\,r\right)}{3}+\frac{\sin \left(  0.0363181223480693\,r\right)}{4}+\frac{\sin \left(  0.04539765293508663\,r\right)}{5} , \sin \left(0.009811584876838586  \,r\right)+\frac{\sin \left(0.01962316975367717\,r\right)}{2}+\frac{  \sin \left(0.02943475463051576\,r\right)}{3}+\frac{\sin \left(  0.03924633950735434\,r\right)}{4}+\frac{\sin \left(  0.04905792438419293\,r\right)}{5} , \sin \left(0.0105816576913495\,r  \right)+\frac{\sin \left(0.021163315382699\,r\right)}{2}+\frac{\sin   \left(0.0317449730740485\,r\right)}{3}+\frac{\sin \left(  0.042326630765398\,r\right)}{4}+\frac{\sin \left(0.0529082884567475  \,r\right)}{5} , \sin \left(0.01139067201557714\,r\right)+\frac{  \sin \left(0.02278134403115428\,r\right)}{2}+\frac{\sin \left(  0.03417201604673142\,r\right)}{3}+\frac{\sin \left(  0.04556268806230857\,r\right)}{4}+\frac{\sin \left(  0.05695336007788571\,r\right)}{5} , \sin \left(0.01223954694042984\,  r\right)+\frac{\sin \left(0.02447909388085967\,r\right)}{2}+\frac{  \sin \left(0.03671864082128951\,r\right)}{3}+\frac{\sin \left(  0.04895818776171934\,r\right)}{4}+\frac{\sin \left(  0.06119773470214918\,r\right)}{5} , \sin \left(0.01312919757078923\,  r\right)+\frac{\sin \left(0.02625839514157846\,r\right)}{2}+\frac{  \sin \left(0.03938759271236769\,r\right)}{3}+\frac{\sin \left(  0.05251679028315692\,r\right)}{4}+\frac{\sin \left(  0.06564598785394615\,r\right)}{5} , \sin \left(0.01406053493400045\,  r\right)+\frac{\sin \left(0.02812106986800089\,r\right)}{2}+\frac{  \sin \left(0.04218160480200134\,r\right)}{3}+\frac{\sin \left(  0.05624213973600178\,r\right)}{4}+\frac{\sin \left(  0.07030267467000223\,r\right)}{5} , \sin \left(0.01503446588876983\,  r\right)+\frac{\sin \left(0.03006893177753966\,r\right)}{2}+\frac{  \sin \left(0.0451033976663095\,r\right)}{3}+\frac{\sin \left(  0.06013786355507933\,r\right)}{4}+\frac{\sin \left(  0.07517232944384916\,r\right)}{5} , \sin \left(0.01605189303448024\,  r\right)+\frac{\sin \left(0.03210378606896047\,r\right)}{2}+\frac{  \sin \left(0.04815567910344071\,r\right)}{3}+\frac{\sin \left(  0.06420757213792094\,r\right)}{4}+\frac{\sin \left(  0.08025946517240118\,r\right)}{5} , \sin \left(0.01711371462093175\,  r\right)+\frac{\sin \left(0.03422742924186351\,r\right)}{2}+\frac{  \sin \left(0.05134114386279526\,r\right)}{3}+\frac{\sin \left(  0.06845485848372701\,r\right)}{4}+\frac{\sin \left(  0.08556857310465876\,r\right)}{5} , \sin \left(0.01822082445851714\,  r\right)+\frac{\sin \left(0.03644164891703428\,r\right)}{2}+\frac{  \sin \left(0.05466247337555141\,r\right)}{3}+\frac{\sin \left(  0.07288329783406855\,r\right)}{4}+\frac{\sin \left(  0.09110412229258569\,r\right)}{5} , \sin \left(0.01937411182884202\,  r\right)+\frac{\sin \left(0.03874822365768404\,r\right)}{2}+\frac{  \sin \left(0.05812233548652607\,r\right)}{3}+\frac{\sin \left(  0.07749644731536809\,r\right)}{4}+\frac{\sin \left(  0.09687055914421011\,r\right)}{5} , \sin \left(0.02057446139579705\,  r\right)+\frac{\sin \left(0.0411489227915941\,r\right)}{2}+\frac{  \sin \left(0.06172338418739115\,r\right)}{3}+\frac{\sin \left(  0.0822978455831882\,r\right)}{4}+\frac{\sin \left(0.1028723069789853  \,r\right)}{5} , \sin \left(0.02182275311709253\,r\right)+\frac{  \sin \left(0.04364550623418506\,r\right)}{2}+\frac{\sin \left(  0.06546825935127759\,r\right)}{3}+\frac{\sin \left(  0.08729101246837012\,r\right)}{4}+\frac{\sin \left(  0.1091137655854627\,r\right)}{5} , \sin \left(0.02311986215626333\,r  \right)+\frac{\sin \left(0.04623972431252665\,r\right)}{2}+\frac{  \sin \left(0.06935958646878998\,r\right)}{3}+\frac{\sin \left(  0.0924794486250533\,r\right)}{4}+\frac{\sin \left(0.1155993107813166  \,r\right)}{5} , \sin \left(0.02446665879515308\,r\right)+\frac{  \sin \left(0.04893331759030617\,r\right)}{2}+\frac{\sin \left(  0.07339997638545925\,r\right)}{3}+\frac{\sin \left(  0.09786663518061234\,r\right)}{4}+\frac{\sin \left(  0.1223332939757654\,r\right)}{5} , \sin \left(0.02586400834688696\,r  \right)+\frac{\sin \left(0.05172801669377391\,r\right)}{2}+\frac{  \sin \left(0.07759202504066087\,r\right)}{3}+\frac{\sin \left(  0.1034560333875478\,r\right)}{4}+\frac{\sin \left(0.1293200417344348  \,r\right)}{5} , \sin \left(0.02731277106934082\,r\right)+\frac{  \sin \left(0.05462554213868165\,r\right)}{2}+\frac{\sin \left(  0.08193831320802247\,r\right)}{3}+\frac{\sin \left(  0.1092510842773633\,r\right)}{4}+\frac{\sin \left(0.1365638553467041  \,r\right)}{5} , \sin \left(0.02881380207911666\,r\right)+\frac{  \sin \left(0.05762760415823331\,r\right)}{2}+\frac{\sin \left(  0.08644140623734997\,r\right)}{3}+\frac{\sin \left(  0.1152552083164666\,r\right)}{4}+\frac{\sin \left(0.1440690103955833  \,r\right)}{5} , \sin \left(0.03036795126603076\,r\right)+\frac{  \sin \left(0.06073590253206151\,r\right)}{2}+\frac{\sin \left(  0.09110385379809227\,r\right)}{3}+\frac{\sin \left(0.121471805064123  \,r\right)}{4}+\frac{\sin \left(0.1518397563301538\,r\right)}{5} ,   \sin \left(0.03197606320812652\,r\right)+\frac{\sin \left(  0.06395212641625303\,r\right)}{2}+\frac{\sin \left(  0.09592818962437955\,r\right)}{3}+\frac{\sin \left(  0.1279042528325061\,r\right)}{4}+\frac{\sin \left(0.1598803160406326  \,r\right)}{5} , \sin \left(0.0336389770872163\,r\right)+\frac{\sin   \left(0.06727795417443261\,r\right)}{2}+\frac{\sin \left(  0.1009169312616489\,r\right)}{3}+\frac{\sin \left(0.1345559083488652  \,r\right)}{4}+\frac{\sin \left(0.1681948854360815\,r\right)}{5} ,   \sin \left(0.03535752660496472\,r\right)+\frac{\sin \left(  0.07071505320992943\,r\right)}{2}+\frac{\sin \left(  0.1060725798148942\,r\right)}{3}+\frac{\sin \left(0.1414301064198589  \,r\right)}{4}+\frac{\sin \left(0.1767876330248236\,r\right)}{5} ,   \sin \left(0.03713253989951881\,r\right)+\frac{\sin \left(  0.07426507979903763\,r\right)}{2}+\frac{\sin \left(  0.1113976196985564\,r\right)}{3}+\frac{\sin \left(0.1485301595980753  \,r\right)}{4}+\frac{\sin \left(0.1856626994975941\,r\right)}{5} ,   \sin \left(0.03896483946269502\,r\right)+\frac{\sin \left(  0.07792967892539004\,r\right)}{2}+\frac{\sin \left(  0.1168945183880851\,r\right)}{3}+\frac{\sin \left(0.1558593578507801  \,r\right)}{4}+\frac{\sin \left(0.1948241973134751\,r\right)}{5} ,   \sin \left(0.0408552420577305\,r\right)+\frac{\sin \left(  0.081710484115461\,r\right)}{2}+\frac{\sin \left(0.1225657261731915  \,r\right)}{3}+\frac{\sin \left(0.163420968230922\,r\right)}{4}+  \frac{\sin \left(0.2042762102886525\,r\right)}{5} , \sin \left(  0.04280455863760801\,r\right)+\frac{\sin \left(0.08560911727521603\,  r\right)}{2}+\frac{\sin \left(0.128413675912824\,r\right)}{3}+\frac{  \sin \left(0.1712182345504321\,r\right)}{4}+\frac{\sin \left(  0.2140227931880401\,r\right)}{5} , \sin \left(0.04481359426396048\,r  \right)+\frac{\sin \left(0.08962718852792095\,r\right)}{2}+\frac{  \sin \left(0.1344407827918814\,r\right)}{3}+\frac{\sin \left(  0.1792543770558419\,r\right)}{4}+\frac{\sin \left(0.2240679713198024  \,r\right)}{5} , \sin \left(0.04688314802656623\,r\right)+\frac{  \sin \left(0.09376629605313247\,r\right)}{2}+\frac{\sin \left(  0.1406494440796987\,r\right)}{3}+\frac{\sin \left(0.1875325921062649  \,r\right)}{4}+\frac{\sin \left(0.2344157401328312\,r\right)}{5} ,   \sin \left(0.04901401296344043\,r\right)+\frac{\sin \left(  0.09802802592688087\,r\right)}{2}+\frac{\sin \left(  0.1470420388903213\,r\right)}{3}+\frac{\sin \left(0.1960560518537617  \,r\right)}{4}+\frac{\sin \left(0.2450700648172022\,r\right)}{5} ,   \sin \left(0.05120697598153157\,r\right)+\frac{\sin \left(  0.1024139519630631\,r\right)}{2}+\frac{\sin \left(0.1536209279445947  \,r\right)}{3}+\frac{\sin \left(0.2048279039261263\,r\right)}{4}+  \frac{\sin \left(0.2560348799076578\,r\right)}{5} , \sin \left(  0.05346281777803219\,r\right)+\frac{\sin \left(0.1069256355560644\,r  \right)}{2}+\frac{\sin \left(0.1603884533340966\,r\right)}{3}+\frac{  \sin \left(0.2138512711121288\,r\right)}{4}+\frac{\sin \left(  0.267314088890161\,r\right)}{5} , \sin \left(0.05578231276230905\,r  \right)+\frac{\sin \left(0.1115646255246181\,r\right)}{2}+\frac{  \sin \left(0.1673469382869271\,r\right)}{3}+\frac{\sin \left(  0.2231292510492362\,r\right)}{4}+\frac{\sin \left(0.2789115638115452  \,r\right)}{5} , \sin \left(0.05816622897846346\,r\right)+\frac{  \sin \left(0.1163324579569269\,r\right)}{2}+\frac{\sin \left(  0.1744986869353904\,r\right)}{3}+\frac{\sin \left(0.2326649159138539  \,r\right)}{4}+\frac{\sin \left(0.2908311448923173\,r\right)}{5} ,   \sin \left(0.06061532802852698\,r\right)+\frac{\sin \left(  0.121230656057054\,r\right)}{2}+\frac{\sin \left(0.1818459840855809  \,r\right)}{3}+\frac{\sin \left(0.2424613121141079\,r\right)}{4}+  \frac{\sin \left(0.3030766401426349\,r\right)}{5} , \sin \left(  0.0631303649963022\,r\right)+\frac{\sin \left(0.1262607299926044\,r  \right)}{2}+\frac{\sin \left(0.1893910949889066\,r\right)}{3}+\frac{  \sin \left(0.2525214599852088\,r\right)}{4}+\frac{\sin \left(  0.315651824981511\,r\right)}{5} , \sin \left(0.06571208837185505\,r  \right)+\frac{\sin \left(0.1314241767437101\,r\right)}{2}+\frac{  \sin \left(0.1971362651155651\,r\right)}{3}+\frac{\sin \left(  0.2628483534874202\,r\right)}{4}+\frac{\sin \left(0.3285604418592752  \,r\right)}{5} , \sin \left(0.06836123997666599\,r\right)+\frac{  \sin \left(0.136722479953332\,r\right)}{2}+\frac{\sin \left(  0.205083719929998\,r\right)}{3}+\frac{\sin \left(0.273444959906664\,  r\right)}{4}+\frac{\sin \left(0.3418061998833299\,r\right)}{5} ,   \sin \left(0.07107855488944881\,r\right)+\frac{\sin \left(  0.1421571097788976\,r\right)}{2}+\frac{\sin \left(0.2132356646683464  \,r\right)}{3}+\frac{\sin \left(0.2843142195577952\,r\right)}{4}+  \frac{\sin \left(0.355392774447244\,r\right)}{5} , \sin \left(  0.07386476137264342\,r\right)+\frac{\sin \left(0.1477295227452868\,r  \right)}{2}+\frac{\sin \left(0.2215942841179303\,r\right)}{3}+\frac{  \sin \left(0.2954590454905737\,r\right)}{4}+\frac{\sin \left(  0.3693238068632171\,r\right)}{5} , \sin \left(0.07672058079958999\,r  \right)+\frac{\sin \left(0.15344116159918\,r\right)}{2}+\frac{\sin   \left(0.23016174239877\,r\right)}{3}+\frac{\sin \left(  0.30688232319836\,r\right)}{4}+\frac{\sin \left(0.3836029039979499\,  r\right)}{5} , \sin \left(0.07964672758239233\,r\right)+\frac{\sin   \left(0.1592934551647847\,r\right)}{2}+\frac{\sin \left(  0.238940182747177\,r\right)}{3}+\frac{\sin \left(0.3185869103295693  \,r\right)}{4}+\frac{\sin \left(0.3982336379119616\,r\right)}{5} ,   \sin \left(0.08264390910047736\,r\right)+\frac{\sin \left(  0.1652878182009547\,r\right)}{2}+\frac{\sin \left(0.2479317273014321  \,r\right)}{3}+\frac{\sin \left(0.3305756364019095\,r\right)}{4}+  \frac{\sin \left(0.4132195455023868\,r\right)}{5} , \sin \left(  0.0857128256298576\,r\right)+\frac{\sin \left(0.1714256512597152\,r  \right)}{2}+\frac{\sin \left(0.2571384768895728\,r\right)}{3}+\frac{  \sin \left(0.3428513025194304\,r\right)}{4}+\frac{\sin \left(  0.428564128149288\,r\right)}{5} , \sin \left(0.08885417027310427\,r  \right)+\frac{\sin \left(0.1777083405462085\,r\right)}{2}+\frac{  \sin \left(0.2665625108193128\,r\right)}{3}+\frac{\sin \left(  0.3554166810924171\,r\right)}{4}+\frac{\sin \left(0.4442708513655214  \,r\right)}{5} , \sin \left(0.09206862889003742\,r\right)+\frac{  \sin \left(0.1841372577800748\,r\right)}{2}+\frac{\sin \left(  0.2762058866701123\,r\right)}{3}+\frac{\sin \left(0.3682745155601497  \,r\right)}{4}+\frac{\sin \left(0.4603431444501871\,r\right)}{5} ,   \sin \left(0.09535688002914089\,r\right)+\frac{\sin \left(  0.1907137600582818\,r\right)}{2}+\frac{\sin \left(0.2860706400874227  \,r\right)}{3}+\frac{\sin \left(0.3814275201165636\,r\right)}{4}+  \frac{\sin \left(0.4767844001457044\,r\right)}{5} , \sin \left(  0.0987195948597075\,r\right)+\frac{\sin \left(0.197439189719415\,r  \right)}{2}+\frac{\sin \left(0.2961587845791225\,r\right)}{3}+\frac{  \sin \left(0.39487837943883\,r\right)}{4}+\frac{\sin \left(  0.4935979742985375\,r\right)}{5} , \sin \left(0.1021574371047232\,r  \right)+\frac{\sin \left(0.2043148742094465\,r\right)}{2}+\frac{  \sin \left(0.3064723113141697\,r\right)}{3}+\frac{\sin \left(  0.408629748418893\,r\right)}{4}+\frac{\sin \left(0.5107871855236162  \,r\right)}{5} , \sin \left(0.1056710629744951\,r\right)+\frac{\sin   \left(0.2113421259489903\,r\right)}{2}+\frac{\sin \left(  0.3170131889234854\,r\right)}{3}+\frac{\sin \left(0.4226842518979805  \,r\right)}{4}+\frac{\sin \left(0.5283553148724757\,r\right)}{5} ,   \sin \left(0.1092611211010309\,r\right)+\frac{\sin \left(  0.2185222422020618\,r\right)}{2}+\frac{\sin \left(0.3277833633030928  \,r\right)}{3}+\frac{\sin \left(0.4370444844041237\,r\right)}{4}+  \frac{\sin \left(0.5463056055051546\,r\right)}{5} , \sin \left(  0.1129282524731764\,r\right)+\frac{\sin \left(0.2258565049463528\,r  \right)}{2}+\frac{\sin \left(0.3387847574195292\,r\right)}{3}+\frac{  \sin \left(0.4517130098927056\,r\right)}{4}+\frac{\sin \left(  0.564641262365882\,r\right)}{5} , \sin \left(0.1166730903725168\,r  \right)+\frac{\sin \left(0.2333461807450337\,r\right)}{2}+\frac{  \sin \left(0.3500192711175505\,r\right)}{3}+\frac{\sin \left(  0.4666923614900673\,r\right)}{4}+\frac{\sin \left(0.5833654518625842  \,r\right)}{5} , \sin \left(0.1204962603100498\,r\right)+\frac{\sin   \left(0.2409925206200996\,r\right)}{2}+\frac{\sin \left(  0.3614887809301494\,r\right)}{3}+\frac{\sin \left(0.4819850412401991  \,r\right)}{4}+\frac{\sin \left(0.6024813015502489\,r\right)}{5} ,   \sin \left(0.1243983799636342\,r\right)+\frac{\sin \left(  0.2487967599272685\,r\right)}{2}+\frac{\sin \left(0.3731951398909027  \,r\right)}{3}+\frac{\sin \left(0.4975935198545369\,r\right)}{4}+  \frac{\sin \left(0.6219918998181712\,r\right)}{5} , \sin \left(  0.1283800591162231\,r\right)+\frac{\sin \left(0.2567601182324462\,r  \right)}{2}+\frac{\sin \left(0.3851401773486692\,r\right)}{3}+\frac{  \sin \left(0.5135202364648923\,r\right)}{4}+\frac{\sin \left(  0.6419002955811154\,r\right)}{5} , \sin \left(0.1324418995948859\,r  \right)+\frac{\sin \left(0.2648837991897719\,r\right)}{2}+\frac{  \sin \left(0.3973256987846578\,r\right)}{3}+\frac{\sin \left(  0.5297675983795438\,r\right)}{4}+\frac{\sin \left(0.6622094979744297  \,r\right)}{5} , \sin \left(0.1365844952106265\,r\right)+\frac{\sin   \left(0.2731689904212531\,r\right)}{2}+\frac{\sin \left(  0.4097534856318796\,r\right)}{3}+\frac{\sin \left(0.5463379808425062  \,r\right)}{4}+\frac{\sin \left(0.6829224760531327\,r\right)}{5} ,   \sin \left(0.140808431699002\,r\right)+\frac{\sin \left(  0.2816168633980041\,r\right)}{2}+\frac{\sin \left(0.4224252950970061  \,r\right)}{3}+\frac{\sin \left(0.5632337267960081\,r\right)}{4}+  \frac{\sin \left(0.7040421584950102\,r\right)}{5} , \sin \left(  0.1451142866615502\,r\right)+\frac{\sin \left(0.2902285733231005\,r  \right)}{2}+\frac{\sin \left(0.4353428599846507\,r\right)}{3}+\frac{  \sin \left(0.580457146646201\,r\right)}{4}+\frac{\sin \left(  0.7255714333077512\,r\right)}{5} , \sin \left(0.1495026295080298\,r  \right)+\frac{\sin \left(0.2990052590160597\,r\right)}{2}+\frac{  \sin \left(0.4485078885240895\,r\right)}{3}+\frac{\sin \left(  0.5980105180321194\,r\right)}{4}+\frac{\sin \left(0.7475131475401492  \,r\right)}{5} , \sin \left(0.1539740213994798\,r\right)+\frac{\sin   \left(0.3079480427989596\,r\right)}{2}+\frac{\sin \left(  0.4619220641984394\,r\right)}{3}+\frac{\sin \left(0.6158960855979192  \,r\right)}{4}+\frac{\sin \left(0.769870106997399\,r\right)}{5}   \right] 
\]
\end{eulerformula}
\begin{eulercomment}
Berikut adalah cara menggambar kurva

\end{eulercomment}
\begin{eulerformula}
\[
y=\sin(x) + \dfrac{\sin 3x}{3} + \dfrac{\sin 5x}{5} + \ldots.
\]
\end{eulerformula}
\begin{eulerprompt}
>plot2d(&sum(sin((2*k+1)*x)/(2*k+1),k,0,20),0,2pi):
\end{eulerprompt}
\begin{euleroutput}
  
  Maxima output too long!
  Error in:
  plot2d(&sum(sin((2*k+1)*x)/(2*k+1),k,0,20),0,2pi): ...
                                            ^
\end{euleroutput}
\begin{eulercomment}
Hal serupa juga dapat dilakukan dengan menggunakan matriks, misalkan
kita akan menggambar kurva

\end{eulercomment}
\begin{eulerformula}
\[
y = \sum_{k=1}^{100} \dfrac{\sin(kx)}{k},\quad 0\le x\le 2\pi.
\]
\end{eulerformula}
\begin{eulercomment}
\end{eulercomment}
\begin{eulerprompt}
>x=linspace(0,2pi,1000); k=1:100; y=sum(sin(k*x')/k)'; plot2d(x,y):
\end{eulerprompt}
\eulerimg{27}{images/Shalih Abdillah_22305144009_ApKom5-252.png}
\eulerheading{Tabel Fungsi}
\begin{eulercomment}
Terdapat cara menarik untuk menghasilkan barisan dengan ekspresi Maxima. Perintah
mxmtable() berguna untuk menampilkan dan menggambar barisan dan menghasilkan barisan sebagai
vektor kolom. 

Sebagai contoh berikut adalah barisan turunan ke-n x\textasciicircum{}x di x=1.
\end{eulercomment}
\begin{eulerprompt}
>mxmtable("diffat(x^x,x=1,n)","n",1,8,frac=1);
\end{eulerprompt}
\begin{euleroutput}
  Maxima said:
  diff: second argument must be a variable; found errexp1
  #0: diffat(expr=[0,1.66665833335744e-7*r,1.33330666692022e-6*r,4.499797504338432e-6*r,1.066581336583994e-5*r,2.08307...,x=[[0,1.66665833335744e-7*r,1.33330666692022e-6*r,4.499797504338432e-6*r,1.066581336583994e-5*r,2.0830...)
   -- an error. To debug this try: debugmode(true);
  
   %mxmevtable:
      return mxm("@expr,@var=@value")();
  Try "trace errors" to inspect local variables after errors.
  mxmtable:
      y[#,1]=%mxmevtable(expr,var,x[#]);
\end{euleroutput}
\begin{eulerprompt}
>$'sum(k, k, 1, n) = factor(ev(sum(k, k, 1, n),simpsum=true)) // simpsum:menghitung deret secara simbolik
\end{eulerprompt}
\begin{eulerformula}
\[
\sum_{k=1}^{n}{k}=\frac{n\,\left(1+n\right)}{2}
\]
\end{eulerformula}
\begin{eulerprompt}
>$'sum(1/(3^k+k), k, 0, inf) = factor(ev(sum(1/(3^k+k), k, 0, inf),simpsum=true))
\end{eulerprompt}
\begin{eulerformula}
\[
\sum_{k=0}^{\infty }{\frac{1}{k+3^{k}}}=\sum_{k=0}^{\infty }{\frac{  1}{k+3^{k}}}
\]
\end{eulerformula}
\begin{eulercomment}
Di sini masih gagal, hasilnya tidak dihitung.
\end{eulercomment}
\begin{eulerprompt}
>$'sum(1/x^2, x, 1, inf)= ev(sum(1/x^2, x, 1, inf),simpsum=true) // ev: menghitung nilai ekspresi
\end{eulerprompt}
\begin{eulerformula}
\[
\sum_{x=1}^{\infty }{\frac{1}{x^2}}=\frac{\pi^2}{6}
\]
\end{eulerformula}
\begin{eulerprompt}
>$'sum((-1)^(k-1)/k, k, 1, inf) = factor(ev(sum((-1)^(x-1)/x, x, 1, inf),simpsum=true))
\end{eulerprompt}
\begin{eulerformula}
\[
\sum_{k=1}^{\infty }{\frac{\left(-1\right)^{-1+k}}{k}}=-\sum_{x=1  }^{\infty }{\frac{\left(-1\right)^{x}}{x}}
\]
\end{eulerformula}
\begin{eulercomment}
Di sini masih gagal, hasilnya tidak dihitung.
\end{eulercomment}
\begin{eulerprompt}
>$'sum((-1)^k/(2*k-1), k, 1, inf) = factor(ev(sum((-1)^k/(2*k-1), k, 1, inf),simpsum=true))
\end{eulerprompt}
\begin{eulerformula}
\[
\sum_{k=1}^{\infty }{\frac{\left(-1\right)^{k}}{-1+2\,k}}=\sum_{k=1  }^{\infty }{\frac{\left(-1\right)^{k}}{-1+2\,k}}
\]
\end{eulerformula}
\begin{eulerprompt}
>$ev(sum(1/n!, n, 0, inf),simpsum=true)
\end{eulerprompt}
\begin{eulerformula}
\[
\sum_{n=0}^{\infty }{\frac{1}{n!}}
\]
\end{eulerformula}
\begin{eulercomment}
Di sini masih gagal, hasilnya tidak dihitung, harusnya hasilnya e.
\end{eulercomment}
\begin{eulerprompt}
>&assume(abs(x)<1); $'sum(a*x^k, k, 0, inf)=ev(sum(a*x^k, k, 0, inf),simpsum=true), &forget(abs(x)<1);
\end{eulerprompt}
\begin{euleroutput}
  Answering "Is -94914474571+15819*r positive, negative or zero?" with "positive"
  Maxima said:
  sum: sum is divergent.
   -- an error. To debug this try: debugmode(true);
  
  Error in:
  ... k, 0, inf)=ev(sum(a*x^k, k, 0, inf),simpsum=true), &forget(abs ...
                                                       ^
\end{euleroutput}
\begin{eulercomment}
Deret geometri tak hingga, dengan asumsi rasional antara -1 dan 1.
\end{eulercomment}
\begin{eulerprompt}
>$'sum(x^k/k!,k,0,inf)=ev(sum(x^k/k!,k,0,inf),simpsum=true)
\end{eulerprompt}
\begin{eulerformula}
\[
\left[ 0 , \sum_{k=0}^{\infty }{\frac{\left(  1.66665833335744 \times 10^{-7}\right)^{k}\,r^{k}}{k!}} , \sum_{k=0  }^{\infty }{\frac{\left(1.33330666692022 \times 10^{-6}\right)^{k}\,  r^{k}}{k!}} , \sum_{k=0}^{\infty }{\frac{\left(  4.499797504338432 \times 10^{-6}\right)^{k}\,r^{k}}{k!}} , \sum_{k=0  }^{\infty }{\frac{\left(1.066581336583994 \times 10^{-5}\right)^{k}  \,r^{k}}{k!}} , \sum_{k=0}^{\infty }{\frac{\left(  2.083072932167196 \times 10^{-5}\right)^{k}\,r^{k}}{k!}} , \sum_{k=0  }^{\infty }{\frac{\left(3.599352055540239 \times 10^{-5}\right)^{k}  \,r^{k}}{k!}} , \sum_{k=0}^{\infty }{\frac{\left(  5.71526624672386 \times 10^{-5}\right)^{k}\,r^{k}}{k!}} , \sum_{k=0  }^{\infty }{\frac{\left(8.530603082730626 \times 10^{-5}\right)^{k}  \,r^{k}}{k!}} , \sum_{k=0}^{\infty }{\frac{\left(  1.214508019889565 \times 10^{-4}\right)^{k}\,r^{k}}{k!}} , \sum_{k=0  }^{\infty }{\frac{\left(1.665833531718508 \times 10^{-4}\right)^{k}  \,r^{k}}{k!}} , \sum_{k=0}^{\infty }{\frac{\left(  2.216991628251896 \times 10^{-4}\right)^{k}\,r^{k}}{k!}} , \sum_{k=0  }^{\infty }{\frac{\left(2.877927110806339 \times 10^{-4}\right)^{k}  \,r^{k}}{k!}} , \sum_{k=0}^{\infty }{\frac{\left(  3.658573803051457 \times 10^{-4}\right)^{k}\,r^{k}}{k!}} , \sum_{k=0  }^{\infty }{\frac{\left(4.568853557635201 \times 10^{-4}\right)^{k}  \,r^{k}}{k!}} , \sum_{k=0}^{\infty }{\frac{\left(  5.618675264007778 \times 10^{-4}\right)^{k}\,r^{k}}{k!}} , \sum_{k=0  }^{\infty }{\frac{\left(6.817933857540259 \times 10^{-4}\right)^{k}  \,r^{k}}{k!}} , \sum_{k=0}^{\infty }{\frac{\left(  8.176509330039827 \times 10^{-4}\right)^{k}\,r^{k}}{k!}} , \sum_{k=0  }^{\infty }{\frac{\left(9.704265741758145 \times 10^{-4}\right)^{k}  \,r^{k}}{k!}} , \sum_{k=0}^{\infty }{\frac{0.001141105023499428^{k}  \,r^{k}}{k!}} , \sum_{k=0}^{\infty }{\frac{0.001330669204938795^{k}  \,r^{k}}{k!}} , \sum_{k=0}^{\infty }{\frac{0.001540100153900437^{k}  \,r^{k}}{k!}} , \sum_{k=0}^{\infty }{\frac{0.001770376919130678^{k}  \,r^{k}}{k!}} , \sum_{k=0}^{\infty }{\frac{0.002022476464811601^{k}  \,r^{k}}{k!}} , \sum_{k=0}^{\infty }{\frac{0.002297373572865413^{k}  \,r^{k}}{k!}} , \sum_{k=0}^{\infty }{\frac{0.002596040745477063^{k}  \,r^{k}}{k!}} , \sum_{k=0}^{\infty }{\frac{0.002919448107844891^{k}  \,r^{k}}{k!}} , \sum_{k=0}^{\infty }{\frac{0.003268563311168871^{k}  \,r^{k}}{k!}} , \sum_{k=0}^{\infty }{\frac{0.003644351435886262^{k}  \,r^{k}}{k!}} , \sum_{k=0}^{\infty }{\frac{0.004047774895164447^{k}  \,r^{k}}{k!}} , \sum_{k=0}^{\infty }{\frac{0.004479793338660443^{k}  \,r^{k}}{k!}} , \sum_{k=0}^{\infty }{\frac{0.0049413635565565^{k}\,r  ^{k}}{k!}} , \sum_{k=0}^{\infty }{\frac{0.005433439383882244^{k}\,r  ^{k}}{k!}} , \sum_{k=0}^{\infty }{\frac{0.005956971605131645^{k}\,r  ^{k}}{k!}} , \sum_{k=0}^{\infty }{\frac{0.006512907859185624^{k}\,r  ^{k}}{k!}} , \sum_{k=0}^{\infty }{\frac{0.007102192544548636^{k}\,r  ^{k}}{k!}} , \sum_{k=0}^{\infty }{\frac{0.007725766724910044^{k}\,r  ^{k}}{k!}} , \sum_{k=0}^{\infty }{\frac{0.00838456803503801^{k}\,r^{  k}}{k!}} , \sum_{k=0}^{\infty }{\frac{0.009079530587017326^{k}\,r^{k  }}{k!}} , \sum_{k=0}^{\infty }{\frac{0.009811584876838586^{k}\,r^{k}  }{k!}} , \sum_{k=0}^{\infty }{\frac{0.0105816576913495^{k}\,r^{k}}{k  !}} , \sum_{k=0}^{\infty }{\frac{0.01139067201557714^{k}\,r^{k}}{k!}  } , \sum_{k=0}^{\infty }{\frac{0.01223954694042984^{k}\,r^{k}}{k!}}   , \sum_{k=0}^{\infty }{\frac{0.01312919757078923^{k}\,r^{k}}{k!}}   , \sum_{k=0}^{\infty }{\frac{0.01406053493400045^{k}\,r^{k}}{k!}}   , \sum_{k=0}^{\infty }{\frac{0.01503446588876983^{k}\,r^{k}}{k!}}   , \sum_{k=0}^{\infty }{\frac{0.01605189303448024^{k}\,r^{k}}{k!}}   , \sum_{k=0}^{\infty }{\frac{0.01711371462093175^{k}\,r^{k}}{k!}}   , \sum_{k=0}^{\infty }{\frac{0.01822082445851714^{k}\,r^{k}}{k!}}   , \sum_{k=0}^{\infty }{\frac{0.01937411182884202^{k}\,r^{k}}{k!}}   , \sum_{k=0}^{\infty }{\frac{0.02057446139579705^{k}\,r^{k}}{k!}}   , \sum_{k=0}^{\infty }{\frac{0.02182275311709253^{k}\,r^{k}}{k!}}   , \sum_{k=0}^{\infty }{\frac{0.02311986215626333^{k}\,r^{k}}{k!}}   , \sum_{k=0}^{\infty }{\frac{0.02446665879515308^{k}\,r^{k}}{k!}}   , \sum_{k=0}^{\infty }{\frac{0.02586400834688696^{k}\,r^{k}}{k!}}   , \sum_{k=0}^{\infty }{\frac{0.02731277106934082^{k}\,r^{k}}{k!}}   , \sum_{k=0}^{\infty }{\frac{0.02881380207911666^{k}\,r^{k}}{k!}}   , \sum_{k=0}^{\infty }{\frac{0.03036795126603076^{k}\,r^{k}}{k!}}   , \sum_{k=0}^{\infty }{\frac{0.03197606320812652^{k}\,r^{k}}{k!}}   , \sum_{k=0}^{\infty }{\frac{0.0336389770872163^{k}\,r^{k}}{k!}} ,   \sum_{k=0}^{\infty }{\frac{0.03535752660496472^{k}\,r^{k}}{k!}} ,   \sum_{k=0}^{\infty }{\frac{0.03713253989951881^{k}\,r^{k}}{k!}} ,   \sum_{k=0}^{\infty }{\frac{0.03896483946269502^{k}\,r^{k}}{k!}} ,   \sum_{k=0}^{\infty }{\frac{0.0408552420577305^{k}\,r^{k}}{k!}} ,   \sum_{k=0}^{\infty }{\frac{0.04280455863760801^{k}\,r^{k}}{k!}} ,   \sum_{k=0}^{\infty }{\frac{0.04481359426396048^{k}\,r^{k}}{k!}} ,   \sum_{k=0}^{\infty }{\frac{0.04688314802656623^{k}\,r^{k}}{k!}} ,   \sum_{k=0}^{\infty }{\frac{0.04901401296344043^{k}\,r^{k}}{k!}} ,   \sum_{k=0}^{\infty }{\frac{0.05120697598153157^{k}\,r^{k}}{k!}} ,   \sum_{k=0}^{\infty }{\frac{0.05346281777803219^{k}\,r^{k}}{k!}} ,   \sum_{k=0}^{\infty }{\frac{0.05578231276230905^{k}\,r^{k}}{k!}} ,   \sum_{k=0}^{\infty }{\frac{0.05816622897846346^{k}\,r^{k}}{k!}} ,   \sum_{k=0}^{\infty }{\frac{0.06061532802852698^{k}\,r^{k}}{k!}} ,   \sum_{k=0}^{\infty }{\frac{0.0631303649963022^{k}\,r^{k}}{k!}} ,   \sum_{k=0}^{\infty }{\frac{0.06571208837185505^{k}\,r^{k}}{k!}} ,   \sum_{k=0}^{\infty }{\frac{0.06836123997666599^{k}\,r^{k}}{k!}} ,   \sum_{k=0}^{\infty }{\frac{0.07107855488944881^{k}\,r^{k}}{k!}} ,   \sum_{k=0}^{\infty }{\frac{0.07386476137264342^{k}\,r^{k}}{k!}} ,   \sum_{k=0}^{\infty }{\frac{0.07672058079958999^{k}\,r^{k}}{k!}} ,   \sum_{k=0}^{\infty }{\frac{0.07964672758239233^{k}\,r^{k}}{k!}} ,   \sum_{k=0}^{\infty }{\frac{0.08264390910047736^{k}\,r^{k}}{k!}} ,   \sum_{k=0}^{\infty }{\frac{0.0857128256298576^{k}\,r^{k}}{k!}} ,   \sum_{k=0}^{\infty }{\frac{0.08885417027310427^{k}\,r^{k}}{k!}} ,   \sum_{k=0}^{\infty }{\frac{0.09206862889003742^{k}\,r^{k}}{k!}} ,   \sum_{k=0}^{\infty }{\frac{0.09535688002914089^{k}\,r^{k}}{k!}} ,   \sum_{k=0}^{\infty }{\frac{0.0987195948597075^{k}\,r^{k}}{k!}} ,   \sum_{k=0}^{\infty }{\frac{0.1021574371047232^{k}\,r^{k}}{k!}} ,   \sum_{k=0}^{\infty }{\frac{0.1056710629744951^{k}\,r^{k}}{k!}} ,   \sum_{k=0}^{\infty }{\frac{0.1092611211010309^{k}\,r^{k}}{k!}} ,   \sum_{k=0}^{\infty }{\frac{0.1129282524731764^{k}\,r^{k}}{k!}} ,   \sum_{k=0}^{\infty }{\frac{0.1166730903725168^{k}\,r^{k}}{k!}} ,   \sum_{k=0}^{\infty }{\frac{0.1204962603100498^{k}\,r^{k}}{k!}} ,   \sum_{k=0}^{\infty }{\frac{0.1243983799636342^{k}\,r^{k}}{k!}} ,   \sum_{k=0}^{\infty }{\frac{0.1283800591162231^{k}\,r^{k}}{k!}} ,   \sum_{k=0}^{\infty }{\frac{0.1324418995948859^{k}\,r^{k}}{k!}} ,   \sum_{k=0}^{\infty }{\frac{0.1365844952106265^{k}\,r^{k}}{k!}} ,   \sum_{k=0}^{\infty }{\frac{0.140808431699002^{k}\,r^{k}}{k!}} ,   \sum_{k=0}^{\infty }{\frac{0.1451142866615502^{k}\,r^{k}}{k!}} ,   \sum_{k=0}^{\infty }{\frac{0.1495026295080298^{k}\,r^{k}}{k!}} ,   \sum_{k=0}^{\infty }{\frac{0.1539740213994798^{k}\,r^{k}}{k!}}   \right] =\left[ 0 , \sum_{k=0}^{\infty }{\frac{\left(  1.66665833335744 \times 10^{-7}\right)^{k}\,r^{k}}{k!}} , \sum_{k=0  }^{\infty }{\frac{\left(1.33330666692022 \times 10^{-6}\right)^{k}\,  r^{k}}{k!}} , \sum_{k=0}^{\infty }{\frac{\left(  4.499797504338432 \times 10^{-6}\right)^{k}\,r^{k}}{k!}} , \sum_{k=0  }^{\infty }{\frac{\left(1.066581336583994 \times 10^{-5}\right)^{k}  \,r^{k}}{k!}} , \sum_{k=0}^{\infty }{\frac{\left(  2.083072932167196 \times 10^{-5}\right)^{k}\,r^{k}}{k!}} , \sum_{k=0  }^{\infty }{\frac{\left(3.599352055540239 \times 10^{-5}\right)^{k}  \,r^{k}}{k!}} , \sum_{k=0}^{\infty }{\frac{\left(  5.71526624672386 \times 10^{-5}\right)^{k}\,r^{k}}{k!}} , \sum_{k=0  }^{\infty }{\frac{\left(8.530603082730626 \times 10^{-5}\right)^{k}  \,r^{k}}{k!}} , \sum_{k=0}^{\infty }{\frac{\left(  1.214508019889565 \times 10^{-4}\right)^{k}\,r^{k}}{k!}} , \sum_{k=0  }^{\infty }{\frac{\left(1.665833531718508 \times 10^{-4}\right)^{k}  \,r^{k}}{k!}} , \sum_{k=0}^{\infty }{\frac{\left(  2.216991628251896 \times 10^{-4}\right)^{k}\,r^{k}}{k!}} , \sum_{k=0  }^{\infty }{\frac{\left(2.877927110806339 \times 10^{-4}\right)^{k}  \,r^{k}}{k!}} , \sum_{k=0}^{\infty }{\frac{\left(  3.658573803051457 \times 10^{-4}\right)^{k}\,r^{k}}{k!}} , \sum_{k=0  }^{\infty }{\frac{\left(4.568853557635201 \times 10^{-4}\right)^{k}  \,r^{k}}{k!}} , \sum_{k=0}^{\infty }{\frac{\left(  5.618675264007778 \times 10^{-4}\right)^{k}\,r^{k}}{k!}} , \sum_{k=0  }^{\infty }{\frac{\left(6.817933857540259 \times 10^{-4}\right)^{k}  \,r^{k}}{k!}} , \sum_{k=0}^{\infty }{\frac{\left(  8.176509330039827 \times 10^{-4}\right)^{k}\,r^{k}}{k!}} , \sum_{k=0  }^{\infty }{\frac{\left(9.704265741758145 \times 10^{-4}\right)^{k}  \,r^{k}}{k!}} , \sum_{k=0}^{\infty }{\frac{0.001141105023499428^{k}  \,r^{k}}{k!}} , \sum_{k=0}^{\infty }{\frac{0.001330669204938795^{k}  \,r^{k}}{k!}} , \sum_{k=0}^{\infty }{\frac{0.001540100153900437^{k}  \,r^{k}}{k!}} , \sum_{k=0}^{\infty }{\frac{0.001770376919130678^{k}  \,r^{k}}{k!}} , \sum_{k=0}^{\infty }{\frac{0.002022476464811601^{k}  \,r^{k}}{k!}} , \sum_{k=0}^{\infty }{\frac{0.002297373572865413^{k}  \,r^{k}}{k!}} , \sum_{k=0}^{\infty }{\frac{0.002596040745477063^{k}  \,r^{k}}{k!}} , \sum_{k=0}^{\infty }{\frac{0.002919448107844891^{k}  \,r^{k}}{k!}} , \sum_{k=0}^{\infty }{\frac{0.003268563311168871^{k}  \,r^{k}}{k!}} , \sum_{k=0}^{\infty }{\frac{0.003644351435886262^{k}  \,r^{k}}{k!}} , \sum_{k=0}^{\infty }{\frac{0.004047774895164447^{k}  \,r^{k}}{k!}} , \sum_{k=0}^{\infty }{\frac{0.004479793338660443^{k}  \,r^{k}}{k!}} , \sum_{k=0}^{\infty }{\frac{0.0049413635565565^{k}\,r  ^{k}}{k!}} , \sum_{k=0}^{\infty }{\frac{0.005433439383882244^{k}\,r  ^{k}}{k!}} , \sum_{k=0}^{\infty }{\frac{0.005956971605131645^{k}\,r  ^{k}}{k!}} , \sum_{k=0}^{\infty }{\frac{0.006512907859185624^{k}\,r  ^{k}}{k!}} , \sum_{k=0}^{\infty }{\frac{0.007102192544548636^{k}\,r  ^{k}}{k!}} , \sum_{k=0}^{\infty }{\frac{0.007725766724910044^{k}\,r  ^{k}}{k!}} , \sum_{k=0}^{\infty }{\frac{0.00838456803503801^{k}\,r^{  k}}{k!}} , \sum_{k=0}^{\infty }{\frac{0.009079530587017326^{k}\,r^{k  }}{k!}} , \sum_{k=0}^{\infty }{\frac{0.009811584876838586^{k}\,r^{k}  }{k!}} , \sum_{k=0}^{\infty }{\frac{0.0105816576913495^{k}\,r^{k}}{k  !}} , \sum_{k=0}^{\infty }{\frac{0.01139067201557714^{k}\,r^{k}}{k!}  } , \sum_{k=0}^{\infty }{\frac{0.01223954694042984^{k}\,r^{k}}{k!}}   , \sum_{k=0}^{\infty }{\frac{0.01312919757078923^{k}\,r^{k}}{k!}}   , \sum_{k=0}^{\infty }{\frac{0.01406053493400045^{k}\,r^{k}}{k!}}   , \sum_{k=0}^{\infty }{\frac{0.01503446588876983^{k}\,r^{k}}{k!}}   , \sum_{k=0}^{\infty }{\frac{0.01605189303448024^{k}\,r^{k}}{k!}}   , \sum_{k=0}^{\infty }{\frac{0.01711371462093175^{k}\,r^{k}}{k!}}   , \sum_{k=0}^{\infty }{\frac{0.01822082445851714^{k}\,r^{k}}{k!}}   , \sum_{k=0}^{\infty }{\frac{0.01937411182884202^{k}\,r^{k}}{k!}}   , \sum_{k=0}^{\infty }{\frac{0.02057446139579705^{k}\,r^{k}}{k!}}   , \sum_{k=0}^{\infty }{\frac{0.02182275311709253^{k}\,r^{k}}{k!}}   , \sum_{k=0}^{\infty }{\frac{0.02311986215626333^{k}\,r^{k}}{k!}}   , \sum_{k=0}^{\infty }{\frac{0.02446665879515308^{k}\,r^{k}}{k!}}   , \sum_{k=0}^{\infty }{\frac{0.02586400834688696^{k}\,r^{k}}{k!}}   , \sum_{k=0}^{\infty }{\frac{0.02731277106934082^{k}\,r^{k}}{k!}}   , \sum_{k=0}^{\infty }{\frac{0.02881380207911666^{k}\,r^{k}}{k!}}   , \sum_{k=0}^{\infty }{\frac{0.03036795126603076^{k}\,r^{k}}{k!}}   , \sum_{k=0}^{\infty }{\frac{0.03197606320812652^{k}\,r^{k}}{k!}}   , \sum_{k=0}^{\infty }{\frac{0.0336389770872163^{k}\,r^{k}}{k!}} ,   \sum_{k=0}^{\infty }{\frac{0.03535752660496472^{k}\,r^{k}}{k!}} ,   \sum_{k=0}^{\infty }{\frac{0.03713253989951881^{k}\,r^{k}}{k!}} ,   \sum_{k=0}^{\infty }{\frac{0.03896483946269502^{k}\,r^{k}}{k!}} ,   \sum_{k=0}^{\infty }{\frac{0.0408552420577305^{k}\,r^{k}}{k!}} ,   \sum_{k=0}^{\infty }{\frac{0.04280455863760801^{k}\,r^{k}}{k!}} ,   \sum_{k=0}^{\infty }{\frac{0.04481359426396048^{k}\,r^{k}}{k!}} ,   \sum_{k=0}^{\infty }{\frac{0.04688314802656623^{k}\,r^{k}}{k!}} ,   \sum_{k=0}^{\infty }{\frac{0.04901401296344043^{k}\,r^{k}}{k!}} ,   \sum_{k=0}^{\infty }{\frac{0.05120697598153157^{k}\,r^{k}}{k!}} ,   \sum_{k=0}^{\infty }{\frac{0.05346281777803219^{k}\,r^{k}}{k!}} ,   \sum_{k=0}^{\infty }{\frac{0.05578231276230905^{k}\,r^{k}}{k!}} ,   \sum_{k=0}^{\infty }{\frac{0.05816622897846346^{k}\,r^{k}}{k!}} ,   \sum_{k=0}^{\infty }{\frac{0.06061532802852698^{k}\,r^{k}}{k!}} ,   \sum_{k=0}^{\infty }{\frac{0.0631303649963022^{k}\,r^{k}}{k!}} ,   \sum_{k=0}^{\infty }{\frac{0.06571208837185505^{k}\,r^{k}}{k!}} ,   \sum_{k=0}^{\infty }{\frac{0.06836123997666599^{k}\,r^{k}}{k!}} ,   \sum_{k=0}^{\infty }{\frac{0.07107855488944881^{k}\,r^{k}}{k!}} ,   \sum_{k=0}^{\infty }{\frac{0.07386476137264342^{k}\,r^{k}}{k!}} ,   \sum_{k=0}^{\infty }{\frac{0.07672058079958999^{k}\,r^{k}}{k!}} ,   \sum_{k=0}^{\infty }{\frac{0.07964672758239233^{k}\,r^{k}}{k!}} ,   \sum_{k=0}^{\infty }{\frac{0.08264390910047736^{k}\,r^{k}}{k!}} ,   \sum_{k=0}^{\infty }{\frac{0.0857128256298576^{k}\,r^{k}}{k!}} ,   \sum_{k=0}^{\infty }{\frac{0.08885417027310427^{k}\,r^{k}}{k!}} ,   \sum_{k=0}^{\infty }{\frac{0.09206862889003742^{k}\,r^{k}}{k!}} ,   \sum_{k=0}^{\infty }{\frac{0.09535688002914089^{k}\,r^{k}}{k!}} ,   \sum_{k=0}^{\infty }{\frac{0.0987195948597075^{k}\,r^{k}}{k!}} ,   \sum_{k=0}^{\infty }{\frac{0.1021574371047232^{k}\,r^{k}}{k!}} ,   \sum_{k=0}^{\infty }{\frac{0.1056710629744951^{k}\,r^{k}}{k!}} ,   \sum_{k=0}^{\infty }{\frac{0.1092611211010309^{k}\,r^{k}}{k!}} ,   \sum_{k=0}^{\infty }{\frac{0.1129282524731764^{k}\,r^{k}}{k!}} ,   \sum_{k=0}^{\infty }{\frac{0.1166730903725168^{k}\,r^{k}}{k!}} ,   \sum_{k=0}^{\infty }{\frac{0.1204962603100498^{k}\,r^{k}}{k!}} ,   \sum_{k=0}^{\infty }{\frac{0.1243983799636342^{k}\,r^{k}}{k!}} ,   \sum_{k=0}^{\infty }{\frac{0.1283800591162231^{k}\,r^{k}}{k!}} ,   \sum_{k=0}^{\infty }{\frac{0.1324418995948859^{k}\,r^{k}}{k!}} ,   \sum_{k=0}^{\infty }{\frac{0.1365844952106265^{k}\,r^{k}}{k!}} ,   \sum_{k=0}^{\infty }{\frac{0.140808431699002^{k}\,r^{k}}{k!}} ,   \sum_{k=0}^{\infty }{\frac{0.1451142866615502^{k}\,r^{k}}{k!}} ,   \sum_{k=0}^{\infty }{\frac{0.1495026295080298^{k}\,r^{k}}{k!}} ,   \sum_{k=0}^{\infty }{\frac{0.1539740213994798^{k}\,r^{k}}{k!}}   \right] 
\]
\end{eulerformula}
\begin{eulerprompt}
>$limit(sum(x^k/k!,k,0,n),n,inf)
\end{eulerprompt}
\begin{eulerformula}
\[
\left[ 0 , \lim_{n\rightarrow \infty }{\sum_{k=0}^{n}{\frac{\left(  1.66665833335744 \times 10^{-7}\right)^{k}\,r^{k}}{k!}}} , \lim_{n  \rightarrow \infty }{\sum_{k=0}^{n}{\frac{\left(  1.33330666692022 \times 10^{-6}\right)^{k}\,r^{k}}{k!}}} , \lim_{n  \rightarrow \infty }{\sum_{k=0}^{n}{\frac{\left(  4.499797504338432 \times 10^{-6}\right)^{k}\,r^{k}}{k!}}} , \lim_{n  \rightarrow \infty }{\sum_{k=0}^{n}{\frac{\left(  1.066581336583994 \times 10^{-5}\right)^{k}\,r^{k}}{k!}}} , \lim_{n  \rightarrow \infty }{\sum_{k=0}^{n}{\frac{\left(  2.083072932167196 \times 10^{-5}\right)^{k}\,r^{k}}{k!}}} , \lim_{n  \rightarrow \infty }{\sum_{k=0}^{n}{\frac{\left(  3.599352055540239 \times 10^{-5}\right)^{k}\,r^{k}}{k!}}} , \lim_{n  \rightarrow \infty }{\sum_{k=0}^{n}{\frac{\left(  5.71526624672386 \times 10^{-5}\right)^{k}\,r^{k}}{k!}}} , \lim_{n  \rightarrow \infty }{\sum_{k=0}^{n}{\frac{\left(  8.530603082730626 \times 10^{-5}\right)^{k}\,r^{k}}{k!}}} , \lim_{n  \rightarrow \infty }{\sum_{k=0}^{n}{\frac{\left(  1.214508019889565 \times 10^{-4}\right)^{k}\,r^{k}}{k!}}} , \lim_{n  \rightarrow \infty }{\sum_{k=0}^{n}{\frac{\left(  1.665833531718508 \times 10^{-4}\right)^{k}\,r^{k}}{k!}}} , \lim_{n  \rightarrow \infty }{\sum_{k=0}^{n}{\frac{\left(  2.216991628251896 \times 10^{-4}\right)^{k}\,r^{k}}{k!}}} , \lim_{n  \rightarrow \infty }{\sum_{k=0}^{n}{\frac{\left(  2.877927110806339 \times 10^{-4}\right)^{k}\,r^{k}}{k!}}} , \lim_{n  \rightarrow \infty }{\sum_{k=0}^{n}{\frac{\left(  3.658573803051457 \times 10^{-4}\right)^{k}\,r^{k}}{k!}}} , \lim_{n  \rightarrow \infty }{\sum_{k=0}^{n}{\frac{\left(  4.568853557635201 \times 10^{-4}\right)^{k}\,r^{k}}{k!}}} , \lim_{n  \rightarrow \infty }{\sum_{k=0}^{n}{\frac{\left(  5.618675264007778 \times 10^{-4}\right)^{k}\,r^{k}}{k!}}} , \lim_{n  \rightarrow \infty }{\sum_{k=0}^{n}{\frac{\left(  6.817933857540259 \times 10^{-4}\right)^{k}\,r^{k}}{k!}}} , \lim_{n  \rightarrow \infty }{\sum_{k=0}^{n}{\frac{\left(  8.176509330039827 \times 10^{-4}\right)^{k}\,r^{k}}{k!}}} , \lim_{n  \rightarrow \infty }{\sum_{k=0}^{n}{\frac{\left(  9.704265741758145 \times 10^{-4}\right)^{k}\,r^{k}}{k!}}} , \lim_{n  \rightarrow \infty }{\sum_{k=0}^{n}{\frac{0.001141105023499428^{k}\,  r^{k}}{k!}}} , \lim_{n\rightarrow \infty }{\sum_{k=0}^{n}{\frac{  0.001330669204938795^{k}\,r^{k}}{k!}}} , \lim_{n\rightarrow \infty   }{\sum_{k=0}^{n}{\frac{0.001540100153900437^{k}\,r^{k}}{k!}}} ,   \lim_{n\rightarrow \infty }{\sum_{k=0}^{n}{\frac{  0.001770376919130678^{k}\,r^{k}}{k!}}} , \lim_{n\rightarrow \infty   }{\sum_{k=0}^{n}{\frac{0.002022476464811601^{k}\,r^{k}}{k!}}} ,   \lim_{n\rightarrow \infty }{\sum_{k=0}^{n}{\frac{  0.002297373572865413^{k}\,r^{k}}{k!}}} , \lim_{n\rightarrow \infty   }{\sum_{k=0}^{n}{\frac{0.002596040745477063^{k}\,r^{k}}{k!}}} ,   \lim_{n\rightarrow \infty }{\sum_{k=0}^{n}{\frac{  0.002919448107844891^{k}\,r^{k}}{k!}}} , \lim_{n\rightarrow \infty   }{\sum_{k=0}^{n}{\frac{0.003268563311168871^{k}\,r^{k}}{k!}}} ,   \lim_{n\rightarrow \infty }{\sum_{k=0}^{n}{\frac{  0.003644351435886262^{k}\,r^{k}}{k!}}} , \lim_{n\rightarrow \infty   }{\sum_{k=0}^{n}{\frac{0.004047774895164447^{k}\,r^{k}}{k!}}} ,   \lim_{n\rightarrow \infty }{\sum_{k=0}^{n}{\frac{  0.004479793338660443^{k}\,r^{k}}{k!}}} , \lim_{n\rightarrow \infty   }{\sum_{k=0}^{n}{\frac{0.0049413635565565^{k}\,r^{k}}{k!}}} , \lim_{  n\rightarrow \infty }{\sum_{k=0}^{n}{\frac{0.005433439383882244^{k}  \,r^{k}}{k!}}} , \lim_{n\rightarrow \infty }{\sum_{k=0}^{n}{\frac{  0.005956971605131645^{k}\,r^{k}}{k!}}} , \lim_{n\rightarrow \infty   }{\sum_{k=0}^{n}{\frac{0.006512907859185624^{k}\,r^{k}}{k!}}} ,   \lim_{n\rightarrow \infty }{\sum_{k=0}^{n}{\frac{  0.007102192544548636^{k}\,r^{k}}{k!}}} , \lim_{n\rightarrow \infty   }{\sum_{k=0}^{n}{\frac{0.007725766724910044^{k}\,r^{k}}{k!}}} ,   \lim_{n\rightarrow \infty }{\sum_{k=0}^{n}{\frac{0.00838456803503801  ^{k}\,r^{k}}{k!}}} , \lim_{n\rightarrow \infty }{\sum_{k=0}^{n}{  \frac{0.009079530587017326^{k}\,r^{k}}{k!}}} , \lim_{n\rightarrow   \infty }{\sum_{k=0}^{n}{\frac{0.009811584876838586^{k}\,r^{k}}{k!}}}   , \lim_{n\rightarrow \infty }{\sum_{k=0}^{n}{\frac{  0.0105816576913495^{k}\,r^{k}}{k!}}} , \lim_{n\rightarrow \infty }{  \sum_{k=0}^{n}{\frac{0.01139067201557714^{k}\,r^{k}}{k!}}} , \lim_{n  \rightarrow \infty }{\sum_{k=0}^{n}{\frac{0.01223954694042984^{k}\,r  ^{k}}{k!}}} , \lim_{n\rightarrow \infty }{\sum_{k=0}^{n}{\frac{  0.01312919757078923^{k}\,r^{k}}{k!}}} , \lim_{n\rightarrow \infty }{  \sum_{k=0}^{n}{\frac{0.01406053493400045^{k}\,r^{k}}{k!}}} , \lim_{n  \rightarrow \infty }{\sum_{k=0}^{n}{\frac{0.01503446588876983^{k}\,r  ^{k}}{k!}}} , \lim_{n\rightarrow \infty }{\sum_{k=0}^{n}{\frac{  0.01605189303448024^{k}\,r^{k}}{k!}}} , \lim_{n\rightarrow \infty }{  \sum_{k=0}^{n}{\frac{0.01711371462093175^{k}\,r^{k}}{k!}}} , \lim_{n  \rightarrow \infty }{\sum_{k=0}^{n}{\frac{0.01822082445851714^{k}\,r  ^{k}}{k!}}} , \lim_{n\rightarrow \infty }{\sum_{k=0}^{n}{\frac{  0.01937411182884202^{k}\,r^{k}}{k!}}} , \lim_{n\rightarrow \infty }{  \sum_{k=0}^{n}{\frac{0.02057446139579705^{k}\,r^{k}}{k!}}} , \lim_{n  \rightarrow \infty }{\sum_{k=0}^{n}{\frac{0.02182275311709253^{k}\,r  ^{k}}{k!}}} , \lim_{n\rightarrow \infty }{\sum_{k=0}^{n}{\frac{  0.02311986215626333^{k}\,r^{k}}{k!}}} , \lim_{n\rightarrow \infty }{  \sum_{k=0}^{n}{\frac{0.02446665879515308^{k}\,r^{k}}{k!}}} , \lim_{n  \rightarrow \infty }{\sum_{k=0}^{n}{\frac{0.02586400834688696^{k}\,r  ^{k}}{k!}}} , \lim_{n\rightarrow \infty }{\sum_{k=0}^{n}{\frac{  0.02731277106934082^{k}\,r^{k}}{k!}}} , \lim_{n\rightarrow \infty }{  \sum_{k=0}^{n}{\frac{0.02881380207911666^{k}\,r^{k}}{k!}}} , \lim_{n  \rightarrow \infty }{\sum_{k=0}^{n}{\frac{0.03036795126603076^{k}\,r  ^{k}}{k!}}} , \lim_{n\rightarrow \infty }{\sum_{k=0}^{n}{\frac{  0.03197606320812652^{k}\,r^{k}}{k!}}} , \lim_{n\rightarrow \infty }{  \sum_{k=0}^{n}{\frac{0.0336389770872163^{k}\,r^{k}}{k!}}} , \lim_{n  \rightarrow \infty }{\sum_{k=0}^{n}{\frac{0.03535752660496472^{k}\,r  ^{k}}{k!}}} , \lim_{n\rightarrow \infty }{\sum_{k=0}^{n}{\frac{  0.03713253989951881^{k}\,r^{k}}{k!}}} , \lim_{n\rightarrow \infty }{  \sum_{k=0}^{n}{\frac{0.03896483946269502^{k}\,r^{k}}{k!}}} , \lim_{n  \rightarrow \infty }{\sum_{k=0}^{n}{\frac{0.0408552420577305^{k}\,r  ^{k}}{k!}}} , \lim_{n\rightarrow \infty }{\sum_{k=0}^{n}{\frac{  0.04280455863760801^{k}\,r^{k}}{k!}}} , \lim_{n\rightarrow \infty }{  \sum_{k=0}^{n}{\frac{0.04481359426396048^{k}\,r^{k}}{k!}}} , \lim_{n  \rightarrow \infty }{\sum_{k=0}^{n}{\frac{0.04688314802656623^{k}\,r  ^{k}}{k!}}} , \lim_{n\rightarrow \infty }{\sum_{k=0}^{n}{\frac{  0.04901401296344043^{k}\,r^{k}}{k!}}} , \lim_{n\rightarrow \infty }{  \sum_{k=0}^{n}{\frac{0.05120697598153157^{k}\,r^{k}}{k!}}} , \lim_{n  \rightarrow \infty }{\sum_{k=0}^{n}{\frac{0.05346281777803219^{k}\,r  ^{k}}{k!}}} , \lim_{n\rightarrow \infty }{\sum_{k=0}^{n}{\frac{  0.05578231276230905^{k}\,r^{k}}{k!}}} , \lim_{n\rightarrow \infty }{  \sum_{k=0}^{n}{\frac{0.05816622897846346^{k}\,r^{k}}{k!}}} , \lim_{n  \rightarrow \infty }{\sum_{k=0}^{n}{\frac{0.06061532802852698^{k}\,r  ^{k}}{k!}}} , \lim_{n\rightarrow \infty }{\sum_{k=0}^{n}{\frac{  0.0631303649963022^{k}\,r^{k}}{k!}}} , \lim_{n\rightarrow \infty }{  \sum_{k=0}^{n}{\frac{0.06571208837185505^{k}\,r^{k}}{k!}}} , \lim_{n  \rightarrow \infty }{\sum_{k=0}^{n}{\frac{0.06836123997666599^{k}\,r  ^{k}}{k!}}} , \lim_{n\rightarrow \infty }{\sum_{k=0}^{n}{\frac{  0.07107855488944881^{k}\,r^{k}}{k!}}} , \lim_{n\rightarrow \infty }{  \sum_{k=0}^{n}{\frac{0.07386476137264342^{k}\,r^{k}}{k!}}} , \lim_{n  \rightarrow \infty }{\sum_{k=0}^{n}{\frac{0.07672058079958999^{k}\,r  ^{k}}{k!}}} , \lim_{n\rightarrow \infty }{\sum_{k=0}^{n}{\frac{  0.07964672758239233^{k}\,r^{k}}{k!}}} , \lim_{n\rightarrow \infty }{  \sum_{k=0}^{n}{\frac{0.08264390910047736^{k}\,r^{k}}{k!}}} , \lim_{n  \rightarrow \infty }{\sum_{k=0}^{n}{\frac{0.0857128256298576^{k}\,r  ^{k}}{k!}}} , \lim_{n\rightarrow \infty }{\sum_{k=0}^{n}{\frac{  0.08885417027310427^{k}\,r^{k}}{k!}}} , \lim_{n\rightarrow \infty }{  \sum_{k=0}^{n}{\frac{0.09206862889003742^{k}\,r^{k}}{k!}}} , \lim_{n  \rightarrow \infty }{\sum_{k=0}^{n}{\frac{0.09535688002914089^{k}\,r  ^{k}}{k!}}} , \lim_{n\rightarrow \infty }{\sum_{k=0}^{n}{\frac{  0.0987195948597075^{k}\,r^{k}}{k!}}} , \lim_{n\rightarrow \infty }{  \sum_{k=0}^{n}{\frac{0.1021574371047232^{k}\,r^{k}}{k!}}} , \lim_{n  \rightarrow \infty }{\sum_{k=0}^{n}{\frac{0.1056710629744951^{k}\,r  ^{k}}{k!}}} , \lim_{n\rightarrow \infty }{\sum_{k=0}^{n}{\frac{  0.1092611211010309^{k}\,r^{k}}{k!}}} , \lim_{n\rightarrow \infty }{  \sum_{k=0}^{n}{\frac{0.1129282524731764^{k}\,r^{k}}{k!}}} , \lim_{n  \rightarrow \infty }{\sum_{k=0}^{n}{\frac{0.1166730903725168^{k}\,r  ^{k}}{k!}}} , \lim_{n\rightarrow \infty }{\sum_{k=0}^{n}{\frac{  0.1204962603100498^{k}\,r^{k}}{k!}}} , \lim_{n\rightarrow \infty }{  \sum_{k=0}^{n}{\frac{0.1243983799636342^{k}\,r^{k}}{k!}}} , \lim_{n  \rightarrow \infty }{\sum_{k=0}^{n}{\frac{0.1283800591162231^{k}\,r  ^{k}}{k!}}} , \lim_{n\rightarrow \infty }{\sum_{k=0}^{n}{\frac{  0.1324418995948859^{k}\,r^{k}}{k!}}} , \lim_{n\rightarrow \infty }{  \sum_{k=0}^{n}{\frac{0.1365844952106265^{k}\,r^{k}}{k!}}} , \lim_{n  \rightarrow \infty }{\sum_{k=0}^{n}{\frac{0.140808431699002^{k}\,r^{  k}}{k!}}} , \lim_{n\rightarrow \infty }{\sum_{k=0}^{n}{\frac{  0.1451142866615502^{k}\,r^{k}}{k!}}} , \lim_{n\rightarrow \infty }{  \sum_{k=0}^{n}{\frac{0.1495026295080298^{k}\,r^{k}}{k!}}} , \lim_{n  \rightarrow \infty }{\sum_{k=0}^{n}{\frac{0.1539740213994798^{k}\,r  ^{k}}{k!}}} \right] 
\]
\end{eulerformula}
\begin{eulerprompt}
>function d(n) &= sum(1/(k^2-k),k,2,n); $'d(n)=d(n)
\end{eulerprompt}
\begin{eulerformula}
\[
d\left(n\right)=\sum_{k=2}^{n}{\frac{1}{-k+k^2}}
\]
\end{eulerformula}
\begin{eulerprompt}
>$d(10)=ev(d(10),simpsum=true)
\end{eulerprompt}
\begin{eulerformula}
\[
\sum_{k=2}^{10}{\frac{1}{-k+k^2}}=\frac{9}{10}
\]
\end{eulerformula}
\begin{eulerprompt}
>$d(100)=ev(d(100),simpsum=true)
\end{eulerprompt}
\begin{eulerformula}
\[
\sum_{k=2}^{100}{\frac{1}{-k+k^2}}=\frac{99}{100}
\]
\end{eulerformula}
\eulerheading{Deret Taylor}
\begin{eulercomment}
Deret Taylor suatu fungsi f yang diferensiabel sampai tak hingga di
sekitar x=a adalah:

\end{eulercomment}
\begin{eulerformula}
\[
f(x) = \sum_{k=0}^\infty \frac{(x-a)^k f^{(k)}(a)}{k!}.
\]
\end{eulerformula}
\begin{eulerprompt}
>$'e^x =taylor(exp(x),x,0,10) // deret Taylor e^x di sekitar x=0, sampai suku ke-11
\end{eulerprompt}
\begin{euleroutput}
  Maxima said:
  taylor: 2.858879991940135*r cannot be a variable.
   -- an error. To debug this try: debugmode(true);
  
  Error in:
   $'e^x =taylor(exp(x),x,0,10) // deret Taylor e^x di sekitar x= ...
                               ^
\end{euleroutput}
\begin{eulerprompt}
>$'log(x)=taylor(log(x),x,1,10)// deret log(x) di sekitar x=1
\end{eulerprompt}
\begin{euleroutput}
  Maxima said:
  log: encountered log(0).
   -- an error. To debug this try: debugmode(true);
  
  Error in:
   $'log(x)=taylor(log(x),x,1,10)// deret log(x) di sekitar x=1 ...
                                ^
\end{euleroutput}
\end{eulernotebook}


\newpage
\chapter{KB Pekan 8: Menggunakan EMT untuk Geometri}
\input{Pertemuan 12/Shalih Abdillah_22305144009_ApKom6}

\newpage
\chapter{KB Pekan 10; Menggunakan EMT untuk Statistika}
\input{Pertemuan 14/Shalih Abdillah_22305144009_ApKom7}

\end{document}